% !TeX program = xelatex
\documentclass[11pt,reqno]{amsart}
\usepackage[margin=25mm]{geometry}
\usepackage{amsmath,amssymb,amsthm,mathtools}
\usepackage{enumitem,booktabs,array,xcolor,needspace}
\usepackage[protrusion=true,expansion=false]{microtype}
\usepackage{hyperref}
\allowdisplaybreaks[2]
\raggedbottom
\setlength{\emergencystretch}{2em}
\setlength{\parindent}{1.5em}
\setlength{\parskip}{0pt}
\setlist[enumerate]{leftmargin=2em,itemsep=3pt,topsep=5pt}
\setcounter{tocdepth}{1}
\numberwithin{equation}{section}
\newtheorem{theorem}{Theorem}[section]
\newtheorem{proposition}[theorem]{Proposition}
\newtheorem{lemma}[theorem]{Lemma}
\newtheorem{corollary}[theorem]{Corollary}
\theoremstyle{definition}
\newtheorem{definition}[theorem]{Definition}
\newtheorem{example}[theorem]{Example}
\theoremstyle{remark}
\newtheorem{remark}[theorem]{Remark}
\newcommand{\R}{\mathbb R}
\newcommand{\C}{\mathbb C}
\newcommand{\Z}{\mathbb Z}
\newcommand{\dd}{\mathrm d}
\newcommand{\g}{\mathfrak g}
\newcommand{\Ia}{I_{\mathrm a}}
\newcommand{\If}{I_{\mathrm f}}
\DeclareMathOperator{\tr}{tr}
\DeclareMathOperator{\ad}{ad}
\DeclareMathOperator{\Vol}{Vol}
\DeclareMathOperator{\supp}{supp}
\title[Yang--Mills: local stationary analysis]
{Four-Dimensional Yang--Mills Theory:\\
Exact Blocking and Local Stationary Analysis}
\author{}
\date{18 November 2026}
\hypersetup{
  pdftitle={Reading Notes on Four-Dimensional Yang--Mills Theory: Rewritten Working Exposition},
  pdfauthor={Zexi He (maintainer)},
  pdfsubject={AI-generated rewritten working exposition; publication copy of 21 November 2026, research snapshot of 20 November 2026},
  pdfkeywords={Yang--Mills theory, reading notes, AI assistance}
}

\begin{document}

% Unnumbered publication information; the research text restarts at page one.
\hypersetup{pageanchor=false}
\pagenumbering{gobble}
\begingroup
\thispagestyle{empty}
\pdfbookmark[0]{Publication information}{publication-information}
\begin{center}
\vspace*{20mm}
{\LARGE\bfseries Reading Notes on\\[4pt]Four-Dimensional Yang--Mills Theory\par}
\vspace{8mm}
{\Large An AI-generated Rewritten Working Exposition\par}
\vspace{3mm}
{\large Exact Blocking and Local Stationary Analysis\par}
\vspace{10mm}
Maintained by Zexi He\par
\medskip
Publication copy --- 21 November 2026\par
\end{center}
\vspace{7mm}
\begin{quote}
\small
\textbf{Status statement.}
This substantially AI-generated rewritten working exposition is maintained
for personal study and discussion. It redevelops arguments from the
companion main exposition \cite{Main}, with an emphasis on motivation,
detailed calculations and explicit proof dependencies. The current
draft develops exact blocking, local stationary analysis and local
integration, and reaches classical matching through quartic order,
including the one-loop quadratic calculation.

This is a work in progress. It has not undergone an exhaustive
independent mathematical review or been mechanically checked against
Lean. The maintainer's study and questions do not constitute a complete
independent verification. It is not an official exposition or an
endorsement by OpenAI or the source contributors, and it does not claim
a proof of a positive mass gap.

This publication copy uses the rewritten research manuscript as
available on 20 November 2026. Its mathematical text, abstract,
numbered statements and printed pagination are retained. The initial
publication information page is unnumbered. The reference [Main] points
to the companion publication of 21 November 2026, based on the
17 November research snapshot. Further provenance and review information
appears in the repository README.
\end{quote}
\endgroup
\clearpage
\pagenumbering{arabic}
\hypersetup{pageanchor=true}

\begin{abstract}
We begin a detailed proof exposition with the finite Wilson measure
and the problem of integrating one scale without changing its
correlations. We construct the regular and axial branches of the
coarse kernel, including their finite probabilities, and prove its
gauge covariance. A normalised coarse kernel transports the density and
every inserted product with the same partition function. Root-fixed
gauge invariance permits an exact tree-gauge reduction. A partition
of unity then puts selected links in exponential coordinates while
retaining the remaining Haar integrals. We derive the complete
Gaussian representation before choosing a stationary centre.
The linear and quadratic terms explain that choice; the resulting
central formula identifies the stationary, determinant and tail
estimates needed next. All identities retain their domains,
branch weights, Jacobians and normalising factors. The local Haar
formula is proved in an appendix. For the specified principal path
action, we derive a volume-uniform stationary graph and two-sided
Hessian bounds from a block Poincar\'e estimate, local derivative
budgets and a weighted inverse estimate. A finite-path reconstruction
then gives a relative tail bound on auxiliary tori, including the full
Gaussian dimension cost. We retain the physical fine-field indicator
and prove its central compatibility near the identity. Rooted
analytic coefficient bounds then localise the stationary action,
loop and Haar terms, and quadratic Gaussian determinant. Their
connected footprints pay the covering weight, and Taylor
subtraction gives a volume-uniform local remainder. A local
gauge construction transfers bounded stencils to physical good
fields. Covariant path records then give compatible extensions
for unbounded footprint families. A bounded-coordinate Gaussian
reference admits a convergent connected logarithm with a full
fourth-order remainder; its difference from the actual central
integral is estimated on auxiliary volumes. A calibrated Gaussian
history gives a fine-curvature gain before choosing the block side.
Comparing its affine centre directly with the stationary graph
proves the central inclusion at a general good auxiliary background.
A selected-window covariance then gives an exact physical
representation with finite reads for bounded Borel incoming
activities. We prove the two masked restorations and their exact
forced merger under explicit absolute-integral and incidence
criteria. Completing the row lift on the full curvature-row
space and retaining the actual positive-likelihood branches give
a physical energy charge per hole seed. The annulus comparison
keeps the negative affine squares with that energy. This proves
a joint Gaussian--Haar packet bound, paying the frozen-link
normalisation and shared incoming occurrences without copying
their variables. Physical Haar variables yield ordered packet and central
derivative estimates. A local affine-Hessian subtraction and
cover compression give sharp incoming regular transport, with
its constant fixed before the block ratio. We identify the first
two ordinary strength coefficients with their retained chart
factors, and compute the classical cubic and quartic jets and
the complete one-loop quadratic jet, including its Haar term.
We match the classical action through degree four by transported
curvatures with path budgets fixed before the block ratio and
bound the resulting local classical remainder.
\end{abstract}
\maketitle
\tableofcontents

\section{The finite problem and the first useful calculation}
\label{sec:finite}

The eventual problem is to obtain continuum correlations from one
family of lattice probability measures, first removing the mesh at
fixed physical side length and then taking a suitable volume
subsequence. Before estimating such a limit, we need an operation
which integrates fine variables while preserving the original
integrals. We shall build this operation and then examine the
integral it actually leaves behind. The broader construction target
is stated in \cite[Theorem~4.1, p.~20]{Main}.

\subsection{What is integrated}

Fix a closed subgroup \(G\subset U(N)\), in its given faithful matrix
representation. The integer \(N\) is fixed throughout. Its real Lie
algebra is
\[
 \g=\{X\in\operatorname{Mat}_N(\C):e^{tX}\in G
                                      \text{ for every }t\in\R\}.
\]
The group data include that this set is a real vector space. We retain
the nonabelian simplicity condition of the final construction: the
bracket \([X,Y]=XY-YX\) is nonzero and its only real Lie ideals are
\(\{0\}\) and \(\g\). Here a Lie ideal is a subspace
\(\mathfrak i\) with \([\g,\mathfrak i]\subset\mathfrak i\).
Simplicity will be used for the flat gauge identities in
Section~\ref{sec:affine-local-reset}.

Since \(G\) is closed in the compact group \(U(N)\), it is compact.
Write \(h_G\) for its probability Haar measure. For every integrable
\(f:G\to\C\) and \(R\in G\),
\begin{equation}
 \int_G f(RU)\,\dd h_G(U)
 =\int_G f(U)\,\dd h_G(U)
 =\int_G f(UR)\,\dd h_G(U).
 \label{eq:haar-invariance}
\end{equation}
The right-invariance follows from uniqueness of left Haar probability:
its right translate is another left-invariant probability measure.

For a positive integer \(\ell\), let
\begin{equation}
 \Lambda_\ell=(\Z/\ell\Z)^4,\qquad
 E_\ell^+=\Lambda_\ell\times\{1,2,3,4\},\qquad
 \mathcal U_\ell=G^{E_\ell^+}.
 \label{eq:lattice}
\end{equation}
The link \((x,i)\) starts at \(x\) and ends at \(x+e_i\), where
\(e_i\) is the \(i\)-th coordinate vector. A configuration
\(U\in\mathcal U_\ell\) assigns a matrix \(U_e=U(x,i)\in G\) to each
positive link \(e=(x,i)\); its reverse is denoted by \(e^{-1}\), with
\(U_{e^{-1}}=U_e^{-1}\). All site
arithmetic is periodic. The reference measure is
\[
 H_\ell=\bigotimes_{e\in E_\ell^+}h_G.
\]
There are \(4\ell^4\) independent positive links.

The smallest closed lattice path is an elementary square. Its ordered
product measures the discrepancy between the two ways of reaching
the opposite corner. This leads to the following energy.

A based oriented path \(\gamma=e_1\cdots e_k\) is an ordered
sequence of composable oriented edges with a specified starting
vertex. Its holonomy is
\[
                  U_\gamma:=U_{e_1}\cdots U_{e_k}.
\]
The empty path has holonomy \(I_N\), and
\(U_{\gamma^{-1}}=U_\gamma^{-1}\).
For a based oriented plaquette \(p\), write \(U_p:=U_{\partial p}\),
where \(\partial p\) is its oriented boundary starting at its base point.
The same conventions apply to any other configuration, giving
\(V_\gamma,V_p\), for example.

\begin{definition}[Plaquette holonomy and Wilson energy]
Let \(\mathcal P_\ell^+=\{(x,i,j):x\in\Lambda_\ell,\ 1\le i<j\le4\}\).
For \(p=(x,i,j)\in\mathcal P_\ell^+\), fix the base point \(x\) and set
\begin{align}
 U_p&=U(x,i)U(x+e_i,j)
                         U(x+e_j,i)^{-1}U(x,j)^{-1},\label{eq:plaquette}\\
 q_G(W)&=1-\frac1N\operatorname{Re}\tr W,\qquad W\in G,\label{eq:q}\\
 W_\ell(U)&=\sum_{p\in\mathcal P_\ell^+}q_G(U_p).
 \label{eq:wilson-energy}
\end{align}
\end{definition}
Changing the starting vertex conjugates the plaquette holonomy.
Its trace and Wilson price are unchanged; an untraced curvature
factor retains the chosen base point.

For a complex matrix \(B\), write
\(\|B\|_{\rm HS}^2=\tr(B^*B)\). Unitarity gives
\begin{align*}
 \|W-I_N\|_{\rm HS}^2
 &=\tr\bigl((W^*-I_N)(W-I_N)\bigr)\\
 &=\tr(2I_N-W-W^*)
   =2N-2\operatorname{Re}\tr W.
\end{align*}
Consequently
\begin{equation}
 q_G(W)=\frac1{2N}\|W-I_N\|_{\rm HS}^2,
 \qquad
 W_\ell(U)=\frac1{2N}\sum_p\|U_p-I_N\|_{\rm HS}^2.
 \label{eq:energy-square}
\end{equation}
Thus zero plaquette energy means that its product is the identity.
Also \(|\tr W|\le N\), so \(0\le q_G\le2\). As there are
\(6\ell^4\) positive plaquettes,
\begin{equation}
                    0\le W_\ell(U)\le12\ell^4.
 \label{eq:energy-bound}
\end{equation}

For \(\beta>0\), the Wilson measure is
\begin{equation}
 \dd\mu_{\ell,\beta}(U)
   =Z_\ell(\beta)^{-1}e^{-\beta W_\ell(U)}\,\dd H_\ell(U),
 \qquad
 Z_\ell(\beta)=\int e^{-\beta W_\ell}\,\dd H_\ell.
 \label{eq:wilson-measure}
\end{equation}
This is a probability measure, since \eqref{eq:energy-bound} implies
\begin{equation}
             e^{-12\beta\ell^4}\le Z_\ell(\beta)\le1.
 \label{eq:partition-bound}
\end{equation}
Continuous observables are integrable because \(\mathcal U_\ell\)
is compact. These are finite-dimensional statements.
The lattice formulation originates with Wilson \cite{Wilson74};
the normalisation here is the one in \cite[Section~2, pp.~13--16]{Main}.

To relate lattice and physical lengths, write
\begin{equation}
              \ell=\frac{M}{\varepsilon},\qquad
              \#\Lambda_\ell=\ell^4,\qquad \Vol=M^4,
 \label{eq:physical-volume}
\end{equation}
with an admissible integer \(\ell\). The physical side length is
\(M\), and one lattice step has length \(\varepsilon\). At terminal
UV level \(n\), these are \(\varepsilon=\varepsilon_n\),
\(\ell_{n,M}=M/\varepsilon_n\), and a coefficient \(\beta_{n,M}\)
still to be selected by the construction. Even at fixed \(M\),
\(\ell_{n,M}^4\) grows when \(\varepsilon_n\downarrow0\).
The lower bound \eqref{eq:partition-bound} therefore gives no
cutoff-uniform control of normalised correlations.

\smallskip\noindent
\emph{Parameter choices.}
For \(a\ge0\) and \(b>0\), we write \(a\ll b\) when \(a/b\) is sufficiently small in
terms of the previously fixed data specified in the estimate.
The threshold is uniform in every variable over which that
estimate is declared uniform. Constants entering a smallness
condition are fixed before the choice that enforces it.
This notation does not assert \(a/b\to0\); a limit is stated
explicitly when vanishing is required. We collect the
constraints before making a choice, and keep prescribed
profile functions distinct from quantities that remain free.

\smallskip\noindent
\emph{Implicit constants.}
For nonnegative quantities \(A,B\), the notation \(A\lesssim B\)
means \(A\le CB\) for a positive constant \(C\).
We write \(A\gtrsim B\) for \(B\lesssim A\), and
\(A\simeq B\) when both comparisons hold.
At a fixed profile these constants may depend on its fixed
data, including the block ratio when it has already been
chosen, but are independent of the physical period and
terminal cutoff. They are uniform in the variables of an
estimate unless a dependence is specified. Any further stated
uniformity remains in force, and relevant extra dependence
is displayed by a subscript. In particular a bound obtained before choosing
the block ratio has a constant independent of that ratio.
Displayed coupling-dependent allowances and budgets remain
explicit.

For self-adjoint matrices these comparisons are in
quadratic-form order: \(H\gtrsim1\) means
\(\langle u,Hu\rangle\gtrsim\|u\|_2^2\) for every \(u\),
with one common positive constant. The notation \(H\simeq1\)
requires both lower and upper bounds. Named gaps and
constants are retained when their size determines a later
radius, decay rate or contraction.

\subsection{Why gauge fixing will enter an integral rather than the measure}

A lattice gauge transformation is a map \(h:\Lambda_\ell\to G\),
acting by
\begin{equation}
             U^h(x,i)=h(x)U(x,i)h(x+e_i)^{-1}.
 \label{eq:gauge}
\end{equation}
An observable \(O\) is gauge invariant when \(O(U^h)=O(U)\) for all
\(h,U\). Along a path, the matrices at consecutive internal vertices
cancel. In particular,
\begin{align*}
 (U^h)_p
 &=h(x)U(x,i)\underline{h(x+e_i)^{-1}h(x+e_i)}
       U(x+e_i,j)\\
 &\quad\cdot\underline{h(x+e_i+e_j)^{-1}h(x+e_i+e_j)}
       U(x+e_j,i)^{-1}\\
 &\quad\cdot\underline{h(x+e_j)^{-1}h(x+e_j)}
       U(x,j)^{-1}h(x)^{-1}\\
 &=h(x)U_ph(x)^{-1}.
\end{align*}
Trace invariance gives \(W_\ell(U^h)=W_\ell(U)\). For a fixed
\(h\), each positive link undergoes a fixed left and right
translation, so \eqref{eq:haar-invariance} gives
\begin{equation}
 \int O(U^h)\,\dd\mu_{\ell,\beta}(U)
                    =\int O(U)\,\dd\mu_{\ell,\beta}(U)
 \label{eq:measure-invariance}
\end{equation}
for every integrable \(O\).

The finite measure is already well defined. Gauge fixing will instead
be used to remove redundant coordinates from a particular blocking
integral. The symmetry that must be preserved there will involve the
coarse variables as well as the Wilson density.

\section{Preserving correlations while integrating one scale}
\label{sec:transport}

Suppose we introduce a coarser connection \(V\). To preserve the
original integrals, it is enough to assign, for each \(U\), a
probability distribution of \(V\). The operation will act on the
unnormalised density; its total integral can then be kept exactly.

Let \((\mathcal U_f,H_f)\) and \((\mathcal U_c,H_c)\) be the fine
and coarse configuration spaces with their product Haar probabilities.

\begin{definition}[Normalised blocking kernel]
A blocking kernel is a nonnegative Borel function \(K(U,V)\) such that
\begin{equation}
                \int_{\mathcal U_c}K(U,V)\,\dd H_c(V)=1
                \quad\text{for every }U.
 \label{eq:kernel-normalisation}
\end{equation}
For \(F\in L^1(H_f)\), its pushed density is
\begin{equation}
         (\mathcal PF)(V)=\int_{\mathcal U_f}K(U,V)F(U)\,\dd H_f(U),
 \label{eq:push}
\end{equation}
defined for almost every \(V\).
\end{definition}

\begin{proposition}[Exact transport of an integral]
\label{prop:transport}
For every complex \(F\in L^1(H_f)\),
\begin{equation}
 \|\mathcal PF\|_{L^1(H_c)}\le\|F\|_{L^1(H_f)},\qquad
 \int\mathcal PF\,\dd H_c=\int F\,\dd H_f.
 \label{eq:mass-preservation}
\end{equation}
\end{proposition}
\begin{proof}
The nonnegative integrand \(K(U,V)|F(U)|\) permits Tonelli's theorem:
\begin{align*}
 \int_{\mathcal U_c}\int_{\mathcal U_f}
              K(U,V)|F(U)|\,\dd H_f(U)\dd H_c(V)
 &=\int_{\mathcal U_f}|F(U)|
            \left(\int_{\mathcal U_c}K(U,V)\,\dd H_c(V)\right)
                                                     \dd H_f(U)\\
 &=\int_{\mathcal U_f}|F(U)|\,\dd H_f(U)<\infty.
\end{align*}
Hence \(\mathcal PF\) exists almost everywhere and the joint
integrand with \(F\) is absolutely integrable. The triangle
inequality inside \eqref{eq:push} proves the norm estimate.
Fubini's theorem then applies to \(KF\) itself, giving
\[
 \int\mathcal PF\,\dd H_c
 =\int F(U)\left(\int K(U,V)\,\dd H_c(V)\right)\dd H_f(U)
 =\int F\,\dd H_f.
\]
\end{proof}

To see what this means for correlations, take \(F\ge0\) with
\(Z=\int F\,\dd H_f>0\), and an observable \(O\) satisfying
\(OF\in L^1(H_f)\). Applying the proposition separately to \(OF\)
and \(F\) gives
\begin{equation}
 \frac{\int\mathcal P(OF)\,\dd H_c}
      {\int\mathcal PF\,\dd H_c}
                  =\frac{\int OF\,\dd H_f}{Z}.
 \label{eq:correlation-preservation}
\end{equation}
There is one denominator, namely the original partition function.
Where \(\mathcal PF(V)>0\), the ratio
\(\mathcal P(OF)(V)/\mathcal PF(V)\) is the conditional expectation
of \(O(U)\) given \(V\) under the joint probability
\(Z^{-1}F(U)K(U,V)\,\dd H_f(U)\dd H_c(V)\).
For \(O=\prod_{j=1}^rO_j\), the transported numerator is
\(\mathcal P(F\prod_jO_j)\). The conditional expectation of a product
need not factor into the product of conditional expectations.

This argument can be repeated. At fixed \(n,M\), set
\[
 F_n=e^{-\beta_{n,M}W_{\ell_{n,M}}},\qquad
 F_m=\mathcal P_mF_{m+1}\quad(m=n-1,\ldots,0).
\]
Here each \(\mathcal P_m\) has a normalised kernel on its own fine
and coarse spaces. Descending induction in
\eqref{eq:mass-preservation} yields
\begin{equation}
          \int F_m\,\dd H_m
             =Z_{\ell_{n,M}}(\beta_{n,M})\quad(0\le m\le n).
 \label{eq:history-mass}
\end{equation}
The same proof transports the full inserted product through the same
history. It remains to choose a useful kernel and obtain local
estimates for these exact densities. Mass preservation alone imposes
no useful local form on \(F_m\).

\section{A coarse kernel whose normalisation is explicit}
\label{sec:kernel}

In one step let the fine side be \(f=Lc\), with \(L\ge3\) odd.
Both \(f,c\) are lattice lengths. A coarse link from \(x\) to
\(x+e_i\) corresponds to fine endpoints \(Lx,Lx+Le_i\).
We want a predicted coarse connection \(Q(U)\), and then allow the
actual \(V\) to fluctuate around it. The mismatch on a coarse link
is the group element \(V_eQ(U)_e^{-1}\). Its Wilson price
\(q_G(V_eQ(U)_e^{-1})\) vanishes precisely at \(V_e=Q(U)_e\).

\subsection{From one prediction to the branch sum}

The simplest example is the straight-path product
\begin{equation}
 Q(U)_i(x)=U(Lx,i)U(Lx+e_i,i)\cdots
                                U(Lx+(L-1)e_i,i).
 \label{eq:straight-predictor}
\end{equation}
It explains both the endpoints of a coarse link and the covariance
required below. The physical averaging maps in the full construction
have further local choices. We retain their branch representation
from the outset; the straight-path map is an example, rather than a
replacement for that averaging procedure.

Fix a finite index set \(\mathcal T\), Borel maps
\(Q_t:\mathcal U_f\to\mathcal U_c\), and Borel weights
\begin{equation}
       \lambda_t(U)\ge0,\qquad \sum_{t\in\mathcal T}\lambda_t(U)=1.
 \label{eq:branch-data}
\end{equation}
These data are fixed independently of the density being pushed.
A branch records one consistent collection of choices over the whole
coarse lattice; shared sites use shared choices. The normalisation
calculation needs no derivatives of these maps or weights.
Section~\ref{sec:actual-branches} constructs the specific maps and
weights used here. Until that construction, the following algebraic
calculation applies to any Borel data satisfying
\eqref{eq:branch-data}.

Fix \(a>0\). For \(\beta>0\), define
\begin{equation}
 z_{\mathrm n}(\beta)=\int_Ge^{-a\beta q_G(W)}\,\dd h_G(W),
 \qquad b_c=\#E_c^+=4c^4.
 \label{eq:noise-normaliser}
\end{equation}
Since \(q_G\in[0,2]\),
\begin{equation}
                     e^{-2a\beta}\le z_{\mathrm n}(\beta)\le1.
 \label{eq:noise-bounds}
\end{equation}
The normalised kernel is
\begin{equation}
 K_\beta(U,V)=z_{\mathrm n}(\beta)^{-b_c}
       \sum_{t\in\mathcal T}\lambda_t(U)
       \exp\left\{-a\beta\sum_{e\in E_c^+}
                         q_G(V_eQ_t(U)_e^{-1})\right\}.
 \label{eq:physical-kernel}
\end{equation}
This is the kernel of \cite[(11.1)--(11.2), p.~36]{Main}.

\begin{proposition}[Normalisation and finite-cutoff integrability]
\label{prop:kernel}
The kernel \eqref{eq:physical-kernel} is Borel, strictly positive,
bounded, and normalised as in \eqref{eq:kernel-normalisation}.
In particular its push is defined at every \(V\) for each
\(F\in L^1(H_f)\).
\end{proposition}
\begin{proof}
All factors are Borel. Each branch exponential is positive and at
least one branch has positive weight. For fixed \(U,t,e\), change
variables \(W=V_eQ_t(U)_e^{-1}\). Right Haar invariance gives
\[
 \int_Ge^{-a\beta q_G(V_eQ_t(U)_e^{-1})}\,\dd h_G(V_e)
                         =z_{\mathrm n}(\beta).
\]
The integral over the \(b_c\) independent coarse links is therefore
\(z_{\mathrm n}(\beta)^{b_c}\). Thus
\[
 \int K_\beta(U,V)\,\dd H_c(V)
 =z_{\mathrm n}(\beta)^{-b_c}
      \sum_t\lambda_t(U)z_{\mathrm n}(\beta)^{b_c}=1.
\]
For boundedness, the branch exponentials are at most one, so
\(K_\beta(U,V)\le z_{\mathrm n}(\beta)^{-b_c}\).
Hence \(\int K_\beta(U,V)|F(U)|\,\dd H_f(U)<\infty\) for every
\(V\).
\end{proof}

The normaliser is outside the complete branch sum. The weights
\(\lambda_t(U)\) remain inside the fine integral. Their sum is one
before the branch-dependent exponential is inserted; this does not
permit removing them from \eqref{eq:physical-kernel}.

\subsection{The action left by a single push}

At the initial Wilson level, inserting \eqref{eq:physical-kernel}
into \eqref{eq:push} gives the exact expression
\begin{align}
 (\mathcal P_\beta e^{-\beta W_f})(V)
 &=z_{\mathrm n}(\beta)^{-b_c}
       \sum_t\int\lambda_t(U)e^{-\beta A_t(V,U)}\,\dd H_f(U),
       \label{eq:initial-block-integral}\\
 A_t(V,U)&=W_f(U)+a\sum_{e\in E_c^+}q_G(V_eQ_t(U)_e^{-1}).
       \label{eq:initial-action}
\end{align}
The notation \(W_f\) here is the Wilson energy at fine side \(f\).
The kernel has added a mismatch penalty to the fine action.
This sum, rather than the Wilson action alone, is the exponent we
will later expand.

Subsequent densities contain more terms. To keep the same calculation
available, write an input density in the form
\begin{equation}
       F(U)=e^{vf^4}\mathcal B(U)e^{-\beta S_f(U)-I_f(U)}.
 \label{eq:general-input}
\end{equation}
Here \(v\in\R\), \(S_f,I_f\) are real Borel functions, and
\(\mathcal B\) is a complex Borel multiplier; the expression is
assumed integrable. This is a factorisation of the given density,
not an assertion about estimates on its factors. Initially
\(v=0\), \(S_f=W_f\), \(I_f=0\), \(\mathcal B=1\).
Direct substitution now gives
\begin{align}
 (\mathcal P_\beta F)(V)
 &=e^{vf^4}z_{\mathrm n}(\beta)^{-b_c}
       \sum_t\int\lambda_t(U)\mathcal B(U)
                    e^{-\beta A_t(V,U)-I_f(U)}\,\dd H_f(U),
       \label{eq:full-block-integral}\\
 A_t(V,U)&=S_f(U)+a\sum_{e\in E_c^+}q_G(V_eQ_t(U)_e^{-1}).
       \label{eq:general-action}
\end{align}
Any cutoff or incoming correction placed in \(\mathcal B\) remains
in this integral. The finite-cutoff integrability of the whole
expression follows from Proposition~\ref{prop:kernel}.

\section{Constructing the actual branch maps}
\label{sec:actual-branches}

The kernel is already normalised for arbitrary branch data.
We now have to choose those data so that the coarse variable reads
an average of nearby fine transports, transforms as a connection,
and remains defined at every configuration. These requirements
explain the two stages of the construction. First choose a common
frame about each coarse root. Then compare the neighbouring frames
using a family of parallel connectors. Where a minimiser is unique,
it supplies the average. Elsewhere a finite family of explicit
transports supplies the fallback branches. All their probabilities
will be part of the original kernel.

\subsection{A matrix distance and the fixed patch}

We use the normalised Frobenius convention of
\cite[(7.1)--(7.3), p.~24]{Main}:
\begin{equation}
 \|A\|_{\rm F}^2=\frac1N\operatorname{Re}\tr(A^*A),\qquad
 \langle X,Y\rangle_{\g}=\frac1N\operatorname{Re}\tr(X^*Y),
 \qquad |X|_{\g}=\|X\|_{\rm F}.
 \label{eq:physical-norm}
\end{equation}
The matrix norm is defined on all complex matrices, regarded as a
real Euclidean space. On \(\g\) the second expression is its inner
product. In terms of the Hilbert--Schmidt norm already defined,
\(\|A\|_{\rm F}=N^{-1/2}\|A\|_{\rm HS}\).
All Lie algebra Lebesgue measures and Haar chart constants below
use the normalised inner product \eqref{eq:physical-norm}.
Define the chord from the group identity by
\begin{equation}
                  d_G(W)=\|W-I_N\|_{\rm F}=\sqrt{2q_G(W)}.
 \label{eq:chord}
\end{equation}
Left and right multiplication by a unitary matrix preserve this
norm. In particular
\begin{equation}
 \begin{gathered}
 \|U-V\|_{\rm F}=d_G(UV^{-1}),\quad
 d_G(RWR^{-1})=d_G(W),\quad d_G(W^{-1})=d_G(W),\\
 d_G(UV)\le d_G(U)+d_G(V),\qquad 0\le d_G(W)\le2.
 \end{gathered}
 \label{eq:chord-properties}
\end{equation}
For the triangle inequality, write \(UV-I=(U-I)V+(V-I)\)
and use right invariance of the norm.

Fix an integer \(J\ge4\), and set
\begin{equation}
       R=(J+3)L,\qquad
       \mathcal Q_R=\{z\in\Z^4:|z|_\infty\le R\},\qquad
       s=|\mathcal Q_R|.
 \label{eq:branch-patch}
\end{equation}
For \(z,z'\in\mathcal Q_R\) and an axis permutation
\(\pi\in S_4\), let \(\gamma^\pi_{z,z'}\) be the path which
changes the four coordinates successively in the order \(\pi\).
Its transport \(U_\gamma\) uses the oriented-path convention above.
Reversing a path gives the inverse transport.
All paths are read on the integer cover with the given periodic
fine connection. Translates of the same prescribed patch are used
at every coarse root.

\subsection{The frame score and the regular test}

A frame on the patch is a collection \(h_z\in G\), with \(h_0=1\).
The framed transport from \(z\) to \(z'\) is
\(h_zU_{\gamma^\pi_{z,z'}}h_{z'}^{-1}\).
To make these transports as close to the identity as possible,
minimise their mean Wilson price:
\begin{equation}
 \mathcal A_U(h)=\frac1{s^2\,4!}
       \sum_{z,z'\in\mathcal Q_R}\sum_{\pi\in S_4}
       q_G\bigl(h_zU_{\gamma^\pi_{z,z'}}h_{z'}^{-1}\bigr).
 \label{eq:actual-frame-score}
\end{equation}
This continuous function has a minimiser on the compact rooted
frame space. Its uniqueness is the issue: an arbitrary choice
among several minimisers need not have the required covariance.

For a reference axial path from the root, put
\(a_z^{0,\pi}=U_{\gamma^\pi_{0,z}}\). The regular frame test is
\begin{equation}
 \max_{\pi,\pi',z,z'}
       d_G\bigl(a_z^{0,\pi}U_{\gamma^{\pi'}_{z,z'}}
                              (a_{z'}^{0,\pi})^{-1}\bigr)<c_{\rm fr}.
 \label{eq:regular-frame-test}
\end{equation}
It says that the route matrices become uniformly close to the
identity after any of the indicated root-based axial changes.
This uniform control will bound every feedback matrix in the
uniqueness proof below; an average score bound alone would not
give that pointwise control at each frame variable.
We shall choose the positive threshold \(c_{\rm fr}\) using \(G\)
alone. These are exactly the score and test in
\cite[(14.2) and (14.4), p.~56]{Main}.

The following elementary nearest-point fact is proved in
Appendix~\ref{app:projection}. The map \(\pi_G(A)\) denotes the
unique group element closest to a matrix \(A\), when that nearest
point is unique. There is a uniform neighbourhood of \(G\) on
which it is real analytic. Near \(I_N\) it satisfies
\begin{equation}
 \pi_G(I_N)=I_N,\qquad D\pi_G(I_N)=\operatorname{pr}_{\g},
 \label{eq:projection-derivative}
\end{equation}
where \(\operatorname{pr}_{\g}\) is orthogonal projection in the
real matrix space. We can choose a ball about \(I_N\) on which
\begin{equation}
 \|\pi_G(A)-\pi_G(B)-\operatorname{pr}_{\g}(A-B)\|_{\rm F}
                              \le\tfrac14\|A-B\|_{\rm F}.
 \label{eq:projection-small-variation}
\end{equation}
The small coefficient will give uniqueness without a bound
depending on the number of frame variables.

\begin{lemma}[Uniqueness on the regular frame test]
\label{lem:regular-frame-unique}
For a sufficiently small \(c_{\rm fr}>0\), depending only on
the fixed matrix group, the minimiser of
\eqref{eq:actual-frame-score} is unique whenever
\eqref{eq:regular-frame-test} holds. It is continuous on that
open set of fine configurations.
\end{lemma}
\begin{proof}
Choose one axis order in \eqref{eq:regular-frame-test} and apply
its root-based axial gauge. If \(a_z\) is that axial transport,
the route matrices in this gauge are
\[
       W_{ij,\pi}=a_iU_{\gamma^\pi_{i,j}}a_j^{-1},
       \qquad \|W_{ij,\pi}-I_N\|_{\rm F}<c_{\rm fr}.
\]
The indices \(i,j\) here denote patch vertices. Changing frames
by \(h_i\mapsto h_i a_i^{-1}\) preserves the score and the
root constraint. It therefore suffices to prove uniqueness
for these near-identity routes. Write \(d=c_{\rm fr}\).
The routes satisfy \(W_{ii,\pi}=I_N\), and reversing a route
gives \(W_{ji,\bar\pi}=W_{ij,\pi}^{-1}\), with \(\bar\pi\)
a permutation of the route labels.

First minimise over all frames, without pinning the root.
Common left multiplication changes every framed route by
conjugation and leaves the score unchanged. Every frame can
therefore be root-normalised without changing its score, so
the unrestricted and rooted minima coincide. Use
\(\operatorname{av}\) for the uniform average over the displayed
finite indices. Comparison with a constant frame gives
\[
 \operatorname{av}_{i,j,\pi}\|h_iW_{ij,\pi}-h_j\|_{\rm F}^2
                                                     \le d^2.
\]
Since \(\|h_i(W_{ij,\pi}-I_N)\|_{\rm F}\le d\),
\begin{align*}
 \operatorname{av}_{i,j}\|h_i-h_j\|_{\rm F}^2
 &\le2\operatorname{av}_{i,j,\pi}
       \bigl(\|h_iW_{ij,\pi}-h_j\|_{\rm F}^2
                      +\|h_i(W_{ij,\pi}-I_N)\|_{\rm F}^2\bigr)\\
 &\le4d^2.
\end{align*}
For \(\bar h=\operatorname{av}_i h_i\), the variance identity
gives \(\operatorname{av}_i\|h_i-\bar h\|_{\rm F}^2\le2d^2\).
Hence \(\operatorname{dist}(\bar h,G)\le\sqrt2d\).
For small enough \(d\), its nearest group point \(b\) is unique.
Replace \(h_i\) by \(b^{-1}h_i\). The resulting minimiser satisfies
\begin{equation}
       \pi_G(\bar h)=I_N,\qquad \|\bar h-I_N\|_{\rm F}\le2d.
 \label{eq:centred-frame-mean}
\end{equation}

For such a minimiser define its feedback matrices by
\[
                 A_i(h)=\operatorname{av}_{j,\pi}
                                      h_jW_{ij,\pi}^{-1}.
\]
Coordinate minimality and route-reversal symmetry imply that
\(h_i\) is a nearest group point to the feedback with \(j\ne i\).
Indeed the two ordered contributions involving \(i\) have the
same real inner product with the proposed new value at \(i\).
Their sum is minimised by minimising squared distance to the
average of \(h_jW_{ij,\pi}^{-1}\) over \(j\ne i\).
Including the diagonal adds the term \(s^{-1}h_i\) to the
feedback. For a competitor \(v\in G\), the difference of its
squared distance and that of \(h_i\) becomes the old nonnegative
difference times \((s-1)/s\), plus
\(s^{-1}\|v-h_i\|_{\rm F}^2\). Thus \(h_i\) is nearest to
\(A_i(h)\) as well. Moreover
\[
    \|A_i(h)-\bar h\|_{\rm F}
       \le\operatorname{av}_{j,\pi}
                  \|h_j(W_{ij,\pi}^{-1}-I_N)\|_{\rm F}\le d.
\]
Consequently \(A_i(h)\) is within \(3d\) of \(I_N\), where the
nearest point is unique, and \(h_i=\pi_G(A_i(h))\).

Centre two unrestricted minimisers \(h,k\) by
\eqref{eq:centred-frame-mean}, and put
\(u=\max_i\|h_i-k_i\|_{\rm F}\).
Nearest-point stationarity at \(I_N\) implies
\(\operatorname{pr}_{\g}(\bar h-\bar k)=0\).
The feedbacks satisfy
\begin{align*}
 \|A_i(h)-A_i(k)\|_{\rm F}&\le u,\\
 A_i(h)-A_i(k)-(\bar h-\bar k)
       &=\operatorname{av}_{j,\pi}
                    (h_j-k_j)(W_{ij,\pi}^{-1}-I_N).
\end{align*}
For the normalised Frobenius norm,
\(\|AB\|_{\rm F}\le\sqrt N\|A\|_{\rm F}\|B\|_{\rm F}\),
by the corresponding Hilbert--Schmidt product bound. Therefore
the norm of the last displayed average is at most \(\sqrt Ndu\).
Applying \eqref{eq:projection-small-variation} gives
\[
          \|h_i-k_i\|_{\rm F}\le(\sqrt N d+\tfrac14)u.
\]
Choose \(d\le(4\sqrt N)^{-1}\) and small enough that the
\(3d\)-ball lies inside the ball of
\eqref{eq:projection-small-variation}. Then \(u\le u/2\), hence
\(u=0\). Every minimiser differs from the same centred minimiser
by a common left multiple. Pinning the root to \(1\) makes that
multiple unique.

Finally the score depends continuously on its finitely many
input links, uniformly over the compact rooted frame space.
For a convergent sequence of regular inputs, each convergent
subsequence of their minimisers minimises the limiting score.
Uniqueness identifies the limit. Compactness then gives continuity
of the minimiser on the regular open set.
\end{proof}

This expands the uniqueness argument in
\cite[Lemma~14.4, pp.~60--61]{Main}, using the projection estimate
\cite[(14.10), p.~59]{Main}. The factor \(\sqrt N\) in the matrix
product bound has been retained explicitly. Since the representation
is fixed, it affects only the allowed group-dependent threshold.
The stronger first-power mean estimate needed for rough-field
packet energies is proved in Appendix~\ref{app:packet-branches};
it concerns these same branches and probabilities.

\subsection{An axial branch at every configuration}

Off the regular test we need explicit frames whose probabilities
remain defined. For a tag \(t=(p,\pi)\in\mathcal Q_R\times S_4\),
put
\begin{equation}
       a_z^t=U_{\gamma^\pi_{p,z}},\qquad
       h_z^t=(a_0^t)^{-1}a_z^t.
 \label{eq:axial-frame}
\end{equation}
The second expression pins the frame to \(h_0^t=1\).
Later the frame must turn plaquette control into link control
before entering a coordinate chart. This suggests bounding the
sum of framed link chords by the patch length times a sum of
plaquette chords. First powers are natural here: changing a link
across one square costs at most that square's chord, by the
triangle inequality. The layer argument below supplies a suitable
base point and proves the bound used in the original kernel.
For an integer box or a coordinate slice \(X\), define
\begin{equation}
 \begin{split}
 \operatorname{edge}_X(U)&=\sum_{z,i:\ z,z+e_i\in X}d_G(U_{z,i}),\\
 \operatorname{face}_X(U)&=\sum_{p:\ p\subset X}d_G(U_p).
 \end{split}
 \label{eq:axial-costs}
\end{equation}
The second sum counts positive plaquettes with all four vertices
in \(X\); it is not the Wilson energy, whose summands are half
the squared chords. Fix \(a_{\rm ax}\ge16\), independently of
\(L,J\), the cutoffs and the period. A tag is eligible when
\begin{equation}
 \operatorname{edge}_{\mathcal Q_R}(U^{a^t})
      \le a_{\rm ax}(2R+1)\operatorname{face}_{\mathcal Q_R}(U).
 \label{eq:axial-eligible}
\end{equation}
Denote the set of eligible tags by \(\mathcal E(U)\).
Changing \(a^t\) to \(h^t\) conjugates every framed link by the
same group element, so it leaves the left side unchanged.

\begin{lemma}[There is always an eligible tag]
\label{lem:eligible-tag}
For each axis order \(\pi\), at least one base point
\(p\in\mathcal Q_R\) satisfies \eqref{eq:axial-eligible}.
In particular \(\mathcal E(U)\ne\varnothing\) for every connection.
\end{lemma}
\begin{proof}
We prove the quantitative statement on a coordinate slice \(X\)
with \(k\) varying axes, each of length \(2R+1\): in any specified
order of those axes there is a base point \(p\) such that
\begin{equation}
       \operatorname{edge}_X(U^{a^{p,\pi}})
           \le4k(2R+1)\operatorname{face}_X(U).
 \label{eq:axial-slice-estimate}
\end{equation}
For \(k=0\) there are no links to estimate.

For the induction, take the last varying axis \(j\), and let
\(X_r\) denote its parallel layers. Choose \(X_{r_0}\) with
least plaquette cost. The within-layer plaquettes of distinct
layers are disjoint subsets of those in \(X\), so
\[
       (2R+1)\operatorname{face}_{X_{r_0}}(U)
                                      \le\operatorname{face}_X(U).
\]
By induction there is a base point in \(X_{r_0}\) whose axial
gauge, in the first \(k-1\) axes, obeys
\[
 \operatorname{edge}_{X_{r_0}}(U^{a})
    \le4(k-1)(2R+1)\operatorname{face}_{X_{r_0}}(U)
    \le4(k-1)\operatorname{face}_X(U).
\]
Extend its paths along the last axis \(j\). In this axial gauge
all internal \(j\)-links are \(1\), and the gauge on the selected
layer is unchanged.

For an internal \(i\)-link, \(i\ne j\), the square between
successive layers is denoted by \(p\), based at \(y\), and has transport
\(U_{y,i}U_{y+e_j,i}^{-1}\), up to inversion and conjugation.
The reverse triangle inequality and
\eqref{eq:chord-properties} therefore give
\[
       |d_G(U_{y+e_j,i})-d_G(U_{y,i})|
                              \le d_G(U_p).
\]
Sum over the within-layer links and telescope between two layers.
Every intervening square is counted at most once for its
direction pair; in particular the harmless bound
\[
 |\operatorname{edge}_{X_r}(U^a)
          -\operatorname{edge}_{X_{r_0}}(U^a)|
                        \le4\operatorname{face}_X(U^a)
                        =4\operatorname{face}_X(U)
\]
holds. The last equality is gauge invariance of plaquette chords.
There is no contribution from \(j\)-links because they are \(1\).
Summing the layer estimates gives
\begin{align*}
 \operatorname{edge}_X(U^a)
 &\le(2R+1)\bigl\{\operatorname{edge}_{X_{r_0}}(U^a)
                                +4\operatorname{face}_X(U)\bigr\}\\
 &\le4k(2R+1)\operatorname{face}_X(U).
\end{align*}
This proves \eqref{eq:axial-slice-estimate}. For \(k=4\), its
coefficient is \(16(2R+1)\), which is paid by
\eqref{eq:axial-eligible}. The proof used no small-field assumption.
\end{proof}

This is the layer argument of
\cite[Lemma~14.2 and (14.8), p.~58]{Main}.
We can now define the local frame law at all inputs:
\begin{equation}
 \begin{split}
 h(U,t)&=
 \begin{cases}
  \text{the rooted minimiser of }\mathcal A_U,&\text{on the regular test},\\
  h^t,&\text{off the regular test},
 \end{cases}\\
 p^{\rm fr}(U,t)&=
 \begin{cases}
  (4!s)^{-1},&\text{on the regular test},\\
  1_{\{t\in\mathcal E(U)\}}/|\mathcal E(U)|,
                                      &\text{off the regular test}.
 \end{cases}
 \end{split}
 \label{eq:frame-law}
\end{equation}
The regular branch gives equal mass to its unused tags; its frame
does not depend on them. Off the test the denominator is positive
by Lemma~\ref{lem:eligible-tag}. Equality at the threshold belongs
to this fallback branch. In either case
\(\sum_t p^{\rm fr}(U,t)=1\).

\subsection{Averaging the connectors between two frames}

At a coarse site \(x\), apply \eqref{eq:frame-law} to the fine
patch translated by \(Lx\), denoted by \(U^{(x)}\). A chosen tag
\(t_x\) gives one frame \(h_z^x\), shared by all links incident
to \(x\). For \(e=(x,i)\), form
\begin{equation}
 Y_{e,z}=h_z^x\,
       U_{\gamma_{Lx+z,L(x+e_i)+z}}\,(h_z^{x+e_i})^{-1},
       \qquad z\in\mathcal Q_R.
 \label{eq:actual-connectors}
\end{equation}
The middle path consists of \(L\) positive \(i\)-steps. These
connectors all compare the same two coarse roots.

The averaging weights are fixed before examining the configuration.
Let \(\nu_L\) be uniform on the centred \(L\)-block:
\begin{equation}
 \nu_L(z)=L^{-4}1_{\{|z|_\infty\le(L-1)/2\}},\qquad
 w=\underbrace{\nu_L*\cdots*\nu_L}_{J\text{ factors}}.
 \label{eq:convolution-weights}
\end{equation}
Here \((f*g)(z)=\sum_yf(y)g(z-y)\) is discrete convolution.
Thus \(w\ge0\), \(\sum_z w(z)=1\), and its support is contained
in \(\{|z|_\infty\le J(L-1)/2\}\subset\mathcal Q_R\).
Each convolution performs one more translated block average.
The fixed number \(J\) belongs to the selected averaging kernel;
the present construction uses only the displayed positivity,
normalisation and support properties.
The score for a predicted link is
\begin{equation}
                 \mathcal A_Y(u)=\sum_z w(z)q_G(Y_zu^{-1}).
 \label{eq:actual-predictor-score}
\end{equation}
This is the weighted group average specified in
\cite[(14.3), p.~56]{Main}.

Put \(A=\sum_z w(z)Y_z\), a matrix which need not lie in \(G\).
Expanding squared distances gives
\begin{equation}
 \mathcal A_Y(u)
     =\tfrac12\sum_z w(z)\|Y_z-A\|_{\rm F}^2
                       +\tfrac12\|A-u\|_{\rm F}^2.
 \label{eq:predictor-variance}
\end{equation}
Hence its minimisers are exactly the nearest group points to \(A\).
Let \(r_*>0\) be a uniform nearest-point radius from
Appendix~\ref{app:projection}, and choose
\(0<c_{\rm pr}<r_*^2/2\). On the open test
\begin{equation}
                        \min_{u\in G}\mathcal A_Y(u)<c_{\rm pr},
 \label{eq:regular-predictor-test}
\end{equation}
equation~\eqref{eq:predictor-variance} gives
\(\operatorname{dist}(A,G)<\sqrt{2c_{\rm pr}}<r_*\).
There is therefore a unique minimiser \(\pi_G(A)\).
The test is open because the minimum of this continuous score
over the fixed compact group is continuous in \(Y\).

The predictor tag is \(z\in\mathcal Q_R\). Its value and mass are
\begin{equation}
 \begin{split}
 Q^{\rm pr}(Y,z)&=
 \begin{cases}
  \pi_G(\sum_y w(y)Y_y),&\text{on }\eqref{eq:regular-predictor-test},\\
  Y_z,&\text{off that test},
 \end{cases}\\
 p^{\rm pr}(Y,z)&=
 \begin{cases}
  s^{-1},&\text{on }\eqref{eq:regular-predictor-test},\\
  w(z),&\text{off that test}.
 \end{cases}
 \end{split}
 \label{eq:predictor-law}
\end{equation}
In particular \(\sum_z p^{\rm pr}(Y,z)=1\). The fallback law uses
the original \(w(z)\), including its zero-weight tags. Replacing
it by a uniform law would change the kernel.

\subsection{The joint law and its covariance}

A complete tag consists of frame tags at sites and predictor tags
at links:
\begin{equation}
       t=((t_x)_{x\in\Lambda_c},(z_e)_{e\in E_c^+})
       \in\mathcal T
       =(\mathcal Q_R\times S_4)^{\Lambda_c}
                                      \times\mathcal Q_R^{E_c^+}.
 \label{eq:complete-tag-space}
\end{equation}
With \(Y_e=Y_e(U,t_x,t_{x+e_i})\) as in
\eqref{eq:actual-connectors}, define
\begin{equation}
 \begin{split}
        Q_t(U)_e&=Q^{\rm pr}(Y_e,z_e),\\
        \lambda_t(U)&=
             \prod_{x\in\Lambda_c}p^{\rm fr}(U^{(x)},t_x)
             \prod_{e\in E_c^+}p^{\rm pr}(Y_e,z_e).
 \end{split}
 \label{eq:actual-joint-law}
\end{equation}
Conditional on the frame tags, each predictor reads the two frames
already chosen at its endpoints. Shared frame tags are not sampled
again for different incident links.

\begin{proposition}[Borel branches and their exact normalisation]
\label{prop:actual-branch-law}
The maps \(Q_t\) and weights \(\lambda_t\) in
\eqref{eq:actual-joint-law} are Borel on the full fine configuration
space, \(\lambda_t\ge0\), and \(\sum_t\lambda_t(U)=1\) for every
\(U\). They use the same fixed finite patches at all periods.
\end{proposition}
\begin{proof}
The regular frame test is open and its minimiser is continuous
there by Lemma~\ref{lem:regular-frame-unique}. Each axial
transport is a finite product of links and inverses, hence
continuous. Both costs in \eqref{eq:axial-eligible} are continuous,
so each eligibility event is closed. Its indicator and the positive
integer \(|\mathcal E(U)|\) are Borel. Formula~\eqref{eq:frame-law}
therefore gives Borel frames and masses on both pieces.
The connector maps are consequently Borel. The regular predictor
test is open in \(Y\), its unique minimiser is continuous, and
the fallback maps are coordinate projections. This proves the
Borel assertion for \eqref{eq:predictor-law} and hence for
\eqref{eq:actual-joint-law}.

To normalise, fix all frame tags and first sum all predictor tags:
\[
     \sum_{(z_e)}\prod_{e\in E_c^+}p^{\rm pr}(Y_e,z_e)
        =\prod_{e\in E_c^+}\left(\sum_{z_e}p^{\rm pr}(Y_e,z_e)\right)=1.
\]
Now sum the frame tags:
\[
       \sum_{(t_x)}\prod_{x\in\Lambda_c}p^{\rm fr}(U^{(x)},t_x)
        =\prod_{x\in\Lambda_c}\left(\sum_{t_x}p^{\rm fr}(U^{(x)},t_x)\right)=1.
\]
All sums are finite. Their order respects the dependence of
predictors on the shared frames and proves the claimed joint
normalisation. The definitions read only the specified translated
patches and connectors.
\end{proof}

This is the branch law in \cite[(14.12), p.~60]{Main}.
The following calculation verifies the symmetry of the actual
choices, rather than adding covariance as a new assumption.

\begin{proposition}[Gauge covariance with every tag fixed]
\label{prop:actual-branch-covariance}
For every fine gauge transformation \(\omega\), put
\(\omega_c(x)=\omega(Lx)\). The actual branches satisfy
\begin{equation}
 \lambda_t(U^\omega)=\lambda_t(U),\qquad
 Q_t(U^\omega)=Q_t(U)^{\omega_c}.
 \label{eq:actual-branch-covariance}
\end{equation}
Consequently the kernel \eqref{eq:physical-kernel} obeys
\begin{equation}
                  K_\beta(U^\omega,V^{\omega_c})=K_\beta(U,V).
 \label{eq:actual-kernel-covariance}
\end{equation}
\end{proposition}
\begin{proof}
Path cancellation gives
\((U^\omega)_{\gamma_{v,v'}}=\omega(v)U_{\gamma_{v,v'}}\omega(v')^{-1}\).
For a patch about \(Lx\), the correspondence of rooted frames is
\begin{equation}
       h'_z=\omega(Lx)h_z\omega(Lx+z)^{-1}.
 \label{eq:frame-transformation}
\end{equation}
It is a bijection of rooted frame spaces. Every framed route is
conjugated by \(\omega(Lx)\), so the score and the regular test
are unchanged. Uniqueness gives \eqref{eq:frame-transformation}
for the regular minimiser.

For an axial tag based at \(Lx+p\),
\[
       (a_z^t)'=\omega(Lx+p)a_z^t\omega(Lx+z)^{-1}.
\]
Substitution into \((a_0^t)^{-1}a_z^t\) gives exactly
\eqref{eq:frame-transformation}. Before root normalisation the
gauged links are conjugated by \(\omega(Lx+p)\). Thus their
edge cost and the original plaquette cost are unchanged, and
the eligible tag set is the same. Both the regular uniform masses
and the eligible-tag masses in \eqref{eq:frame-law} are unchanged.

Substituting the frame transformation on both ends of a connector
cancels its interior gauge values and yields
\begin{equation}
             Y'_{e,z}=\omega_c(x)Y_{e,z}\omega_c(x+e_i)^{-1}.
 \label{eq:connector-transformation}
\end{equation}
The change \(u'=\omega_c(x)u\omega_c(x+e_i)^{-1}\) is a bijection
of \(G\), and \(Y'_{e,z}(u')^{-1}\) is the conjugate of
\(Y_{e,z}u^{-1}\) by \(\omega_c(x)\). Hence the predictor score,
its minimum and its regular test are invariant. The unique regular
minimiser transforms as \(u'\), and the fallback value \(Y_{e,z}\)
transforms in the same way. The masses \(s^{-1}\) and \(w(z)\)
are unchanged. Formula~\eqref{eq:actual-joint-law} now proves
\eqref{eq:actual-branch-covariance}.

Finally
\[
        (V^{\omega_c})_eQ_t(U^\omega)_e^{-1}
         =\omega_c(x)V_eQ_t(U)_e^{-1}\omega_c(x)^{-1}.
\]
Each mismatch price and each branch likelihood in the kernel
is unchanged. Summing the branches proves
\eqref{eq:actual-kernel-covariance}.
\end{proof}

\begin{corollary}[Gauge invariance survives the push]
\label{cor:pushed-gauge-invariance}
If \(F\in L^1(H_f)\) is gauge invariant, then
\(\mathcal P_\beta F\) is gauge invariant on the coarse lattice.
The exact density history starting at the Wilson density therefore
remains gauge invariant at every step using these kernels.
\end{corollary}
\begin{proof}
For any coarse gauge \(\eta\), choose a fine gauge \(\omega\)
with \(\omega(Lx)=\eta(x)\); values at nonroots may be \(1\).
The fixed gauge change \(U\mapsto U^\omega\) preserves \(H_f\).
Consequently
\begin{align*}
 (\mathcal P_\beta F)(V^\eta)
   &=\int K_\beta(U^\omega,V^\eta)F(U^\omega)\,\dd H_f(U)\\
   &=\int K_\beta(U,V)F(U)\,\dd H_f(U)
     =(\mathcal P_\beta F)(V).
\end{align*}
Induction gives the assertion for the density history. The same
argument applies to \(OF\) when the inserted product \(O\) is
gauge invariant and \(OF\) is integrable.
\end{proof}

\subsection{Analytic branch data near the flat connection}

The global definition and covariance above are sufficient for tree
integration. The later stationary equation needs more: the actual
branch action must be smooth on a common local region. We now
verify the first such local input, with its dependence on the fixed
patch stated explicitly.

\begin{lemma}[A regular analytic branch on a fixed patch]
\label{lem:actual-analytic-branch}
For fixed \(G,L,J\) and the chosen branch thresholds, there is
\(r_{\rm br}>0\) such that the following holds. If the fine links
in the read set of a frame or a predicted link have logarithms
\(X_e\) with \(\max_e|X_e|_{\g}<r_{\rm br}\), the corresponding
choices are on their regular branches and do not depend on their
unused tags. Their frames and predicted links are real analytic
functions of these logarithms. On a smaller closed input ball,
every fixed order of their derivatives is bounded by a constant
depending on \(G,L,J\), the thresholds and the derivative order,
independently of the number of blocks and the period.
The local probabilities are the fixed uniform masses on their
unused tags.
\end{lemma}
\begin{proof}
At the flat connection \(U=\mathbf1\), every route is \(1\).
The rooted frame score vanishes precisely when every \(h_z\)
is the same, hence precisely when \(h_z=1\) for all \(z\).
Write a rooted frame variation as \(h_z(t)=e^{t\xi_z}\),
with \(\xi_0=0\). Since
\(h_z(t)h_{z'}(t)^{-1}=1+t(\xi_z-\xi_{z'})+O(t^2)\),
\eqref{eq:chord} gives its Hessian as
\begin{align}
 D_h^2\mathcal A_{\mathbf1}(\mathbf1)[\xi,\xi]
 &=\frac1{s^2}\sum_{z,z'}|\xi_z-\xi_{z'}|_{\g}^2\notag\\
 &=\frac2s\sum_z|\xi_z|_{\g}^2
       -\frac2{s^2}\left|\sum_z\xi_z\right|_{\g}^2
 \ge\frac2{s^2}\sum_z|\xi_z|_{\g}^2.
 \label{eq:frame-flat-hessian}
\end{align}
The inequality uses \(\xi_0=0\) and Cauchy--Schwarz on the
remaining \(s-1\) vectors. Thus the finite rooted Hessian is
invertible. The score is real analytic in the fine link and frame
logarithms, so the analytic implicit function theorem provides
a unique nearby analytic critical frame.

This is the actual minimiser for sufficiently small input links.
Indeed compactness and continuity imply that any limit of
minimisers as \(U\to\mathbf1\) minimises the flat score, whose
rooted minimiser is unique. All minimisers therefore enter the
implicit-function neighbourhood. The frame test is strict at
\(\mathbf1\), so shrinking the input neighbourhood makes its
unique global minimiser the analytic one just constructed.

For the predictor, the flat frames and all connectors are \(1\).
Then \(\mathcal A_Y(u)=q_G(u^{-1})\), whose Hessian at \(1\)
is \(|\xi|_{\g}^2\), by \eqref{eq:chord}.
The same analytic implicit-function argument applies, or one
can use the analytic nearest-point expression in
\eqref{eq:predictor-law}. Its regular test is also strict at
the flat input. Continuity keeps both stages on their regular
branches after one further reduction of the input radius.

Every frame and connector reads a fixed finite number of links
in the prescribed translated patches. The same radius therefore
works at every coarse site and link, independently of how many
such patches occur. On a smaller closed coordinate ball the
analytic functions and each fixed order of their derivatives are
bounded by compactness. The number of input coordinates depends
only on \(G,L,J\), so converting these derivative bounds between
fixed finite-dimensional coordinate norms does not introduce a
volume or cutoff factor. Finally the probabilities on both regular
branches are constant by \eqref{eq:frame-law} and
\eqref{eq:predictor-law}.
\end{proof}

This proves the fixed-patch input of
\cite[Lemma~14.1, p.~57]{Main} for the branch maps actually used
in our kernel. The explicit lower bound in
\eqref{eq:frame-flat-hessian} may depend on the patch cardinality;
the lemma fixes \(L,J\) before the volume varies. It is not an
estimate uniform as \(L\to\infty\), nor a Hessian bound for the
full lattice action. Those are different analytic questions.
The proof of Lemma~\ref{lem:uniform-action-derivatives} below
expands the first three implicit derivatives of these same
frames, followed by the original connector average and
nearest-point projection defining the actual predictor.

\section{Removing the redundant fine coordinates}
\label{sec:tree}

The exponent in \eqref{eq:initial-action} still contains gauge
directions. Before using ordinary coordinates near a minimiser, we
should remove the directions that a root-fixed gauge transformation
can change freely. We first identify the symmetry needed for this
particular integral.

\subsection{Which gauge transformations may keep \(V\) fixed?}

For a fine gauge transformation \(h\), define its coarse restriction
by \(h_c(x)=h(Lx)\). Along the straight path in
\eqref{eq:straight-predictor}, cancellation at its internal vertices
gives
\begin{equation}
 Q(U^h)_i(x)=h_c(x)Q(U)_i(x)h_c(x+e_i)^{-1}.
 \label{eq:predictor-covariance}
\end{equation}
For the actual averaging procedure, this covariance has now been
proved, with every tag fixed, in
Proposition~\ref{prop:actual-branch-covariance}. In particular
\begin{equation}
                K_\beta(U^h,V^{h_c})=K_\beta(U,V).
 \label{eq:kernel-covariance}
\end{equation}
This uses both the covariance of the predictors and invariance of
their actual likelihoods, not covariance of the predictors alone.
The tree reduction below requires only the resulting identity for
the full kernel and no differentiability in the fine field.

If \(h(Lx)=1\) at every embedded coarse site, then \(h_c=1\), so
\eqref{eq:kernel-covariance} reads
\begin{equation}
                     K_\beta(U^h,V)=K_\beta(U,V).
 \label{eq:root-kernel}
\end{equation}
Consequently \(K_\beta(U,V)F(U)\) is invariant under these
transformations whenever \(F\) is root-fixed gauge invariant.
Gauge invariance of \(F\) alone would not give this conclusion for
an arbitrary kernel.

\subsection{One tree in each block}

The fine sites split into blocks
\[
                  Lx+\{0,\ldots,L-1\}^4,\qquad x\in\Lambda_c.
\]
In each block choose a spanning tree rooted at \(Lx\): a connected
collection of links containing every block vertex and no cycle.
For an explicit choice, write a vertex as \(Lx+z\). For \(z\ne0\),
select its first nonzero coordinate \(j\) and join it to its parent
\(Lx+z-e_j\). Repeatedly taking the parent decreases
\(z_1+\cdots+z_4\), so every vertex reaches \(Lx\).
There is one edge per nonroot vertex, hence \(L^4-1\) edges;
connectedness and this edge count show that this is a tree.
Its root paths have length at most \(4(L-1)\).

Let \(T\subset E_f^+\) be the union of the tree links, expressed
with positive orientation, and set \(I=E_f^+\setminus T\).
Tree links are internal to disjoint blocks, and
\begin{equation}
       |T|=f^4-c^4,\qquad |I|=3f^4+c^4.
 \label{eq:tree-count}
\end{equation}
Any fixed choice of block trees gives the integration identity below.

\begin{definition}[Tree gauge]
A configuration is in tree gauge if \(U_e=1\) for every \(e\in T\).
For \(u\in G^I\), denote by \(U_T(u)\) the connection equal to \(1\)
on \(T\) and to \(u_e\) on \(I\).
\end{definition}

\begin{lemma}[The unique root-fixed tree transformation]
\label{lem:tree-transformation}
Every \(U\in\mathcal U_f\) has a unique gauge transformation
\(h_U\), equal to \(1\) at every root, such that \(U^{h_U}\)
is in tree gauge. It depends only on the tree links of \(U\).
\end{lemma}
\begin{proof}
Put \(h_U=1\) at each root. If \(z\) is a parent and \(z'\) its
child, let \(U_{zz'}\) denote the oriented link along that tree
edge, with an inverse when its orientation is negative. Define
\[
                           h_U(z')=h_U(z)U_{zz'}.
\]
Then
\[
 U^{h_U}_{zz'}=h_U(z)U_{zz'}h_U(z')^{-1}
                =h_U(z)U_{zz'}(h_U(z)U_{zz'})^{-1}=1.
\]
The recursion reaches every vertex and reads only tree links.
Conversely, the requirement that this transformed link equal \(1\)
forces the same value at each child. Starting from the fixed root
values proves uniqueness.
\end{proof}

\begin{proposition}[Exact tree-gauge integration]
\label{prop:tree-integration}
If a Borel function \(J\in L^1(H_f)\) is invariant under all gauge
transformations equal to \(1\) at the roots, then
\begin{equation}
        \int_{\mathcal U_f}J(U)\,\dd H_f(U)
          =\int_{G^I}J(U_T(u))\,\dd h_G^I(u).
 \label{eq:tree-integration}
\end{equation}
\end{proposition}
\begin{proof}
First fix the tree links. By Lemma~\ref{lem:tree-transformation},
\(h_U\) is now fixed, independently of every remaining link.
For each \(e\in I\), make the change of variables
\begin{equation}
                 u_e=h_U(e_-)U_eh_U(e_+)^{-1}.
 \label{eq:free-link-change}
\end{equation}
It is invertible, with
\(U_e=h_U(e_-)^{-1}u_eh_U(e_+)\), and preserves \(h_G\) by
left and right invariance. The product measure on \(G^I\) is
therefore unchanged. Invariance of \(J\) gives
\[
                         J(U)=J(U^{h_U})=J(U_T(u)).
\]
For a nonnegative \(J\), Tonelli's theorem now yields
\begin{align*}
 \int_{G^T}\int_{G^I}J(U)\,\dd h_G^I(U_I)\dd h_G^T(U_T)
 &=\int_{G^T}\int_{G^I}J(U_T(u))\,\dd h_G^I(u)\dd h_G^T(U_T)\\
 &=\int_{G^I}J(U_T(u))\,\dd h_G^I(u).
\end{align*}
The last equality uses that \(h_G^T\) has mass one. Applying this
to \(|J|\) establishes absolute integrability of the reduced
integrand. Fubini and the same changes of variables then prove the
formula for complex \(J\).
\end{proof}

For the actual kernel, \eqref{eq:kernel-covariance} has been proved.
Take an integrable input \(F\) which is root-fixed gauge invariant.
This property holds at each step of the Wilson density history by
Corollary~\ref{cor:pushed-gauge-invariance}; it also holds for an
integrable gauge-invariant inserted product times that density.
Apply the proposition to the full
integrand \(J(U)=K_\beta(U,V)F(U)\), and only then expand the
branch sum. Equation~\eqref{eq:full-block-integral} becomes
\begin{equation}
 \begin{split}
 (\mathcal P_\beta F)(V)
 ={}&e^{vf^4}z_{\mathrm n}(\beta)^{-b_c}
       \sum_t\int_{G^I}\lambda_t(U_T(u))\mathcal B(U_T(u))\\
 &\hspace{22mm}\cdot
       e^{-\beta A_t(V,U_T(u))-I_f(U_T(u))}\,\dd h_G^I(u).
 \end{split}
 \label{eq:tree-block-integral}
\end{equation}
This does not require each branch term to be invariant on its own.
The calculation is the compact-group identity in
\cite[Lemma~11.2 and (11.5), p.~37]{Main}; the proof above fixes the
tree links before using Haar invariance, so no Jacobian for gauge
variables has to be computed.

\section{Putting only the selected links into a chart}
\label{sec:sectors}

The free variables in \eqref{eq:tree-block-integral} still range over
\(G\). A Taylor expansion requires coordinates, but an exponential
chart covers only a neighbourhood of the identity component. We
therefore split each link integral before making any coordinate
change. The part outside the chart keeps its original Haar measure.

\subsection{The local Haar factor}

Use the fixed normalised inner product \eqref{eq:physical-norm}
on \(\g\), with the Lebesgue measure determined by an orthonormal basis.
Changing that basis by an orthogonal matrix leaves this measure
unchanged. The precise Haar formula needed below is proved in
Appendix~\ref{app:haar}.

\begin{lemma}[Haar measure in a translated exponential chart]
\label{lem:haar-chart}
There are \(r_G,\zeta_G>0\) and a smooth positive function \(j_G\)
on \(\{|X|_{\g}<r_G\}\), with \(j_G(0)=1\), such that exponential
is one-to-one on this ball and
\begin{equation}
       \dd h_G(e^XR)=\zeta_Gj_G(X)\,\dd X
       \quad (R\in G,\ |X|_{\g}<r_G).
 \label{eq:haar-chart}
\end{equation}
Its Jacobian is
\begin{equation}
 j_G(X)=\left|\det_{\g}\int_0^1e^{-t\ad X}\,\dd t\right|,
 \qquad (\ad X)Y=[X,Y].
 \label{eq:haar-jacobian}
\end{equation}
Every fixed number of derivatives of \(j_G,j_G^{-1},\log j_G\)
is bounded on \(\{|X|_{\g}\le r_G/2\}\).
\end{lemma}

The factor \(\zeta_G\) is tied to probability Haar measure on the
whole group, including all components. It is not replaced by one.
Right translation by \(R\) leaves it and \(j_G\) unchanged.

Choose a smooth \(\chi:G\to[0,1]\) such that
\begin{equation}
 \chi(e^X)=1\quad\text{if }|X|_{\g}\le r_G/4,
 \qquad
 \supp\chi\subset\{e^X:|X|_{\g}\le r_G/2\}.
 \label{eq:chart-cutoff}
\end{equation}
For instance, use a smooth radial cutoff in this chart and extend
it by zero. Its support is a compact subset of the chart, so this
extension is smooth on \(G\).
For every fixed reference link \(R\in G\),
\begin{align}
 \int_G f(U)\,\dd h_G(U)
 &=\int_G\chi(UR^{-1})f(U)\,\dd h_G(U)\notag\\
 &\quad+\int_G[1-\chi(UR^{-1})]f(U)\,\dd h_G(U),
 \label{eq:one-link-split}\\
 \int_G\chi(UR^{-1})f(U)\,\dd h_G(U)
 &=\zeta_G\int_{|y|_{\g}<r_G}
               \chi(e^y)f(e^yR)j_G(y)\,\dd y.
 \label{eq:one-link-chart}
\end{align}
Both identities are exact for integrable \(f\). We have not
discarded the second term of \eqref{eq:one-link-split}.

\subsection{A reference connection and the exact sector expansion}

The coordinates should record the discrepancy from the fixed coarse
connection \(V\). For a free fine link \(e=(Lx+z,i)\), choose the
blockwise reference
\begin{equation}
 R_e(V)=
 \begin{cases}
  1,&z_i<L-1,\\
  V(x,i),&z_i=L-1.
 \end{cases}
 \label{eq:reference-link}
\end{equation}
Thus internal links have reference \(1\), while a link crossing the
positive face of a block carries the corresponding coarse link.
In particular \(R_e(\mathbf1)=1\). This is the reference of
\cite[Section~11, just before (11.8), p.~38]{Main}.
The coordinate identity itself would hold for any fixed
\(R_e(V)\in G\); its later estimates do depend on the choice.

For each \(e\in I\), insert
\[
       1=\chi(U_eR_e(V)^{-1})+[1-\chi(U_eR_e(V)^{-1})].
\]
Multiplying these identities gives
\begin{equation}
 1=\sum_{\Ia\subset I}
       \prod_{e\in\Ia}\chi(U_eR_e(V)^{-1})
       \prod_{e\in\If}[1-\chi(U_eR_e(V)^{-1})],
 \qquad \If=I\setminus\Ia.
 \label{eq:sector-unity}
\end{equation}
Each summand is called a sector. The active set \(\Ia\) consists
of the links put into exponential coordinates; the frozen set
\(\If\) consists of links kept in Haar coordinates. Frozen variables
will be held fixed during an active integral and then integrated.
These weights form a partition of unity and need not be disjoint
indicators.

For fixed \(V\) and a sector, let
\begin{equation}
 \Omega=\{y\in\g^{\Ia}:|y_e|_{\g}<r_G\text{ for all }e\in\Ia\},
 \qquad Y\in G^{\If}.
 \label{eq:chart-domain}
\end{equation}
Define the full fine connection by
\begin{equation}
 U(V,Y,y)_e=
 \begin{cases}
  1,&e\in T,\\
  e^{y_e}R_e(V),&e\in\Ia,\\
  Y_e,&e\in\If.
 \end{cases}
 \label{eq:sector-connection}
\end{equation}
All products in the active measure are finite, and
\begin{equation}
      \dd h_G^{\Ia}(U_{\Ia})
         =\zeta_G^{|\Ia|}\prod_{e\in\Ia}j_G(y_e)\,\dd y,
      \qquad \dd y=\prod_{e\in\Ia}\dd y_e.
 \label{eq:product-haar}
\end{equation}
The frozen measure remains \(\dd h_G^{\If}(Y)\).

For the sector weight write
\begin{equation}
 \Theta(V,Y,y)=\prod_{e\in\Ia}\chi(e^{y_e})
          \prod_{e\in\If}[1-\chi(Y_eR_e(V)^{-1})].
 \label{eq:sector-weight}
\end{equation}
Pull back the branch action, and place the Jacobian in the exponent:
\begin{align}
 \varphi_t(V,Y,y)&=A_t(V,U(V,Y,y)),\label{eq:sector-action}\\
 \psi(V,Y,y)&=I_f(U(V,Y,y))
                              -\sum_{e\in\Ia}\log j_G(y_e).
                              \label{eq:haar-correction}
\end{align}
Finally retain the factors which have not been put into this exponent:
\begin{equation}
 \theta_t(V,Y,y)=\Theta(V,Y,y)\lambda_t(U(V,Y,y))
                                      \mathcal B(U(V,Y,y)).
 \label{eq:full-sector-multiplier}
\end{equation}
This multiplier can be merely Borel and can be complex. The
representation does not require differentiating it.

\begin{proposition}[The complete chart formula]
\label{prop:sector-formula}
For the actual kernel and an integrable root-fixed gauge-invariant
input \(F\) as in \eqref{eq:general-input},
\begin{equation}
 \begin{split}
 (\mathcal P_\beta F)(V)
 ={}&e^{vf^4}z_{\mathrm n}(\beta)^{-b_c}
       \sum_{t\in\mathcal T}\sum_{\Ia\subset I}\zeta_G^{|\Ia|}\\
 &\quad\cdot\int_{G^{\If}}\int_\Omega
       \theta_t(V,Y,y)e^{-\beta\varphi_t(V,Y,y)-\psi(V,Y,y)}
                                   \,\dd y\,\dd h_G^{\If}(Y).
 \end{split}
 \label{eq:full-sector-formula}
\end{equation}
\end{proposition}
\begin{proof}
Insert \eqref{eq:sector-unity} into \eqref{eq:tree-block-integral}.
For an active link, \(\chi(U_eR_e(V)^{-1})\) vanishes outside its
translated chart. Equation~\eqref{eq:one-link-chart} therefore
changes that integral to the unique coordinate \(y_e\), with factor
\(\zeta_Gj_G(y_e)\). Repeat over all \(e\in\Ia\).
No change is made to a frozen coordinate. The product of the
Jacobian factors satisfies
\[
 e^{-I_f(U(V,Y,y))}\prod_{e\in\Ia}j_G(y_e)
       =\exp\left\{-I_f(U(V,Y,y))
                         +\sum_{e\in\Ia}\log j_G(y_e)\right\}
       =e^{-\psi(V,Y,y)}.
\]
The remaining factors are precisely
\eqref{eq:sector-action} and \eqref{eq:full-sector-multiplier}.
This proves the formula for nonnegative inputs. For a complex
multiplier apply the same changes to the absolute value first;
the finite partition and the original integrability justify Fubini.
\end{proof}

This gives the full version of \cite[(11.8)--(11.9), p.~38]{Main}.
The factor \(\zeta_G\) occurs once per active link, not once per
real component. The dimension used by a subsequent Gaussian
integral will instead be
\begin{equation}
                         d=|\Ia|\dim_{\R}\g.
 \label{eq:active-dimension}
\end{equation}
If \(\Ia=\varnothing\), the active space is \(\R^0\), its integral
is evaluation at its unique point, and no chart factor is introduced.
In particular \eqref{eq:full-sector-formula} retains configurations
outside the identity component through their frozen links.

\section{The Gaussian representation before choosing a centre}
\label{sec:gaussian}

Fix \(V,Y,t,\Ia\) and temporarily suppress them from the notation.
The remaining integral is
\begin{equation}
                    \int_\Omega\theta(y)e^{-\beta\varphi(y)-\psi(y)}
                    \,\dd y.
 \label{eq:active-integral}
\end{equation}
Choose an orthonormal basis of \(\g\) and identify its active product
with \(\R^d\), using the Euclidean norm
\( |y|^2=\sum_{e\in\Ia}|y_e|_{\g}^2\).
The large coefficient \(\beta\) suggests a fluctuation scale
\(g=\beta^{-1/2}\). Before looking for a preferred centre, we can
make this rescaling around any \(m\in\R^d\).

For a real symmetric positive definite matrix \(C\), let
\(\mu_C\) be the centred Gaussian probability with density
\begin{equation}
 \dd\mu_C(z)=(2\pi)^{-d/2}(\det C)^{-1/2}
                e^{-z^{\mathsf T}C^{-1}z/2}\,\dd z.
 \label{eq:gaussian-density}
\end{equation}
Its covariance is \(C\). The normalisation follows by orthogonal
diagonalisation of \(C\) and the one-dimensional Gaussian integral.

\begin{lemma}[An exact Gaussian change of variables]
\label{lem:gaussian-change}
Let \(g>0\), \(m\in\R^d\), and \(C>0\). For every
\(B\in L^1(\R^d)\),
\begin{equation}
 \int_{\R^d}B(y)\,\dd y
  =(2\pi)^{d/2}g^d(\det C)^{1/2}
       \int_{\R^d}e^{z^{\mathsf T}C^{-1}z/2}B(m+gz)
                                                     \,\dd\mu_C(z).
 \label{eq:gaussian-change}
\end{equation}
Absolute integrability on the two sides is equivalent.
\end{lemma}
\begin{proof}
The affine substitution \(y=m+gz\) has Jacobian \(g^d\), so
\(\int B(y)\,\dd y=g^d\int B(m+gz)\,\dd z\).
Solving \eqref{eq:gaussian-density} for Lebesgue measure gives
\[
 \dd z=(2\pi)^{d/2}(\det C)^{1/2}
                   e^{z^{\mathsf T}C^{-1}z/2}\,\dd\mu_C(z).
\]
Substitution proves the identity. Repeating the calculation with
\(|B|\) proves the assertion about absolute integrability.
\end{proof}

Apply the lemma to the function equal to
\(\theta(y)e^{-g^{-2}\varphi(y)-\psi(y)}\) on \(\Omega\) and zero
outside. Its image domain is
\begin{equation}
                   \Omega_{m,g}=\{z:m+gz\in\Omega\}.
 \label{eq:rescaled-domain}
\end{equation}
Thus the exact sector identity is
\begin{equation}
 \begin{split}
 &\int_\Omega\theta(y)e^{-g^{-2}\varphi(y)-\psi(y)}\,\dd y\\
 &\quad=(2\pi)^{d/2}g^d(\det C)^{1/2}
      \int_{\Omega_{m,g}}\theta(m+gz)\\
 &\hspace{30mm}\cdot\exp\left\{-g^{-2}\varphi(m+gz)
             -\psi(m+gz)+\tfrac12z^{\mathsf T}C^{-1}z\right\}
                                                    \,\dd\mu_C(z).
 \end{split}
 \label{eq:sector-gaussian}
\end{equation}
This is \cite[Lemma~11.3 and (11.10), p.~38]{Main}, with the domain
displayed explicitly. No stationary-point hypothesis or smoothness
of \(\theta\) was needed. The compensating positive quadratic
exponential comes from rewriting Lebesgue measure; its cancellation
with a part of the action remains to be arranged.

The identity is pointwise in the fixed frozen variables whenever
the active integral is absolutely integrable. This holds for almost
every \(Y\) by Proposition~\ref{prop:sector-formula}. Measurable
choices of \(m(V,Y)\) and \(C(V,Y)>0\) allow the same formula
to be integrated against \(h_G^{\If}\). In particular the choices
\(m=0\) and \(C=1\) always give such a representation.
For the regular all-active reference contribution,
Theorem~\ref{thm:uniform-stationary} supplies an analytic, hence
measurable, stationary choice. The exact formula in other sectors
does not require such a choice.

The domains and dimensions in \eqref{eq:sector-gaussian} have distinct
roles. The chart domain is a product of link balls; it is generally
not a Euclidean ball. Translation and rescaling do not remove its
indicator. Only the \(d\) active real coordinates contribute
\(g^d\) and \(\det C\). The frozen Haar variables contribute neither.

\section{Why stationarity and the inverse Hessian are the useful choices}
\label{sec:stationary}

Formula~\eqref{eq:sector-gaussian} leaves the shift and covariance
arbitrary. We now inspect its exponent to determine what choices
would make its remaining exponential small. For this calculation
fix a branch and a region on which the pulled-back action is smooth.
The Borel multiplier remains separate.

Suppose \(\varphi\) is \(C^2\) on a neighbourhood of the segment
from \(m\) to \(m+gz\). Apply the fundamental theorem of
calculus twice to \(t\mapsto\varphi(m+tgz)\). It gives
\begin{align}
 \varphi(m+gz)
 &=\varphi(m)+gD\varphi(m)[z]+g^2V_g(z),
       \label{eq:taylor-with-linear}\\
 V_g(z)&=\int_0^1(1-t)
                 D^2\varphi(m+tgz)[z,z]\,\dd t.
       \label{eq:taylor-integral}
\end{align}
The term \(gD\varphi(m)[z]\) in the action would leave
\(-g^{-1}D\varphi(m)[z]\) in the exponent. Also
\(V_0(z)=\tfrac12D^2\varphi(m)[z,z]\), while the compensating
term in \eqref{eq:sector-gaussian} is
\(\tfrac12z^{\mathsf T}C^{-1}z\). Thus the desired cancellations are
\begin{equation}
        D\varphi(m)=0,\qquad
        C^{-1}=D^2\varphi(m)=:H>0.
 \label{eq:desired-stationary-data}
\end{equation}
The first condition removes the negative power \(g^{-1}\); the
second removes the quadratic term of order one. These conditions
are the reason to seek a stationary centre and a positive Hessian.
They have not followed merely from a change of variables.

The separation into a stationary action, a Gaussian determinant
and a fluctuation integral is also the organisation in
\cite[Section~2, (2.10)--(2.13), pp.~267--268]{Bal87}.
That paper uses constrained block averages and its own gauge
construction. Here we prove the finite-dimensional identities
for the normalised kernel above; no analytic estimate is transferred
from that different constrained problem.

\subsection{A quadratic calculation that closes exactly}

Before using a nonlinear action, consider
\(\varphi(y)=\tfrac12|Dy+\alpha|^2\), where
\(D:\R^d\to\R^k\) is a real matrix and \(\alpha\in\R^k\).
Let \(H=D^{\mathsf T}D\). At an arbitrary shift,
\begin{equation}
 |D(m+gz)+\alpha|^2
  =|Dm+\alpha|^2
    +2g\langle Hm+D^{\mathsf T}\alpha,z\rangle
    +g^2\langle z,H z\rangle.
 \label{eq:quadratic-expansion}
\end{equation}
If \(H>0\), take
\begin{equation}
             m=-H^{-1}D^{\mathsf T}\alpha,
             \qquad C=H^{-1}.
 \label{eq:quadratic-choice}
\end{equation}
The linear term in \eqref{eq:quadratic-expansion} vanishes and the
quadratic term cancels the compensating exponential in
\eqref{eq:gaussian-change}. Consequently, for every Borel
multiplier \(B\) for which the integral is absolutely convergent,
\begin{equation}
 \int B(y)e^{-|Dy+\alpha|^2/(2g^2)}\,\dd y
  =\frac{(2\pi)^{d/2}g^d}{\sqrt{\det H}}
    e^{-|Dm+\alpha|^2/(2g^2)}
             \int B(m+gz)\,\dd\mu_{H^{-1}}(z).
 \label{eq:quadratic-identity}
\end{equation}
Indicators and incoming factors may be included in \(B\) without
changing this calculation. For a nonlinear action, its deviation
from the quadratic expression will remain inside the expectation.

\subsection{The nonlinear central integral}

Suppose now that a point \(m\) satisfying
\eqref{eq:desired-stationary-data} has been obtained, and let
\(D>0\) be such that the closed ball
\(\{|y-m|\le gD\}\) is inside the smooth branch region.
This is a finite-dimensional conditional calculation. Its
stationary and Hessian inputs for the actual regular principal
action are proved in Theorem~\ref{thm:uniform-stationary}.

Write
\begin{align}
 P(g,z)&=\tfrac12z^{\mathsf T}H z-V_g(z)
                           +\psi(m)-\psi(m+gz),
                           \label{eq:central-exponent}\\
 N_g&=\frac{(2\pi)^{d/2}g^d}{\sqrt{\det H}}
                 e^{-g^{-2}\varphi(m)-\psi(m)},
                           \label{eq:gaussian-normalisation}\\
 Z_D(g)&=\int_{|z|\le D}e^{P(g,z)}\,\dd\mu_{H^{-1}}(z).
                           \label{eq:central-Z}
\end{align}
The factor \(N_g\) extracts exactly the stationary value,
the Haar correction at the centre, the rescaling and the Gaussian
determinant. The remaining \(Z_D\) still contains the nonlinear
Taylor remainder and the change in the Haar correction.

\begin{proposition}[Exact central identity]
\label{prop:central}
If \(\varphi\) is \(C^2\) and \(\psi\) is continuous on a
neighbourhood of the central ball, then
\begin{equation}
 I_D(g):=\int_{|y-m|\le gD}
                       e^{-g^{-2}\varphi(y)-\psi(y)}\,\dd y
                         =N_gZ_D(g).
 \label{eq:central-identity}
\end{equation}
If these functions and \(m,H\) are fixed as \(g\)
varies near zero, \(Z_D\) has an even continuous extension and
\begin{equation}
          Z_D(-g)=Z_D(g),\qquad
          Z_D(0)=\mu_{H^{-1}}\{|z|\le D\}.
 \label{eq:central-parity}
\end{equation}
\end{proposition}
\begin{proof}
By \eqref{eq:taylor-with-linear} and stationarity,
\[
          \varphi(m+gz)=\varphi(m)+g^2V_g(z).
\]
Substitute \(y=m+gz\) in \(I_D\), insert the Gaussian
density, and obtain
\begin{align*}
 I_D(g)
 &=g^de^{-g^{-2}\varphi(m)-\psi(m)}
       \int_{|z|\le D}
               e^{-V_g(z)+\psi(m)-\psi(m+gz)}\,\dd z\\
 &=\frac{(2\pi)^{d/2}g^d}{\sqrt{\det H}}
       e^{-g^{-2}\varphi(m)-\psi(m)}
       \int_{|z|\le D}
         e^{\frac12z^{\mathsf T}H z-V_g(z)
                           +\psi(m)-\psi(m+gz)}
                                       \,\dd\mu_{H^{-1}}(z).
\end{align*}
This proves \eqref{eq:central-identity}.

The integral formula \eqref{eq:taylor-integral} also makes sense
for small negative \(g\), and
\(V_{-g}(-z)=V_g(z)\). Therefore
\(P(-g,-z)=P(g,z)\). The ball and the centred Gaussian measure
are invariant under \(z\mapsto-z\), proving evenness.
At zero, \(V_0(z)=\tfrac12z^{\mathsf T}H z\), so
\(P(0,z)=0\). Continuity and boundedness on a fixed compact
central region justify dominated convergence and give the last
formula.
\end{proof}

This is the identity in \cite[Proposition~12.1, (12.7)--(12.8),
p.~40]{Main}. In particular \(Z_D(0)\) is a Gaussian probability
of a finite ball. It is not automatically one. For positive active
dimension it is strictly between zero and one. The empty active
space uses the convention \(\det H=1\) and Gaussian mass one.

The role of \(\psi\) is now visible. Its variation appears as
\(\psi(m)-\psi(m+gz)\), without the large coefficient
\(g^{-2}\). When it is differentiable, its linear variation is of
order \(g\), whereas a nonstationary linear term of \(\varphi\)
would have order \(g^{-1}\). This explains the present choice of
centre: we first cancel the terms magnified by the coupling and
keep the remaining Haar and lower-order terms in the fluctuation
integral. It does not assert that their derivatives vanish.

\subsection{Keeping the Gaussian measure fixed when \(V\) varies}

To compare coarse configurations, we will need derivatives of the
integral. A covariance which varies with the coarse data also
varies its Gaussian measure. We can put that variation into an
explicit density instead.

Let \(x\) be real coordinates for a coarse connection near a
reference value. Here \(x\) denotes a coordinate vector, not a
lattice site. Suppress the frozen variables, and suppose a smooth
stationary graph \(m(x)\) is available, with
\(H(x)=D_y^2\varphi(x,m(x))>0\).
Comparing the two Gaussian densities gives
\begin{equation}
 \frac{\dd\mu_{H(x)^{-1}}}{\dd\mu_{H(0)^{-1}}}(z)
 =\sqrt{\frac{\det H(x)}{\det H(0)}}
        e^{-\frac12z^{\mathsf T}(H(x)-H(0))z}.
 \label{eq:fixed-gaussian-density}
\end{equation}
Using \eqref{eq:central-exponent} inside this identity yields
\begin{align}
 Z_D(g,x)&=\int_{|z|\le D}e^{\widetilde P(g,x,z)}
                                    \,\dd\mu_{H(0)^{-1}}(z),
                                    \label{eq:fixed-Gaussian-Z}\\
 \widetilde P(g,x,z)
   &=\tfrac12\log\frac{\det H(x)}{\det H(0)}
          +\tfrac12z^{\mathsf T}H(0)z-V_g(x,z)\notag\\
   &\quad+\psi(x,m(x))-\psi(x,m(x)+gz).
                                    \label{eq:fixed-exponent}
\end{align}
Here \(V_g(x,z)\) is \eqref{eq:taylor-integral} with the coarse
coordinate fixed at \(x\). The Gaussian measure in
\eqref{eq:fixed-Gaussian-Z} is now independent of \(x\).

The determinant terms in \eqref{eq:gaussian-normalisation} and
\eqref{eq:fixed-exponent} are both necessary. The first normalises
the original quadratic integral; the second changes the reference
Gaussian to a fixed one. Neither term may be dropped when
differentiating in \(x\).

For example, differentiation in any coarse direction gives
\begin{equation}
 \partial\log\det H
        =\tr(H^{-1}\partial H),\qquad
 \partial H^{-1}
        =-H^{-1}(\partial H)H^{-1}.
 \label{eq:matrix-derivatives}
\end{equation}
The first follows by expanding
\(\det(H+t\partial H)
 =\det H\det(1+tH^{-1}\partial H)\)
to first order, using \(\det(1+tA)=1+t\tr A+O(t^2)\).
The second follows by differentiating
\(HH^{-1}=1\).
Also differentiation of stationarity gives
\begin{equation}
        H(x)Dm(x)[u]
                  =-D_xD_y\varphi(x,m(x))[u].
 \label{eq:stationary-derivative}
\end{equation}
In this formula \(u\) is a coarse coordinate direction and
\(D_y\varphi\) is identified with its Euclidean gradient.
Thus inverse-Hessian bounds arise directly in derivatives of
the centre and of the Gaussian density. A finite-dimensional
implicit function theorem supplies no lattice-uniform constants
by itself; the needed inverse bounds and common neighbourhood must
come from the lattice action.

\section{Returning to the actual sector and identifying the next estimates}
\label{sec:next}

The auxiliary integral \(I_D\) isolates a smooth exponent. The
actual sector in \eqref{eq:active-integral} still has the multiplier
\(\theta\), the chart domain, and its complement outside the
central ball. We now put these pieces back before asking for an
estimate. This prevents a bound on \(Z_D\) from being mistaken
for a bound on the full blocked density.

\subsection{When does the central ball lie in the sector chart?}

For the active coordinates write
\(\|y\|_\infty=\max_{e\in\Ia}|y_e|_{\g}\).
Recall that \(|y|\) in the Gaussian calculation is the Euclidean
norm of all active components. For \(|z|\le D\),
\[
       |m_e+gz_e|_{\g}
          \le|m_e|_{\g}+g|z_e|_{\g}
          \le\|m\|_\infty+gD.
\]
Thus the explicit sufficient condition
\begin{equation}
                     \|m\|_\infty+gD\le r_G/4
 \label{eq:central-compatibility}
\end{equation}
implies both
\begin{equation}
 \{|y-m|\le gD\}\subset\Omega,
 \qquad \chi(e^{y_e})=1\quad(e\in\Ia)
                         \text{ on this ball}.
 \label{eq:compatible-ball}
\end{equation}
Smoothness of the branch action on that ball is an additional
requirement, separate from chart membership. The purpose of
\eqref{eq:central-compatibility} is now precise: it controls both
the location of the centre and the size of the fluctuation region.
Small \(gD\) alone is insufficient if the centre is near the edge
of the cutoff.

Assume \eqref{eq:compatible-ball} and the stationary conditions of
Section~\ref{sec:stationary}. Split the actual integral into
\begin{align}
 C_D(g)&=\int_{|y-m|\le gD}
                      \theta(y)e^{-g^{-2}\varphi(y)-\psi(y)}\,\dd y,
                      \label{eq:actual-central}\\
 T_D(g)&=\int_{\Omega\cap\{|y-m|>gD\}}
                      \theta(y)e^{-g^{-2}\varphi(y)-\psi(y)}\,\dd y.
                      \label{eq:actual-tail}
\end{align}
The central calculation, now keeping the multiplier, gives
\begin{equation}
 C_D(g)=N_g\int_{|z|\le D}
                       \theta(m+gz)e^{P(g,z)}
                                        \,\dd\mu_{H^{-1}}(z).
 \label{eq:actual-central-gaussian}
\end{equation}
The full sector is exactly \(C_D(g)+T_D(g)\).
Although the active chart cutoffs equal one here, the remaining
multiplier is
\begin{equation}
 \theta_t(V,Y,m+gz)
  =\left(\prod_{e\in\If}[1-\chi(Y_eR_e(V)^{-1})]\right)
     \lambda_t(U(V,Y,m+gz))\mathcal B(U(V,Y,m+gz)).
 \label{eq:remaining-multiplier}
\end{equation}
The frozen product is constant in \(z\), while the branch likelihood
and the input multiplier generally are not. Compatibility of the
chart cutoff therefore does not remove them. In general
\eqref{eq:actual-central-gaussian} cannot be replaced by
\(N_gZ_D(g)\).

If the sector is nonnegative and \(C_D(g)>0\), positivity yields
the elementary comparison
\begin{equation}
 0\le\log(C_D(g)+T_D(g))-\log C_D(g)
       =\log\left(1+\frac{T_D(g)}{C_D(g)}\right)
       \le\frac{T_D(g)}{C_D(g)}.
 \label{eq:positive-tail-comparison}
\end{equation}
For complex multipliers the additive decomposition remains valid,
but this positive logarithmic comparison does not apply.

\subsection{Why the tail must be measured relative to the Gaussian scale}

In the positive smooth reference calculation, an unweighted central
integral is \(N_gZ_D(g)\). To compare its tail \(T\) to
this integral, the first relevant ratio is
\begin{equation}
 \frac{T}{N_g}
   =(2\pi)^{-d/2}g^{-d}\sqrt{\det H}\,
       e^{g^{-2}\varphi(m)+\psi(m)}T.
 \label{eq:normalised-tail-factors}
\end{equation}
An absolute tail estimate must pay all the factors on this right
side. For example, a bound \(H\le K_{\mathrm H}1\)
would give
\begin{equation}
                    \sqrt{\det H}\le K_{\mathrm H}^{d/2}.
 \label{eq:determinant-price}
\end{equation}
Indeed every positive eigenvalue of \(H\) is at most
\(K_{\mathrm H}\), and their product is its determinant.
This upper bound concerns the fine active-coordinate Hessian.
An estimate on derivatives in the coarse coordinates concerns a
different matrix and does not supply \eqref{eq:determinant-price}.
Both bounds are derived for the actual regular principal action
in \eqref{eq:model-determinant-upper} and
\eqref{eq:model-coarse-row-bound}. They are not additional
hypotheses on the final Wilson construction.

There is also a dimension cost even for the reference Gaussian.
If \(H\ge\eta1\), then
\begin{equation}
 \mu_{H^{-1}}\{|z|>D\}
                         \le2^{d/2}e^{-\eta D^2/4}.
 \label{eq:gaussian-tail}
\end{equation}
To verify this, diagonalise \(H\). Its eigenvalues
\(\lambda_j\) satisfy \(\lambda_j\ge\eta\), and the independent
Gaussian coordinates have variances \(\lambda_j^{-1}\). Hence
\[
 \int e^{\eta|z|^2/4}\,\dd\mu_{H^{-1}}(z)
       =\prod_{j=1}^d\left(1-\frac{\eta}{2\lambda_j}\right)^{-1/2}
       \le2^{d/2}.
\]
On \(\{|z|>D\}\) the integrand is at least
\(e^{\eta D^2/4}\), proving \eqref{eq:gaussian-tail}.

For a requested error \(0<\delta<1\), this particular bound is
at most \(\delta\) if
\begin{equation}
                D^2\ge\frac4\eta
                        \left(\frac d2\log2+\log\delta^{-1}\right).
 \label{eq:Gaussian-radius-request}
\end{equation}
At a centre with \(\|m\|_\infty\ll r_G\), the Gaussian
error and chart requests can be collected as
\[
 \left(\frac{d+\log\delta^{-1}}{\eta}\right)^{1/2}
                         \ll D\ll\frac{r_G}{g}.
\]
The constants implicit in these comparisons are chosen to
give \eqref{eq:Gaussian-radius-request} and
\eqref{eq:central-compatibility}, with their displayed
numerical bounds. This range is available once
\[
        \frac g{r_G}
          \left(\frac{d+\log\delta^{-1}}{\eta}\right)^{1/2}
                                                     \ll1.
\]
Thus first fix the requested error and the other data,
check this compatibility, and then choose \(D\).
Chart inclusion needs a fixed upper bound on \(gD\),
not a limit \(gD\to0\). In the all-active sector,
\eqref{eq:tree-count} gives
\(d=(3f^4+c^4)\dim_{\R}\g\), so a finite-dimensional calculation
with its full dimension dependence is not yet a uniform estimate
as either cutoff is removed. The Gaussian tail
\eqref{eq:gaussian-tail} also does not estimate the tail of the
nonlinear action by itself; that requires a relative energy bound
for the actual action.

\subsection{The output of the exact calculation}

We can already write the complete Gaussian representation without
assuming a stationary point in every sector. In each fixed sector
choose measurable \(m=m(V,Y)\) and \(C=C(V,Y)>0\), possibly
depending on \(t,\Ia\). Substitute \eqref{eq:sector-gaussian} into
\eqref{eq:full-sector-formula}. With the suppressed dependence of
\(\theta_t,\varphi_t,\psi\) restored through their argument
\((V,Y,m+gz)\), the result is
\begin{equation}
 \begin{split}
 (\mathcal P_\beta F)(V)
 ={}&e^{vf^4}z_{\mathrm n}(\beta)^{-b_c}
     \sum_t\sum_{\Ia\subset I}\zeta_G^{|\Ia|}(2\pi)^{d/2}g^d\\
 &\cdot\int_{G^{\If}}\sqrt{\det C}
       \int_{\Omega_{m,g}}\theta_t(V,Y,m+gz)\\
 &\qquad\cdot\exp\left\{-g^{-2}\varphi_t(V,Y,m+gz)
             -\psi(V,Y,m+gz)+\tfrac12z^{\mathsf T}C^{-1}z\right\}
                     \,\dd\mu_C(z)\dd h_G^{\If}(Y).
 \end{split}
 \label{eq:complete-Gaussian}
\end{equation}
All sums are finite at fixed cutoff and volume. Absolute
integrability follows from the same identity for absolute values.
The choices \(m=0,C=1\) make the representation immediately
available, whereas the stationary choices are useful where the
model analysis proves them.

The factors and variables in this formula have the following origins:
\begin{center}
\small
\begin{tabular}{@{}ll@{}}
\toprule
Origin & Retained factor or integral\\
\midrule
Input scalar & \(e^{vf^4}\)\\
Normalised coarse noise & \(z_{\mathrm n}(\beta)^{-b_c}\)\\
Branch choices & \(\sum_t\), with \(\lambda_t\) inside \(\theta_t\)\\
Active Haar charts & \(\zeta_G^{|\Ia|}\), with \(j_G\) inside \(\psi\)\\
Active Gaussian coordinates & \((2\pi)^{d/2}g^d\sqrt{\det C}\)\\
Frozen links & \(\dd h_G^{\If}(Y)\)\\
Chart domain and incoming factors & \(\Omega_{m,g}\) and \(\theta_t\)\\
\bottomrule
\end{tabular}
\end{center}
For an integrable gauge-invariant inserted product times the density,
the same operations act on its full numerator,
with the same fixed kernel and the same denominator
\eqref{eq:history-mass}. A new sector normalisation has not been
introduced.

The analytic task is now specific. We must construct a stationary
graph in a common regular chart and control its Hessian before
comparing the central integral with its physical complement.
Sections~\ref{sec:smooth-principal}--\ref{sec:uniform-stationary}
carry out the stationary and Hessian analysis for the principal
actions used in the construction.
Section~\ref{sec:physical-tail} proves the reference tail comparison
on auxiliary tori. The subsequent localisation and input-remainder
steps retain all other branch, cutoff and input factors in the
original integral. The corresponding source
development is \cite[Sections~14--19, pp.~55--100]{Main}.

\section{The smooth principal action to which stationarity applies}
\label{sec:smooth-principal}

The decomposition \eqref{eq:general-input} allowed an arbitrary
Borel \(S_f\). Such a function has no reason to have a positive
Hessian. We therefore need the actual structure of the smooth
principal action, rather than a further abstract positivity
assumption. The Wilson starting point and the actions retained by
the RG construction share a positive direct plaquette term,
nonnegative quadratic row squares, and higher curvature terms.
The last terms have no quadratic variation at the flat connection.
This is the structure we will use to prove the gap.

\subsection{Curvature slots and the path budgets}

For a group element \(W\), define
\begin{equation}
       \mathcal C(W)=\operatorname{pr}_{\g}
                                      \bigl((W-W^{-1})/2\bigr).
 \label{eq:principal-curvature}
\end{equation}
Here the projection is for \eqref{eq:physical-norm}; its argument
is skew Hermitian. Projection commutes with conjugation by \(G\),
because both its subspace and its inner product are invariant.
Consequently \(\mathcal C\) is a smooth covariant curvature map,
and
\begin{equation}
 \mathcal C(1)=0,\qquad
 D(\mathcal C\circ\exp)(0)=1,\qquad
 |\mathcal C(W)|_{\g}\le d_G(W).
 \label{eq:curvature-properties}
\end{equation}
The derivative follows by expanding \(e^{tX}-e^{-tX}=2tX+O(t^3)\).
The norm bound follows from contraction of orthogonal projection
and \(\|W-W^{-1}\|_{\rm F}\le2d_G(W)\).

Let \(\mathcal F=\{(i,j):1\le i<j\le4\}\) be the six plaquette
orientations. A path template \(\gamma\) is based at the origin;
write \(\operatorname{end}\gamma\) for its displacement and
\(U_\gamma(o)\) for the holonomy of its translate based at \(o\).
For \(\alpha=(i,j)\in\mathcal F\), let
\(p=(o+\operatorname{end}\gamma,i,j)\) be the terminal plaquette.
A transported curvature slot is
\begin{equation}
 Z_{\gamma,\alpha}(U;o)
   =U_\gamma(o)\,\mathcal C(U_p)\,U_\gamma(o)^{-1}.
 \label{eq:principal-slot}
\end{equation}
It transforms by conjugation at \(o\). Its length is
\(l=|\gamma|+4\); its footprint contains the path and the
terminal plaquette. We count a repeated link each time it occurs
in a word. Periodic identification can reduce the number of
distinct variables read, but does not increase that word length.

At each anchor take the six row functions
\begin{equation}
       \ell_b(U;o)=\sum_j r_{bj}
                          Z_{\gamma_{bj},\alpha_{bj}}(U;o),
       \qquad b\in\mathcal F,
       \qquad
       \sum_j|r_{bj}|e^{\sigma(|\gamma_{bj}|+4)}\le C_{\rm p},
 \label{eq:stationary-row-budget}
\end{equation}
where the coefficients \(r_{bj}\) are real and
\(\sigma,C_{\rm p}>0\) are fixed path data.
For the higher terms let
\(p_\nu(U;o)=a_\nu T_\nu(Z_1,\ldots,Z_{k_\nu})\), with
\(a_\nu\in\R\) and \(k_\nu\in\{3,4\}\).
Each \(Z_j\) is a slot based at \(o\)
and \(T_\nu\) is a real multilinear form invariant under
simultaneous conjugation. If \(l_\nu\) is their total length,
assume the stated principal path budget
\begin{equation}
                   \sum_\nu |a_\nu|\,\|T_\nu\|e^{\sigma l_\nu}
                                                         \le C_{\rm p}.
 \label{eq:stationary-higher-budget}
\end{equation}
The norm of a multilinear form is its supremum on arguments of
\(\g\)-norm at most one. Equivalently its norm may be absorbed
into the coefficient magnitude. Component coefficient budgets
give this budget with a group-dependent constant at these fixed
degrees. These are the path data in
\cite[(13.15)--(13.16), p.~50]{Main}.

For completeness, the signed symmetry average below means the
arithmetic average over all signed permutations of the four axes.
Such a map sends an oriented edge to an oriented edge; its
transport is read using an inverse if the resulting orientation
is negative. Pull back a configuration by this rule and average
the resulting scalar functions. This finite average is denoted
by \(\langle\cdot\rangle_{\rm sym}\).
It preserves positivity, word lengths and all path budgets up
to the same fixed counting conventions.

The unmasked principal action is
\begin{equation}
 S_f^{\rm un}(U)=\lambda W_f(U)
       +\left\langle\frac12\sum_{o\in\Lambda_f}\sum_{b\in\mathcal F}
                                      |\ell_b(U;o)|_{\g}^2
                    +\sum_{o\in\Lambda_f}\sum_\nu p_\nu(U;o)
         \right\rangle_{\rm sym},\qquad \lambda>0.
 \label{eq:unmasked-principal-action}
\end{equation}
The direct Wilson term is itself invariant under this average.
Formula~\eqref{eq:unmasked-principal-action} is the smooth
unmasked part of \cite[(13.18), p.~51]{Main}.
Initially \(\lambda=1\), all rows and higher terms vanish,
and \(S_f^{\rm un}=W_f\). The RG step must subsequently maintain
the path budgets for the same family of measures. Proving that
propagation is a later induction task; the estimates below apply
to every profile with these specified path data.

The unmasked analytic extension \(I_f^{\rm un}\) of the
\(\beta\)-independent loop term has the same based-path
structure with two curvature slots and its own finite exponential
coefficient budget \(\Lambda\). More explicitly, it is a signed
symmetry average of sums of invariant bilinear forms with real coefficients
\(b_\nu T_\nu^{(2)}(Z_{\nu,1},Z_{\nu,2})\), with
\(\sum_\nu|b_\nu|\|T_\nu^{(2)}\|e^{\sigma(l_{\nu,1}+l_{\nu,2})}
\le\Lambda\) per anchor. Since every slot has norm at most two,
\[
       |I_f^{\rm un}(U)|\le4\Lambda f^4,
       \qquad
       |I_f^{\rm un}(U)|\le\Lambda f^4h^2
              \quad\text{when every plaquette chord is at most }h.
\]
The original physical loop action \(I_f\) multiplies each such
path term by its prescribed fine mask \(\chi_{Y_\nu}\), where
\(Y_\nu\) contains that term's full slot footprint. Thus
\(I_f=I_f^{\rm un}\) on the fine all-good set, whereas its
physical value price is \(C\Lambda h^2\) per anchor at every
configuration. The general physical formulas retain \(I_f\);
the reference, stationary and analytic formulas below use
\(I_f^{\rm un}\). This distinction leaves the reference splitting
unchanged because its indicator imposes the all-good set.
Likewise \(|\ell_b|\le2C_{\rm p}\), the higher terms have
absolute sum at most \(16C_{\rm p}\) per anchor, and hence
\[
       |S_f^{\rm un}(U)|
          \le(12\lambda+12C_{\rm p}^2+16C_{\rm p})f^4.
\]
These bounds make the regular reference below integrable at fixed
cutoff and volume. This specifies the analytic loop extension
used in the reference calculation; its physical masks stay in
the original density. The stationary
equation does not contain \(I_f^{\rm un}\), so its constants below do not
require a bound on this separate loop allowance.

All constants in this stationary analysis may depend on
\(G,L,J,a,\lambda,\sigma,C_{\rm p}\) and the fixed branch parameters.
At a fixed such profile they will be independent of the period,
the number of blocks and the terminal UV cutoff. This does not
declare the constants uniform as the block ratio varies.
The fine-curvature gain needed to choose that ratio will be
proved separately for the prescribed Gaussian history profiles
in Section~\ref{sec:calibrated-rows}.

\subsection{The regular reference integral and the chart action}

The action \eqref{eq:unmasked-principal-action} is used in a
regular reference contribution. For a fine threshold \(h>0\),
let
\[
       \chi_{\rm sf}(U)
          =1_{\{d_G(U_p)\le h\text{ for every fine plaquette }p\}},
       \qquad
       F_{{\rm sf},\beta}
          =\chi_{\rm sf}e^{-\beta S_f^{\rm un}-I_f^{\rm un}}.
\]
Every slot transforms by conjugation at its anchor; the row norms
and invariant multilinear forms therefore make \(S_f^{\rm un}\)
and \(I_f^{\rm un}\) gauge invariant. The indicator is gauge invariant
as well, so the reference splitting retains the symmetry needed
for tree integration.
The prescribed clipped and masked principal action agrees with
\(S_f^{\rm un}\) on this all-good set. For any original density
\(F\), the splitting
\begin{equation}
       F=e^{vf^4}F_{{\rm sf},\beta}
                         +\bigl(F-e^{vf^4}F_{{\rm sf},\beta}\bigr)
 \label{eq:reference-exact-splitting}
\end{equation}
is exact, and the blocking operator acts on both terms. The
stationary calculation concerns the smooth exponent of the first
term. Its indicator, and every multiplier in the second term,
remain in the physical integral. This is the role of the reference
splitting in \cite[(14.1), p.~56]{Main}.

We now take the all-active tree chart. Write \(V(x)_e=e^{x_e}\)
and give free fine links the values
\(U(x,y)_e=e^{y_e}R_e(V(x))\), keeping the tree links equal to
\(1\). Define
\begin{equation}
 \varphi(x,y)=S_f^{\rm un}(U(x,y))
       +a\sum_{e\in E_c^+}
                  q_G\bigl(V(x)_eQ(U(x,y))_e^{-1}\bigr).
 \label{eq:stationary-model-action}
\end{equation}
On a sufficiently small common chart the actual branches all
use the same regular predictor \(Q\), by
Lemma~\ref{lem:actual-analytic-branch}. The function in
\eqref{eq:stationary-model-action} is therefore independent of
the unused branch tags there. It is exactly
\cite[(14.19), p.~64]{Main}. The term \(I_f^{\rm un}\) and the Haar
Jacobian remain in \(\psi\), not in this \(\beta\)-multiplied
stationary equation.

\section{Why the actual flat Hessian has a gap}
\label{sec:stationary-coercivity}

The Hessian should measure the fine fluctuations that remain after
fixing the coarse connection. A Wilson quadratic form controls
plaquette curls, but a curl alone misses gradient directions.
Tree gauge removes gradients inside a block. The coarse mismatch
penalty must control the remaining differences between block roots.
We now verify this reconstruction quantitatively.

\subsection{The linear predictor and a local reconstruction}

For a \(\g\)-valued fine link array \(A\), zero on the tree links,
define the linearised regular predictor by
\begin{equation}
          Q(e^{tA})_b=\exp\{t(Q_0A)_b+O(t^2)\}.
 \label{eq:linear-predictor}
\end{equation}
This derivative exists by Lemma~\ref{lem:actual-analytic-branch}.
It is a fixed finite-range linear operator. Its local coefficient
row sums, and the number of its rows reading a given fine link,
are bounded using \(G,L,J\) alone. The linearised plaquette curl is
\begin{equation}
       (\dd A)_{ij}(z)
           =A_j(z+e_i)-A_j(z)-A_i(z+e_j)+A_i(z).
 \label{eq:linear-curl}
\end{equation}
Both arrays use counting measure and the fixed \(\g\)-norm.

For the quantitative local arguments we use the same admissible
period regime as the stationary analysis in \cite[Section~14]{Main}:
the fixed enlarged boxes needed by a predictor and its block-tree
reconstruction have chosen nonwrapping lifts to the integer cover.
Their radii are fixed from \(L,J\) before the two cutoffs vary.

\begin{lemma}[Block Poincar\'e estimate]
\label{lem:block-poincare}
There is \(C_{\rm P}<\infty\), depending only on the fixed group,
block and branch data, such that
\begin{equation}
       \sum_{e\in I}|A_e|_{\g}^2
          \le C_{\rm P}
              \left(\sum_p|(\dd A)_p|_{\g}^2
                          +\sum_b|(Q_0A)_b|_{\g}^2\right).
 \label{eq:block-poincare}
\end{equation}
The constant is independent of the number of blocks and of the
terminal UV cutoff.
\end{lemma}
\begin{proof}
Fix a block and enlarge a box about it to contain its fine links,
all links read by their predictor rows, and the trees of every
block met by those reads. In this fixed box choose a root \(r\)
and coordinate-ordered paths \(\gamma_z:r\to z\). Set
\[
       \phi(z)=\sum_{e\in\gamma_z}A_e,\qquad
       (d_0\phi)_e=\phi(e_+)-\phi(e_-),\qquad
       B_e=A_e-(d_0\phi)_e.
\]
An oriented link contributes its negative when traversed backwards.
The difference defining \(B_e\) is the sum of \(A\) around the
closed path \(\gamma_{e_-}\,e\,\gamma_{e_+}^{-1}\).
Move its coordinate steps into a common order and cancel
backtracks. Commuting an \(i\)-step with a \(j\)-step changes
the path sum by exactly the intervening plaquette curl, with
the corresponding sign. The paths have uniformly bounded
length in this fixed box. Thus the procedure gives integers
\(c_{ep}\) with
\begin{equation}
              B_e=\sum_{p\text{ in the box}}c_{ep}(\dd A)_p,
              \qquad\sup_e\sum_p|c_{ep}|\le N_*<\infty.
 \label{eq:local-additive-reconstruction}
\end{equation}
One can bound the number of commutations by the square of the
maximum path length. This bound depends on the chosen box, not
on the period. All intermediate squares stay in a fixed further
enlargement, which we include in that box.

For a block with root \(Lv\), put \(c_v=\phi(Lv)\) and extend
this coarse array as a block-constant fine potential. Along a
tree link \(A_e=0\), hence \(d_0\phi=-B\). Summing along the
tree path from \(Lv\) to a vertex \(z\) gives
\[
       |\phi(z)-c_v|_{\g}
             \le(L^4-1)N_*\sum_{p\text{ in the box}}|(\dd A)_p|_{\g}.
\]
Subtract the block-constant gradient from \(A\), writing
\(A=d_0c+\widetilde B\). By the preceding bound at both endpoints,
\begin{equation}
       |\widetilde B_e|_{\g}
           \le C_B\sum_{p\text{ in the box}}|(\dd A)_p|_{\g},
       \qquad C_B=\{1+2(L^4-1)\}N_*.
 \label{eq:block-gradient-remainder}
\end{equation}
The estimate holds on the entire fine read of the relevant
predictor, by the initial enlargement of the box.

It remains to recover the gradient between block roots.
Apply the already proved predictor covariance to the fine pure
gauge with values \(\omega_t(z)=e^{-tc_v}\) on block \(v\),
and differentiate at zero. The first variation of
\(\mathbf1^{\omega_t}\) is \(d_0c\), and its coarse first
variation is \((d_{0,c}c)_{(v,i)}=c_{v+e_i}-c_v\). Hence
\begin{equation}
                           Q_0d_0c=d_{0,c}c.
 \label{eq:predictor-gradient-identity}
\end{equation}
The locally constructed values \(c_v\) can be extended arbitrarily
outside the enlarged nonwrapping box when applying this identity;
the chosen predictor row reads only inside it. No periodic
potential for the original full array \(A\) has been assumed.

For an internal fine link, \(d_0c=0\), so
\eqref{eq:block-gradient-remainder} already controls \(A_e\).
For a crossing link with coarse bond \(b=(v,i)\),
\eqref{eq:predictor-gradient-identity} gives
\[
       A_e=(Q_0A)_b+\widetilde B_e-(Q_0\widetilde B)_b.
\]
Let \(Q_*\) bound the sum of the operator norms in a local row
of \(Q_0\). Thus, on a fixed enlarged plaquette set
\(\mathcal P_v\), both cases give
\begin{equation}
       |A_e|_{\g}\le |(Q_0A)_b|_{\g}
                 +(1+Q_*)C_B\sum_{p\in\mathcal P_v}|(\dd A)_p|_{\g},
 \label{eq:link-reconstruction-bound}
\end{equation}
with the noise term absent for internal links.

The sets \(\mathcal P_v\) have a common cardinality bound \(p_*\)
and a common overlap bound \(M_*\). There are at most \(4L^4\)
fine links per block and at most \(L^3\) crossing fine links for
each coarse bond. Square \eqref{eq:link-reconstruction-bound},
use \((u+v)^2\le2u^2+2v^2\) and Cauchy--Schwarz for its
plaquette sum, and sum over blocks. A permissible constant is
\[
       C_{\rm P}=1+8L^4(1+Q_*)^2C_B^2p_*M_*+2L^3.
\]
All four quantities \(Q_*,C_B,p_*,M_*\) were fixed locally,
so this proves \eqref{eq:block-poincare} with the required
volume independence.
\end{proof}

This is the reconstruction in
\cite[Lemma~14.11, (14.20)--(14.23), pp.~65--66]{Main}.
The two terms in its right side have separate roles: the curl
controls the internal block fluctuation, and the actual averaged
noise controls the remaining block-root gradient.

\subsection{The quadratic form of the specified action}

Let \(H_0=D_y^2\varphi(0,0)\), for
\eqref{eq:stationary-model-action}. Its quadratic form can now
be read directly from the action, with no assumed positivity.
At \(x=0\), a variation \(y=tA\) gives the configuration
\(U_e(t)=e^{tA_e}\), with \(A_e=0\) on tree links.
For this flat variation extend \(A\) to reverse links by
\(A_{e^{-1}}=-A_e\), and write \(U_p(t)\) for its plaquette holonomy.
Multiplication around the oriented boundary gives
\[
       U_p(t)=1+t(\dd A)_p+O(t^2),\qquad
       q_G(U_p(t))=\tfrac12t^2|(\dd A)_p|_{\g}^2+O(t^3).
\]
Every transported curvature slot is zero at the flat connection.
Its first variation is the curl at its terminal plaquette;
the transporting path contributes nothing at first order because
the curvature itself is zero. Therefore a row square has
nonnegative second variation. A term with three or four slots
has zero second variation. Finally \eqref{eq:linear-predictor}
gives the noise second variation
\(a\sum_b|(Q_0A)_b|_{\g}^2\). Signed symmetry averaging preserves
nonnegativity and reindexes the direct plaquette sum. Hence
\begin{equation}
       \langle A,H_0A\rangle
          \ge\lambda\sum_p|(\dd A)_p|_{\g}^2
                         +a\sum_b|(Q_0A)_b|_{\g}^2.
 \label{eq:actual-flat-quadratic}
\end{equation}
Combining this with Lemma~\ref{lem:block-poincare} proves
\begin{equation}
                  H_0\ge\eta_0\,1,\qquad
                  \eta_0=\frac{\min(\lambda,a)}{C_{\rm P}}>0.
 \label{eq:proved-flat-gap}
\end{equation}
This inequality uses the counting Euclidean norm
\(\|A\|_2^2=\sum_{e\in I}|A_e|_{\g}^2\).
Also \(\varphi(0,0)=0\) and \(D\varphi(0,0)=0\), since all
curvatures and noise mismatches vanish at the flat connection.
This proves the model estimate of
\cite[Proposition~14.12 and (14.24), p.~66]{Main}.

\section{Local derivative bounds and an inverse without a volume factor}
\label{sec:stationary-local-bounds}

The gap \eqref{eq:proved-flat-gap} bounds the inverse in the
counting \(\ell^2\) norm. The chart, however, bounds each link
separately: for a fine or coarse array write
\[
                   \|u\|_\infty=\max_e|u_e|_{\g}.
\]
Converting a whole-array \(\ell^2\) bound to this norm by a
dimension factor would make the coordinate neighbourhood shrink
as the lattice grows. We instead use the spatial decay of the
Hessian entries. That decay follows from the specified path
budgets and the finite predictor read.

\subsection{The local matrix bound supplied by the action}

Attach a fine coordinate to the coarse block containing the
initial site of its link, and a coarse coordinate to its initial
coarse site. Let \(d(i,j)\) be the periodic maximum distance
between their attached sites. There are at most
\(m_0=4(L^4+1)\) link-coordinate blocks at one attached site,
including fine and coarse coordinates. A matrix block
\(A_{ij}:\g\to\g\) is measured in its operator norm. Define
\begin{equation}
 \|A\|_\theta=
       \max\left\{\sup_i\sum_j e^{\theta d(i,j)}\|A_{ij}\|_{\rm op},
                  \sup_j\sum_i e^{\theta d(i,j)}\|A_{ij}\|_{\rm op}
           \right\}.
 \label{eq:weighted-local-matrix}
\end{equation}
The operator norm uses the fixed \(\g\)-inner product. Writing
these blocks in an orthonormal basis gives the scalar-coordinate
Schur estimates of \cite[(14.25), p.~67]{Main}, with equivalence
constants depending only on the fixed group. Block notation
keeps \(\|u\|_\infty\) identical to the link norm already used
in the chart and cutoff conditions.

Two elementary properties will be useful. If the largest absolute
row and column sums of block norms are \(R,C\), then
\begin{align*}
 \|Au\|_2^2
 &\le\sum_i\left(\sum_j\|A_{ij}\|_{\rm op}|u_j|_{\g}\right)^2\\
 &\le R\sum_{i,j}\|A_{ij}\|_{\rm op}|u_j|_{\g}^2
 \le RC\|u\|_2^2.
\end{align*}
Thus \(\|A\|_{2\to2}\le\|A\|_\theta\). Also the triangle
inequality for \(d\) in a matrix product gives
\begin{equation}
                        \|AB\|_\theta\le\|A\|_\theta\|B\|_\theta.
 \label{eq:weighted-product}
\end{equation}
These estimates use counting measure, not a volume-normalised
Euclidean norm.

The next perturbation bound is in this weighted row and column
norm. A quadratic-form estimate supplies an operator bound,
but the weighted entry sums also have to be proved. We will
keep the first derivative coordinate fixed before summing
translated path terms and actual predictor reads.

\begin{lemma}[Uniform local derivatives of the actual action]
\label{lem:uniform-action-derivatives}
Fix the profile data in Section~\ref{sec:smooth-principal}.
There are \(r_0,\theta,K>0\), independent of the period and
terminal UV cutoff, such that \(\varphi\) is real analytic on
\(\max(\|x\|_\infty,\|y\|_\infty)<2r_0\) and
\begin{align}
 \|D_y^2\varphi(x,y)\|_\theta&\le K,\label{eq:proved-local-hessian}\\
 \|D_y^2\varphi(x,y)-H_0\|_\theta
                    &\le K\max(\|x\|_\infty,\|y\|_\infty),
                    \label{eq:proved-hessian-perturbation}\\
 \|D_y\varphi(x,0)\|_\infty&\le K\|x\|_\infty.
                    \label{eq:proved-gradient-forcing}
\end{align}
Here gradients are identified using the counting product inner
product. For each fixed \(k\ge1\), the joint gradient has
block derivative bounds
\begin{equation}
 \sup_i\sum_{j_1,\ldots,j_k}
       \bigl\|D_{j_1}\cdots D_{j_k}\nabla_i\varphi(x,y)\bigr\|_{\rm op}
                                                            \le C_k,
 \label{eq:joint-gradient-sums}
\end{equation}
with the same volume and cutoff independence. The indices may
be fine or coarse link-coordinate blocks. The norm of each
multilinear block is its supremum on unit \(\g\)-vectors.
\end{lemma}
\begin{proof}
All derivatives in this proof are ordinary derivatives of the
exponential-chart coordinates. Write \(z=(x,y)\), and extend a
fine direction by zero on tree links. Constants may depend on
the fixed data \(G,L,J,a,\lambda,\sigma,C_{\rm p}\) and the
branch thresholds. They will not depend on the number of
translated blocks or the terminal cutoff. We first obtain real
derivative bounds, then justify analyticity and termwise
differentiation on a common complex neighbourhood.

\smallskip\noindent
\emph{Exponential factors and their derivative placements.}
For \(Y,Z_1,\ldots,Z_r\in\g\), the iterated integral formula is
\begin{equation}
 \begin{split}
 D^r\exp_Y[Z_1,\ldots,Z_r]
  =\sum_{\pi\in S_r}\int_{0<t_1<\cdots<t_r<1}
       e^{t_1Y}Z_{\pi(1)}e^{(t_2-t_1)Y}
          \cdots Z_{\pi(r)}e^{(1-t_r)Y}\,
                                \dd t_1\cdots\dd t_r .
 \end{split}
 \label{eq:action-exponential-simplex}
\end{equation}
Here \(S_r\) is the set of permutations of \(r\) labels.
For \(r=1\), differentiate the equation
\(\partial_t e^{tY}=Ye^{tY}\) in the direction \(Z_1\)
and solve the resulting inhomogeneous linear equation;
it gives \(D\exp_Y[Z_1]=\int_0^1e^{tY}Z_1e^{(1-t)Y}\,\dd t\).
For higher \(r\), differentiate this integral repeatedly.
Inserting the next direction in each exponential interval
partitions the integration region according to the order of
the insertion times. Its pieces are exactly the permutations
and ordered simplices in \eqref{eq:action-exponential-simplex}.
This also proves the formula inductively, including repeated
directions.

All exponential factors in this formula are unitary. In the
normalised Frobenius convention \eqref{eq:physical-norm},
\[
 \|AB\|_{\rm F}\le\|A\|_{\rm op}\|B\|_{\rm F},\qquad
                       \|A\|_{\rm op}\le\sqrt N\|A\|_{\rm F}.
\]
Keep one inserted matrix in Frobenius norm and the remaining
\(r-1\) in operator norm. Each integral has volume \(1/r!\),
which cancels the \(r!\) permutations. Therefore
\begin{equation}
 \|D^r\exp_Y[Z_1,\ldots,Z_r]\|_{\rm F}
       \le N^{(r-1)/2}\prod_{\nu=1}^r|Z_\nu|_{\g},
                                      \qquad r\ge1.
 \label{eq:action-real-exponential-price}
\end{equation}
No factor depending on \(|Y|\) is needed for these real
derivatives. The same bound applies to \(e^{-Y}\), with
the signs of its directions reversed.

After the replica substitution, a positive link is
\(e^{y_e}\) or \(e^{y_e}e^{x_b}\); its reverse is
\(e^{-y_e}\) or \(e^{-x_b}e^{-y_e}\). Tree factors are one.
Thus a path word can be written as
\[
       W(z)=\prod_{v=1}^{m}e^{\epsilon_vz_{i(v)}},
              \qquad \epsilon_v\in\{1,-1\}.
\]
The map \(i(v)\) records the fine or coarse coordinate used
at the \(v\)-th occurrence; it can take the same value at
several occurrences. For joint directions
\(h_1,\ldots,h_r\), apply the product rule by assigning
each labelled derivative to one of the \(m\) occurrences.
The derivatives assigned to a common occurrence are evaluated
by \eqref{eq:action-real-exponential-price}. There are no
norm costs from the undifferentiated unitary factors. Hence
\begin{equation}
 \|D^rW(z)[h_1,\ldots,h_r]\|_{\rm F}
   \lesssim_r\prod_{\nu=1}^r
                     \left(\sum_{v=1}^{m}|h_{\nu,i(v)}|_{\g}\right).
 \label{eq:action-word-direction-price}
\end{equation}
The implicit constants in these exponential and word bounds
depend only on \(G,r\), independently of the block ratio,
period and cutoff. Summing the coordinate derivative block
norms over all ordered input indices therefore costs
\(\lesssim_r m^r\). If an input index occurs more than once,
the occurrences are counted before they are grouped into
that coordinate block. The triangle inequality justifies
the grouping.

\smallskip\noindent
\emph{A complete Wilson calculation.}
For a positive plaquette \(p\), put
\(f_p(x,y)=q_G(U(x,y)_p)\) and
\(\delta=\max(\|x\|_\infty,\|y\|_\infty)\).
There are at most eight exponential occurrences in its word.
For real group-valued inputs,
\(q_G(W)=1-N^{-1}\operatorname{Re}\tr W\), and
\(N^{-1}|\operatorname{Re}\tr A|\le\|A\|_{\rm F}\).
Taking one background direction and two fine directions in
\eqref{eq:action-word-direction-price} therefore gives
\begin{align}
 &\left|D^3f_p(tx,ty)[(x,y),(0,u),(0,u)]\right|\notag\\
 &\qquad\lesssim\delta
              \left(\sum_{e\in\partial p}|u_e|_{\g}\right)^2
       \lesssim\delta\sum_{e\in\partial p}|u_e|_{\g}^2,
                       \qquad 0\le t\le1.
 \label{eq:action-Wilson-third-price}
\end{align}
The implicit constants depend only on \(G\) and include the
fixed factor eight for the background occurrences; a fine
direction occurs only in the four fine exponential factors. The fundamental theorem of
calculus now gives the exact identity
\begin{equation}
 \begin{split}
 &D_y^2f_p(x,y)[u,u]-D_y^2f_p(0,0)[u,u]\\
 &\quad=\int_0^1
       D^3f_p(tx,ty)[(x,y),(0,u),(0,u)]\,\dd t.
 \end{split}
 \label{eq:action-Wilson-Hessian-integral}
\end{equation}
Each positive fine link has six plaquette-boundary occurrences
in dimension four, counting multiplicities if a periodic
identification occurs. Thus
\[
 \sum_{p\in\mathcal P_f^+}\sum_{e\in\partial p}|u_e|_{\g}^2
                              =6\sum_{e\in I}|u_e|_{\g}^2,
\]
where \(u=0\) on trees. Consequently
\begin{equation}
 \lambda\left|D_y^2(W_f\circ U)(x,y)[u,u]
           -D_y^2(W_f\circ U)(0,0)[u,u]\right|
                    \lesssim\lambda\delta\sum_{e\in I}|u_e|_{\g}^2.
 \label{eq:action-Wilson-volume-bound}
\end{equation}
This calculation explains the maximum background norm and
the local square sum. The weighted matrix estimate requires
the corresponding rooted entry sums, which we prove next.

\smallskip\noindent
\emph{Root the path term before summing its translates.}
Consider one path-polynomial template \(\alpha\), translated
to a fine anchor \(v\), with total length \(l_\alpha\) and
coefficient magnitude \(b_\alpha\), including the norm of
its multilinear form. A slot contains two transports along
its path and one four-link plaquette. After replacing a
fine link by at most two exponentials, the total number of
occurrences is at most \(C(1+l_\alpha)\). The curvature map
\(\mathcal C(W)=\operatorname{pr}_{\g}((W-W^{-1})/2)\)
has a fixed linear projection and two word contributions.
There are at most four slots in a term. Their zero-order
norms are at most two, and the product rule and
\eqref{eq:action-word-direction-price} therefore give the
same placement estimate, with coefficient \(C_{G,r}b_\alpha\).
The additional factors depend on the fixed degree, not on
the number of links in its paths.

Fix the first derivative coordinate \(i\), rather than
first summing over all anchors. For a fine coordinate,
choosing its occurrence in the template determines the
translated anchor modulo the fine period. For a coarse
coordinate, its occurrence is in a replica link carrying
that coarse bond; there are at most \(C_L\) choices of the
fine anchor, since only the finitely many residues in its
block are possible. Thus the sum over the translated
anchor and the first differentiated occurrence costs at
most \(C_L(1+l_\alpha)\). Each remaining derivative has
at most \(C(1+l_\alpha)\) occurrences to choose. This
counts the first derivative and translation together;
there is no factor equal to the number of lattice anchors.

The projected paths join every read site to the anchor.
For two derivative coordinates read by this term,
\(d(i,j)\le C_L(1+l_\alpha)\). Consequently, for \(r\ge2\),
the sum over all translates of this one template satisfies
\begin{equation}
 \begin{split}
 &\sup_i\sum_{j_2,\ldots,j_r}e^{\theta d(i,j_2)}
   \left\|D_iD_{j_2}\cdots D_{j_r}
                   \sum_v f_{\alpha,v}(z)\right\|_{\rm op}\\
 &\qquad\lesssim_{L,r}b_\alpha(1+l_\alpha)^r
                                      e^{C_L\theta(1+l_\alpha)}.
 \end{split}
 \label{eq:action-rooted-template-price}
\end{equation}
Here \(f_{\alpha,v}\) includes its coefficient.
The implicit prefactor depends only on \(G,L,r\);
the geometric constant \(C_L\) stays explicit because it
determines the later choice of \(\theta\).
All indices may be fine or coarse. The multilinear norm is the supremum
on unit \(\g\)-vectors, as in the statement of the lemma.
Periodic identification can merge occurrences and translated
anchors; the preceding absolute occurrence count remains
an upper bound. It never replaces native path length by
a shorter periodic distance in a coefficient budget.

The path budget for a row square is explicitly
\begin{align}
 &\sum_{j,k}|r_{bj}r_{bk}|
       e^{\sigma\{(|\gamma_{bj}|+4)+(|\gamma_{bk}|+4)\}}
                                                        \notag\\
 &\qquad=
 \left(\sum_j|r_{bj}|e^{\sigma(|\gamma_{bj}|+4)}\right)^2
                                                \le C_{\rm p}^2.
 \label{eq:action-double-row-budget}
\end{align}
Indeed expand \(\frac12|\ell_b|_{\g}^2\) as its double
slot sum before differentiating. The factor one half is
harmless, and there are six rows per fine anchor. Higher
terms have the budget \eqref{eq:stationary-higher-budget}.
Wilson terms have a fixed length and coefficient \(\lambda\).
Thus all principal templates have a total exponential
coefficient budget bounded by
\(C_G(\lambda+C_{\rm p}^2+C_{\rm p})\). The normalised
signed symmetry average preserves this bound.

Choose \(\theta>0\) so that \(C_L\theta<\sigma/2\).
For each fixed \(r\),
\[
 \sup_{l\ge0}(1+l)^r e^{-\sigma l/2}<\infty.
\]
Summing \eqref{eq:action-rooted-template-price} with these
budgets gives, for the actual principal action,
\begin{equation}
 \sup_i\sum_{j_2,\ldots,j_r}e^{\theta d(i,j_2)}
       \left\|D_iD_{j_2}\cdots D_{j_r}
                      S_f^{\rm un}(U(x,y))\right\|_{\rm op}
                                                \le C_r^{\rm p}.
 \label{eq:action-principal-rooted-sums}
\end{equation}
For example its right side is at most
\(C_{G,L,r,\sigma}(\lambda+C_{\rm p}^2+C_{\rm p})\).
This proves the weighted entry estimate that the quadratic
form calculation alone did not supply.

\smallskip\noindent
\emph{Derivatives of the actual frame.}
We next treat the noise part of
\eqref{eq:stationary-model-action} with the actual predictor
of \eqref{eq:predictor-law}. Work on one of its fixed
patches, regarding the finitely many input link coordinates
as independent before any periodic identification. Write
\(h_\zeta=e^{\xi_\zeta}\), with \(\xi_0=0\), and define
the finite vector equation
\[
 F(z,\xi)=\nabla_\xi
                \mathcal A_{U(z)}((e^{\xi_\zeta})_\zeta)=0.
\]
The gradient and the norm on this rooted frame space use
the counting product of \(\g\)-inner products. Its input
directions use the maximum norm on the finite read set.
By \eqref{eq:frame-flat-hessian},
\(F_\xi(0,0)\ge2s^{-2}1\).
Lemma~\ref{lem:actual-analytic-branch} identifies the unique
nearby solution with the actual regular frame. On a
sufficiently small closed input neighbourhood it satisfies
\begin{equation}
             F_\xi(z,\xi(z))\ge s^{-2}1,
                 \qquad \|F_\xi(z,\xi(z))^{-1}\|_{2\to2}\le s^2.
 \label{eq:action-frame-inverse-price}
\end{equation}
This is a fixed-patch bound; \(s=|\mathcal Q_R|\) is fixed
when the number of blocks varies.

Here are the derivative equations needed for the predictor.
Subscripts \(1,2,3\) denote partial derivatives in three
fixed joint input directions, holding \(\xi\) fixed when
applied to \(F\); write \(\xi_1,\xi_{12},\xi_{123}\)
for the total derivatives of the frame solution.
All partials below are evaluated at \((z,\xi(z))\), and
\(A=F_\xi(z,\xi(z))\). Differentiating the same equation
once, twice, and three times gives
\begin{align}
 A\xi_1={}&-F_1,\notag\\
 A\xi_{12}={}&-F_{12}-F_{1\xi}\xi_2-F_{2\xi}\xi_1
                                      -F_{\xi\xi}[\xi_1,\xi_2],
 \label{eq:action-frame-first-second}\\
 A\xi_{123}={}&-F_{123}
       -\sum_{\rm cyc}F_{12\xi}\xi_3
       -\sum_{\rm cyc}F_{1\xi\xi}[\xi_2,\xi_3]
       -F_{\xi\xi\xi}[\xi_1,\xi_2,\xi_3]\notag\\
 &\quad-\sum_{\rm cyc}F_{1\xi}\xi_{23}
       -\sum_{\rm cyc}F_{\xi\xi}[\xi_1,\xi_{23}].
 \label{eq:action-frame-third}
\end{align}
Each cyclic sum has the three cyclic permutations of
\((1,2,3)\). In particular the last two sums retain the
second derivatives of the frame; they cannot be omitted
in a third derivative estimate.

Each route in the frame score has length at most \(8R\).
Its two frame factors and replica substitution give at
most \(16R+2\) exponential occurrences. To estimate a
partial derivative of \(F\), pair its output with a frame
vector of counting norm one and differentiate the score
once in that vector. Each of its components has norm at
most one. The word bound through order four, and the fact
that the score weights sum to one, therefore bound all
partials of \(F\) through order three by
\(C_F=C_G(1+R)^4\ge1\), in the mixed input maximum and
frame counting norms just specified. There is no sum over
the output components after this duality estimate. Thus
\(C_F\) does not depend on how many patches occur. Inserting
\eqref{eq:action-frame-inverse-price} in the displayed
identities gives the explicit choices
\begin{align}
 B_1&=s^2C_F,\qquad
 B_2=s^2C_F(1+2B_1+B_1^2),\notag\\
 B_3&=s^2C_F(1+3B_1+3B_1^2+B_1^3+3B_2+3B_1B_2),
 \label{eq:action-frame-jet-prices}
\end{align}
for bounds on \(\|\xi_1\|_2,\|\xi_{12}\|_2,
\|\xi_{123}\|_2\) when the three input directions have
maximum norm at most one. Applying the chain rule for
\(e^{\xi_\zeta}\), with
\eqref{eq:action-real-exponential-price}, bounds the first
three derivatives of every frame matrix and its inverse
by fixed constants. For its third derivative the terms
are \(D\exp[\xi_{123}]\), the three terms
\(D^2\exp[\xi_{12},\xi_3]\), and
\(D^3\exp[\xi_1,\xi_2,\xi_3]\). The inverse is
\(e^{-\xi_\zeta}\), so the same estimates apply.

\smallskip\noindent
\emph{Connectors, averaging, projection, and the mismatch.}
For a coarse link \(e=(v,i)\), let \(\gamma_{e,\zeta}\)
denote the original straight connector from \(Lv+\zeta\)
to \(L(v+e_i)+\zeta\). Retain the full connector
\(Y_{e,\zeta}\) of \eqref{eq:actual-connectors}. Its middle
transport has \(L\) fine links and at most \(2L\)
exponential factors. Its \(r\)-th directional derivative
therefore has norm at most \(C_{G,r}(2L)^r\) for unit
maximum-norm input directions. For the full connector,
the precise product rule is
\[
 D_1\cdots D_rY_{e,\zeta}
   =\sum_{A_1\,\dot\cup\,A_2\,\dot\cup\,A_3=\{1,\ldots,r\}}
       (D_{A_1}h_\zeta^v)
       (D_{A_2}U_{\gamma_{e,\zeta}})
       (D_{A_3}(h_\zeta^{v+e_i})^{-1}).
\]
The three subsets are ordered and may be empty;
\(D_\varnothing\) leaves its factor unchanged. There are
\(3^r\) terms. The frame bounds and the preceding word
bound give \(\|D_1\cdots D_rY_{e,\zeta}\|_{\rm F}\le C_{Y,r}\)
for \(r\le3\), independently of \(\zeta\), the period,
and the number of blocks.

Set \(A_e(z)=\sum_\zeta w(\zeta)Y_{e,\zeta}(z)\).
Positivity and normalisation of the original weights give
\begin{equation}
 \|D_1\cdots D_rA_e\|_{\rm F}
       \le\sum_\zeta w(\zeta)C_{Y,r}=C_{Y,r},\qquad r\le3.
 \label{eq:action-connector-average-price}
\end{equation}
At zero, \(A_e=I_N\). Shrink the common input radius so
that all these matrices stay in one closed ball strictly
inside the nearest-point neighbourhood of
Appendix~\ref{app:projection}. On this ball the first three
derivatives of \(\pi_G\) have fixed bounds. For the actual
regular prediction \(Q_e=\pi_G(A_e)\), the full chain rule is
\begin{align}
 (Q_e)_1={}&D\pi_G(A_e)[(A_e)_1],\notag\\
 (Q_e)_{12}={}&D^2\pi_G(A_e)[(A_e)_1,(A_e)_2]
                          +D\pi_G(A_e)[(A_e)_{12}],\notag\\
 (Q_e)_{123}={}&D^3\pi_G(A_e)[(A_e)_1,(A_e)_2,(A_e)_3]
        +\sum_{\rm cyc}D^2\pi_G(A_e)[(A_e)_{12},(A_e)_3]
                          +D\pi_G(A_e)[(A_e)_{123}].
 \label{eq:action-predictor-chain-rule}
\end{align}
If \(C_{\pi,r}\) bounds the \(r\)-th projection derivative
on that closed ball, an explicit third-order choice is
\[
 C_{Q,3}=C_{\pi,3}C_{Y,1}^3
           +3C_{\pi,2}C_{Y,1}C_{Y,2}+C_{\pi,1}C_{Y,3}.
\]
The first two lines give the corresponding \(C_{Q,1}\)
and \(C_{Q,2}\). These bounds include
both frames, the original connector average, and its
projection. Lemma~\ref{lem:actual-analytic-branch} also
keeps both regular tests strict on this chart, so all
unused-tag probabilities are constant here.

Write the noise summand as
\[
 n_e(z)=a q_G(V(x)_eQ_e(z)^{-1})
       =a\{1-N^{-1}\operatorname{Re}\tr(e^{x_e}Q_e(z)^*)\}.
\]
The second expression is used only for real inputs. Matrix
adjoint is real linear and preserves Frobenius norm. Thus
for \(r\ge1\) the exact derivative formula is
\begin{equation}
 D_1\cdots D_rn_e
   =-\frac aN\operatorname{Re}\tr
       \sum_{A\subset\{1,\ldots,r\}}
             (D_Ae^{x_e})(D_{A^c}Q_e)^*.
 \label{eq:action-noise-product-rule}
\end{equation}
The exponential and predictor bounds, and its \(2^r\)
terms, bound every such coordinate derivative through
order three by \(C_{n,r}\). For example normalised
trace Cauchy--Schwarz and
\eqref{eq:action-real-exponential-price} give
\[
 C_{n,3}=a\{C_{Q,3}+3C_{Q,2}
                                    +3\sqrt N\,C_{Q,1}+N\}.
\]
No path replacement of the
predictor has entered this calculation.

For higher fixed orders the same differentiation of
\(F(z,\xi(z))=0\) has just one highest response,
\(F_\xi\xi_{1\cdots r}\), and finitely many products
of lower responses and partial derivatives of \(F\).
Its inverse bound and induction give finite constants
at each order. The product and projection rules then
give \(C_{Q,r},C_{n,r}\) for every fixed \(r\). This
proves all the fixed-order assertions needed in
\eqref{eq:joint-gradient-sums}, without a new global
regularity assumption.

\smallskip\noindent
\emph{The actual predictor's weighted read sums.}
Let \(T_e\) be the joint coordinate read set of \(n_e\),
including both frames, all original connectors and replica
coordinates, and its explicit coarse coordinate \(x_e\).
There are constants \(D_Q,N_Q,R_Q\), fixed by \(L,J\),
such that
\[
 |T_e|\le D_Q,\qquad
 \#\{e:i\in T_e\}\le N_Q,\qquad
                d(i,j)\le R_Q\quad(i,j\in T_e).
\]
Indeed each read lies in the union of the two fixed translated
patches and their length-\(L\) connectors. This union has
finitely many fine and coarse coordinate positions. For a
fixed coordinate, the possible coarse origins of a patch
containing it belong to a fixed box about its parent block.
For example these constants may be chosen to satisfy
\(D_Q\le C(1+J)^4L^4\), \(N_Q\le C(1+J)^4\), and
\(R_Q\le C(1+J)\), after fixed geometric enlargements.
Periodic identification
can identify positions; counting their original occurrences
gives the same bounds on every compatible period. Directional
derivatives in identified variables are obtained by identifying
their directions before the fixed-patch calculation, so the
maximum input norm does not increase.

For \(r\ge2\), only summands with \(i\in T_e\) contribute
to a rooted derivative, and all remaining derivative indices
must belong to that same \(T_e\). The explicit estimate is
\begin{equation}
 \begin{split}
 &\sup_i\sum_{j_2,\ldots,j_r}e^{\theta d(i,j_2)}
       \left\|D_iD_{j_2}\cdots D_{j_r}\sum_e n_e(z)\right\|_{\rm op}\\
 &\qquad\le
 \sup_i\sum_{e:i\in T_e}e^{\theta R_Q}D_Q^{r-1}C_{n,r}
                 \le N_Qe^{\theta R_Q}D_Q^{r-1}C_{n,r}.
 \end{split}
 \label{eq:action-noise-rooted-sums}
\end{equation}
This is the required fixed-patch summation for the actual
predictor. Its constants can depend on the block side, which
the lemma fixes, and are independent of the number of blocks.

\smallskip\noindent
\emph{Analyticity and the exchanges of infinite path sums.}
The preceding real estimates used unitarity. For complexified
coordinate blocks on a ball of radius \(2r_0\), instead use
\(\|e^Y\|_{\rm op}\le e^{\|Y\|_{\rm op}}\).
The simplex formula is still valid and contributes this
exponential factor on each occurrence. Consequently, for
one coefficient-weighted path term and every fixed \(r\),
\begin{equation}
 \sum_{i_1,\ldots,i_r}
        \|D_{i_1}\cdots D_{i_r}f_{\alpha,v}(z)\|_{\rm op}
     \lesssim_r b_\alpha(1+l_\alpha)^r e^{C_Gr_0l_\alpha}.
 \label{eq:word-local-derivative-cost}
\end{equation}
The curvature projection and real multilinear forms are
extended complex linearly. The row inner product is extended
complex bilinearly. For the Wilson price use the holomorphic
extension \(1-\tr(W+W^{-1})/(2N)\). All agree with the
original functions on real group-valued inputs.

The finite-coordinate implicit frame and projection maps
are real analytic on their regular branch, and have
holomorphic extensions to a smaller common patch polydisc.
Extend the resulting actual \(Q_e(z)\) in these input
coordinates. It stays invertible near zero, and the noise
term uses \(Q_e(z)^{-1}\) in its complex extension, rather
than the real-adjoint formula used for
\eqref{eq:action-noise-product-rule}. The same finitely
many patch templates work on every period.

The group and geometry constants \(C_G,C_L\) and the local
analytic neighbourhoods have now been fixed from the profile.
First choose \(\theta>0\) with \(C_L\theta\ll\sigma\),
then choose \(r_0>0\) with \(C_Gr_0\ll\sigma\) and with
the closed real \(2r_0\)-ball in all those neighbourhoods.
In particular we arrange
\(C_L\theta<\sigma/4\) and \(C_Gr_0<\sigma/4\), which give
\(C_Gr_0+C_L\theta<\sigma/2\).
The derivative budgets below are fixed after these choices.
The length polynomials
in \eqref{eq:word-local-derivative-cost} are absorbed by
the remaining path-width reserve, exactly as in the sum
of \eqref{eq:action-rooted-template-price}. At each fixed
period there are finitely many anchors, and their
holomorphic word series converge locally uniformly on
the common polydisc. This proves analyticity and permits
termwise derivatives. Rooting before the anchor sum
gives the uniform derivative majorants; analyticity of
the extensive scalar action does not give or require a
volume-independent bound on its total value.

\smallskip\noindent
\emph{The weighted Schur bounds and the forcing estimate.}
Combine \eqref{eq:action-principal-rooted-sums} and
\eqref{eq:action-noise-rooted-sums}. For every fixed
\(r\ge2\) there is \(K_r\), uniform in period and cutoff,
such that the full action satisfies
\begin{equation}
 \sup_i\sum_{j_2,\ldots,j_r}e^{\theta d(i,j_2)}
                \|D_iD_{j_2}\cdots D_{j_r}\varphi(z)\|_{\rm op}
                                                   \le K_r.
 \label{eq:action-full-rooted-sums}
\end{equation}
The case \(r=2\), restricted to two fine indices, bounds
the weighted Hessian row sums. Symmetry of the scalar
action's Hessian gives the same column sums, proving
\eqref{eq:proved-local-hessian}. Dropping the weight and
taking \(r=k+1\) proves \eqref{eq:joint-gradient-sums},
with the counting inner product identifying the first
derivative slot with the gradient block.

For the difference estimate let
\(H_{ij}(z)=D_{y_i}D_{y_j}\varphi(z)\), regarded as a
bilinear block, or as its associated linear map on
\(\g\). Its operator norm is the same under this
identification. The exact block identity is
\begin{equation}
 H_{ij}(x,y)-H_{ij}(0,0)
       =\int_0^1\sum_k(D_kH_{ij})(tx,ty)[z_k],\dd t,
 \label{eq:action-block-Hessian-integral}
\end{equation}
where \(k\) ranges over both fine and coarse coordinate
blocks. Thus \eqref{eq:action-full-rooted-sums} with
\(r=3\) gives
\begin{align}
 &\sup_i\sum_j e^{\theta d(i,j)}
                 \|H_{ij}(x,y)-H_{ij}(0,0)\|_{\rm op}\notag\\
 &\quad\le\delta\int_0^1\sup_i\sum_{j,k}e^{\theta d(i,j)}
                        \|D_kH_{ij}(tx,ty)\|_{\rm op}\,\dd t
        \le K_3\delta .
 \label{eq:action-weighted-Hessian-difference}
\end{align}
The identical column bound follows from Hessian symmetry.
Since \(H(0,0)=H_0\), this proves
\eqref{eq:proved-hessian-perturbation} in its stated
weighted norm, rather than only in counting \(\ell^2\).

Finally \(D_y\varphi(0,0)=0\). For a Wilson or noise
summand this follows from \(Dq_G(1)=0\). Every curvature
slot is zero at the flat field, so row squares and
degree-three or degree-four terms also have zero first
derivative there. For a fine coordinate \(i\), integrate
the mixed second derivatives along \((tx,0)\):
\[
 (\nabla_y\varphi(x,0))_i
     =\int_0^1\sum_{b\ {\rm coarse}}
             (D_{x_b}\nabla_{y_i}\varphi)(tx,0)[x_b],\dd t.
\]
Taking norms, then the supremum over \(i\), and using
the joint \(r=2\) bound gives
\(\|D_y\varphi(x,0)\|_\infty\le K_2\|x\|_\infty\).
This is \eqref{eq:proved-gradient-forcing}. Enlarge
\(K\) to the maximum of the two required fixed-order
constants. All conclusions hold on the same chart and
with the same period and cutoff uniformity.
\end{proof}

This proves the bounds needed here from the specified action,
as in \cite[Lemma~14.13, pp.~67--68]{Main}, whose current
PDF locators were rechecked on 17 November 2026.
Only the derivative sums are uniform. The value of the entire
action is an extensive scalar and can grow with the number
of sites.

\subsection{From a gap and spatial decay to a maximum-norm inverse}

The weighted matrix bound also controls powers of a scaled
Hessian. Combining this control with their spectral decay
makes a direct Neumann series exponentially summable.
We use the weighted row and column norm already defined in
\eqref{eq:weighted-local-matrix}. For comparison,
\cite[Section~2.4, Lemma~2.12, p.~14]{DST23} proves inverse
decay for scalar lattice Green functions by weighted conjugation.
The following Neumann-series proof uses our block norm directly.

\begin{lemma}[A volume-uniform weighted inverse]
\label{lem:uniform-inverse}
Suppose \(H=H^*\ge\eta\,1\) in counting \(\ell^2\), and
\(\|H\|_\theta\le K\), with at most \(m_0\) coordinate blocks
attached to a coarse site. Set
\begin{equation}
               \mu=\min\{\theta/2,\eta\theta/(8K)\}.
 \label{eq:inverse-rate}
\end{equation}
Then
\begin{equation}
       \|(H^{-1})_{ij}\|_{\rm op}\le\frac2\eta e^{-\mu d(i,j)},
       \qquad
       \|H^{-1}\|_{\mu/2}
              \le C m_0\eta^{-1}(1+\mu^{-1})^4.
 \label{eq:proved-weighted-inverse}
\end{equation}
The numerical constant \(C\) is independent of the torus side
and the number of coordinates. The same bounds hold for a
principal restriction to a fixed subset of the coordinate blocks.
\end{lemma}
\begin{proof}
An empty coordinate set is immediate. Otherwise the Schur
estimate gives \(0<\eta\le\|H\|_{2\to2}\le K\). Define
\[
              A=1-\frac{H}{2K},\qquad
              \alpha=\frac{\eta}{4K}\in(0,1/4].
\]
The spectrum of \(A\) lies in \([1/2,1-2\alpha]\), so
\(\|A\|_{2\to2}<1\) and
\[
                H^{-1}=\frac1{2K}\sum_{n=0}^{\infty}A^n
\]
in counting \(\ell^2\) operator norm. We show that the same
series has exponential entry decay. The spectral bound
and submultiplicativity of the weighted matrix norm give,
for every \(n\ge0\),
\[
 \|A^n\|_{2\to2}\le(1-2\alpha)^n,\qquad
 \|A\|_\theta\le1+\frac{\|H\|_\theta}{2K}\le\frac32,
 \qquad \|A^n\|_\theta\le(3/2)^n.
\]
Consequently every block obeys both
\[
 \|(A^n)_{ij}\|_{\rm op}\le(1-2\alpha)^n,\qquad
 \|(A^n)_{ij}\|_{\rm op}\le(3/2)^n e^{-\theta d(i,j)}.
\]
If a nonnegative number is bounded by \(a\) and \(b\),
it is bounded by \(a^{1-\alpha}b^\alpha\). Applying this
to the two entry bounds yields
\[
 \|(A^n)_{ij}\|_{\rm op}
 \le\bigl[(1-2\alpha)^{1-\alpha}(3/2)^\alpha\bigr]^n
                                      e^{-\alpha\theta d(i,j)}.
\]
The weighted arithmetic--geometric mean inequality and
\(0<\alpha\le1/4\) give
\[
 \begin{split}
 (1-2\alpha)^{1-\alpha}(3/2)^\alpha
 &\le(1-\alpha)(1-2\alpha)+\tfrac32\alpha\\
 &=1-\tfrac32\alpha+2\alpha^2
 \le1-\alpha.
 \end{split}
\]
The entry series is therefore absolutely convergent, with
\[
 \begin{split}
 \|(H^{-1})_{ij}\|_{\rm op}
 &\le\frac1{2K}\sum_{n\ge0}(1-\alpha)^n
                                      e^{-\alpha\theta d(i,j)}\\
 &=\frac1{2K\alpha}e^{-\alpha\theta d(i,j)}
   =\frac2\eta e^{-\alpha\theta d(i,j)}.
 \end{split}
\]
Since \(\mu\le\eta\theta/(8K)=\alpha\theta/2\), this
proves the first estimate in
\eqref{eq:proved-weighted-inverse}. The argument uses
exponential summability, with no finite-range condition.

On any four-dimensional torus the number of sites at periodic
maximum distance \(t\in\Z_{\ge0}\) is at most
\(C(1+t)^3\), uniformly in its side. Therefore
\[
       \sup_i\sum_l e^{-\mu d(i,l)/2}
          \le C m_0\sum_{t\ge0}(1+t)^3e^{-\mu t/2}
          \le C m_0(1+\mu^{-1})^4.
\]
Combining this with the entry estimate proves the weighted
row and column bound. A principal restriction retains the
gap by extending test vectors by zero. It decreases the
weighted sums and shell counts, so the same argument applies.
\end{proof}

The estimate is \cite[Lemma~14.14 and (14.30), p.~68]{Main}.
Its row bound gives in particular
\(\|H^{-1}\|_{\infty\to\infty}\le\|H^{-1}\|_{\mu/2}\).
For the actual \(H_0\), the proved gap, upper bound and
locality estimate can now be summarised at a fixed profile as
\[
                H_0\simeq1,\qquad
                \|H_0^{-1}\|_{\infty\to\infty}\lesssim1.
\]
These constants may depend on the already fixed block ratio.
The first comparison concerns quadratic forms; the second
also uses spatial locality. The quantitative dependence on
\(\eta,\theta,K,m_0\) remains visible in the lemma whenever
a later choice uses the gap or decay rate. There is no
factor equal to the full active dimension.

\section{Constructing the stationary graph on a common chart}
\label{sec:uniform-stationary}

We now have both inputs which the finite-dimensional central
calculation left open: an actual positive flat Hessian, and local
bounds controlling its variation. The weighted inverse estimate
lets us solve the stationary equation in the same maximum norm
as the coordinate chart. The radii will be chosen from these
proved bounds, not from the total number of variables.

\begin{theorem}[The stationary graph and its Hessian bounds]
\label{thm:uniform-stationary}
Fix the profile of Section~\ref{sec:smooth-principal}, in the
admissible period regime of Lemma~\ref{lem:block-poincare}.
There are positive constants \(r_x,r_y,\eta,K_{\rm H},\nu\),
fixed before the period and terminal UV cutoff vary,
with the following properties.
For \(\|x\|_\infty<r_x\), the equation
\begin{equation}
                   D_y\varphi(x,y)=0,\qquad \|y\|_\infty\le r_y
 \label{eq:model-stationary-equation}
\end{equation}
has a unique solution \(y=m(x)\), and
\begin{equation}
       m(0)=0,\qquad
       \|m(x)\|_\infty\lesssim\|x\|_\infty,
       \qquad \|m(x)\|_\infty\le r_y/2.
 \label{eq:model-centre-size}
\end{equation}
On the whole common chart,
\begin{equation}
       \eta\,1\le D_y^2\varphi(x,y)\le K_{\rm H}1
             \quad(\|x\|_\infty\le r_x,\ \|y\|_\infty\le r_y).
 \label{eq:model-two-sided-hessian}
\end{equation}
The graph is real analytic. For each fixed derivative order,
its derivatives as multilinear maps between maximum-norm
coordinate spaces have volume- and cutoff-independent bounds.
Moreover
\begin{equation}
       H(x):=D_y^2\varphi(x,m(x)),\qquad
       \|H(x)^{-1}\|_\nu\lesssim1.
 \label{eq:model-moving-inverse}
\end{equation}
The same graph serves every unused branch tag on this chart.
\end{theorem}
\begin{proof}
Put \(\mathcal G(x,y)=D_y\varphi(x,y)\). Apply
Lemma~\ref{lem:uniform-inverse} to the actual \(H_0\), using
\eqref{eq:proved-flat-gap} and \eqref{eq:proved-local-hessian}.
It gives a constant \(C_0\) with
\(\|H_0^{-1}\|_{\infty\to\infty}\le C_0\), independent of
the period. For fixed \(x\), consider
\begin{equation}
                         \mathcal T_x(y)=y-H_0^{-1}\mathcal G(x,y).
 \label{eq:model-contraction}
\end{equation}
On a proposed chart of radii \(r_x,r_y\), the proved action
bounds give
\begin{align*}
 \|D_y\mathcal T_x(y)\|_{\infty\to\infty}
   &=\|H_0^{-1}(H_0-D_y\mathcal G(x,y))\|_{\infty\to\infty}
       \le C_0K\max(r_x,r_y),\\
 \|\mathcal T_x(0)\|_\infty
   &\le C_0K\|x\|_\infty.
\end{align*}
For contraction and invariance of the closed \(r_y\)-ball it
therefore suffices to arrange
\begin{equation}
       C_0K\max(r_x,r_y)\le\tfrac14,
       \qquad C_0Kr_x\le r_y/4.
 \label{eq:derived-contraction-conditions}
\end{equation}
For positivity on that same ball, the Schur estimate and
\eqref{eq:proved-hessian-perturbation} show that it suffices to
arrange
\begin{equation}
                         K\max(r_x,r_y)\le\eta_0/2.
 \label{eq:derived-gap-condition}
\end{equation}
Collecting the chart, contraction and gap requests gives the
following order of choices:
\[
 0<r_y\ll\min\{r_0,r_G,(C_0K)^{-1},\eta_0/K\},
 \qquad
 0<r_x\ll\min\{r_y,r_y/(C_0K)\}.
\]
First fix the profile and the bounds \(r_0,\theta,K,\eta_0,C_0\)
already obtained. Choose \(r_y\), and then \(r_x\), in these
ranges. The smallness thresholds are chosen so that
\(r_y<r_G/8\) and
\eqref{eq:derived-contraction-conditions}--%
\eqref{eq:derived-gap-condition} hold with their displayed
constants. This is a fixed-profile choice: the radii may
depend on \(L\), but not on the period or terminal cutoff.
No dimension-dependent norm conversion enters.

The map \(\mathcal T_x\) has Lipschitz constant at most \(1/4\)
and maps the closed \(r_y\)-ball into the \(r_y/2\)-ball.
Starting from \(y^{(0)}=0\), let
\(y^{(j+1)}=\mathcal T_x(y^{(j)})\). Then
\[
       \|y^{(j+1)}-y^{(j)}\|_\infty
                    \le4^{-j}\|\mathcal T_x(0)\|_\infty.
\]
The geometric sum gives a limit \(m(x)\) in the ball.
Continuity makes it a fixed point, hence a solution of
\eqref{eq:model-stationary-equation}. The same contraction
inequality proves uniqueness, and
\[
       \|m(x)\|_\infty
            \le(1-1/4)^{-1}C_0K\|x\|_\infty.
\]
The image-ball bound gives
\(\|m(x)\|_\infty\le r_y/2\). At \(x=0\),
\(\mathcal G(0,0)=0\), so uniqueness gives \(m(0)=0\).

For all points of the common chart,
\eqref{eq:derived-gap-condition} gives
\(D_y^2\varphi\ge\eta_0/2\); put \(\eta=\eta_0/2\).
The upper bound follows from
\(\|D_y^2\varphi\|_{2\to2}\le\|D_y^2\varphi\|_\theta\le K\).
Take \(K_{\rm H}=\max(K,1)\).
Apply Lemma~\ref{lem:uniform-inverse} with this common lower
gap \(\eta\), the same \(\theta,K,m_0\), and put \(\nu=\mu/2\)
for its resulting \(\mu\). This proves
\eqref{eq:model-moving-inverse}, even for Hessians throughout
the common chart.

At each point of the graph the analytic implicit function
theorem applies, since \(D_y\mathcal G\) is invertible.
It gives a local analytic solution. Uniqueness in the common
ball identifies these local solutions with \(m\), proving
analyticity on the entire coarse ball.

We verify the claimed uniform derivative bounds separately.
For a coarse direction \(u\),
\[
       Dm(x)[u]
          =-H(x)^{-1}D_x\mathcal G(x,m(x))[u].
\]
The inverse bound and \eqref{eq:joint-gradient-sums} control this
in the maximum norm. For higher orders set
\(Z(x)=(x,m(x))\) and differentiate
\(\mathcal G(Z(x))=0\). For \(k\ge2\), the chain rule gives
\begin{equation}
 \begin{split}
 H(x)D^km(x)[u_1,\ldots,u_k]
   =-\sum_{\substack{\pi\text{ a partition of }\{1,\ldots,k\}\\
                      |\pi|\ge2}}
       D^{|\pi|}\mathcal G(Z(x))
          \bigl[D^{|B|}Z(x)[u_B]:B\in\pi\bigr].
 \end{split}
 \label{eq:model-graph-derivative-recursion}
\end{equation}
Here a partition is a collection of nonempty disjoint blocks
whose union is \(\{1,\ldots,k\}\), and \(u_B\) denotes the
directions indexed by that block. The omitted one-block term
is exactly \(H D^km\), because the first coordinate
of \(Z\) is linear. Every block in the remaining sum has size
at most \(k-1\). Its factors are therefore bounded inductively
using lower derivatives of \(m\), the uniform gradient
derivative bounds and \(H^{-1}\). At fixed \(k\) the
number of partitions depends only on \(k\), so this induction
has no volume factor.

Finally the chart was chosen within the common regular branch
region. All unused tags give the same \(Q\) there, hence the
same \(\varphi\), stationary equation and unique graph.
\end{proof}

This is the stationary conclusion of
\cite[Theorem~14.17 and (14.39)--(14.40), p.~73]{Main}.
The present proof supplies its maximum-norm construction and
fixed-order maximum-norm derivative bounds.
Theorem~\ref{thm:local-analytic-graph} below strengthens them to
rooted spatial bounds by retaining the path weights in the
same stationary equation.

\subsection{The two matrix inputs now follow from one action}

For the specified fixed profile, the preceding conclusions
can be written concisely as
\[
       H(x)\simeq1,\qquad
                   \|H(x)^{-1}\|_\nu\lesssim1.
\]
The weighted inverse estimate is separate from the
quadratic-form comparison. In a determinant price the
dimension must still be kept explicit: in the all-active
real dimension \(d=|I|\dim_{\R}\g\),
\begin{equation}
       \eta\,1\le H(x)\le K_{\rm H}1,
       \qquad \sqrt{\det H(x)}\le K_{\rm H}^{d/2}.
 \label{eq:model-determinant-upper}
\end{equation}
The dimension in this determinant price is retained explicitly.

For the centre-energy comparison we need a different matrix,
the Hessian of \(x\mapsto\varphi(x,0)\). Equation
\eqref{eq:joint-gradient-sums} gives the uniform row bound
\begin{equation}
       \sup_x\sup_i\sum_j
            \|D_{x_i}D_{x_j}\varphi(x,0)\|_{\rm op}\lesssim1
               \quad(\|x\|_\infty\le r_x).
 \label{eq:model-coarse-row-bound}
\end{equation}
The indices here are coarse-link blocks, whereas the matrix in
\eqref{eq:model-determinant-upper} acts on free fine coordinates.
Both bounds come from the same profile budgets, but they
control different derivatives.

To see the scalar consequence, apply Taylor's formula to
\(t\mapsto\varphi(tx,0)\). Its value and first derivative at
zero vanish, so symmetry and the Schur bound give
\begin{equation}
       |\varphi(x,0)|
            \lesssim\sum_{e\in E_c^+}|x_e|_{\g}^2
            \lesssim c^4\|x\|_\infty^2.
 \label{eq:model-reference-energy}
\end{equation}
In particular this is an extensive energy estimate, not a
volume-independent bound on the total scalar action.

\subsection{Local coercivity and the central chart}

For fixed \(x\), the whole maximum-norm \(r_y\)-ball is
convex. Taylor's integral formula on the segment from \(m(x)\)
to \(y\), using stationarity and
\eqref{eq:model-two-sided-hessian}, proves
\begin{equation}
       \varphi(x,y)-\varphi(x,m(x))
              \ge\frac\eta2\sum_{e\in I}|y_e-m_e(x)|_{\g}^2
              \qquad(\|y\|_\infty\le r_y).
 \label{eq:model-local-coercivity}
\end{equation}
Thus \(\varphi(x,m(x))\le\varphi(x,0)\), and
\eqref{eq:model-reference-energy} bounds the stationary value
from above. Also \(|y-m(x)|>gD\) inside this chart has
relative action price at least \(\eta D^2/2\) after multiplication
by \(g^{-2}\).

The central coordinate condition can now be checked using
the constructed graph. If \(gD\le r_y/2\), then for \(|z|\le D\),
\[
       \|m(x)+gz\|_\infty
            \le r_y/2+gD\le r_y<r_G/8.
\]
The central ball lies in the common smooth chart and all its
active exponential cutoffs equal one. The degree-two path
argument for \(I_f^{\rm un}\), and the smooth positive Haar density,
make \(\psi\) smooth there. The stationary and chart inputs
for the exact central identities are therefore available for
this regular all-active reference contribution.
The physical all-good indicator \(\chi_{\rm sf}\) and the
remaining input multipliers still have their explicit place in
\(\theta\), as in \eqref{eq:actual-central-gaussian}.

For the positive reference contribution, the remaining comparison
concerns configurations outside this common chart, including the
other sectors of its Haar decomposition.
Section~\ref{sec:physical-tail} controls that physical complement
relative to the full Gaussian normalisation on auxiliary tori.
The incoming multipliers remain in the separate remainder of
\eqref{eq:reference-exact-splitting}; their later estimates are
distinct from this reference comparison.

\section{Comparing the physical reference tail with the Gaussian scale}
\label{sec:physical-tail}

The complement of the central ball contains configurations which
need not lie in the common coordinate chart. Local convexity cannot
estimate that whole complement. The actual positive reference
integral supplies two other pieces of information: its fine
plaquettes are small, and its kernel penalises the discrepancy
between the predicted and prescribed coarse links. We first show
that small values of both quantities force entry into the common
chart. A configuration which fails that test pays a definite noise
price instead. The two cases will give a bound relative to the
complete Gaussian prefactor, on the auxiliary volumes where its
dimension cost can be paid.

Throughout this section the principal profile is fixed as in
Section~\ref{sec:smooth-principal}. All local geometric arguments
use the same admissible nonwrapping period regime as
Lemma~\ref{lem:block-poincare}. Constants may depend on that profile
and the fixed branch parameters; they do not depend on the number
of blocks or the terminal UV cutoff.

\subsection{Stability on the set actually retained by the indicator}

The reference density contains \(\chi_{\rm sf}\). We therefore
need positivity of its principal action on this good set, rather
than a global convexity assertion. The nonnegative row squares
already have the correct sign. Each higher term contains at least
three curvatures, so its extra factors can be paid using the fine
plaquette threshold.

\begin{lemma}[Good-set stability]
\label{lem:tail-good-stability}
There exists \(h_0>0\), determined by the fixed principal path data,
such that
\begin{equation}
       S_f^{\rm un}(U)\ge\frac\lambda2 W_f(U)\ge0
          \quad\text{if }\chi_{\rm sf}(U)=1,\quad h\le h_0.
 \label{eq:tail-good-stability}
\end{equation}
On the same set the loop allowance gives
\begin{equation}
                         |I_f^{\rm un}(U)|\le f^4\Lambda h^2.
 \label{eq:tail-good-loop-price}
\end{equation}
\end{lemma}
\begin{proof}
Let \(0\le a_1,\ldots,a_k\le h\), with \(k\ge2\).
Putting \(b_j=a_j/h\) when \(h>0\), we have
\[
       \prod_j b_j\le\left(\prod_j b_j\right)^{2/k}
                         \le\frac1k\sum_jb_j^2.
\]
The first inequality uses \(\prod_jb_j\le1\), and the second
is the arithmetic--geometric mean inequality. Hence
\begin{equation}
                 \prod_{j=1}^k a_j
                     \le\frac{h^{k-2}}k\sum_{j=1}^k a_j^2.
 \label{eq:bounded-curvature-product}
\end{equation}
The case \(h=0\) follows directly.

Each transported curvature has the norm of its terminal
curvature, bounded by that plaquette chord. For a fixed path
term of degree \(k\), let \(p_j(o)\) denote the terminal plaquette
of its \(j\)-th slot at anchor \(o\). Summing \eqref{eq:bounded-curvature-product}
over its anchors gives
\[
 \begin{split}
       \sum_o|p_\nu(U;o)|
       &\le |a_\nu|\,\|T_\nu\|\frac{h^{k-2}}k
             \sum_{j=1}^k\sum_o
                 d_G(U_{p_j(o)})^2\\
       &\le 2|a_\nu|\,\|T_\nu\|h^{k-2}W_f(U).
 \end{split}
\]
For the last inequality, translation permutes the anchors of
each orientation and \(d_G(P)^2=2q_G(P)\). No number of
anchors is introduced. Summing the coefficient budget for
degrees three and four yields
\[
       \left|\sum_{o,\nu}p_\nu(U;o)\right|
                     \le2C_{\rm p}(h+h^2)W_f(U).
\]
Signed symmetry averaging preserves this bound and the
nonnegativity of the row squares. Choose \(h_0\) so that
\(2C_{\rm p}(h_0+h_0^2)\le\lambda/2\). This proves
\eqref{eq:tail-good-stability}. For \(I_f^{\rm un}\), both slots
have norm at most \(h\); its bilinear coefficient budget is
at most \(\Lambda\) per anchor. Summing over the \(f^4\)
anchors gives \eqref{eq:tail-good-loop-price}.
\end{proof}

These are the stability estimates of
\cite[Lemma~15.2, (15.2)--(15.3), p.~74]{Main}.

\subsection{Small plaquettes and small mismatch force chart entry}

The additive block reconstruction used square moves to turn a
closed path sum into a sum of curls. Its group-valued version
uses the same moves: changing the order of two perpendicular
steps inserts a conjugate plaquette holonomy. The chord triangle
inequality replaces the additive triangle inequality.

\begin{lemma}[Finite-path comparison]
\label{lem:tail-path-comparison}
Suppose two paths with the same endpoints in a fixed integer
box are joined by a sequence of backtrack cancellations and
elementary square moves. Then
\begin{equation}
       d_G\bigl(U_{\gamma_1}U_{\gamma_2}^{-1}\bigr)
               \le\sum_{\text{square moves }p}d_G(U_p).
 \label{eq:tail-path-comparison}
\end{equation}
The finite collection of frame, connector and block-tree path
comparisons needed at one block admits a common filling bound,
depending only on \(L,J\).
\end{lemma}
\begin{proof}
If \(\gamma_1=\alpha p_1\omega\) and
\(\gamma_2=\alpha p_2\omega\) differ by replacing two sides
of one square, then
\[
       U_{\gamma_1}U_{\gamma_2}^{-1}
           =U_\alpha\,U_{p_1p_2^{-1}}\,U_\alpha^{-1}.
\]
The middle loop is an oriented plaquette. Conjugation and
orientation do not change its chord. A backtrack cancels
exactly. Applying the chord triangle inequality along the
sequence proves \eqref{eq:tail-path-comparison}.

To construct the required moves, choose coordinate-ordered
paths from a box root to all its vertices. For one link,
compare the root path to its first endpoint followed by
that link with the root path to its second endpoint.
Commute the extra step through the later coordinate segments
and cancel backtracks. All squares stay in the box, and at
most \(C(1+R_*)\) moves are needed if its radius is \(R_*\).
Applying this to the links of a path reduces it to the
coordinate path with the same endpoints. A length-\(l\)
path therefore costs at most \(C(1+R_*)l\) squares.
The relevant coordinate paths and connectors have fixed
lengths, and each tree path has length at most \(L^4-1\).
Their common filling bound is finite and independent of the
total period. Repeated uses of a square are counted with
multiplicity.
\end{proof}

For a coarse block \(v\), choose one fixed enlarged box which
contains its links, the two frame patches and all connectors
for each adjacent coarse bond, and the necessary block-tree
paths and fillings. Let \(\mathcal P_v\) be its plaquettes,
and put
\[
       \epsilon_v(U)=\sum_{p\in\mathcal P_v}d_G(U_p),
       \qquad |\mathcal P_v|\le m_*.
\]
The bound \(m_*\) depends on \(L,J\); there is no summation
over all blocks in this local definition.

\begin{proposition}[Nonlinear block reconstruction]
\label{prop:tail-chart-entry}
There exists \(\delta_{\rm g}>0\), depending only on the fixed
local data, with the following property. If every fine plaquette chord is at most
\(\delta_{\rm g}\), all actual branches use the same regular
predictor \(Q(U)\). For a tree-gauge free link \(e\) in block
\(v\),
\begin{equation}
 d_G\bigl(U_eR_e(V)^{-1}\bigr)
       \lesssim\epsilon_v(U)
       +\begin{cases}
           0,&e\text{ is internal},\\
           d_G\bigl(Q(U)_{(v,i)}V_{(v,i)}^{-1}\bigr),
                                      &e\text{ crosses in direction }i.
         \end{cases}
 \label{eq:tail-nonlinear-reconstruction}
\end{equation}
The claim allows arbitrary prescribed coarse links \(V\).
\end{proposition}
\begin{proof}
Fix the enlarged box at \(v\), and abbreviate
\(\epsilon=\epsilon_v(U)\). For a frame rooted at a coarse
root \(Lw\), let \(a_{w,z}\) be its coordinate-ordered
transport from \(Lw\) to \(Lw+z\). The closed paths in
the regular frame test have chords at most \(C\epsilon\)
by Lemma~\ref{lem:tail-path-comparison}. Here and below
the finite filling multiplicities are included in \(C\).
Choosing \(\delta_{\rm g}\) so that
\(Cm_*\delta_{\rm g}<c_{\rm fr}\) makes every such test
strict. Thus the frame is the unique minimiser of
\eqref{eq:actual-frame-score} for every tag.

We need its quantitative proximity to the axial frame.
Testing the score with \(h_z=a_{w,z}\) gives a value at
most \(C\epsilon^2\). The minimising score is no larger.
Every term of its finite nonnegative average is consequently
at most \(s^2 4! C\epsilon^2\). Use the term with endpoints
\(0,z\) and the order defining \(a_{w,z}\). Since \(h_0=1\),
\(q_G=d_G^2/2\) gives
\begin{equation}
                 d_G(h_z a_{w,z}^{-1})\lesssim\epsilon.
 \label{eq:tail-frame-axial-comparison}
\end{equation}
The factor \(s^2 4!\) is part of a fixed patch constant,
not a factor depending on the number of blocks.

Now fix a coarse bond \(b=(w,i)\). Choose a path from
\(Lw\) to \(L(w+e_i)\) which follows the first block tree,
crosses one fixed face link and returns along the second
block tree. Denote its transport by \(W_b\). For each
connector, replacing its two frames by the axial frames
changes its transport by at most \(C\epsilon\), by
\eqref{eq:tail-frame-axial-comparison} and bi-invariance
of the matrix norm. The resulting transport is a path
between those same coarse roots. Finite-path comparison
with \(W_b\) therefore gives
\begin{equation}
                              d_G(Y_zW_b^{-1})\lesssim\epsilon
                                    \quad\text{for every connector }z.
 \label{eq:tail-connectors-common-path}
\end{equation}
Testing the predictor score at \(u=W_b\) bounds its
minimum by \(C\epsilon^2\). Decrease \(\delta_{\rm g}\)
so that this is below \(c_{\rm pr}\). The predictor
therefore uses its regular branch for every tag.

Its distance to \(W_b\) has an especially direct estimate.
Write \(A_b=\sum_z w(z)Y_z\). The already proved
predictor formula says \(Q_b=\pi_G(A_b)\). Since
\(W_b\in G\) is an admissible nearest-point competitor,
\[
       \|Q_b-W_b\|_{\rm F}
          \le\|Q_b-A_b\|_{\rm F}+\|A_b-W_b\|_{\rm F}
          \le2\|A_b-W_b\|_{\rm F}\lesssim\epsilon.
\]
This uses only \(\sum_z w(z)=1\), without a reciprocal
of a smallest positive weight.

For an internal free link \(e\), its two tree paths
close it to a loop. The loop transport is \(U_e\)
in tree gauge, so \(d_G(U_e)\le C\epsilon\).
For a crossing link, the tree path through that link
has transport \(U_e\). Comparing it with the path
defining \(W_b\) gives
\(d_G(U_eW_b^{-1})\le C\epsilon\); the preceding predictor
bound then gives \(d_G(U_eQ_b^{-1})\le C\epsilon\).
The chord triangle inequality with \(V_b\), and the
definition of \(R_e(V)\), prove
\eqref{eq:tail-nonlinear-reconstruction}.
\end{proof}

\begin{corollary}[Entry into a prescribed common chart]
\label{cor:tail-common-chart}
For every sufficiently small fixed \(r_*>0\), there is
\(\delta_*>0\), with the same local parameter dependence,
such that the following holds. If all fine plaquette chords
and all blocking mismatches
\(d_G(Q(U)_bV_b^{-1})\) are at most \(\delta_*\), a
tree-gauge configuration has a unique representation
\(U=U(V,y)\) with \(\|y\|_\infty\le r_*\).
\end{corollary}
\begin{proof}
Take \(0<\delta_*\le\delta_{\rm g}\) with
\(\delta_*\ll r_*\), in terms of the fixed local data.
Each right side of \eqref{eq:tail-nonlinear-reconstruction}
is \(\lesssim(m_*+1)\delta_*\lesssim\delta_*\), since
\(m_*\) is fixed by those data. On a fixed neighbourhood
of the identity the exponential chart has a unique
logarithm and
\( |\log W|_{\g}\lesssim_G d_G(W)\).
Indeed its smooth inverse fixes the identity and has
bounded derivative on a smaller neighbourhood.
The threshold in \(\delta_*\ll r_*\) is chosen to put
these chords in that neighbourhood and to ensure that their
logarithms have norm at most \(r_*\).
Set \(y_e=\log(U_eR_e(V)^{-1})\) for all free links.
Injectivity of the chart gives uniqueness.
\end{proof}

The reconstruction and its chart consequence are
\cite[Lemma~14.8, Proposition~14.9, (14.18), and
Corollary~14.10, pp.~63--64]{Main}. The use of the nearest
group point above replaces the individual positive-weight
estimate for the predictor by one estimate on its matrix mean.


\subsection{The two action prices outside the central ball}

Choose \(\delta_*\) from Corollary~\ref{cor:tail-common-chart}
with \(r_*=r_y\), and require \(h\le\min(h_0,\delta_*)\).
For \(g=\beta^{-1/2}\) and \(gD\le r_y/2\), put
\begin{equation}
       \mathcal C_x
         =\{U(x,m(x)+gz):|z|\le D\}.
 \label{eq:tail-central-set}
\end{equation}
This is a subset of the common chart. Its definition does
not yet assert that the physical good-field indicator is one.
For a branch tag \(t\), use the actual reference action
\[
       A_t(V,U)=S_f^{\rm un}(U)
                     +a\sum_{b\in E_c^+}q_G(V_bQ_t(U)_b^{-1}).
\]
On the good set under consideration, the preceding construction
already makes \(Q_t(U)=Q(U)\) independent of \(t\).

If a good tree-gauge configuration lies outside
\(\mathcal C_x\), then either all its mismatches are at
most \(\delta_*\), or one is larger. In the first case
Corollary~\ref{cor:tail-common-chart} puts it in the form
\(U(x,y)\) with \(\|y\|_\infty\le r_y\). Uniqueness
of the coordinates and \(U\notin\mathcal C_x\) imply
\(|y-m(x)|>gD\). The already proved local coercivity
therefore gives
\begin{equation}
       \beta\{A_t(V(x),U)-\varphi(x,m(x))\}
                                   \ge\eta D^2/2.
 \label{eq:tail-small-mismatch-price}
\end{equation}
In the second case, stability and nonnegativity of every
noise summand give
\begin{equation}
                             A_t(V(x),U)\ge a\delta_*^2/2.
 \label{eq:tail-large-mismatch-price}
\end{equation}
These are \cite[Proposition~15.3, (15.4)--(15.5), p.~75]{Main}.
The first price is already relative to the stationary value.
For the second we must still pay that value. The proved
coarse Taylor bound supplies exactly the needed upper bound:
\begin{equation}
       \varphi(x,m(x))\le\varphi(x,0)
                   \lesssim c^4\|x\|_\infty^2
                   \lesssim f^4\|x\|_\infty^2.
 \label{eq:tail-centre-energy-price}
\end{equation}
Only an upper bound is used here. There is no need to bound
the absolute stationary value by another joint Hessian estimate.

\subsection{The tail divided by its full stationary prefactor}

Let \(n_f=f^4\) and \(d=|I|\dim_{\R}\g\). Work directly
with the tree Haar integral, so that the tail includes all
sectors outside the central coordinate set. Define
\begin{equation}
 \begin{split}
 T_\beta(x)={}&z_{\mathrm n}(\beta)^{-b_c}
       \sum_t\int_{G^I\setminus\mathcal C_x}
          \chi_{\rm sf}(U)\lambda_t(U)\\
       &\hspace{12mm}\cdot
          e^{-\beta A_t(V(x),U)-I_f^{\rm un}(U)}\,\dd h_G^I(U).
 \end{split}
 \label{eq:tail-full-physical-integral}
\end{equation}
The trees are fixed to one in this integral. It is finite and
nonnegative by the reference bounds and the actual kernel law.
Inserting the sector partition of unity would give the same
quantity: its multipliers sum to one pointwise. In particular
no factor counting the number of sectors is needed.

For the all-active stationary chart restore the coarse
parameter in \(\psi\) and define
\begin{align}
       \psi(x,y)&=I_f^{\rm un}(U(x,y))
                            -\sum_{e\in I}\log j_G(y_e),\nonumber\\
       \psi_*(x)&=\psi(x,m(x)),\nonumber\\
       N_g^{\mathrm{phys}}(x)&=z_{\mathrm n}(\beta)^{-b_c}
          \zeta_G^{|I|}(2\pi)^{d/2}g^d
          \frac{e^{-\beta\varphi(x,m(x))-\psi_*(x)}}
               {\sqrt{\det H(x)}}.
 \label{eq:tail-full-prefactor}
\end{align}
Thus \(N_g\) is the Gaussian normalisation of
\eqref{eq:gaussian-normalisation}, while
\(N_g^{\mathrm{phys}}=z_{\mathrm n}^{-b_c}\zeta_G^{|I|}N_g\)
is the full physical prefactor. The additional factors are the
blocking-noise normaliser and the active Haar chart constant.
Both prefactors are strictly positive.

The factor \(g^{-d}\) in the ratio \(T_\beta/N_g^{\mathrm{phys}}\)
has logarithmic cost of order \(d\log\beta\). We now
specify the auxiliary size regime in which the central
radius pays it. For \(t>e^e\), set
\begin{equation}
                         P(t)=\exp\{(\log\log t)^2\}.
 \label{eq:tail-subpower-scale}
\end{equation}
This is the scale used in \cite[Section~15.4, pp.~76--77]{Main}.
For each fixed \(\nu>0\) and \(m>0\),
\begin{equation}
       P(t)t^{-\nu}\longrightarrow0,\qquad
       \frac{P(t)}{(1+\log t)^m}\longrightarrow\infty.
 \label{eq:tail-scale-comparisons}
\end{equation}
The first follows by comparing \((\log\log t)^2\) with
\(\nu\log t\); the logarithm of the second ratio is
\((\log\log t)^2-m\log(1+\log t)\), which tends to
infinity. These two elementary comparisons will check all
the parameter requests together.

Here \(\widehat\beta>0\) is a prescribed reference coupling for the
auxiliary scale choices, whereas \(\beta\) is the coefficient of
the actual action and \(g=\beta^{-1/2}\). We fix the auxiliary
quantities as functions of \(\widehat\beta\), then ask for estimates
uniform over the displayed band of actual couplings. This lets
one family of cutoff and size choices serve every input in that
band. The functions prescribed below are not retuned as \(\beta\)
varies; after fixing their profile data, we only increase the
common entry threshold.

\begin{theorem}[Relative reference tail on auxiliary tori]
\label{thm:tail-auxiliary}
Fix the principal and branch profile, a number \(d_0>0\),
and constants \(C_1,m_1>0\). As \(\widehat\beta\to\infty\),
suppose
\begin{equation}
 \begin{gathered}
       \widehat\beta/2\le\beta\le2\widehat\beta,\qquad
       h=\widehat\beta^{-1/2}P(\widehat\beta),\qquad
       D=d_0P(\widehat\beta),\\
       \|x\|_\infty\le\widehat\beta^{-7/16},\qquad
       n_f\le C_1(1+\log\widehat\beta)^{m_1}.
 \end{gathered}
 \label{eq:tail-auxiliary-regime}
\end{equation}
Assume the loop allowances satisfy, for every \(\nu>0\),
\(\Lambda\le C_\nu\widehat\beta^\nu\) eventually,
uniformly over this family of auxiliary tori. Then there
exists \(c>0\), fixed by the local profile, such that
\begin{equation}
                0\le\frac{T_\beta(x)}{N_g^{\mathrm{phys}}(x)}
                                      \le e^{-cD^2}
 \label{eq:tail-auxiliary-result}
\end{equation}
for all sufficiently large \(\widehat\beta\), uniformly
in that family. The ratio is smaller than every inverse
power of \(\widehat\beta\).
The lower threshold may also depend on
\(d_0,C_1,m_1\) and the stated loop bounds.
\end{theorem}
\begin{proof}
First collect the applicability conditions:
\[
 h\ll\min(h_0,\delta_*),\qquad
 gD\ll r_y,\qquad \|x\|_\infty\ll r_x.
\]
The profile and \(d_0,C_1,m_1\) are fixed first. The functions
\(h,D\) are prescribed by \eqref{eq:tail-auxiliary-regime},
and the actual coupling varies in its stated comparable band.
In particular \(gD\le\sqrt2\,d_0h\).
Equation~\eqref{eq:tail-scale-comparisons} lets us increase the
entry coupling until
\(h\le\min(h_0,\delta_*)\), \(gD\le r_y/2\) and
\(\|x\|_\infty<r_x\), uniformly over the auxiliary tori.
These fixed inequalities are the chart requirements.

If all mismatches are small,
\eqref{eq:tail-small-mismatch-price} supplies the relative
loss \(\eta D^2/2\). If some mismatch exceeds
\(\delta_*\), the centre-energy estimate gives
\[
       \varphi(x,m(x))
           \le C n_f\widehat\beta^{-7/8}
           \le a\delta_*^2/4
\]
eventually. Consequently
\[
       \beta\{A_t(V(x),U)-\varphi(x,m(x))\}
                               \ge\beta a\delta_*^2/4.
\]
For this to be at least \(\eta D^2/2\), the additional
condition is precisely
\[
                    \eta(gD)^2\le a\delta_*^2/2.
\]
It follows by requiring \(gD\ll\delta_*\sqrt{a/\eta}\).
The fixed profile and the prescribed scales meet this
condition after increasing entry; no limit of \(gD\)
is needed in this comparison.
Both cases therefore have a common relative loss
\(c_0D^2\), with \(c_0=\eta/2\).

On the retained good set,
\(-I_f^{\rm un}(U)\le n_f\Lambda h^2\).
Factor this bound and the common action loss out of
\eqref{eq:tail-full-physical-integral}. The remaining
integral is at most
\[
       \int_{G^I}\sum_t\lambda_t(U)\,\dd h_G^I(U)=1.
\]
Here both the exact finite branch law and the probability
normalisation of Haar measure are used. We obtain
\begin{equation}
       T_\beta(x)\le z_{\mathrm n}(\beta)^{-b_c}
          e^{-\beta\varphi(x,m(x))+n_f\Lambda h^2-c_0D^2}.
 \label{eq:tail-before-normalisation}
\end{equation}
The possibly large noise normaliser cancels exactly on
division by \(N_g^{\mathrm{phys}}\); it is not estimated separately.

We record every other cost in that division. At the
stationary configuration, \eqref{eq:model-centre-size}
and the definition of the replica links give
\[
       \max_e d_G(U(x,m(x))_e)
                            \lesssim\|x\|_\infty.
\]
Indeed \(d_G(e^X)\le|X|_{\g}\), and a crossing link is
\(e^{m_e}e^{x_b}\). Telescoping four factors therefore
bounds each stationary plaquette chord by
\(C\|x\|_\infty\). The bilinear loop budget implies
\[
       |I_f^{\rm un}(U(x,m(x)))|
                       \lesssim n_f\Lambda\|x\|_\infty^2.
\]
The function \(\log j_G\) is smooth on a fixed smaller
chart and vanishes at zero. Bounded first derivatives
give \(|\log j_G(m_e)|\lesssim|m_e|_{\g}\), hence
\begin{equation}
       |\psi_*(x)|\lesssim n_f
                    \{\Lambda\|x\|_\infty^2+\|x\|_\infty\}.
 \label{eq:tail-centre-loop-haar-price}
\end{equation}
Also \(\sqrt{\det H}\le K_{\rm H}^{d/2}\) by
\eqref{eq:model-determinant-upper}, and
\(|I|\le4n_f\), \(d\le4n_f\dim_{\R}\g\).

For \(T_\beta>0\), division of
\eqref{eq:tail-before-normalisation} by
\eqref{eq:tail-full-prefactor} gives explicitly
\begin{align*}
 \log\frac{T_\beta}{N_g^{\mathrm{phys}}}
       \le{}&-c_0D^2+n_f\Lambda h^2+\psi_*(x)
                  -|I|\log\zeta_G-\frac d2\log(2\pi)\\
             &\hspace{17mm}-d\log g
                              +\frac12\log\det H(x)\\
       \le{}&-c_0D^2+Cn_f
          \{1+\log\widehat\beta+\Lambda h^2
                +\Lambda\widehat\beta^{-7/8}
                +\widehat\beta^{-7/16}\}.
\end{align*}
If \(T_\beta=0\), the desired conclusion is immediate.
The displayed calculation keeps the Haar chart constant
per free link and the Gaussian factor per real component.

To bound the last line, use, for example,
\(\Lambda\le C_{1/4}\widehat\beta^{1/4}\) and
\(P(\widehat\beta)^2\le\widehat\beta^{1/4}\).
Then \(\Lambda h^2\lesssim\widehat\beta^{-1/2}\) and
\(\Lambda\widehat\beta^{-7/8}\lesssim\widehat\beta^{-5/8}\).
After entry these and the remaining negative-power term
are at most one. All positive costs are therefore at most
\(Cn_f(1+\log\widehat\beta)\).

The conditions used in the proof can now be read together:
\[
 \begin{gathered}
 \left(\frac{n_f(1+\log\widehat\beta)}{\eta}\right)^{1/2}
              \ll D,\qquad
 gD\ll\min\{r_y,\delta_*\sqrt{a/\eta}\},\\
 \|x\|_\infty\ll
           \min\{r_x,\delta_*\sqrt{a/n_f}\},\qquad
 h\ll\min\{h_0,\delta_*\}.
 \end{gathered}
\]
The first condition pays the full normalisation cost.
The other three ensure chart applicability and the two
relative-energy cases. Every constant entering these
conditions has now been fixed. With
\(n_f\le C_1(1+\log\widehat\beta)^{m_1}\),
\(D=d_0P(\widehat\beta)\), \(gD\le\sqrt2\,d_0h\) and
\(\|x\|_\infty\le\widehat\beta^{-7/16}\),
\eqref{eq:tail-scale-comparisons} verifies all four
conditions simultaneously by increasing the same entry
threshold. No further parameter is chosen.
In particular the total positive cost is at most
\(c_0D^2/2\). This proves
\eqref{eq:tail-auxiliary-result} with \(c=c_0/2\).
The stronger consequence
\(D^2/\log\widehat\beta\to\infty\) is used to conclude
decay faster than every inverse power; it is not an
additional tuning condition.
\end{proof}

This is the conclusion of
\cite[Theorem~15.4 and (15.7), pp.~76--77]{Main}.
The auxiliary volume bound has a concrete role: it pays
the full \(d\log\beta\) cost in a comparison of total
integrals. No assertion that this cost is uniformly
bounded on arbitrary physical volumes is used.

\subsection{Keeping the physical indicator in the central formula}

The tail theorem does not require every point of
\(\mathcal C_x\) to be good. Restoring the coarse parameter
as \(P(g,z;x)\) in \eqref{eq:central-exponent}, the exact
central change of variables instead gives
\begin{equation}
 \begin{split}
       (\mathcal P_\beta F_{{\rm sf},\beta})(V(x))
       ={}&N_g^{\mathrm{phys}}(x)\int_{|z|\le D}
             \chi_{\rm sf}(U(x,m(x)+gz))\\
          &\hspace{12mm}\cdot e^{P(g,z;x)}
                           \,\dd\mu_{H(x)^{-1}}(z)
                     +T_\beta(x).
 \end{split}
 \label{eq:tail-exact-masked-central}
\end{equation}
On this ball all chart cutoffs are one and all predictor
branches coincide. Their likelihoods sum to one, which
justifies removing their sum here. The physical indicator
has been retained explicitly. We are considering precisely
\(F_{{\rm sf},\beta}\), without an additional incoming
multiplier; the remainder in
\eqref{eq:reference-exact-splitting} is still a separate
term of the original blocked density.

To replace the integral in
\eqref{eq:tail-exact-masked-central} by the unmasked
\(Z_D(g,x)\), we need a quantitative plaquette bound on
the whole central ball. The estimate
\(\|m(x)\|_\infty\lesssim\|x\|_\infty\) alone
does not provide it on the full auxiliary coarse domain:
\begin{equation}
       \frac{\widehat\beta^{-7/16}}h
                 =\frac{\widehat\beta^{1/16}}{P(\widehat\beta)}
                         \longrightarrow\infty.
 \label{eq:tail-two-coarse-widths}
\end{equation}
This identifies why physical cutoff compatibility must
be checked separately from coordinate compatibility.

A single central configuration with a reserve in its
plaquette bounds is enough to make this check locally.
On the common stationary chart,
Theorem~\ref{thm:uniform-stationary} gives
\[
       \sup_{\|x\|_\infty<r_x}
              \|Dm(x)\|_{\infty\to\infty}\lesssim1.
\]
The implicit bound is fixed by the profile, independently
of the period and terminal cutoff.

\begin{lemma}[Transporting a good central configuration]
\label{lem:tail-central-reserve}
Suppose \(U(x_0,m(x_0)+gz_0)\), with \(|z_0|\le D\),
has every plaquette chord at most \(h/2\).
Every configuration in \(\mathcal C_x\) is good if
\begin{equation}
       \|x-x_0\|_\infty\ll h,\qquad gD\ll h.
 \label{eq:tail-central-reserve}
\end{equation}
Both central sets are required to lie in the common chart.
The smallness thresholds depend only on the fixed profile,
independently of the period and terminal cutoff.
\end{lemma}
\begin{proof}
For \(X,Y\in\g\), which are skew Hermitian,
\[
       e^X-e^Y=\int_0^1e^{(1-t)X}(X-Y)e^{tY}\,\dd t,
       \qquad \|e^X-e^Y\|_{\rm F}\le|X-Y|_{\g}.
\]
The identity follows by differentiating
\(e^{(1-t)X}e^{tY}\), and its norm bound uses unitarity.
For two free links in our replica coordinates, the chord
distance is thus at most the change of the fine coordinate
plus, for a crossing link, the change of its coarse coordinate.
Integration of \(Dm\) along the coarse segment gives
\[
       \|m(x)-m(x_0)\|_\infty
                              \lesssim\|x-x_0\|_\infty.
\]
For \(|z|,|z_0|\le D\),
\(\|z-z_0\|_\infty\le2D\). Hence the maximum fine-link
distance between the two configurations is
\(\lesssim\|x-x_0\|_\infty+gD\).
Telescoping the four unitary factors of a plaquette
bounds the change of its matrix by four times this
distance, and the change of its chord by the same quantity.
Choose the thresholds in \eqref{eq:tail-central-reserve}
so that this change is at most \(h/2\). Adding the original
\(h/2\) bound proves the assertion.
\end{proof}

This is \cite[Lemma~15.5 and (15.8), p.~77]{Main}, using
the volume-uniform Lipschitz bound of the stationary graph.
At the flat coarse field the required starting
configuration is available: \(x_0=0,z_0=0\) gives
\(U=1\), because \(m(0)=0\).

\begin{corollary}[Unmasked central comparison near the identity]
\label{cor:tail-near-identity}
Under the auxiliary bounds of
Theorem~\ref{thm:tail-auxiliary}, there exist
\(c_{\rm in},d_{\rm in}>0\), fixed by the profile,
such that, if \(0<d_0\le d_{\rm in}\) and
\(\|x\|_\infty\le c_{\rm in}h\), then eventually
the entire central ball is good and
\begin{equation}
       0\le (\mathcal P_\beta F_{{\rm sf},\beta})(V(x))
                          -N_g^{\mathrm{phys}}(x)Z_D(g,x)
                           \le N_g^{\mathrm{phys}}(x)e^{-cD^2}.
 \label{eq:tail-near-identity}
\end{equation}
Here \(Z_D(g,x)\) is exactly
\eqref{eq:central-Z} with its coarse parameter restored.
\end{corollary}
\begin{proof}
The bound \(gD/h\le\sqrt2\,d_0\) follows from
\(\beta\ge\widehat\beta/2\). After fixing the profile, choose
\(0<c_{\rm in},d_{\rm in}\ll1\). Then
\[
       \frac{\|x\|_\infty}{h}\le c_{\rm in},\qquad
       \frac{gD}{h}\le\sqrt2\,d_0\le\sqrt2\,d_{\rm in},
\]
so both conditions in \eqref{eq:tail-central-reserve} hold.
With the flat starting configuration,
Lemma~\ref{lem:tail-central-reserve} proves that the
indicator in \eqref{eq:tail-exact-masked-central} is
identically one. The coarse restriction also fits the
auxiliary domain eventually, since
\(c_{\rm in}h/\widehat\beta^{-7/16}
=c_{\rm in}\widehat\beta^{-1/16}P(\widehat\beta)\to0\).
The exact masked formula therefore becomes
\(\mathcal P_\beta F_{{\rm sf},\beta}=N_g^{\mathrm{phys}} Z_D+T_\beta\).
Apply \eqref{eq:tail-auxiliary-result}.
\end{proof}

This is the near-identity conclusion of
\cite[Corollary~15.6 and (15.9), pp.~77--78]{Main}.
For a general good auxiliary background, calibration will supply
the starting reserve in Section~\ref{sec:general-central-comparison}.
The direct comparison there also controls every nearby central
configuration. Local gauge and extension operations remain
necessary when transferring auxiliary chart statements to
unrestricted physical volumes.
We next localise the stationary and quadratic Gaussian
terms of this same prefactor. The distinction between
the common coordinate chart and physical good fields
will remain explicit.


\section{Keeping spatial information in the stationary equation}
\label{sec:rooted-stationary}

The central calculation has produced a positive factor
\(N_g^{\mathrm{phys}}(x)\), but its logarithm is a sum over the whole
torus. A bound on that total scalar cannot supply a local
RG estimate independent of the volume.
The useful question is how much can be assigned to one
anchor, after summing over all the field coordinates it reads.

Write
\[
        \Psi(x)=\varphi(x,m(x)),\qquad
        \mathcal L(x)=\frac12
                      \log\frac{\det H(x)}{\det H_0}.
\]
All matrices here are on the same actual finite torus.
Since \(m(0)=0\), the action and loop vanish at the
identity, and \(j_G(0)=1\), the exact prefactor
\eqref{eq:tail-full-prefactor} gives
\begin{equation}
       -\log\frac{N_g^{\mathrm{phys}}(x)}{N_g^{\mathrm{phys}}(0)}
                =\beta\Psi(x)+\psi_*(x)+\mathcal L(x).
 \label{eq:local-prefactor-target}
\end{equation}
The field-independent factor \(N_g^{\mathrm{phys}}(0)\) keeps the
vacuum normalisation, including \(\det H_0\).
We shall expand precisely the three functions on the
right, without changing the original normalisation.
Their separation may be compared with
\cite[Sect.~1, (1.3)--(1.7), printed pp.~260--261]{Bal87}.
The local estimates below are proved for the present action.

\subsection{A norm that survives the contractions}

Keep the coordinate attachment and periodic distance
\(d(i,j)\) of \eqref{eq:weighted-local-matrix}.
An index denotes a link-coordinate block with values
in \(\g\), rather than an individual colour component.
For a list of such indices, let
\(\mathsf t(i_0,\ldots,i_k)\) be the smallest total
length of a tree containing their attached sites;
edge length is the periodic maximum distance.
Additional lattice sites may be vertices of the tree.
Repeated sites add no length.
This convention matters when two kernels are joined at
an index which is then summed out.

For a multilinear kernel
\(T_{i;j_1,\ldots,j_k}:\g^k\to\g\), define
\begin{equation}
 |T|_{\theta,k}
    =\sup_i\sum_{j_1,\ldots,j_k}
        e^{\theta\mathsf t(i,j_1,\ldots,j_k)}
                       \|T_{i;j_1,\ldots,j_k}\|_{\rm op}.
 \label{eq:local-rooted-tensor}
\end{equation}
The operator norm is the supremum on unit input vectors.
The word rooted means that \(i\) is held fixed while
all input indices are summed. At order zero the norm
is \(\sup_i|T_i|_{\g}\).
The same definition applies to scalar kernels assigned
to an anchor. Rectangular kernels may have fine and
coarse indices of different types.

\begin{lemma}[Joining rooted kernels]
\label{lem:local-rooted-joining}
Substitute rooted kernels into the input slots of another
rooted kernel and sum the joined indices.
The norm \eqref{eq:local-rooted-tensor} of the resulting
kernel is at most the product of the input norms, at
the same weight \(\theta\).
For matrices this includes the weighted row estimate
for a product; the column estimate follows by reversing
the product.
\end{lemma}
\begin{proof}
For one summand choose a connecting tree for every
factor. The trees meet at the joined indices.
Their union is connected and contains every remaining
input and the output root. Deleting cycle edges
therefore gives
\[
 \mathsf t(\text{remaining indices})
       \le \mathsf t(\text{outer indices})
            +\sum_{\text{insertions}}
                          \mathsf t(\text{inserted indices}).
\]
Additional vertices are allowed, so a joined index
may remain a vertex after it ceases to be an external
index. Exponentiating pays the desired weight.
Take operator norms and sum first over the external
inputs of each inserted kernel. Its rooted norm bounds
this sum uniformly in the joined index.
The remaining sum is bounded by the outer rooted norm.
All bounds use nonnegative sums, so their order may
be changed. This also proves the rectangular version.
\end{proof}

A graph arising in this calculation can have cycles.
Only its connecting length is bounded by a tree.
No one-to-one correspondence with labelled trees is
used. This is Lemma~16.1, p.~78, of \cite{Main}.

The eventual local functions require all Taylor
degrees, not just a fixed number of derivatives.
For an analytic rooted function write its ordered-index
Taylor series as
\[
 F_i(z)=\sum_{k\ge0}\sum_{j_1,\ldots,j_k}
       F^{[k]}_{i;j_1,\ldots,j_k}
                         [z_{j_1},\ldots,z_{j_k}],
 \qquad
 F^{[k]}=\frac1{k!}D^kF(0).
\]
Here \(z\) can be \(x\), \(y\), or their joint array.
The bracket means evaluation of a multilinear coefficient.
Define the absolute analytic norm
\begin{equation}
             \mathcal A_{r,\theta}(F)
                  =\sum_{k\ge0}r^k|F^{[k]}|_{\theta,k}.
 \label{eq:local-analytic-norm}
\end{equation}
For an analytic matrix \(K(x)\), its analogous norm is
\begin{equation}
 \mathcal A_{r,\theta}(K)
  =\sum_{k\ge0}r^k\sup_i
     \sum_{j,j_1,\ldots,j_k}
       e^{\theta\mathsf t(i,j,j_1,\ldots,j_k)}
                    \|K^{[k]}_{ij;j_1,\ldots,j_k}\|_{\rm op}.
 \label{eq:local-analytic-matrix}
\end{equation}
Thus the second matrix index is summed along with the
field inputs. For a constant matrix this is its
weighted row norm.
Expanding a product degree by degree and applying
Lemma~\ref{lem:local-rooted-joining} gives
\begin{equation}
 \mathcal A_{r,\theta}(K_1K_2)
        \le\mathcal A_{r,\theta}(K_1)
                             \mathcal A_{r,\theta}(K_2).
 \label{eq:local-analytic-product}
\end{equation}
Symmetrising the field inputs by averaging their
permutations does not increase any of these bounds.

\subsection{The path budget supplies an absolute analytic bound}

We now retain the spatial information in
Lemma~\ref{lem:uniform-action-derivatives}.
A useful stronger bookkeeping convention is to keep
the original path labels in a coefficient, even when
two labels yield the same monomial.
Each label comes with a connecting tree. Its length,
and later the sum of these lengths under substitution,
will pay for an actual connected footprint.

\begin{lemma}[An absolute bound from the actual local reads]
\label{lem:local-path-majorant}
For the fixed principal profile there exist positive
\(r_*,\theta_*\), independent of the period and terminal
UV cutoff, such that the full joint gradient satisfies
\[
                   \mathcal A_{r_*,\theta_*}
                                   (\nabla\varphi)<\infty
\]
with a uniform bound. The same assertion holds for
its Hessian on a smaller fixed radius.
The coefficients can be kept with connecting-tree
labels so that the absolute sum weighted by each
label's own length has the same uniform bounds.
The gradient has zero constant term.
\end{lemma}
\begin{proof}
Consider one based path term with total word length
\(l\) and coefficient magnitude \(b\), including its
multilinear form norm. A link in the replica coordinates
is a product of at most two coordinate exponentials.
For a matrix exponential the degree-\(k\) term has
operator norm at most \(C_G^k/k!\) on unit Lie algebra
inputs. Multiplication of the word series therefore
has absolute coefficient sum at radius \(r\) bounded
by \(C e^{C_G r(1+l)}\).
One gradient derivative introduces at most a fixed
multiple of \(1+l\) additional choices. It does not
introduce a factor depending on the torus volume.

For a prescribed output link, only translated anchors
whose read contains that link contribute.
There are at most a fixed polynomial in \(1+l\)
such translations. All their read sites, including
the output, can be joined by the projected based
paths, of total length at most \(C_L(1+l)\).
The resulting bound per output is consequently of
the form
\begin{equation}
       Cb(1+l)^C
           \exp\{C_G r(1+l)+C_L\theta(1+l)\}.
 \label{eq:local-word-analytic-price}
\end{equation}
This estimate sums all degrees before summing the
path labels. Repetitions of a coordinate in a word
are already counted by the exponential series.

Choose \(r,\theta>0\) so small that
\(C_G r+C_L\theta<\sigma/2\).
The remaining fixed polynomial is paid by the
unused exponential path width.
Equations~\eqref{eq:stationary-row-budget} and
\eqref{eq:stationary-higher-budget} now sum
\eqref{eq:local-word-analytic-price}.
A row square uses two path labels; its absolute
sum is bounded by the product of the two row
budgets. The direct Wilson terms have length four.
The finite signed symmetry average preserves
these bounds.

The regular predictor in a noise term reads a
fixed finite patch. Lemma~\ref{lem:actual-analytic-branch}
provides a complex analytic extension on a
neighbourhood fixed by that patch and the branch data.
On a smaller complex polydisc its Taylor coefficients
have a finite absolute sum.
More explicitly, for a function bounded by \(K\)
on a polydisc of radius \(R\) in \(v\) scalar
coordinates, the Cauchy coefficient estimate gives
\[
       |c_\alpha|\le KR^{-|\alpha|},\qquad
       \sum_{\alpha\in\mathbb N_0^v}
          |c_\alpha|r^{|\alpha|}
                       \le K(1-r/R)^{-v}\quad(0\le r<R).
\]
Here \(\alpha\) is a multi-index, \(c_\alpha\)
its Taylor coefficient, and \(v\) is the fixed
number of scalar patch coordinates.
Thus this estimate is applied on one patch,
whose dimension is fixed by the profile.
Each patch has a connecting tree of bounded length,
and any output belongs to only a fixed number of
translated patches. The same rooted estimate follows
for all noise terms.

This proves the gradient bound while retaining the
specified trees. To differentiate its series once
more, use a smaller radius \(r'<r\). A coefficient
of degree \(k\) costs the factor
\(k(r'/r)^{k-1}/r\), whose supremum is finite.
The connecting tree already contains the newly
distinguished index. This proves the Hessian
assertion without a loss of spatial weight.
Finally every curvature and noise mismatch vanishes
at the flat configuration. Each term of the action
thus has zero first derivative there.
\end{proof}

The constants depend on
\(G,L,J,a,\lambda,\sigma,C_{\rm p}\) and the fixed
branch data. They do not depend on \(\Lambda\), since
the separate loop \(I_f^{\rm un}\) does not enter \(\varphi\).
This develops the absolute-majorant assertion used
in \cite[Lemma~14.13, p.~67, and Section~16.3,
pp.~80--81]{Main}.

\subsection{One analytic recursion gives all local responses}

The fixed-point map \eqref{eq:model-contraction}
has derivative zero in \(y\) at \((0,0)\).
Its Taylor expansion can therefore be written
\begin{equation}
       y=Dm(0)x+\mathcal R(x,y),\qquad
       Dm(0)=-H_0^{-1}\varphi_{yx}(0,0),
 \label{eq:local-series-equation}
\end{equation}
where the newly defined remainder \(\mathcal R\)
has no terms of total degree zero or one.
The same inverse used to construct the stationary
graph will now be used in the absolute analytic norm.

\begin{theorem}[A local analytic stationary graph]
\label{thm:local-analytic-graph}
There exist \(r_{\rm a},\theta>0\), fixed
before the period and cutoff vary, such that
\begin{equation}
 \mathcal A_{r,\theta}(m)\lesssim r,\qquad
 \mathcal A_{r,\theta}(H-H_0)\lesssim r
                   \quad(0\le r\le r_{\rm a}).
 \label{eq:local-graph-hessian-majorants}
\end{equation}
For every fixed \(k\ge1\), on
\(\|x\|_\infty\le r_{\rm a}/2\),
\begin{equation}
 \sup_i\sum_{j_1,\ldots,j_k}
 e^{\theta\mathsf t(i,j_1,\ldots,j_k)}
       \|D_{j_1}\cdots D_{j_k}m_i(x)\|_{\rm op}
                                                     \le C_k.
 \label{eq:local-graph-response}
\end{equation}
The constants have the same profile dependence as
Lemma~\ref{lem:local-path-majorant}, and are independent
of the volume and terminal UV cutoff.
Both series in \eqref{eq:local-graph-hessian-majorants}
may retain their connecting-tree labels.
\end{theorem}
\begin{proof}
Reduce \(\theta\) below \(\theta_*\) and the inverse
weight \(\nu\) of \eqref{eq:model-moving-inverse}.
Use the radius \(r_*\) of Lemma~\ref{lem:local-path-majorant},
decreased once if necessary so that both its gradient and Hessian
majorants converge there. These data are fixed before choosing
\(r_{\rm a}\).
The fixed matrix \(H_0^{-1}\) then has a uniform
weighted row bound. By
Lemma~\ref{lem:local-rooted-joining} and
Lemma~\ref{lem:local-path-majorant}, constants
\(c_1,c_2\ge1\) can be chosen so that
\[
 |Dm(0)|_{\theta,1}\le c_1,\qquad
 \sum_{k\ge2}u^k|\mathcal R^{[k]}|_{\theta,k}
               \le\frac{c_2u^2}{1-c_2u}
                     \quad(0\le u<c_2^{-1}).
\]
Indeed a finite absolute majorant at a fixed radius
bounds each coefficient geometrically; enlarging
\(c_2\) gives the displayed series bound.

Consider the scalar iteration
\begin{equation}
 t_0(r)=0,\qquad
 t_{n+1}(r)=c_1r+
       \frac{c_2(r+t_n(r))^2}{1-c_2(r+t_n(r))}.
 \label{eq:local-scalar-majorant}
\end{equation}
Its power-series coefficients are nonnegative and
increase with \(n\).
For this iteration to preserve \(0\le t\le2c_1r\), it suffices that
\[
 c_2(1+2c_1)r_{\rm a}\le\tfrac12,\qquad
 2c_2(1+2c_1)^2r_{\rm a}\le c_1.
\]
The stationary uniqueness chart and the later Hessian substitution
also require
\[
 r_{\rm a}<r_x,\qquad 2c_1r_{\rm a}<r_y/2,\qquad
                         (1+2c_1)r_{\rm a}<r_*.
\]
Collect all these requests and choose once
\[
 0<r_{\rm a}\ll
 \min\left\{r_x,\frac{r_y}{1+2c_1},
          \frac{r_*}{1+2c_1},\frac1{c_2(1+2c_1)}\right\}.
\]
All data on the right are already fixed. The smallness threshold
is chosen to give the preceding numerical inequalities; this is
possible since \(c_1\ge1\) and \((1+2c_1)/c_1\le3\).
For \(0\le r\le r_{\rm a}\), \(0\le t\le2c_1r\), they give
\[
 c_2(r+t)\le\tfrac12,\qquad
 \frac{c_2(r+t)^2}{1-c_2(r+t)}
 \le2c_2(1+2c_1)^2r^2\le c_1r.
\]
Thus \(t_n(r)\le2c_1r\) for every \(n\).
Monotone convergence of the coefficients at
\(r=r_{\rm a}\) gives a convergent limiting series
with value at most \(2c_1r_{\rm a}\).

Iterate \eqref{eq:local-series-equation} as a vector
power series, starting with \(y^{(0)}=0\).
The joining lemma shows inductively that its
absolute rooted coefficient series is dominated
by \(t_n\). In making this comparison it is enough
to bound the radii for both joint inputs \(x,y\)
by \(r+t_n(r)\).
Each fixed degree stabilises: if two iterates
first differ at degree \(k\), the nonlinear
remainder makes their next difference start at
degree at least \(k+1\).
The limiting vector series consequently converges
absolutely and solves \eqref{eq:local-series-equation}.
By the chosen radius, the solution lies in the uniqueness ball of
Theorem~\ref{thm:uniform-stationary}, so it is the
same \(m\), with
\(\mathcal A_{r,\theta}(m)\le2c_1r\).

The Hessian has a joint analytic majorant by
Lemma~\ref{lem:local-path-majorant}.
Substitute \((x,m(x))\). Its majorant converges at
\((1+2c_1)r_{\rm a}<r_*\) by the same radius choice.
The constant coefficient is \(H_0\).
The remaining absolute series has no constant term,
so its value at radius \(r\le r_{\rm a}\) is
bounded by a fixed multiple of \(r\).
This proves the second bound in
\eqref{eq:local-graph-hessian-majorants}.

Termwise differentiation now gives, for fixed \(k\),
\[
 |D^km(x)|_{\theta,k}
  \le2c_1r_{\rm a}^{\,1-k}
        \sup_{l\ge k}(l)_k\,2^{-(l-k)} ,
 \qquad
 (l)_k=l(l-1)\cdots(l-k+1).
\]
In this estimate the unassigned field coordinates
are bounded by \(r_{\rm a}/2\).
Removing them from a connecting tree cannot
increase its minimum length.
The displayed supremum is finite, proving
\eqref{eq:local-graph-response}.

Finally, keep the separate trees in every substitution.
An entry of \(H_0^{-1}\) is assigned a path joining
its two sites, of length \(d(i,j)\).
The proof of the joining lemma bounds the stronger
sum weighted by the sum of all constituent lengths.
The scalar majorant therefore applies to these
labelled coefficients as well.
\end{proof}

This uses the same equation as the earlier
maximum-norm construction and avoids a second
fixed-order induction.
It proves the conclusions of
\cite[Proposition~16.2 and Lemma~16.3,
pp.~79--81]{Main}.
In particular, if two coarse fields agree within
distance \(R\) of the site attached to \(i\), and
their segment stays in the smaller chart, then
\begin{equation}
       |m_i(x^1)-m_i(x^0)|_{\g}
            \le C_1e^{-\theta R}
                         \|x^1-x^0\|_\infty .
 \label{eq:local-distant-response}
\end{equation}
Integrate the first derivative along that segment
and sum only over its distant inputs.
This is a comparison of the same finite-period
solution at two fields; no infinite-volume
stationary solution has been substituted.


\section{Localising the three stationary contributions}
\label{sec:local-stationary-prefactor}

We now apply the same rooted coefficient bounds to
each term of \eqref{eq:local-prefactor-target}.
Their total values may be extensive.
All uniform statements below concern the absolute
coefficient sum assigned to one root.

\subsection{The classical action and its quadratic response}

\begin{proposition}[The local stationary action]
\label{prop:local-classical-action}
On a smaller fixed chart, \(\Psi\) has an absolutely
convergent rooted expansion of degrees at least two.
Its analytic coefficient sum at radius \(r\) is
bounded by \(Cr^2\) per coarse root.
Moreover
\begin{equation}
 D^2\Psi
       =\varphi_{xx}
          -\varphi_{xy}H^{-1}\varphi_{yx}
                     \quad\hbox{at }(x,m(x)).
 \label{eq:local-classical-schur}
\end{equation}
These bounds are uniform in the period and terminal
UV cutoff at the fixed principal profile.
\end{proposition}
\begin{proof}
Stationarity cancels the derivative of the fine
coordinate:
\[
       D\Psi(x)[u]=
             D_x\varphi(x,m(x))[u].
\]
Let \(F_i(x)=\nabla_{x_i}\varphi(x,m(x))\),
using the fixed Lie algebra inner product.
Lemma~\ref{lem:local-path-majorant},
Theorem~\ref{thm:local-analytic-graph}, and the
joining lemma give a rooted analytic majorant
for \(F\). Its constant term is zero.
After decreasing the radius, this yields
\(\mathcal A_{r,\theta}(F)\le Cr\).

Since \(\Psi(0)=0\), radial integration gives the
exact identity
\begin{equation}
 \Psi(x)=\sum_i\sum_{k\ge1}
       \sum_{j_1,\ldots,j_k}\frac1{k+1}
       \big\langle x_i,
          F^{[k]}_{i;j_1,\ldots,j_k}
                        [x_{j_1},\ldots,x_{j_k}]
       \big\rangle_{\g}.
 \label{eq:local-classical-radial}
\end{equation}
Indeed \(F_i(tx)\) contributes \(t^k\), whose
integral on \([0,1]\) is \(1/(k+1)\).
At fixed period the root sum is finite, and the
absolute analytic bound justifies integrating
the coefficient series.
Multiplication by \(x_i\) adds at most \(r\)
to its coefficient bound and no connecting length,
because \(i\) was already the root.
The resulting bound is \(Cr^2\).

Differentiating \(D\Psi=\varphi_x\) once more,
and using
\[
             Dm(x)
                =-H(x)^{-1}
                               \varphi_{yx}(x,m(x)),
\]
proves \eqref{eq:local-classical-schur}.
Every contraction in that formula has the rooted
spatial bounds already established.
\end{proof}

This is \cite[Proposition~16.4 and (16.9)--(16.10),
p.~82]{Main}. In particular the quadratic coarse
response comes from eliminating the fine variables
in the actual action. It is not obtained by
replacing the fine Hessian with a scalar gap.

\subsection{The stationary loop and Haar terms}

Recall
\[
 \psi_*(x)=I_f^{\rm un}(U(x,m(x)))
                  -\sum_{e\in I}\log j_G(m_e(x)).
\]
For the loop, two curvature slots vanish at the
identity. Its based-path analytic coefficient
bound thus starts at degree two and carries
the explicit allowance \(\Lambda\).
The proof of Lemma~\ref{lem:local-path-majorant},
now applied to these bilinear terms, followed by
substitution of the graph, gives a rooted scalar
bound \(C\Lambda r^2\) per coarse site.
There are at most a fixed \(L^4\)-multiple of fine
anchors in a coarse block.

The Haar term also starts at degree two.
The formula proved in Appendix~\ref{app:haar} is
\[
 j_G(X)=\det_{\g}\left(\int_0^1
                         e^{-t\,\ad X}\,\dd t\right).
\]
Near zero the determinant is positive, so
\(\log j_G\) is real analytic.
Its linear term is
\[
       D\log j_G(0)[X]=-\tfrac12
                                  \operatorname{Tr}_{\g}(\ad X)=0 .
\]
Here \(\operatorname{Tr}_{\g}\) is the real trace
on the finite-dimensional Lie algebra.
The last equality follows because \(\ad X\)
is skew adjoint for the invariant inner product.
Consequently \(\log j_G(X)=O(|X|_{\g}^2)\)
with an absolute analytic majorant on a fixed ball.
Substitution of \(m_e\), whose constant
coefficient is zero, preserves this starting degree.
The bounded number of fine links in one coarse block
therefore gives
\begin{equation}
   \text{rooted coefficient sum for }\psi_*
                         \ \lesssim(1+\Lambda)r^2
                         \quad\hbox{per coarse site}.
 \label{eq:local-loop-haar-majorant}
\end{equation}
This notation denotes the sum of the labelled scalar
coefficient norms, not a bound on the whole scalar
\(\psi_*\) as the volume grows.
It is the estimate \cite[(16.12), p.~83]{Main}.

\subsection{The determinant as closed matrix chains}

Define
\[
                 A(x)=H_0^{-1}(H(x)-H_0).
\]
By \eqref{eq:local-graph-hessian-majorants} and
\eqref{eq:local-analytic-product}, a smaller fixed
analytic radius can be chosen so that
\begin{equation}
       \mathcal A_{r,\theta}(A)\lesssim r,\qquad
       \mathcal A_{r_{\rm a},\theta}(A)\le\tfrac12 .
 \label{eq:local-det-smallness}
\end{equation}
Whenever a radius is decreased below, all earlier
bounds are retained by restricting their domains.

\begin{proposition}[A local determinant expansion]
\label{prop:local-determinant}
On this chart,
\begin{equation}
 \mathcal L(x)=\sum_{i\in I}p_i(x),\qquad
 p_i(x)=\frac12\sum_{n\ge1}\frac{(-1)^{n+1}}n
                         \operatorname{Tr}_{\g}[(A(x)^n)_{ii}].
 \label{eq:local-det-chain}
\end{equation}
Each \(p_i\) has an absolutely convergent rooted
coefficient expansion of degrees at least one.
Its analytic coefficient sum is bounded by \(Cr\)
per fine link root, with the same volume and cutoff
independence as before.
Internal matrix indices may be retained as coefficient
labels, together with their connecting lengths.
\end{proposition}
\begin{proof}
For fixed real \(x\) in the chart put
\(E=H(x)-H_0\).
The matrix \(H_0+tE\) is positive definite for
\(0\le t\le1\), being a convex combination of
the two positive Hessians.
Jacobi's formula gives
\[
 \frac{\dd}{\dd t}\log\det(H_0+tE)
       =\operatorname{Tr}\bigl((1+tA)^{-1}A\bigr).
\]
The trace here is the full real trace on
\(\g^I\). The absolute row norm of \(A(x)\)
is at most \(1/2\), so
\[
          (1+tA)^{-1}A
                =\sum_{n\ge1}(-t)^{n-1}A^n
\]
converges uniformly on \([0,1]\).
Integration and the finite trace prove
\eqref{eq:local-det-chain}.
This proof does not require \(A\) to be symmetric.

Let \(a(r)=\mathcal A_{r,\theta}(A)\).
The product estimate bounds the rooted matrix
coefficient sum of \(A^n\) by \(a(r)^n\).
Restricting its second matrix index to its root
decreases that absolute sum, and taking the
block trace costs at most \(\dim_{\R}\g\).
Hence the rooted coefficient sum of \(p_i\) is
bounded by
\[
       \frac{\dim_{\R}\g}{2}
                 \sum_{n\ge1}\frac{a(r)^n}{n}
                     \le(\dim_{\R}\g)a(r)\lesssim r .
\]

More explicitly,
\[
 (A^n)_{ii}=\sum_{i_1,\ldots,i_{n-1}}
          A_{ii_1}A_{i_1i_2}\cdots A_{i_{n-1}i}.
\]
Each factor has its rooted analytic expansion.
Keep the intermediate matrix indices, its path
labels, and all allocated field inputs.
For each such label the joining proof pays the
sum of the individual connecting lengths before
any matrix indices are combined.
The same bound \(a(r)^n\) therefore controls this
more detailed series.
It describes a closed chain with attached trees,
whose union is connected.
Finally \(A(0)=0\), so no constant coefficient occurs.
\end{proof}

This proves \cite[Proposition~16.5 and (16.11),
pp.~82--83]{Main} for Lie algebra blocks.
The conversion from fine roots to coarse-site roots
costs only the fixed number of fine links in a block.
The scalar \(\frac12\log\det H_0\) stays in
\(N_g^{\mathrm{phys}}(0)\); it has not been assigned a
volume-independent scalar bound.

\subsection{From connecting lengths to connected footprints}

A coefficient must also pay for the size of the
local set on which it will be used.
Fix an integer \(B\ge2\). For a nonempty set \(X\)
of coarse sites, define \(n_B(X)\) to be the smallest
number of periodic maximum-metric balls of radius
\(B\) whose union is connected and contains \(X\).
The centres need not be in \(X\).
Sites are adjacent when their periodic maximum
distance is at most one. Put \(n_B(\varnothing)=0\).
This is the covering number used for the activity
weights in \cite[Section~8.1, pp.~26--27]{Main};
\(B\) is a covering radius, distinct from the block
ratio \(L\).

For each scalar coefficient label above, retain
all its connecting trees and internal indices.
Let \(l\) be the sum of their metric edge lengths.
Replace every edge by a shortest coordinate path
on the torus; such a path has at most four times
the maximum-metric length.
Their union is connected. Define \(X\) as its
distance-one neighbourhood.
It contains the root and both endpoints of every
coarse link read by the monomial.

\begin{lemma}[A covering price for a path footprint]
\label{lem:local-footprint-cover}
Every footprint constructed in this way satisfies
\begin{equation}
                    n_B(X)\le2+\frac{16l}{B}.
 \label{eq:local-footprint-cover}
\end{equation}
Thus, for \(\kappa>0\) and
\(B\ge\max(2,64\kappa/\theta)\),
\[
                 e^{\kappa n_B(X)}
                         \le e^{2\kappa}e^{\theta l/4}.
\]
\end{lemma}
\begin{proof}
The coordinate paths have at most \(4l\) edges
in total. Starting at the root, traverse a spanning
tree of their connected union out and back.
This gives a walk of at most \(8l\) edges visiting
every site of that union.
Place centres along the walk every \(B-1\) steps,
including its initial site and, if necessary,
one final centre.
At most \(1+\lceil8l/(B-1)\rceil\) centres suffice.
Every site on the walk is within \(B-1\) of one
of them. Its distance-one neighbourhood is thus
covered by the radius-\(B\) balls.
Consecutive balls intersect, so their union is
connected. Since \(B-1\ge B/2\), their number
is at most \(2+16l/B\).
If \(l=0\), one ball covers the distance-one
neighbourhood of the single site.
\end{proof}

The argument is \cite[Lemma~16.6 and (16.13),
p.~83]{Main}. It works with the actual periodic
paths, including those that wrap.

Fix integers \(q\ge6\), \(s\ge2\).
For a scalar \(f\) assigned to root \(o\) and
reading coarse links in \(X\), define
\begin{equation}
 \|f\|^{\rm chart}_{q,r,s;o,X}
  =\max_{0\le j\le q}\ 
     \sup_{\|x\|_\infty\le r}\ 
     \sup_{\substack{|u_{\nu,e}|_{\g}\le
                          r(1+d(e,o))^s\\1\le\nu\le j}}
                  |D^jf(x)[u_1,\ldots,u_j]|.
 \label{eq:local-chart-activity-norm}
\end{equation}
Here \(d(e,o)\) is the periodic maximum distance
from the initial site of the coarse link \(e\)
to \(o\). At order zero the expression is
\(|f(x)|\). These are ordinary coordinate
derivatives. The domain requires each coarse
link logarithm itself to be small.

\begin{theorem}[An exact local family for the stationary prefactor]
\label{thm:local-stationary-family}
For fixed \(q,s,\kappa\) and
\(B\ge\max(2,64\kappa/\theta)\), there is a
smaller fixed chart on which
\begin{equation}
       \beta\Psi(x)+\psi_*(x)+\mathcal L(x)
                            =\sum_\alpha f_\alpha(x).
 \label{eq:local-stationary-family}
\end{equation}
Each summand is assigned a root \(o_\alpha\),
reads a connected footprint \(X_\alpha\), and
vanishes at the flat field. Moreover
\begin{equation}
 \sup_o\sum_{\alpha:o_\alpha=o}
      e^{\kappa n_B(X_\alpha)}
      \|f_\alpha\|^{\rm chart}_{q,r,s;o_\alpha,X_\alpha}
       \lesssim_{q,s,\kappa}
                 \{\beta r^2+(1+\Lambda)r^2+r\}.
 \label{eq:local-stationary-family-bound}
\end{equation}
The constant is independent of the period, terminal
UV cutoff, and total number of roots.
Its dependence on the fixed principal profile remains
as specified above.
\end{theorem}
\begin{proof}
Take the labelled series of
\eqref{eq:local-classical-radial},
\eqref{eq:local-loop-haar-majorant}, and
\eqref{eq:local-det-chain}, and assign the
footprints just constructed.
Different labels may yield the same monomial;
keep their separate summands.

For a scalar monomial of degree \(k\), coefficient
norm \(b\), and label length \(l\), every read
coordinate is at distance at most \(4l\) from its
root. An order-\(j\) derivative chooses at most
\((k)_j\) ordered factors.
The remaining field factors give \(r^{k-j}\);
the allowed directions give at most
\(r^j(1+4l)^{sj}\).
Thus its chart norm is bounded by
\[
                    C_qb(1+k)^q r^k(2+4l)^{sq}.
\]
The covering price costs at most
\(e^{2\kappa}e^{\theta l/4}\).
For fixed \(q,s,\theta\),
\[
        (2+4l)^{sq}\le C_{q,s,\theta}e^{\theta l/4},
        \qquad (1+k)^q\le C_q2^k .
\]
Consequently the sum is paid by the labelled
analytic majorants at radius \(2r\), retaining
at least half of their exponential spatial width.
Choose the chart so that those majorants converge.
Their classical, loop--Haar, and determinant
parts cost respectively
\[
                  C\beta r^2,\qquad
                  C(1+\Lambda)r^2,\qquad Cr .
\]
At most \(m_0=4(L^4+1)\) possible link roots
are attached to one coarse site.
This gives \eqref{eq:local-stationary-family-bound}
and absolute convergence of all differentiated
series through order \(q\).
At each fixed period the number of roots is
finite, so the series identity is preserved as well.
\end{proof}

This is \cite[Proposition~16.7 and (16.14)--(16.15),
pp.~84--85]{Main}. The norm and the parameter
dependence are the same; grouping repeated labels
is optional and can only decrease the bound.

\subsection{The smallness obtained after Taylor subtraction}

For an analytic function let \(T_{\le d}f\) be its
Taylor polynomial at zero through total degree
\(d\). For an extensive scalar we use the rooted
series just constructed, truncating degree by degree
on its finite torus. Define
\begin{equation}
 \mathcal R_{\rm st}(x)
       =\beta(1-T_{\le4})\Psi(x)
                  +(1-T_{\le2})(\psi_*(x)+\mathcal L(x)).
 \label{eq:local-stationary-remainder}
\end{equation}
The symbol \(1-T_{\le d}\) denotes the Taylor
remainder operation.
The classical factor carries \(\beta\), which is
why retaining only its quadratic polynomial would
not give the remainder size needed here.

\begin{corollary}[The local stationary Taylor remainder]
\label{cor:local-stationary-remainder}
The remainder \(\mathcal R_{\rm st}\) has an exact
expansion \(\sum_\alpha r_\alpha\) with the same
type of roots and connected footprints, satisfying
\begin{equation}
 \sup_o\sum_{\alpha:o_\alpha=o}
 e^{\kappa n_B(X_\alpha)}
 \|r_\alpha\|^{\rm chart}_{q,r,s;o_\alpha,X_\alpha}
       \lesssim_{q,s,\kappa}
                    \{\beta r^5+(1+\Lambda)r^3\}.
 \label{eq:local-stationary-remainder-bound}
\end{equation}
If \(\beta\le2\widehat\beta\),
\(\Lambda\le C_\nu\widehat\beta^\nu\) eventually
for every \(\nu>0\), and
\(r=w=\widehat\beta^{-7/16}\), the right side
is at most
\begin{equation}
                     C_\nu\widehat\beta^{-19/16+\nu}.
 \label{eq:local-stationary-remainder-exponent}
\end{equation}
There is no factor depending on physical volume.
\end{corollary}
\begin{proof}
Remove the classical coefficient labels of degrees
at most four and the other labels of degrees at most
two. This preserves their exact series identities,
roots, and footprint assignments.
The remaining degrees are at least five and three,
respectively.

Choose a fixed radius \(r_*>0\) at which the
labelled analytic majorants are finite.
For \(2r\le r_*\) and \(k\ge d\),
\[
                     (2r)^k
                         \le(2r/r_*)^d r_*^k.
\]
In the preceding proof, the degree polynomial
costs radius \(2r\), and the spatial polynomial
costs part of the exponential width.
Use the displayed inequality with \(d=5\)
for the classical series and \(d=3\) for
the remaining series.
Their fixed-radius bounds retain precisely the
explicit factors \(\beta\) and \(1+\Lambda\).
This proves \eqref{eq:local-stationary-remainder-bound}.

At the stated width,
\[
       \beta w^5\le2\widehat\beta^{-19/16},\qquad
       (1+\Lambda)w^3
                   \le C_\nu\widehat\beta^{-21/16+\nu}.
\]
The second term has stronger decay, proving
\eqref{eq:local-stationary-remainder-exponent}.
\end{proof}

This is \cite[Corollary~16.8 and (16.16)--(16.17),
p.~85]{Main}. It is a bound after Taylor
subtraction in the chart norm.
The retained polynomials have not yet been
identified with the prescribed coarse principal
action or coupling counterterm.
Those matching identities require the later
gauge invariant extraction.


\section{A local gauge bridge to physical good fields}
\label{sec:local-physical-bridge}

Small plaquettes do not put every link of a periodic
configuration near the identity.
Thus the chart norm just proved must be related to
physical fields by a local gauge construction.
For a bounded set of links this can be done directly,
retaining the size of its read.
The path mechanism is the one used in
\cite[Sect.~B, Lemma~1 and (1.24)--(1.26),
printed pp.~79--80]{Bal85g}.
That reference uses its own background and averaging
conditions. We prove the estimates for the paths and
norms needed here.

\subsection{The physical derivative norm}

For a coarse field \(V\) and a Lie algebra array
\(A=(A_e)_e\), define its left derivative by
\[
       D_A f(V)=\left.\frac{\dd}{\dd t}
                        f((e^{tA_e}V_e)_e)\right|_{t=0}.
\]
Products of these derivatives are ordered
compositions; their order is retained.
For a site footprint \(X\), take the good set on
which every plaquette having a vertex in \(X\)
has chord at most \(h\).
The footprints below include the box of all read
paths and a fixed plaquette halo.
Define
\begin{equation}
 \|f\|_{q,h,s;o,X}
   =\max_{0\le j\le q}
       \sup_{\substack{V\ {\rm good\ on}\ X\\
                 |A_{\nu,e}|_{\g}\le h(1+d(e,o))^s}}
               |D_{A_j}\cdots D_{A_1}f(V)|.
 \label{eq:local-physical-activity-norm}
\end{equation}
At order zero the expression is \(|f(V)|\).
The unmasked function must be \(C^q\) on a
neighbourhood of each good configuration, including
the boundary of its good set.
No derivative of the good-field indicator is part
of this norm.
These are the ordered derivatives and good-set
norm of \cite[(8.4)--(8.6), pp.~27--28]{Main}.

Let \(\mathcal Q\) be a coarse box of radius
\(R\ge1\), centred at \(o\), which together with
its halo has a lift not crossing the periodic seam.
For each site \(z\) choose the coordinate-ordered
path \(T_z\) from \(o\) to \(z\) in this box.
For a link \(e=(z,z+e_i)\) in the box set
\begin{equation}
       W_e(V)=V_{T_z}V_eV_{T_{z+e_i}}^{-1},\qquad
       \kappa_e(V)=\log W_e(V).
 \label{eq:local-physical-gauge-map}
\end{equation}
The logarithm uses the fixed small exponential chart.
The notation \(\kappa_e\) denotes this link
coordinate; the unindexed \(\kappa>0\) above
continues to denote the support exponent.

\begin{lemma}[The size and derivatives of the local gauge map]
\label{lem:local-physical-gauge}
If all plaquettes in the box have chord at most
\(h\), and \(Ch(1+R)^2\) lies within the fixed
logarithm chart, then
\begin{align}
       \max_e|\kappa_e(V)|_{\g}
                         &\le Ch(1+R)^2,
                         \label{eq:local-gauge-value}\\
       \max_e|D_{A_j}\cdots D_{A_1}\kappa_e(V)|_{\g}
                         &\lesssim_{G,q,s}
                                 [h(1+R)^{s+1}]^j
                      \quad(1\le j\le q).
                         \label{eq:local-gauge-ordered-derivatives}
\end{align}
Under a gauge transformation \(u\),
\begin{equation}
                \kappa_e(u\cdot V)
                         =\operatorname{Ad}_{u_o}\kappa_e(V).
 \label{eq:local-gauge-conjugation}
\end{equation}
The constants are independent of the surrounding
period and terminal UV cutoff.
\end{lemma}
\begin{proof}
The word for \(W_e\) is a loop of two
coordinate-ordered paths and one link.
Moving that link through the later coordinate
segments fills the loop by at most \(C(1+R)^2\)
elementary plaquettes, all in the box.
The successive square moves used in
Lemma~\ref{lem:tail-path-comparison} change a
path transport by at most the chord of the
corresponding plaquette.
Telescoping therefore bounds \(d_G(W_e)\) by
\(C(1+R)^2h\).
The fixed-chart logarithm is Lipschitz on a
smaller neighbourhood of the identity, giving
\eqref{eq:local-gauge-value}.

For a word of length \(n\), each ordered
derivative inserts its direction into one
of the \(n\) factors:
\[
                  D_AV_e=A_eV_e,\qquad
                  D_AV_e^{-1}=-V_e^{-1}A_e .
\]
After \(j\) derivatives there are at most \(n^j\)
insertion choices, counting repeated occurrences
of the same link separately.
The undifferentiated factors are unitary.
Keeping the insertion order gives the bound
\begin{equation}
 \|D_{A_j}\cdots D_{A_1}
           (V_{e_1}^{\epsilon_1}\cdots V_{e_n}^{\epsilon_n})\|_{\rm F}
   \le C_G^j n^j
                    \prod_{\nu=1}^j
                          \max_{e\ {\rm read}}|A_{\nu,e}|_{\g}.
 \label{eq:local-ordered-word-price}
\end{equation}
The word for \(W_e\) has length at most
\(C(1+R)\), and its allowed directions have
size at most \(C_s h(1+R)^s\).
This proves the stated bound for derivatives
of \(W_e\).
All derivatives of the logarithm through order
\(q\) are bounded on its smaller fixed chart.
In the iterated chain rule the orders of the
derivatives of \(W_e\) sum to \(j\).
Their product therefore retains exactly the
factor \([h(1+R)^{s+1}]^j\).
The number of terms depends only on \(q\).
This proves
\eqref{eq:local-gauge-ordered-derivatives}.

Finally \(V_{T_z}\) transforms to
\(u_oV_{T_z}u_z^{-1}\).
All endpoint factors cancel in the loop, leaving
\(W_e\mapsto u_oW_eu_o^{-1}\).
The fixed small logarithm commutes with
conjugation, proving
\eqref{eq:local-gauge-conjugation}.
\end{proof}

This proves the scalar chart-map part of
\cite[Lemma~16.9, pp.~86--87]{Main}.
The construction reads only the specified box.

\subsection{Pulling back a bounded stencil}

A stencil here means a specified finite collection
of read links.
To pull back a scalar chart function we need its
invariance under simultaneous conjugation of those
link coordinates.
This property is available for the stationary
families without altering their bounds.
Indeed the action obeys
\[
       \varphi(\operatorname{Ad}_u x,
                            \operatorname{Ad}_u y)=\varphi(x,y).
\]
The maximum link norm and the uniqueness domains
are invariant under \(\operatorname{Ad}_u\).
Uniqueness of the stationary graph gives
\(m(\operatorname{Ad}_u x)=\operatorname{Ad}_um(x)\).
Thus \(\Psi,\psi_*\) and \(\mathcal L\) are invariant.
In their coefficient families, group the colour
contractions at each fixed list of link labels.
The adjoint action does not mix links, so these
scalar polynomials are invariant separately.
For determinant chains the adjacent orthogonal
colour matrices cancel under the block trace.
The footprints depend on link labels alone, and
grouping cannot increase an absolute coefficient
bound. The same grouping applies after Taylor
subtraction.

\begin{proposition}[A physical local function from a bounded stencil]
\label{prop:local-physical-pullback}
Let \(f\) be a scalar chart function invariant
under simultaneous conjugation and reading links
in \(\mathcal Q\).
Suppose it is defined on a neighbourhood of the
closed maximum-link \(2w\)-ball, where \(0<w\le1\)
and the \(4w\)-ball is in the fixed group chart.
If
\begin{equation}
                   h(1+R)^{s+1}\ll w
 \label{eq:local-physical-radius-request}
\end{equation}
with the smallness threshold fixed by \(G,q,s\), then
there is a gauge invariant local function
\(\widetilde f\) reading the box and its halo.
It is \(C^q\) near every configuration good
on that footprint, equals \(f(\kappa(V))\)
there, and satisfies
\begin{equation}
              \|\widetilde f\|_{q,h,s;o,X}
                      \lesssim_{G,q,s}
                           \|f\|^{\rm chart}_{q,w,s;o,X}.
 \label{eq:local-physical-pullback-bound}
\end{equation}
The constant is independent of the surrounding
period and both cutoffs.
\end{proposition}
\begin{proof}
Since \(s\ge2\), the radius request also pays
the value bound \eqref{eq:local-gauge-value}.
Every good configuration maps into the closed
\(w\)-ball.
On the open region where the logarithms exist
and \(\max_e|\kappa_e|_{\g}<2w\), define
\(\widetilde f=f\circ\kappa\), and set it to
zero elsewhere.
This gives a Borel local function.
The defining region is invariant under gauge
transformations by
\eqref{eq:local-gauge-conjugation}; the assumed
conjugation invariance of \(f\) makes the
extension gauge invariant.
Every good configuration has a neighbourhood
inside the defining region, so the extension
is \(C^q\) there, including at the boundary
of the physical good set.

For an order-\(j\) chain-rule term let \(p\)
be the derivative order falling on \(f\), and
let \(j_1,\ldots,j_p\ge1\) be the orders
falling on its chart arguments; their sum is
\(j\).
Lemma~\ref{lem:local-physical-gauge} bounds
these argument vectors in maximum norm by a fixed multiple of
\[
                          [h(1+R)^{s+1}]^{j_\nu}.
\]
Directions with maximum norm at most \(w\)
are among the allowed directions of the chart
norm. Its bound on \(D^pf\) therefore gives a fixed
multiple of
\[
       w^{-p}[h(1+R)^{s+1}]^j
                       \le w^{j-p}\le1.
\]
These fixed multiples depend only on \(G,q,s\).
There are finitely many chain-rule terms
through order \(q\).
The order-zero bound follows by evaluating
on the \(w\)-ball.
This proves \eqref{eq:local-physical-pullback-bound}.
Only the unmasked function has been differentiated.
\end{proof}

This is \cite[Proposition~16.11 and (16.26)--(16.27),
pp.~88--89]{Main}, with the local extension
region and its reserve specified.
In the auxiliary regime already used for the tail,
\[
       h=\widehat\beta^{-1/2}P(\widehat\beta),
       \qquad w=\widehat\beta^{-7/16},
\]
and for every fixed \(C_1,m_1>0\) and power bound
\(R\le C_1(1+\log\widehat\beta)^{m_1}\), one has
\begin{equation}
       \frac{h(1+R)^{s+1}}w
        =\widehat\beta^{-1/16}
             P(\widehat\beta)(1+R)^{s+1}
                         \longrightarrow0 .
 \label{eq:local-physical-request-compatible}
\end{equation}
Equation~\eqref{eq:tail-scale-comparisons}
pays this request uniformly over those boxes.
Thus a bounded nonwrapping stencil can be
transferred to physical good fields with its
local derivative estimate.

\subsection{What this step supplies to the next one}

Equation~\eqref{eq:local-stationary-family}
is an exact expansion of the stationary
prefactor on its actual finite-period chart.
It includes the classical action, the stationary
loop and Haar factors, and the quadratic Gaussian
determinant.
The covering weight and the per-anchor summation
are paid explicitly.
Corollary~\ref{cor:local-stationary-remainder}
records the small remainder after Taylor
subtraction; it does not perform principal-action
matching.

Proposition~\ref{prop:local-physical-pullback}
supplies the physical norm estimate for each
specified bounded stencil.
For an entire family containing arbitrarily
long footprints, its box-size request cannot
simply be imposed uniformly.
Compatible physical extensions must retain
the connected-path covering price rather than
replace every footprint by a large filled box.
Section~\ref{sec:compatible-physical-kernels}
constructs the compatible extensions with
these path prices. Their exact masked
reassembly remains a subsequent step.

The factor \(Z_D(g,x)\) in the central integral
also remains separate.
The determinant series expands the quadratic
Gaussian normalisation; the nonquadratic
factor is expanded in
Section~\ref{sec:connected-Gaussian-logarithm},
with its original central cutoff compared in
Section~\ref{sec:connected-original-central}.
The incoming multipliers in
\eqref{eq:reference-exact-splitting} keep their
original place as well.
We next construct compatible physical records and
expand this nonquadratic factor. The construction
will use covariant kernels with their own declared
read sets. The microscopic calibration needed for
the general-background central comparison remains
a separate task; it does not follow from the
local gauge pullback alone.


\section{Calculus for connected local families}
\label{sec:connected-local-calculus}

The remaining expansions contain products of arbitrarily
many local factors. Applying a two-factor derivative bound
repeatedly would attach an unnecessary constant to every
vertex. We first keep the derivative orders separately.
This also makes the cost of changing the anchor visible.

\subsection{Factorial weights and one anchor change}

For either the chart norm or the physical good-set norm,
write \(M_j(f)\) for its order-\(j\) supremum, at a
specified anchor and domain. Introduce
\begin{equation}
               \|f\|^{\rm alg}_q
                      =\sum_{j=0}^q\frac{M_j(f)}{j!}.
 \label{eq:connected-algebra-norm}
\end{equation}
The suppressed parameters are those of the original norm.
They are held fixed during multiplication.
This auxiliary seminorm satisfies
\begin{equation}
 \|f\|_q\le q!\|f\|^{\rm alg}_q
       \le e q!\|f\|_q,\qquad
 \|fg\|^{\rm alg}_q\le
                  \|f\|^{\rm alg}_q\|g\|^{\rm alg}_q .
 \label{eq:connected-algebra-product}
\end{equation}
Indeed an ordered derivative of order \(j\) of \(fg\)
chooses a subset of its positions for \(f\); the
remaining positions act on \(g\), in their original
order. Hence
\[
        M_j(fg)\le
          \sum_{a+b=j}\frac{j!}{a!b!}M_a(f)M_b(g).
\]
Divide by \(j!\) and sum. The equivalence follows from
\(j!\le q!\) and \(\sum_{j=0}^q1/j!\le e\).
The same calculation gives
\begin{equation}
              \|e^f-1\|^{\rm alg}_q
                         \le e^{\|f\|^{\rm alg}_q}-1 .
 \label{eq:connected-algebra-exponential}
\end{equation}
These arguments apply to ordinary chart derivatives
and to ordered physical left derivatives.
No derivatives are commuted.

Suppose \(f_i\) is local on \(X_i\), with anchor
\(o_i\in X_i\), and the overlap graph of these
footprints is connected. Put
\[
                X=\bigcup_iX_i,\qquad
                K=\sum_i n_B(X_i).
\]
The connected covers of the \(X_i\) form a connected
cover of \(X\) with at most \(K\) balls. Its periodic
diameter is at most \(3BK\).
For any \(o\in X\),
\[
 1+d(e,o)\le
             (1+d(o,o_i))(1+d(e,o_i)).
\]
Thus multiplying a direction allowed at \(o\) by
\(\tau=(1+3BK)^{-s}\) makes it allowed at every \(o_i\).
Applying the preceding factorial convolution at a
field in the common domain yields
\begin{equation}
 \left\|\prod_i f_i\right\|_{q;o,X}
       \le q!(1+3BK)^{sq}
                         \prod_i\|f_i\|^{\rm alg}_{q;o_i,X_i}.
 \label{eq:connected-whole-product}
\end{equation}
For an algebraic output seminorm the \(q!\) is
omitted. In the physical case, goodness on \(X\)
implies goodness on every \(X_i\).
The threshold is unchanged when directions are
rescaled.

For each fixed \(\eta>0\), the inequality
\[
        \sup_{K\ge1}(1+3BK)^{sq}e^{-\eta K}
                            \lesssim_{q,s,\eta}(1+B)^{sq}
\]
and \(n_B(X)\le K\) give
\begin{equation}
 e^{\kappa n_B(X)}
       \left\|\prod_i f_i\right\|_{q;o,X}
 \lesssim_{q,s,\eta}(1+B)^{sq}
       \prod_i e^{(\kappa+\eta)n_B(X_i)}
                         \|f_i\|^{\rm alg}_{q;o_i,X_i}.
 \label{eq:connected-one-reserve}
\end{equation}
The reserve pays the anchor change once for the
whole product, independently of its number of
factors. This is the calculus of
\cite[Lemmas~8.2--8.3, pp.~29--30]{Main},
proved here directly for the two norms in use.

\subsection{A positive tree majorant}

For a nonempty footprint \(X_i\), abbreviate
\(k_i=n_B(X_i)\) and put \(D_B=(2B+1)^4\).
Every covering ball has at most \(D_B\) sites, so
\begin{equation}
                         |X_i|\le D_B k_i .
 \label{eq:connected-cover-cardinality}
\end{equation}
Two labels are called incompatible when their
footprints intersect; a label is incompatible
with itself.

\begin{lemma}[A tree sum with one fixed root]
\label{lem:connected-positive-tree}
Let \(u_i\ge0\) and suppose
\[
       \sup_x\sum_{i:x\in X_i}u_i e^{k_i}
                                  \le\epsilon,\qquad
                                  D_B\epsilon\le1.
\]
The sum of rooted labelled trees with prescribed
root \(i\), omitting its own price and weighting
each ordered list of \(m\) children by \(1/m!\),
is at most \(e^{k_i}\).
Repeated labels are allowed in this tree sum.
Consequently the sum of distinct nonempty
connected label sets whose union contains \(x\),
with weights \(\prod_i u_i\), is at most
\(\epsilon\).
\end{lemma}
\begin{proof}
The root sums of height at most \(n\) obey
\[
 \Phi_i^{(0)}=1,\qquad
 \Phi_i^{(n+1)}
       =\exp\left\{\sum_{j:X_j\cap X_i\ne\varnothing}
                                  u_j\Phi_j^{(n)}\right\}.
\]
The exponential supplies the child factorials.
If \(\Phi_j^{(n)}\le e^{k_j}\), the exponent is
at most \(|X_i|\epsilon\le k_i\).
Induction and monotone convergence prove the
root bound, first for finite alphabets and then
for countable ones.

For a distinct connected set choose a label
whose footprint contains \(x\) as its root.
Summing all such choices and all spanning
trees only increases its price.
The tree sum permits additional repeated
labels and is therefore an upper bound.
Summing the root prices gives
\[
          \sum_{i:x\in X_i}u_i\Phi_i
                   \le\sum_{i:x\in X_i}u_i e^{k_i}
                   \le\epsilon .
\]
For distinct nonroot labels the exponential
child factorials agree with the ordered
nonroot sum divided by its factorial.
Thus no extra multiplicity is required.
\end{proof}

This is the rooted part of
\cite[Lemma~9.1, pp.~31--32]{Main}.
The distinction between distinct sets and
repeated labels is essential.
The underlying pinned-tree organisation is
the one in \cite[Section~2, (2.6)--(2.7),
preprint pp.~3--4]{FP07}.

\subsection{One shifted resolvent estimate}

\begin{lemma}[The same spatial bound for every positive shift]
\label{lem:connected-shifted-resolvent}
If \(H=H^*\ge\eta I\) and \(\|H\|_\theta\le K\)
in the weighted block matrix norm, then for
some \(\sigma>0\) and \(C<\infty\),
\begin{equation}
                  \|(H+tI)^{-1}\|_\sigma
                             \le\frac{C}{\eta+t}
                                      \quad(t\ge0).
 \label{eq:connected-shifted-resolvent}
\end{equation}
The constants depend only on the gap, weighted
budget, and the fixed number of indices at
each site. In particular,
\begin{equation}
 H^{-1/2}=\frac1\pi\int_0^\infty
                 t^{-1/2}(H+tI)^{-1}\,\dd t,\qquad
                  \|H^{-1/2}\|_\sigma\le C\eta^{-1/2}.
 \label{eq:connected-square-root-integral}
\end{equation}
All bounds are uniform in the finite period.
\end{lemma}
\begin{proof}
The empty coordinate set is immediate.
Otherwise the Schur bound gives \(\eta\le K\).
For every \(t\ge0\),
\[
 H+tI\ge(\eta+t)I,\qquad
 \|H+tI\|_\theta\le K+t,\qquad
 \frac{\eta+t}{K+t}\ge\frac{\eta}{K}.
\]
Apply Lemma~\ref{lem:uniform-inverse} with gap
\(\eta+t\) and weighted bound \(K+t\).
The last inequality makes its rate in
\eqref{eq:inverse-rate} at least
\(\mu=\min\{\theta/2,\eta\theta/(8K)\}\).
Consequently \eqref{eq:proved-weighted-inverse}
at the common width \(\sigma=\mu/2\) gives
\[
 \|(H+tI)^{-1}\|_\sigma
 \le C m_0(\eta+t)^{-1}(1+\mu^{-1})^4,
\]
where \(m_0\) bounds the number of coordinate
blocks per site as in Lemma~\ref{lem:uniform-inverse}.
This proves \eqref{eq:connected-shifted-resolvent}.
For a scalar eigenvalue \(\lambda>0\),
substitution \(t=\lambda u^2\) gives
\[
     \frac1\pi\int_0^\infty
                \frac{t^{-1/2}}{\lambda+t}\,\dd t
                                          =\lambda^{-1/2}.
\]
Diagonalisation proves the matrix identity.
The displayed weighted norm bound makes its
integral absolutely convergent at both
endpoints.
\end{proof}

For the moving chart Hessian the inverse
resolvent also has an absolute analytic
majorant with the same \(C/(\eta+t)\)
bound. To see this put \(E(x)=H(x)-H_0\)
and \(R_t^0=(H_0+tI)^{-1}\).
On a smaller fixed radius,
\[
 (H(x)+tI)^{-1}
       =R_t^0\sum_{n\ge0}(-E(x)R_t^0)^n,\qquad
       \mathcal A_{r,\sigma}(E)
                                \frac C\eta\le\tfrac12 .
\]
Use \eqref{eq:local-analytic-product} and
integrate its absolute coefficients in
\eqref{eq:connected-square-root-integral}.
Thus \(W(x)=H(x)^{-1/2}\), as well as
\(C(x)=H(x)^{-1}\), has rooted
analytic labels with connecting paths.
This proves the local whitening estimate of
\cite[Lemma~23.1, pp.~115--116]{Main} by integrating
the positive-shift inverses. The locators were checked on
20 November 2026.


\section{Compatible kernels on physical good backgrounds}
\label{sec:compatible-physical-kernels}

We now replace separate local coordinate choices
by physical records defined from the original links.
A record means one specified path expression,
predictor patch, or finite contraction of such
expressions, together with its endpoints and read set.
It is a function of those links before any
background outside its footprint is chosen.
Local gauges will estimate these functions; they
will not define different versions on different
overlaps.

\subsection{Parameters, records, and the family norm}

Fix the principal profile throughout this section.
For the coarse nominal coupling choose
\begin{equation}
 B=\left\lceil(\log\widehat\beta)^2\right\rceil,\qquad
 h_{\rm c}=\widehat\beta^{-1/2}P(\widehat\beta).
 \label{eq:compatible-physical-parameters}
\end{equation}
The original fine-field indicator and its threshold
remain those of the physical integral.
Here \(h_{\rm c}\) tests the coarse background.
Assume the coarse period \(c\) satisfies
\begin{equation}
                         c\ge C_{\rm geom}B .
 \label{eq:compatible-period-reserve}
\end{equation}
The fixed geometric constant contains the frame
patches and the short filling boxes below.
This is the period reserve in
\cite[(25.1), p.~126]{Main}.
We verify joint parameter compatibility at the
end of this construction.

The notation \(C_B\) denotes a fixed polynomial
allowance in \(1+B\). Its coefficients and degree
may depend on
\(G,L,J,a,\lambda,\sigma,C_{\rm p},q,s,\kappa\)
and the chosen positive length reserves.
Its degree can be increased finitely many times while deriving
an estimate. All such increases are completed before its entry
threshold is selected. The resulting coefficients and degrees
are fixed as the coupling, period and cutoff vary.
It never denotes a volume allowance. Since \(B\) varies
with the coupling, this polynomial growth is not absorbed
into an implicit constant. The same rule applies to the
displayed path and loop budgets.

The parent block attaches a coarse site \(d(a)\)
to each free fine link \(a\).
A matrix record \(K_\alpha\), with fine output
\(a_\alpha\) and input \(b_\alpha\), is covariant
at these two attached sites.
Its footprint \(X_\alpha\) contains both sites
and all its branch reads.
Its native length \(l_\alpha\ge1\) is the sum
of its declared cover-path lengths plus one
for each primitive vertex.
Repeated reads are counted with multiplicity,
before periodic identification.
For \(\theta>0\), use the family norm
\begin{equation}
 \begin{split}
 \|\mathcal K\|_{\kappa,\theta}^{\rm phys}
 =\max\Bigl\{&
 \sup_a\sum_{\alpha:a_\alpha=a}
   e^{\kappa n_B(X_\alpha)+\theta l_\alpha}
       \|K_\alpha\|_{q,h_{\rm c},s;d(a),X_\alpha},\\
 &\sup_b\sum_{\alpha:b_\alpha=b}
   e^{\kappa n_B(X_\alpha)+\theta l_\alpha}
       \|K_\alpha\|_{q,h_{\rm c},s;d(b),X_\alpha}
                                                        \Bigr\}.
 \end{split}
 \label{eq:compatible-kernel-family-norm}
\end{equation}
The local norm uses the operator norm of the
Lie algebra block. Tensors sum all fine input
indices at one output, as in
\eqref{eq:local-rooted-tensor}.
Vectors have only the output index.
For anchored scalar families write
\begin{equation}
 \|f\|_{\mathcal E,\kappa}
       =\sup_o\sum_{\alpha:o_\alpha=o}
           e^{\kappa n_B(X_\alpha)}
              \|f_\alpha\|_{q,h_{\rm c},s;o,X_\alpha}.
 \label{eq:compatible-scalar-family-norm}
\end{equation}
Whenever a native reserve is stated, the same
sum also includes \(e^{\theta l_\alpha}\).
Scalar normalisation and quadratic extraction
are separate from the definition of this norm.

Let \(V\) be a coarse field good at threshold
\(h_{\rm c}\). In the same block-tree coordinates
as before, write \(U(V,y)_e=e^{y_e}R_e(V)\)
on free links, with the trees fixed to one.
The function \(\varphi(V,y)\) is the actual
unmasked principal phase of
Section~\ref{sec:smooth-principal}.
On each fixed predictor patch a local gauge
puts the replica links within \(C_Lh_{\rm c}\)
of the identity. Lemma~\ref{lem:actual-analytic-branch}
therefore supplies its regular branch on a
common small \(y\)-ball.
The path sums converge there: a complex
perturbation \(y\) costs \(e^{C\|y\|_\infty l}\)
in a word of length \(l\), paid by a fixed
part of the path width.

Under a coarse gauge transformation \(u\),
the tangent coordinates transform by
\[
          (\mathcal O_u y)_a
                      =\operatorname{Ad}_{u_{d(a)}}y_a,
          \qquad
          \varphi(u\cdot V,\mathcal O_u y)=\varphi(V,y).
\]
The replica and path transports telescope
their endpoint matrices; uniqueness gives
the same covariance for each regular branch.
Thus these are identities of the original
physical functions.

\subsection{Covariantising a flat kernel on the integer cover}

Put
\[
             A_0=H_0=\varphi_{yy}(1,0),\qquad
             A(V)=\varphi_{yy}(V,0),\qquad
             b(V)=\varphi_y(V,0).
\]
The flat quadratic formula has scalar colour
entries: its cochain maps act identically
on each Lie algebra component.
On the integer coarse cover, write
\(A_{0,ab}(z)\) for its scalar coefficients.
Here \(a,b\) are the finitely many free-link
residue types in a block and \(z\in\mathbb Z^4\)
is the coarse displacement.
Their exponentially weighted sum is finite
by the original path budget.
Absolute periodisation gives the flat matrix
on each admissible finite torus.

We need the same cover coefficients for
\(P_t^0=(A_0+tI)^{-1}\), uniformly in \(t\ge0\).
Here is a direct verification.
The matrix Fourier symbol
\[
             A_0(p)=\sum_{z\in\mathbb Z^4}
                              A_0(z)e^{-ip\cdot z},
                    \qquad p\in[-\pi,\pi]^4,
\]
extends holomorphically to a strip by its
exponential coefficient bound.
The finite-period flat gap implies
\(A_0(p)\ge\eta I\) on every admissible Fourier
grid. These grids become dense as the periods
grow, so the same inequality holds for every
real \(p\).
On a sufficiently small complex strip,
\(\|A_0(p+iu)-A_0(p)\|\le\eta/4\).
The resolvent identity then gives
\[
             \|(A_0(p+iu)+tI)^{-1}\|
                                    \le\frac2{\eta+t}.
\]
Shift each Fourier integration contour within
this strip. At a smaller width \(\sigma>0\),
the inverse coefficients satisfy
\begin{equation}
       |P_t^0|_\sigma:=
       \max\left\{\sup_a\sum_{b,z},
                  \sup_b\sum_{a,z}\right\}
            e^{\sigma|z|_1}|P_{t,ab}^0(z)|
                                 \le\frac C{\eta+t}.
 \label{eq:compatible-flat-cover-resolvent}
\end{equation}
The finite residue multiplicity is part of \(C\).
Absolute periodisation of these coefficients
is the actual flat finite-period inverse,
by multiplication of their Fourier symbols.

For each ordered pair of cover endpoints
choose a coordinate path \(\gamma_{ab,z}\)
of length at most \(C_L(1+|z|_1)\), using the
reverse path for the reversed pair.
Choose the empty path when the endpoints
are identical.
For a scalar kernel \(K\) define
\begin{equation}
               K[V]_{ab}(z)
                    =K_{ab}(z)
                          \operatorname{Ad}_{V_{\gamma_{ab,z}}}.
 \label{eq:compatible-covariant-copy}
\end{equation}
Covariantise on the cover first and periodise
afterwards. Symmetric kernels remain symmetric.
Adjoint transport is orthogonal, so the
undifferentiated row and column bounds have
the original constant, independent of \(B\).
Ordered word insertions give the differentiated
family bound \(C_B|K|_\sigma\), at smaller
length width.

\subsection{Short comparisons and long path tails}

The obstruction to multiplying these copies
is explicit:
\[
 K[V]L[V]-(KL)[V]
  =\sum_{i,z,w}K_{ai}(z)L_{ib}(w)
       \left\{\operatorname{Ad}_{V_{\gamma_{ai,z}\gamma_{ib,w}}}
                 -\operatorname{Ad}_{V_{\gamma_{ab,z+w}}}\right\}.
\]
The two paths have the same cover endpoints.
Their quotient has zero cover displacement
and length at most \(C_Ll\), where
\(l=1+|z|_1+|w|_1\).
Square moves fill it in a box of radius
\(C_Ll\) with at most \(C_Ll^2\) plaquettes.

For \(l\le B\) include this box and its fixed
halo in the footprint.
It is nonwrapping by
\eqref{eq:compatible-period-reserve}, and is
covered by a fixed number of \(B\)-balls.
Its value defect is \(O(l^2h_{\rm c})\).
The following calculation pays all the
physical derivatives as well.

In the local gauge, let \(x(V)\) be its
link logarithms and \(g_z(V)=V_{T_z}\) the
endpoint transports.
Lemma~\ref{lem:local-physical-gauge} and the
word estimate give, for \(1\le j\le q\),
\[
 |x(V)|_\infty\le C(1+l)^2h_{\rm c},\qquad
 |D_1\cdots D_jx(V)|_\infty
       +\max_z\|D_1\cdots D_j\operatorname{Ad}_{g_z(V)}\|
              \le C_{q,s}[h_{\rm c}(1+l)^{s+1}]^j .
\]
If a record minus its transported flat
record is denoted by \(K_\alpha\), its
local coordinate expression vanishes at
\(x=0\).
Its fixed-order coordinate derivative sums
cost at most its coefficient magnitude
times a fixed polynomial in \(1+l\).
This follows by word insertions; predictor
patches use their complete analytic branch
derivatives.
Integrating its first coordinate derivative
along the segment from zero gives a value
factor \(h_{\rm c}(1+l)^2\).
For a positive-order physical derivative,
the chain rule supplies \(h_{\rm c}^j\),
or the value factor remains while derivatives
fall on endpoint transports.
Since \(h_{\rm c}\le1\), every term is bounded
by
\begin{equation}
 \|K_\alpha\|_{\rm physical}
          \le C|c_\alpha|h_{\rm c}(1+l)^{C_{q,s}} .
 \label{eq:compatible-short-record-price}
\end{equation}
Here \(c_\alpha\) is the retained scalar
coefficient. The notation on the left is
the local norm at either endpoint anchor.
There are finitely many derivative allocations
at fixed \(q\); their number is independent
of \(l\) and the period.

For \(l>B\), keep only the actual paths
and the fixed branch halos as footprint.
No filling box is declared.
Bound the two transported terms separately.
The same insertion calculation gives
\begin{equation}
       \|K_\alpha\|_{\rm physical}
                    \le C|c_\alpha|(1+l)^{C_{q,s}},
             \qquad n_B(X_\alpha)\le C_L(1+l/B).
 \label{eq:compatible-long-record-price}
\end{equation}
Choose a final native width \(\theta>0\)
strictly below the available coefficient
width. For sufficiently large \(B\),
the covering factor \(e^{C_L\kappa l/B}\)
uses only part of that width.
Fixed length polynomials, including endpoint
rebasing and incidence counts, use another
fixed part. The long coefficient sum is
then at most \(Ce^{-cB}\).
This proves
\begin{equation}
 \|K[V]L[V]-(KL)[V]\|_{\kappa,\theta}^{\rm phys}
       \le C_B(h_{\rm c}+e^{-cB})|K|_\sigma|L|_\sigma .
 \label{eq:compatible-product-defect}
\end{equation}
All comparisons were made on the integer
cover. A winding image after periodisation
does not change their validity.
Coincident images are summed after taking
absolute cover prices.
This is \cite[Lemma~25.1, p.~128]{Main}.

\subsection{Comparison with the actual gradient and Hessian}

Expand the row squares of
\eqref{eq:unmasked-principal-action} into
their two original path labels.
Differentiate each principal or noise
record at \(y=0\), retaining those labels.
For a Hessian record subtract its flat
scalar coefficient with the prescribed
transport between its two tangent endpoints.
For a gradient record its flat value is zero.
These subtractions sum exactly to
\(A(V)-A_0[V]\) and \(b(V)\).
They retain the complete predictor derivative,
including the derivative of its frame.

Each short subtracted record vanishes in
its flat local chart, so
\eqref{eq:compatible-short-record-price}
applies. The transport in a product of
two row derivatives is compared with
the single prescribed endpoint transport
by \eqref{eq:compatible-product-defect}.
Each long record has
\eqref{eq:compatible-long-record-price}.
The exponential path budgets and their
products sum the remaining coefficients
per row and per column. Consequently
\begin{equation}
       \|b\|_{\kappa,\theta}^{\rm phys}
       +\|A-A_0[V]\|_{\kappa,\theta}^{\rm phys}
                         \le C_B(h_{\rm c}+e^{-cB}).
 \label{eq:compatible-actual-comparison}
\end{equation}
This is the estimate of
\cite[Lemma~25.2, pp.~128--132]{Main}.
The preceding record argument proves it
in its own good-footprint norm.

For later value bounds there is also
\begin{equation}
                         \|b(V)\|_\infty\lesssim h_{\rm c}
 \label{eq:compatible-gradient-value}
\end{equation}
on globally good fields, with an implicit constant
independent of \(B\).
A direct plaquette derivative retains
one small curvature; a row-square derivative
retains one row; a higher curvature term
retains at least two curvature slots.
The noise mismatch at the replica is
\(O(h_{\rm c})\) on its fixed patch, so
its first derivative has the same bound.
Rooted path sums pay their length factors.
Thus no filled long box is needed for
this value estimate.

The same word majorants give a fixed fine
analytic radius \(r_0>0\).
At pinned output, the order-\(k\) fine
derivative prices are
\(C_B k!r_0^{-k}\) for the phase and
\(C_B(1+\Lambda)k!r_0^{-k}\) for the
loop--Haar term, including the prescribed
coarse derivatives through \(q\).
Coarse derivative insertions cost only
fixed degree polynomials in \(k\), paid
by one further reduction of \(r_0\).


\subsection{One parametrix gives positivity and the exact inverse}

Set \(P_t(V)=P_t^0[V]\), and define its defect by
\[
                  E_t(V)=I-P_t(V)(A(V)+tI).
\]
An approximate inverse together with such a
defect is called a parametrix.
Using \(P_t^0(A_0+tI)=I\) on the cover gives
\[
 E_t=P_t(V)\{A_0[V]-A(V)\}
                  +\{I-P_t(V)(A_0[V]+tI)\}.
\]
The second bracket is the product defect
of \eqref{eq:compatible-product-defect};
the identity part of \(tI\) has no defect.
Hence, at suitably reduced fixed reserves,
\begin{equation}
       \|E_t\|_{\kappa,\theta}^{\rm phys}
         \le\frac{C_B(h_{\rm c}+e^{-c_0B})}{\eta+t},
       \qquad
       \|P_t\|_{\kappa,\theta}^{\rm phys}
                            \le\frac{C_B}{\eta+t}.
 \label{eq:compatible-parametrix-price}
\end{equation}
Here and below \(c_0>0\) is a fixed decay
constant; it is unrelated to the period \(c\).

\begin{lemma}[Positivity from the shifted parametrix]
\label{lem:compatible-positive-Hessian}
For sufficiently large nominal coupling,
\(A(V)\) has a positive lower bound independent
of \(B\), cutoff, and period on globally good
fields. Its inverse has the exact local expansion
\begin{equation}
              (A(V)+tI)^{-1}
                         =\sum_{n\ge0}E_t(V)^nP_t(V).
 \label{eq:compatible-Neumann-inverse}
\end{equation}
\end{lemma}
\begin{proof}
Increase the common threshold until the
right side of the first bound in
\eqref{eq:compatible-parametrix-price}
is at most \(1/4\) for all \(t\ge0\).
Its order-zero row and column bounds imply
\(\|E_t\|_{2\to2}\le1/4\).
Therefore \(I-E_t=P_t(A+tI)\) is invertible.
On the finite coordinate space this makes
\(A+tI\) invertible for every \(t\ge0\),
and proves the Neumann identity.

The matrix \(A\) is self-adjoint.
If it had an eigenvalue \(\lambda\le0\),
then \(A+(-\lambda)I\) would be singular,
a contradiction.
Thus \(A>0\).
For the quantitative gap use the
undifferentiated transport bound in
\eqref{eq:compatible-flat-cover-resolvent}:
\[
          \|P_0(V)\|_{2\to2}\le C_0/\eta ,
          \qquad
          \|A(V)^{-1}\|_{2\to2}
                           \le\frac{4C_0}{3\eta}.
\]
The constant \(C_0\) is independent of \(B\),
since adjoint transports are orthogonal.
This proves \(A(V)\ge\eta_*I\) with
\(\eta_*=3\eta/(4C_0)>0\).
Its upper bound is also independent of
\(B\): the covariant flat row budget is
fixed, and the comparison price can be
made at most one.

For local functions keep the intermediate
indices and record labels in each power
of \(E_t\). Their footprints meet at the
joined indices and have connected unions.
Use stronger input covering and native
weights, and apply
\eqref{eq:connected-whole-product} to the
whole chain. Through derivative order \(q\)
the extra cost is polynomial in its total
cover size and length.
A fixed part of the input reserves pays
this cost once, while the geometric ratio
is at most \(1/4\).
The chain and its fixed-order differentiated
series therefore converge absolutely in
the final family norm.
\end{proof}

The inputs are those of
\cite[Section~25.4, pp.~132--134]{Main}.
The shifted-invertibility argument above
deduces positivity from the same parametrix;
it does not assume positivity of a
covariantised flat copy.

\subsection{The stationary graph as a sum of physical trees}

Expand the actual gradient at \(y=0\):
\[
       \varphi_y(V,y)=b(V)+A(V)y+\mathcal R_{\rm nl}(V,y),
       \qquad
       \mathcal R_{\rm nl}(V,0)
          =D_y\mathcal R_{\rm nl}(V,0)=0 .
\]
With \(P_0=P_0(V)\), its stationary equation is
\begin{equation}
       y=u+E_0y+\sum_{k\ge2}F_k[y,\ldots,y],
       \quad
       u=-P_0b,\qquad
       F_k=-\frac1{k!}P_0D_y^k\mathcal R_{\rm nl}(V,0).
 \label{eq:compatible-physical-tree-equation}
\end{equation}
The \(1/k!\) is the original Taylor factor.
The fine input slots of \(F_k\) are ordered.

\begin{proposition}[Compatible stationary and Gaussian families]
\label{prop:compatible-stationary-data}
For sufficiently large nominal coupling,
every globally good \(V\) has a unique small
stationary point \(m(V)\), and
\[
                         \|m(V)\|_\infty\le Ch_{\rm c}.
\]
The actual stationary Hessian
\(H(V)=\varphi_{yy}(V,m(V))\)
has a uniform positive gap.
The graph, its inverse Hessian, and
\(W(V)=H(V)^{-1/2}\) have absolutely
convergent local covariant family expansions.
Their physical prices are \(C_Bh_{\rm c}\),
\(C_B\), and \(C_B\), respectively.
Each retains a positive native length reserve.
On the intersection with the earlier chart,
these are the same stationary point and
the same operators.
\end{proposition}
\begin{proof}
Choose input weights at covering exponent
\(\kappa+8\) and native exponent \(8\theta\),
with \(\theta\) small enough for all primitive
path budgets. The tree, Hessian substitution,
and resolvent steps use only finitely many
reductions, ending at \(\kappa,\theta\).
No reduction is made at every tree height.
Use the algebraic derivative seminorm.

The preceding estimates give a polynomial
\(K_B\ge1\) and a fixed \(r_0>0\) with
\[
 \|u\|_{\rm in}\le K_Bh_{\rm c},\qquad
 \|E_0\|_{\rm in}\le K_B(h_{\rm c}+e^{-c_0B}),\qquad
 \|F_k\|_{\rm in}\le K_Br_0^{\,1-k}\quad(k\ge2).
\]
The input norm here sums the fine slots at
the pinned output and uses the stronger
weights just specified.
The bound on \(u\) uses
\(e^{-c_0B}\le h_{\rm c}\), an eventual
consequence of
\eqref{eq:compatible-physical-parameters}.
For \(F_k\), fine analytic derivatives,
division by \(k!\), and one fixed radius
loss pay all degree factors.

A leaf of an expansion tree carries \(u\);
a vertex with one child carries \(E_0\);
a vertex with \(k\ge2\) ordered children
carries \(F_k\).
At an edge sum the joined coordinate.
Its attached site belongs to both read
sets, so their union is connected.
For a tree \(T\) and \(j\) ordered derivatives,
the full Leibniz expansion is
\[
 D_1\cdots D_j[T]
  =\sum_{\omega:\{1,\ldots,j\}\to V(T)}
       \operatorname{Contr}_T
            \bigotimes_{v\in V(T)}
                          D_{\omega^{-1}(v)}K_v .
\]
The subsequence at each vertex retains
its order. Factorial convolution bounds
these products by the algebraic vertex
prices. Rebase to the output only after
forming the whole product.
If \(K\) is the sum of its vertex covering
numbers, this costs at most
\((1+3BK)^{sq}\), paid by one fixed covering
reserve. The native reserves pay the
analogous fixed length polynomials.
Thus the final norm is bounded by one
polynomial allowance times the positive
tree majorant of the input prices.

Put
\[
 \delta=K_Bh_{\rm c},\qquad
 \epsilon=K_B(h_{\rm c}+e^{-c_0B}),\qquad
 \mathcal M(t)=\delta+\epsilon t
                              +K_B\frac{t^2}{r_0-t}.
\]
Choose the threshold so that
\[
                    \epsilon\le\tfrac18,\qquad
                    \delta\le r_0/(64K_B).
\]
For \(t_*=2\delta\le r_0/4\),
\[
       \mathcal M(t_*)\le2\delta,\qquad
       \sup_{0\le t\le t_*}\mathcal M'(t)
                    \le\epsilon+4K_Bt_*/r_0\le\tfrac14 .
\]
The nonnegative iteration from zero
therefore increases to a finite limit.
The tree iteration of
\eqref{eq:compatible-physical-tree-equation}
is dominated by this scalar iteration.
Summing all finite trees gives an
absolutely convergent physical family,
of price \(C_Bh_{\rm c}\), which solves
the original stationary equation.
Repeated records in different slots are
distinct occurrences; their Taylor
factorials are kept exactly once.

The third fine derivative has a uniform
order-zero rooted budget independent
of \(B\). Consequently
\[
       \|\varphi_{yy}(V,y)-A(V)\|_{2\to2}
                                      \le C\|y\|_\infty .
\]
On a fixed small ball this preserves
half the positive gap of
Lemma~\ref{lem:compatible-positive-Hessian}.
Integrating that Hessian along a segment
proves uniqueness there.
The counting gap and the undifferentiated
weighted row bound for \(A(V)\) also give
a maximum-norm inverse constant independent
of \(B\), by
Lemma~\ref{lem:uniform-inverse}.
Using \eqref{eq:compatible-gradient-value}
in the stationary equation gives
\[
       \|m(V)\|_\infty
           \le Ch_{\rm c}+C\|m(V)\|_\infty^2 .
\]
Decrease the fixed uniqueness ball and
absorb its last term. This proves the
stated sharper value bound.

Taylor expansion of \(\varphi_{yy}\) at
zero and substitution of the tree family
now give \(H-A\) a physical family
price \(C_Bh_{\rm c}\).
Its operator value is \(O(h_{\rm c})\),
so \(H\) retains a fixed gap.
Define
\[
      \widetilde E_t
           =I-P_t(H+tI)
           =E_t-P_t(H-A).
\]
It obeys the same small defect bound
as \eqref{eq:compatible-parametrix-price}.
Thus
\begin{equation}
       (H+tI)^{-1}
             =\sum_{n\ge0}\widetilde E_t^{\,n}P_t,\qquad
       \|(H+tI)^{-1}\|_{\rm family}
                                      \le\frac{C_B}{\eta_*+t}.
 \label{eq:compatible-stationary-resolvent}
\end{equation}
The whole-chain estimate above applies
again. Integrate in
\eqref{eq:connected-square-root-integral}
to obtain \(W\); at \(t=0\) obtain the
inverse Hessian.
For a fixed discrete label, \(t\)
changes scalar coefficients, not its
read footprint. Perform that integrable
coefficient integral before declaring
the resulting local record.

Every finite product is covariant at
its external slots, since orthogonal
transformations cancel at internal slots.
On the chart the stationary equations
are identical in the same uniqueness
ball. Their graphs therefore agree;
the Hessian and its positive square
root then agree as well.
\end{proof}

This proves \cite[Proposition~25.3,
pp.~134--137]{Main} with the same physical
norm and parameter dependence.

\subsection{Raw local definitions and scalar compatibility}

We make the extensions used above precise.
Choose a fixed small branch reserve
\(h_*>0\). A primitive record is given
by its original smooth, gauge-covariant
formula where its fixed patch plaquettes
are less than \(h_*\) and its regular
branch tests have their fixed reserve.
Set it to zero elsewhere.
The defining region is open and gauge
invariant. This gives a Borel covariant
local function.
For \(2h_{\rm c}<h_*\), the record agrees
with its smooth formula near its own
closed good set.
The short filling box is used to
estimate that formula, not to choose
an unrelated gauge-dependent value.
Long records retain their actual paths.
Products use the union of their reads.

All tree and chain records above use
these same primitive definitions.
On a globally good field their sums
are the exact stationary and inverse
equations. On a footprint good set
their estimates use only that footprint.
Replacing \(h_{\rm c}\) by \(2h_{\rm c}\)
in the estimates still gives the same
eventual smallness conditions. At each
fixed period the differentiated sums
thus converge on a neighbourhood of
globally good fields, including the
boundary at the nominal threshold.
No guard or hard mask is differentiated.

The stationary action and loop--Haar
factor also have compatible physical
scalar families. Expand them in \(y\)
at zero and substitute the tree family.
At the replica every curvature and
noise mismatch is \(O(h_{\rm c})\).
The phase and loop values start
quadratically; their first fine
derivatives retain one such factor.
Terms with two or more fine derivatives
contain at least two graph factors
after substitution.
The Haar logarithm starts quadratically
in \(m\).
The rooted majorants therefore give
the physical prices
\begin{equation}
       \|\Psi\|_{\rm local}^{\rm phys}
                         \le C_Bh_{\rm c}^2,\qquad
       \|\psi_*\|_{\rm local}^{\rm phys}
                         \le C_B(1+\Lambda)h_{\rm c}^2 .
 \label{eq:compatible-stationary-scalar-prices}
\end{equation}
The notation means these scalar
families, not their extensive total
values. Positive-order physical
derivatives preserve the two small
factors: removing a curvature supplies
its allowed \(h_{\rm c}\)-scaled
direction instead.

For completeness the determinant has
a compatible scalar expansion as well.
At fixed period the spectral identity is
\begin{equation}
 \frac12\log\frac{\det H(V)}{\det H_0}
   =\frac12\int_0^\infty
        \operatorname{Tr}
          \{(H_0+tI)^{-1}-(H(V)+tI)^{-1}\}\,\dd t .
 \label{eq:compatible-physical-logdet}
\end{equation}
Insert \eqref{eq:compatible-stationary-resolvent}.
Every term with at least one
\(\widetilde E_t\) has family bound
\[
             \frac{C_B(h_{\rm c}+e^{-c_0B})}
                                      {(\eta_*+t)^2},
\]
which is integrable. The trace closes
its endpoint path, giving a gauge
invariant scalar record.

For its zero-defect term compare the
flat periodisation of \(P_t^0\) with
the diagonal of \(P_t^0[V]\).
The zero-displacement diagonal path
is empty, so it cancels exactly.
Any remaining diagonal image has a
nonzero cover displacement equal to
a period multiple. Its length is at
least \(c\); it is kept as a winding
path, with an \(e^{-c_0B}\) tail.
To check integrability of this leading
difference at large \(t\), use
\[
 P_t^0-\frac1{1+t}I
           =\frac1{1+t}(I-A_0)P_t^0 .
\]
The nonzero-displacement coefficients
are therefore bounded by
\(C/((1+t)(\eta+t))\) in the cover
norm. This is integrable in \(t\).
The label-by-label integral defines
local scalars as above. Their price is
\begin{equation}
                 C_B(h_{\rm c}+e^{-c_0B})
 \label{eq:compatible-physical-logdet-price}
\end{equation}
per attached root, with a positive
native reserve.
All terms vanish at \(V=1\).
The winding terms have been retained,
including on flat backgrounds with
nontrivial holonomy.

Equations~\eqref{eq:compatible-stationary-scalar-prices}
and \eqref{eq:compatible-physical-logdet}
provide physical families for the same
stationary prefactor as
\eqref{eq:local-prefactor-target}.
Their sums agree with the chart sums
on the common domain by the exact
stationary and spectral identities.
This compatibility does not require
individual coordinate monomials in
unrelated gauges to agree.

Finally collect the requests of this construction.
After fixing the profile, native and covering reserves,
the chart and tree estimates require
\[
 B^kh_{\rm c}\ll1,\qquad B^ke^{-c_0B}\ll1,\qquad
                                      e^{-c_0B}\le h_{\rm c}
\]
for the finitely many powers \(k\) occurring above.
The prescribed definitions
\eqref{eq:compatible-physical-parameters} verify these
requests simultaneously. Indeed, for every fixed finite power,
\[
       B^k h_{\rm c}\longrightarrow0,\qquad
       B^k e^{-c_0B}\longrightarrow0,\qquad
                         e^{-c_0B}\le h_{\rm c}
                         \quad\hbox{eventually}.
\]
They pay the short chart requests,
all tree smallness inequalities, and
all fixed reserve reductions at one
profile threshold.
Periods satisfying
\eqref{eq:compatible-period-reserve}
can simultaneously be chosen of order
\(B\). Then \(f^4=(Lc)^4\le C_LB^4\),
within the auxiliary size regime of
Theorem~\ref{thm:tail-auxiliary} with
a fixed sufficiently large logarithmic
power. The estimates themselves are
uniform on all larger admissible
periods satisfying the reserve.
Its placement in one Wilson calendar
is the separate choice described in
\cite[Section~26.4, p.~142]{Main};
no volume subsequence has been used
to construct these local functions.


\section{The nonquadratic Gaussian logarithm}
\label{sec:connected-Gaussian-logarithm}

The determinant accounts for the quadratic integral.
The remaining \(Z_D\) contains correlated fluctuations
of the nonlinear phase and of the loop--Haar factor.
We first identify its lowest connected coefficient,
then prove a convergent local logarithm.
The hard-core expansion and tree organisation used
below are those of
\cite[Section~2, (2.2)--(2.7), preprint pp.~3--4;
Section~4.1, Proposition~5, pp.~10--11]{FP07}.
The gauge-theory separation into local connected
contributions may be compared with
\cite[Section~2, (2.11)--(2.13), printed p.~14]{Bal88a}.
Their model-specific interaction estimates are not
being substituted for the estimates proved here.

\subsection{Whitening retains the central domain}

For fixed coarse \(x\) let
\[
             C(x)=H(x)^{-1},\qquad
             W(x)=H(x)^{-1/2}.
\]
Let \(\xi_e\), one for each free fine link,
be independent standard Gaussian vectors in the
fixed Lie algebra inner product.
Then \(z=W(x)\xi\) has covariance \(C(x)\).
This change of variables is called whitening:
the integration variables \(\xi\) are independent,
while the field \(z\) remains correlated.
The central domain becomes the ellipsoid
\[
                            |W(x)\xi|_2\le D .
\]
It does not factor over \(\xi_e\).

Restore \(x\) in the existing interaction
\eqref{eq:central-exponent}:
\begin{equation}
 \begin{split}
 P(g,z;x)={}&-g^{-2}
     \{\varphi(x,m(x)+gz)-\varphi(x,m(x))
                         -\tfrac{g^2}{2}\langle z,H(x)z\rangle\}\\
           &-\{\psi(x,m(x)+gz)-\psi(x,m(x))\}.
 \end{split}
 \label{eq:connected-actual-interaction}
\end{equation}
Thus, with all original normalisers already in
\(N_g^{\mathrm{phys}}\),
\[
           Z_D(g,x)=\int_{|z|_2\le D}
                                e^{P(g,z;x)}\,\dd\mu_{C(x)}(z).
\]
The apparent singularity at \(g=0\) is removable:
stationarity removes the linear phase term and
the displayed quadratic term removes its Hessian.
When \(g\) is complex below, the profile, coarse
field, nominal coupling, \(\Lambda,D\), and all
Gaussian cutoffs are held fixed. The physical prefactor
\(N_g^{\mathrm{phys}}\) is evaluated at the actual positive
coupling; this auxiliary complex variation of the interaction
does not vary that prefactor.

At \((x,m(x))\) put
\[
               T_k=D_y^k\varphi(x,m(x)),\qquad
               P_k=D_y^k\psi(x,m(x)).
\]
They are symmetric multilinear forms on the free
coordinate space. In particular
\[
 P(g,z;x)=g u_1(x,z)+g^2u_2(x,z)+O(g^3),
 \quad
 u_1=-\tfrac16T_3[z^3]-P_1[z],\qquad
 u_2=-\tfrac1{24}T_4[z^4]-\tfrac12P_2[z^2].
\]
For this formula alone the remainder is a
finite-domain Taylor remainder.
Let \(p_D(x)=\mu_{C(x)}\{|z|_2\le D\}\).
Parity under \((g,z)\mapsto(-g,-z)\) gives
\begin{equation}
 \log Z_D(g,x)=\log p_D(x)+g^2L_{2,D}(x)+O(g^4),
 \quad
 L_{2,D}=\frac1{p_D}
       \int_{|z|_2\le D}(u_2+\tfrac12u_1^2)\,\dd\mu_C .
 \label{eq:connected-finite-second-coefficient}
\end{equation}
This keeps \(Z_D(0,x)=p_D(x)\).
It is not yet a volume-uniform remainder estimate.

\subsection{The full Gaussian second coefficient}

The polynomial coefficient without the central
cutoff is well defined:
\[
                    L_2(x)=\mathbb E_{\mu_C}
                                      (u_2+\tfrac12u_1^2).
\]
For a centred Gaussian,
\(\mathbb E(z_af)=\sum_bC_{ab}\mathbb E(\partial_bf)\).
Applying this integration-by-parts identity
inductively to monomials proves Wick's rule:
even moments sum over pairings, with one
covariance factor per pair, and odd moments vanish.
In an orthonormal real basis, all free indices
including colours are summed in
\begin{equation}
 \begin{split}
 L_2={}&-\tfrac18T_{abcd}C_{ab}C_{cd}
                  -\tfrac12P_{ab}C_{ab}\\
       &+\tfrac18T_{abc}T_{def}C_{ab}C_{de}C_{cf}
                  +\tfrac1{12}T_{abc}T_{def}C_{ad}C_{be}C_{cf}\\
       &+\tfrac12T_{abc}P_dC_{ab}C_{cd}
                  +\tfrac12P_aP_bC_{ab}.
 \end{split}
 \label{eq:connected-explicit-L2}
\end{equation}
For two cubic vertices, nine of the fifteen
pairings have one edge between vertices,
and six have three; their common coefficient
\(1/72\) gives \(1/8\) and \(1/12\).
The cubic--linear term has three pairings
and coefficient \(1/6\), giving \(1/2\).
All displayed graphs have one vertex or two
vertices joined by at least one covariance.
This is connectedness at this order.

Each of the six terms has an exact connected
local family, of anchored price
\begin{equation}
                              C_B(1+\Lambda)^2.
 \label{eq:connected-L2-family-price}
\end{equation}
Here and below the same statement is available
in the chart norm and in the physical norm;
the family norm uses the corresponding
single-function norm.
To verify the bound, keep the rooted labels
of \(T_3,T_4,P_1,P_2,C\) and every covariance
endpoint. A vertex has price \(C_B\), or
\(C_B(1+\Lambda)\) for a loop--Haar vertex.
For one vertex pin an index and sum its
remaining indices by its rooted bound;
internal covariance pairs use their
weighted suprema.
For two vertices choose an edge joining
them. Pin the first endpoint, sum the
first vertex, then the chosen covariance
by its row bound, then the second vertex
rooted at its endpoint.
Remaining covariance entries use weighted
suprema. This gives no volume factor.
The union of the vertex trees and
covariance paths is connected.
The footprint and derivative estimates of
Section~\ref{sec:local-stationary-prefactor}
pay the covering weight and fixed-order
coarse derivatives.
At most two loop--Haar vertices occur,
which gives the stated allowance.
This is the calculation of
\cite[Section~23.3, pp.~117--119]{Main}.

On the chart, discarding degrees at most
two in the same analytic families gives
\begin{equation}
             \|g^2(1-T_{\le2})L_2\|_{\rm local}^{\rm chart}
                         \le C_Bg^2(1+\Lambda)^2r^3 .
 \label{eq:connected-L2-chart-subtraction}
\end{equation}
The radius-ratio proof is exactly that of
Corollary~\ref{cor:local-stationary-remainder}.
The retained jets remain part of the
later exact scalar and coupling extraction.

\subsection{Bounded independent coordinates}

For \(R\ge1\), project each independent vector
radially to its radius-\(R\) ball:
\[
 \xi_e^R=
 \begin{cases}
          \xi_e,&|\xi_e|_{\g}\le R,\\
          R\xi_e/|\xi_e|_{\g},&|\xi_e|_{\g}>R.
 \end{cases}
 \qquad
 z^R=W\xi^R,\qquad
 \mathcal Z_R(g,x)=\mathbb E e^{P(g,z^R;x)} .
\]
No derivative is taken in the Gaussian
variables. The radial projection is odd
and commutes with orthogonal transformations,
including \(\operatorname{Ad}(G)\).
The same definition applies to physical
good \(V\), using the data of
Proposition~\ref{prop:compatible-stationary-data}.

At the stationary point the exact series is
\begin{equation}
 P(g,W\xi^R)
   =-\sum_{k\ge3}\frac{g^{k-2}}{k!}T_k[(W\xi^R)^k]
     -\sum_{k\ge1}\frac{g^k}{k!}P_k[(W\xi^R)^k].
 \label{eq:connected-bounded-interaction-series}
\end{equation}
Keep every vertex, kernel, and Gaussian endpoint
label in this expansion.
The resulting primitive scalar letters
\(v_\alpha(g,\cdot,\xi)\) have finite connected
footprints \(X_\alpha\), including the Gaussian
variables they read.
They satisfy \(v_\alpha(0)=0\) and
\[
             v_\alpha(-g,\cdot,\xi)
                         =v_\alpha(g,\cdot,-\xi).
\]
Use an incidence norm, charging every site
of a footprint rather than just its anchor:
\begin{equation}
 \sup_o\sum_{\alpha:o\in X_\alpha}
       e^{(\kappa+6)n_B(X_\alpha)}
            \sup_\xi\|v_\alpha(g,\cdot,\xi)\|^{\rm alg}_{q;o_\alpha,X_\alpha}
                   \le C_B|g|H_R,\qquad
                 H_R=R^3+(1+\Lambda)R .
 \label{eq:connected-primitive-incidence}
\end{equation}
The local norm on the left is the chosen
chart or physical norm.

Here are the quantitative sums.
The vertex analytic bounds give \(C_B k!a^{-k}\)
and \(C_B(1+\Lambda)k!a^{-k}\) at a fixed
smaller fine radius \(a>0\).
Kernel labels pay each Gaussian endpoint
sum. Each endpoint contributes at most \(R\).
For \(|g|C_BR\) sufficiently small, the two
degree sums are bounded by
\[
 C_B\sum_{k\ge3}|g|^{k-2}(C_BR)^k
                      \le C_B|g|R^3,\qquad
 C_B(1+\Lambda)\sum_{k\ge1}(|g|C_BR)^k
                      \le C_B(1+\Lambda)|g|R .
\]
Finite powers of the polynomial allowance
are included in the final \(C_B\).
Choosing the fixed small ratio before
summing is essential.
The spatial joining lemma pays the union
of the vertex trees and kernel paths.
To pass from anchors to incidence, a
specified site in that union is at
distance at most its native length
from the anchor.
Reserve half the native exponential
weight; the sum of possible anchor
positions is bounded by
\(\sum_{z\in\mathbb Z^4}(1+|z|)^m e^{-c|z|}<\infty\)
for the fixed derivative and incidence
powers. Short filling boxes introduce
only polynomial \(B\)-costs.
The covering and whole-product estimates
then prove
\eqref{eq:connected-primitive-incidence}.
This supplies the actual interaction
bound of \cite[Lemma~24.1, pp.~120--121]{Main}.

\subsection{From primitive letters to a hard-core gas}

Put \(a_\alpha=e^{v_\alpha}-1\).
It has the same read footprint and finite
length certificate as its primitive letter.
If \(t_\alpha\) bounds the algebraic norm
of \(v_\alpha\), then
\[
          \|a_\alpha\|^{\rm alg}_q
                   \le e^{t_\alpha}-1\le2t_\alpha
                            \quad(t_\alpha\le1).
\]
Absolute summability at fixed period
justifies
\[
           e^{\sum_\alpha v_\alpha}
                =\sum_{A\ {\rm finite}}\prod_{\alpha\in A}a_\alpha.
\]
Each primitive address is selected at
most once in this identity.
Call a nonempty connected finite set
of addresses a polymer, and define
\[
       X(A)=\bigcup_{\alpha\in A}X_\alpha,\qquad
       \zeta(A)=\mathbb E\prod_{\alpha\in A}a_\alpha .
\]
Different overlap components have disjoint
declared footprints, hence read disjoint
independent Gaussian vectors.
Their expectations factor exactly.
It follows that
\begin{equation}
          \mathcal Z_R
              =\sum_{\{A_1,\ldots,A_n\}\ {\rm disjoint}}
                                      \prod_{j=1}^n\zeta(A_j).
 \label{eq:connected-polymer-gas}
\end{equation}
Disjointness refers to footprints.
An enlarged declared read can join extra
addresses; disjoint components still
factor, which is all this identity uses.

Take products and expectations in the
algebraic derivative norm.
Their Gaussian law is fixed, so the
norm of an expectation is at most the
expectation of its norm.
The whole-product anchor estimate costs
one covering reserve and a polynomial
allowance.
Apply Lemma~\ref{lem:connected-positive-tree}
to the primitive prices at the remaining
stronger covering exponent.
For sufficiently small \(C_B|g|H_R\) this
gives
\begin{equation}
       \sup_o\sum_{A:o\in X(A)}
          e^{(\kappa+3)n_B(X(A))}
                         \|\zeta(A)\|^{\rm alg}_q
                                   \le C_B|g|H_R .
 \label{eq:connected-polymer-incidence}
\end{equation}
The finite factors \(D_B\) are included
in the polynomial smallness allowance.
Repeated species will occur in the
logarithm, even though they did not
occur in a selected primitive subset.

\begin{lemma}[A convergent rooted logarithm]
\label{lem:connected-log-bound}
Suppose the left side of
\eqref{eq:connected-polymer-incidence}
is at most \(\epsilon\), with
\(D_B\epsilon\le1/2\).
The gas has a nonvanishing analytic
continuation in the smallness region.
Its logarithm is an absolutely convergent
family of connected footprints, with
anchored and incidence norm at covering
exponent \(\kappa\) at most \(C_B\epsilon\).
All estimates are uniform in the finite
volume.
\end{lemma}
\begin{proof}
For an ordered tuple of polymers define
\[
 \mathcal U(A_1,\ldots,A_n)
   =\sum_{G\ {\rm connected\ on}\ \{1,\ldots,n\}}
       \prod_{\{i,j\}\in E(G)}
                  (-\mathbf1_{\{X(A_i)\cap X(A_j)\ne\varnothing\}}),
\]
with value one for \(n=1\).
The logarithm series is
\begin{equation}
        \sum_{n\ge1}\frac1{n!}
          \sum_{A_1,\ldots,A_n}
               \mathcal U(A_1,\ldots,A_n)\prod_{i=1}^n\zeta(A_i).
 \label{eq:connected-Ursell-log}
\end{equation}
Repeated polymers are allowed in these
ordered tuples, with their factorial.

For a fixed incompatibility graph,
\(|\mathcal U|\) is bounded by its
number of spanning trees.
To see the cancellations, root a connected
spanning subgraph at vertex one, use its
distance levels, and choose the
smallest-labelled predecessor as parent.
For the resulting tree, optional edges
join vertices at the same level or
adjacent levels without replacing that
chosen parent.
The connected subgraphs giving the
same tree form precisely an interval
between the tree and its allowed
completion. Summing signs over the
optional edges gives zero or a number
of absolute value one.
The intervals partition the connected
subgraphs, proving the tree bound.
This is the partition argument in
\cite[Section~4.1, Proposition~5 and (4.4),
preprint p.~10]{FP07}.

Write \(k_A=n_B(X(A))\) and
\[
              u_A=e^{(\kappa+1)k_A}\|\zeta(A)\|^{\rm alg}_q .
\]
The hypothesis gives
\(\sup_o\sum_{A:o\in X(A)}u_Ae^{k_A}\le\epsilon\).
The positive-tree lemma therefore bounds
the price-free root sum of logarithmic
clusters by \(e^{k_A}\).
The whole-product estimate
\eqref{eq:connected-one-reserve} bounds
each cluster's differentiated output,
including its final covering weight,
by \(C_B\prod_i u_{A_i}\).
This pays the anchor change and the
factor \(q!\) once, for the whole cluster.

For a fixed site \(o\), bound the
indicator that its union contains \(o\)
by the number of tuple entries whose
footprint contains \(o\).
Selecting this entry changes \(1/n!\)
to \(1/(n-1)!\).
Its total absolute price is at most
\[
                 C_B\sum_{A:o\in X(A)}u_Ae^{k_A}
                                  \le C_B\epsilon .
\]
This is the incidence bound and also
controls any prescribed anchored
assignment, without a sum over all sites.

For finite alphabets the gas can also
be written
\[
 \sum_{n\ge0}\frac1{n!}\sum_{A_1,\ldots,A_n}
        \prod_i\zeta(A_i)
        \prod_{i<j}
          (1-\mathbf1_{\{X(A_i)\cap X(A_j)\ne\varnothing\}}).
\]
Self-incompatibility kills tuples
with repeated species in this gas.
Expand the last product into graphs.
Decomposing a graph into its connected
components gives precisely the
exponential of
\eqref{eq:connected-Ursell-log},
with the same factorials, near zero
fugacity.
The absolute bounds continue the
identity throughout the stated
smallness region. Countable alphabets
follow by their uniform absolute
majorants at each finite period.
The exponential identity implies
nonvanishing and fixes the branch of
the logarithm which is zero at \(g=0\).
\end{proof}

This is the norm and convergence
argument of \cite[Lemma~24.2,
pp.~122--123]{Main}.
It has not assumed independence of
the correlated field \(z\).

\subsection{The entire fourth-order remainder}

After enlarging a fixed polynomial
allowance, there is a common complex
radius
\[
                         r_g=\frac{c_1}{C_BH_R}
\]
on which the preceding logarithm and
all its absolute local majorants
converge. This also pays
\(|g|C_BR\) in the primitive expansion.
Here \(c_1>0\) is a fixed profile constant.

\begin{theorem}[The full connected interacting remainder]
\label{thm:connected-fourth-order}
For \(|g|\le r_g/2\), the bounded-coordinate
reference has an exact connected local
family expansion
\begin{equation}
        \log\mathcal Z_R(g)
                    =g^2L_2+\sum_\alpha r_\alpha(g),
 \label{eq:connected-full-reference-identity}
\end{equation}
whose anchored norm is at most
\begin{equation}
       C_B\left\{
          |g|^4H_R^4+
          |g|^2(1+\Lambda)^2(1+R)^6e^{-c_1R^2}
                                                        \right\}.
 \label{eq:connected-fourth-order-price}
\end{equation}
The statement holds in the chart norm,
and in the physical norm for the
compatible families just constructed.
In the physical case the functions
are gauge invariant raw local scalars,
smooth near their own good sets.
The bounds are independent of cutoff
and physical volume within the fixed
profile and period reserve, and retain
a positive native length reserve.
\end{theorem}
\begin{proof}
Parity of each primitive letter and
the change \(\xi\mapsto-\xi\) make each
polymer weight even in \(g\).
Every term of the logarithm is therefore
even and vanishes at zero.
Its second coefficient sums to
\[
            L_2^R=\mathbb E(u_2(W\xi^R)
                                      +\tfrac12u_1(W\xi^R)^2),
\]
by differentiating the bounded-coordinate
integral twice.
Apply the Cauchy estimate to each
connected label on a fixed smaller
circle inside \(r_g\), keeping the
absolute family majorant.
Summing its even degrees at least
four gives \(C_B|g|^4H_R^4\).
All finite powers of the polynomial
allowance caused by \(r_g^{-4}\)
are included in \(C_B\).

To replace \(L_2^R\) by \(L_2\),
expand its polynomial into the same
vertex and \(W\)-entry labels.
Each monomial has Gaussian degree
at most six and reads at most six
independent Gaussian vectors, allowing
repetitions.
It agrees with its unprojected version
when all these vectors have norm
at most \(R\).
Gaussian tails and Cauchy--Schwarz
therefore bound its expectation
difference by
\[
                       C_G(1+R)^6e^{-c_GR^2}
\]
times its coefficient norm.
The coefficient majorants and their
coarse derivatives sum by
\eqref{eq:connected-L2-family-price}.
This proves the second term in
\eqref{eq:connected-fourth-order-price}.
There was no union bound over the
whole torus in this local estimate.

For physical functions the tensors,
kernel entries, and primitive guards
are the compatible records of
Section~\ref{sec:compatible-physical-kernels}.
Internal adjoint transformations cancel
in contractions. The remaining Gaussian
endpoints transform orthogonally;
their independent radial-projected
law is invariant.
Thus each contracted polymer and each
logarithmic cluster is a gauge
invariant raw local scalar.
Expectations use a fixed Gaussian law,
so physical derivatives differentiate
only the displayed coefficients and
the actual \(W\)-entries.
The same product and rooted estimates
apply on the cluster's own good set.

Choose the final native exponent
below the finitely many input widths.
Length is additive under tree
substitution, matrix multiplication,
and formation of a cluster.
The same convergent majorants then
include that final exponential weight.
Short boxes keep their bounded
covering price, while long records
keep their path covers.
\end{proof}

This proves the connected-reference
conclusions of \cite[Theorem~24.4,
pp.~123--124, and Theorem~26.1,
pp.~139--140]{Main}.
The leading coefficient is the full
Gaussian polynomial \(L_2\);
the remainder is an estimate for
the whole convergent logarithm.
It remains to subtract the physical
vacuum and compare with the original
central cutoff.


\section{Vacuum subtraction and the original central cutoff}
\label{sec:connected-original-central}

Two operations are still needed before the connected
reference estimate can be used.
Its leading physical coefficient has a flat scalar
value, which belongs in the vacuum normalisation.
The original \(Z_D\) also retains its finite
ellipsoidal domain.
We handle these operations separately, keeping
their exact signs and errors.

\subsection{A small-curvature gain after physical flat subtraction}

\begin{lemma}[The physical leading coefficient minus its flat value]
\label{lem:connected-physical-flat-subtraction}
The physical family \(L_2=\sum_\alpha l_\alpha\)
can be relabelled, preserving its sum, so that
\begin{equation}
       \|(l_\alpha-l_\alpha(1))_\alpha\|_{\mathcal E,\kappa}
             \le C_B(1+\Lambda)^2(h_{\rm c}+e^{-c_0B}).
 \label{eq:connected-physical-flat-gain}
\end{equation}
Only short footprints are enlarged to boxes;
long footprints retain their connected paths.
\end{lemma}
\begin{proof}
Split the retained native labels at
\(l_\alpha\le B/C_*\), with a sufficiently
large fixed geometric constant \(C_*\).
A short label, including all its paths,
primitive short boxes and halos, fits in
a nonwrapping rooted box of radius
\(C_L(1+l_\alpha)\).
This box has bounded \(B\)-covering number.

We record the coordinate radius needed
for the flat comparison.
The unguarded primitive formulas are
holomorphic near the identity.
On a complex coarse link ball of radius
\(r_{x,B}\), their seed and linear defect
prices in the stationary tree equation
are bounded by \(K_Br_{x,B}\) and
\(K_B(r_{x,B}+e^{-c_0B})\).
Word growth uses only a fixed part of
the native path reserve.
For example one may choose
\begin{equation}
       r_{x,B}
            =\min\{r_{\rm br}/4,\ r_0/(256K_B^2)\},
 \label{eq:connected-physical-complex-radius}
\end{equation}
where \(r_{\rm br}>0\) is a fixed branch
and word radius and \(K_B\) includes
the preceding finitely many constants.
The scalar tree majorant used in
Proposition~\ref{prop:compatible-stationary-data}
then converges on this ball.
The resolvent series does as well.
Thus the unguarded \(L_2\) labels have
an absolute complex coefficient majorant
bounded by \(C_B(1+\Lambda)^2\), with
a remaining native width.
This concerns analytic formulas near
the identity, not analyticity across
an extension guard.

At a smaller radius the Cauchy estimates
through order \(q\) cost only powers
of \(r_{x,B}^{-1}\), hence polynomials
in \(B\).
For a short footprint its number of
link-colour variables is also polynomial
in \(B\). Consequently the sums of its
fixed-order coordinate derivative norms,
including their coefficient labels,
have total price \(C_B(1+\Lambda)^2\).

Put the physical field into the local
gauge on its enlarged short box.
Its link logarithms have maximum
size \(C(1+l_\alpha)^2h_{\rm c}\),
inside the preceding chart eventually,
since every fixed polynomial in \(B\)
times \(h_{\rm c}\) tends to zero.
The scalar \(l_\alpha(V)-l_\alpha(1)\)
is gauge invariant and vanishes at the
flat local chart.
The value and ordered derivative
calculation of
\eqref{eq:compatible-short-record-price}
therefore gives its price as
\[
                   Cp_\alpha h_{\rm c}
                                      (1+l_\alpha)^{C_{q,s}},
\]
where the native-weighted rooted sum
of \(p_\alpha\) is at most
\(C_B(1+\Lambda)^2\).
Polynomial length factors and the
incidence of the enlarged boxes are
paid by the remaining native width.
This proves the short part of
\eqref{eq:connected-physical-flat-gain}.

For long labels make no enlargement.
Their original physical prices, and
their flat values, have sum at most
\(C_B(1+\Lambda)^2e^{-c_0B}\).
The native reserve gives this bound
directly. Possible winding contributions
are included; no flat-holonomy
cancellation is assumed for them.
\end{proof}

This is the flat-subtraction estimate
of \cite[Lemma~26.2, pp.~140--141]{Main}.
The polynomial radius in
\eqref{eq:connected-physical-complex-radius}
is sufficient; a radius independent
of \(B\) is not needed for this step.

At real physical \(g=\beta^{-1/2}\),
define the signed vacuum scalar
\[
                         c_{\rm int}=-\log\mathcal Z_R(g,1).
\]
Subtract each local family's flat
value in \eqref{eq:connected-full-reference-identity}.
This gives the exact physical good-field
identity
\begin{equation}
          -\log\mathcal Z_R(g,V)
                   =c_{\rm int}+\sum_\alpha f_\alpha(V),
                   \qquad f_\alpha(1)=0 ,
 \label{eq:connected-physical-vacuum-identity}
\end{equation}
with the local price
\begin{equation}
 \begin{split}
 H_{\rm int}\le C_B\{&
     g^2(1+\Lambda)^2(h_{\rm c}+e^{-c_0B})
             +g^4H_R^4\\
     &+g^2(1+\Lambda)^2(1+R)^6e^{-c_1R^2}\}.
 \end{split}
 \label{eq:connected-physical-nonconstant-price}
\end{equation}
If \(n_c=c^4\) is the number of coarse
sites, the extensive scalar satisfies
\[
       |c_{\rm int}|/n_c
                \le C_B\{g^2(1+\Lambda)^2+g^4H_R^4\}.
\]
It is estimated per site after the
local bounds, not as a uniformly
bounded total scalar.

\subsection{A common projected radius on auxiliary volumes}

We now compare to the original central
integral on precisely the auxiliary
size regime of
Theorem~\ref{thm:tail-auxiliary}.
Let \(d_*(\widehat\beta)\) be a fixed
polynomial in \(1+\log\widehat\beta\)
bounding the number of free real
coordinates on every allowed auxiliary
torus. Choose
\begin{equation}
       R=\frac{D}{4C_W\sqrt{d_*}},\qquad
       D=d_0P(\widehat\beta),\qquad
       C_W\ge\max\{1,\sup_x\|W(x)\|_{2\to2}\}.
 \label{eq:connected-common-white-radius}
\end{equation}
The gap makes \(C_W\) a fixed profile
constant. The choice uses the common
auxiliary size bound, not the eventual
physical volume.
It gives \(R\ge1\) eventually.
Moreover \(R,H_R,C_B\) are subpower
functions of \(\widehat\beta\), whereas
\(R^2\) dominates every fixed power
of \(\log\widehat\beta\).
In particular the physical
\(|g|\le\sqrt2\widehat\beta^{-1/2}\)
eventually lies below \(r_g/2\).

\begin{lemma}[The original central logarithm and the bounded reference]
\label{lem:connected-central-comparison}
On these auxiliary volumes, for every fixed
derivative order \(q\), the chart difference
\[
                    e_D(x)=\log Z_D(g,x)-\log\mathcal Z_R(g,x)
\]
has \(C_w^q\) norm at most \(Ce^{-c_1R^2}\)
at physical \(g=\beta^{-1/2}\).
Here \(w=\widehat\beta^{-7/16}\), and
\begin{equation}
 \|F\|_{C_w^q}
   =\max_{0\le j\le q}
      \sup_{\|x\|_\infty\le w}
      \sup_{\|u_\nu\|_\infty\le w}
                         |D^jF(x)[u_1,\ldots,u_j]|.
 \label{eq:connected-global-chart-norm}
\end{equation}
The constants are uniform over the stated
auxiliary tori.
\end{lemma}
\begin{proof}
On the event
\(\max_e|\xi_e|_{\g}\le R\), projection
does nothing, and
\[
              |W\xi|_2\le C_W\sqrt{d_*}R\le D/4 .
\]
The original central integrand and the
bounded reference agree there.
The complementary probability is at
most \(C_Gd_*e^{-c_GR^2}\).
On the central domain, the rooted
interaction bounds summed over the
auxiliary sites give
\[
         |P(g,z;x)|
             \le C_Ln_c|g|\{D^3+(1+\Lambda)D\}=o(1).
\]
The same holds for the projected
integrand, using \(|z^R|_2\le D/4\).
Every factor apart from \(g\) is
subpower and \(n_c\) is polylogarithmic.
Thus the value difference is at most
\(Ce^{-c_1R^2}\), and both positive
integrals are bounded below by \(1/2\)
eventually. Their logarithms have the
same value-error bound, also on the
slightly larger chart
\(\|x\|_\infty\le2w\).

We give the derivative argument without
differentiating a moving ellipsoid.
For \(Z_D\) use its fixed-\(z\) ball
and its moving Gaussian density with
respect to Lebesgue measure.
Its logarithmic derivative is
\[
       D\log\left(\frac{\dd\mu_{C(x)}}{\dd z}\right)(z)
          =\tfrac12\operatorname{Tr}(C\,DH)
                              -\tfrac12\langle z,DH z\rangle .
\]
Repeated derivatives give polynomials
in \(z\), derivatives of \(H\),
and its inverse. The trace costs at
most a fixed multiple of the dimension
per derivative. The interaction
derivatives have the same rooted
analytic bounds, summed over the
auxiliary sites.
For every fixed order their integrals
are bounded by a fixed polynomial
in \(d_*,D,\Lambda,g^{-1}\).
For \(\mathcal Z_R\), differentiate
its fixed independent Gaussian
expectation; the bounded variables
and derivatives of \(W\) give the
same type of polynomial bounds.
The lower bounds just proved give
them for the logarithms as well.
At every fixed order these prices
are at most \(C\widehat\beta^K\)
for some fixed \(K\).

For an order \(j\le q\) mixed derivative,
use finite differences in its allowed
directions with dimensionless step
\(t=e^{-aR^2}\).
The difference divided by \(t^j\)
differs from the derivative by at most
\(Ct\) times the order-\((j+1)\)
bound. The value-error contribution
is at most \(Ct^{-j}e^{-c_1R^2}\).
Choose \(a>0\) with \(qa<c_1/2\).
Since \(R^2/\log\widehat\beta\to\infty\),
both errors are \(Ce^{-c_2R^2}\).
All difference arguments stay in
the larger \(2w\)-chart.
This proves the claimed fixed-order
bound, after renaming the decay
constant.
\end{proof}

This proves \cite[Lemma~24.5,
p.~124]{Main} with the derivative
interpolation made explicit.
The difference includes the fact
that \(Z_D(0,x)=p_D(x)\), whereas
\(\mathcal Z_R(0,x)=1\).

\subsection{The resulting estimates and their scope}

Combining the convergent family with
the central comparison gives the
exact auxiliary identity
\begin{equation}
             \log Z_D(g,x)
                    =g^2L_2(x)+\sum_\alpha r_\alpha(g,x)+e_D(x).
 \label{eq:connected-original-central-identity}
\end{equation}
The family price is
\eqref{eq:connected-fourth-order-price}.
After retaining the degree-at-most-two
coarse Taylor polynomial of \(L_2\),
the local chart error has price
\begin{equation}
 C_B\{g^4H_R^4+g^2(1+\Lambda)^2w^3\}
                      +Ce^{-c_1R^2}
                     \le C_\nu\widehat\beta^{-2+\nu}.
 \label{eq:connected-original-central-rate}
\end{equation}
The clipped coefficient tail is
included in the exponential term.
The two polynomial orders are
\(\widehat\beta^{-2+\nu}\) and
\(\widehat\beta^{-37/16+\nu}\).
To assign \(e_D\) its full auxiliary
footprint, changing to the spatially
weighted chart norm costs only
polylogarithmic diameter factors.
Its covering price is at most
\(\exp\{C(\log\widehat\beta)^m\}\)
for one fixed finite \(m\).
Both are absorbed by the Gaussian
error since \(R^2\) dominates every
such logarithmic power.
This is the full auxiliary estimate
of \cite[Corollary~24.6,
p.~125]{Main}.

For the physical bounded-coordinate
reference, under the same comparable
couplings and uniform subpower loop
allowances, flat subtraction gives
instead
\begin{equation}
                         H_{\rm int}
                            \le C_\nu\widehat\beta^{-3/2+\nu}.
 \label{eq:connected-physical-final-rate}
\end{equation}
Indeed the leading nonconstant term
is \(g^2(1+\Lambda)^2h_{\rm c}\);
the analytic fourth-order term has
order \(\widehat\beta^{-2+\nu}\);
the two exponential tails are
smaller than every inverse power.
These are physical family estimates
on unrestricted volumes satisfying
the period reserve, for the reference
\(\mathcal Z_R\).

The comparison of a single central
ball with that reference was proved
only on auxiliary volumes.
On an unrestricted physical volume
the Gaussian-coordinate failure
events must remain in the exact
large-field reassembly.
Likewise the fine indicator in
\eqref{eq:tail-exact-masked-central}
can be replaced by one only on its
proved compatibility domain.
The present expansion does not
enlarge that domain.
Incoming activities keep their
original terms in
\eqref{eq:reference-exact-splitting}.
Sections~\ref{sec:calibrated-rows}--\ref{sec:general-central-comparison}
now construct the calibrated fine-curvature reserve and prove
the indicator criterion at a general good auxiliary background.
Sections~\ref{sec:physical-selected-Gaussian}--\ref{sec:exact-masked-reassembly}
then return to the exact physical integral and give the masked
reassembly with its explicit convergence input. Principal-action
matching and the prescribed quadratic extraction follow after
the remaining physical packet estimates.


\section{Calibration before the block side is chosen}
\label{sec:calibrated-rows}

The gap and path budgets proved above give a stationary point for
every admissible principal profile. They do not, by themselves,
say that a coarse curvature of size \(h_{\rm c}\) produces a fine
curvature smaller than \(h\). The two thresholds are asymptotically
equal. We therefore need the scaling of the Gaussian minimiser,
with its constant chosen before the block side \(L\). We first
construct that datum for the actual Gaussian histories. The
subsequent nonlinear argument will use only two row identities
and their absolute coefficient prices.

An \(r\)-cochain is an array on oriented elementary \(r\)-cells,
with a sign change on reversing orientation. Write \(d_r\) for
the oriented coboundary: \(d_0\) is the gradient and \(d_1\) is
the curl. Cancellation of opposite faces gives
\(d_{r+1}d_r=0\). Their real inner products are counting sums,
including cell types and colour components.
In particular, a superscript \(T\) below is the transpose for
counting measure.

For a translation kernel from coarse rows to fine rows define
\[
 |k|_\sigma=\max_{r,\alpha}
       \sum_{u\in\Z^4,\gamma}e^{\sigma|u|_1}
                             |k_{r\alpha,\gamma}(u)|.
\]
Here \(r\) is one fine residue in a centred block, \(\alpha\)
is an output row type, \(\gamma\) is an input row type, and
\(u\) is a displacement in coarse units. For a coarse-to-coarse
kernel omit \(r\). The maximum over fine residues is essential:
this norm does not sum over all fine output sites.

\subsection{Which principal rows are being calibrated}

At smoothing index \(j\ge4\), take \(J=j\) in
\eqref{eq:convolution-weights}. A finite Gaussian history starts
with the Wilson quadratic form and appends odd block sides;
at each step its smoothing index stays fixed or decreases by
one. The repeated averaging and the noise covariance \(a^{-1}\)
determine its Gaussian Schur complement. Appendix~\ref{app:calibration}
constructs this history directly, including the zero momentum.
It produces positive curvature metrics \(\widetilde H_v\) and
row factors \(S_v\) such that
\begin{equation}
 S_v^{\mathsf T}S_v=\widetilde H_v,\qquad
 K_v=d_1^{\mathsf T}\widetilde H_vd_1.
 \label{eq:calibrated-row-metric}
\end{equation}
Their exponential coefficient bounds depend on the final index
and its fixed profile prices, never on history depth or period.
At the empty history \(\widetilde H_v=S_v=1\).
At a nonempty history we use
\(S_v=(\sqrt{\lambda_j}\,1,R_v)^{\mathsf T}\), where
\(R_v^{\mathsf T}R_v=\widetilde H_v-\lambda_j1>0\).
The positive \(\lambda_j\) is chosen before all block sides.
The additional positive-extension price is denoted by \(l_j\)
and is defined in \eqref{eq:calibration-positive-extension}.

These are precisely rows of the form
\eqref{eq:stationary-row-budget}: transport a terminal curvature
along a fixed coordinate path for each cover coefficient of
\(R_v\), then periodise. Finite symmetry copies carry their
square-root averaging weights. The empty history is the literal
Wilson action, with no extra nonlinear row action.
The stationary estimates above allowed a larger class of row
budgets. The calibration below concerns these prescribed Gaussian
profiles; it is not inferred for arbitrary rows from their budget.
Propagation of the complete nonlinear principal action remains
the principal-action matching task.

Collect the energetic rows of \(\varphi\) in a vector
\(\mathbf F(x,y)\): its fine rows are the direct
\(\sqrt\lambda\,\mathcal C(U_p)\) and the principal row sums;
its noise rows are
\(\sqrt a\,\mathcal C(V_bQ(U)_b^{-1})\).
Include the actual square-root symmetry weights. Then exactly
\begin{equation}
          \varphi(x,y)=\tfrac12\|\mathbf F(x,y)\|_2^2
                                      +\mathcal R(x,y).
 \label{eq:calibrated-square-decomposition}
\end{equation}
The remainder contains the cubic and quartic slot polynomials
and the direct and noise defects
\[
 \tfrac\lambda2\{d_G(W)^2-|\mathcal C(W)|_{\g}^2\},
 \qquad
 \tfrac a2\{d_G(W)^2-|\mathcal C(W)|_{\g}^2\}.
\]
For \(W=e^X\) these defects are \(O(|X|^4)\), since
\(\mathcal C(e^X)=X+O(|X|^3)\) and
\(d_G(e^X)^2=|X|^2+O(|X|^4)\).
Thus \(\mathcal R\) has no quadratic term at \((0,0)\).
On a fixed smaller group chart we also have
\begin{equation}
                        d_G(W)\le2|\mathcal C(W)|_{\g}.
 \label{eq:calibrated-curvature-chord}
\end{equation}
This follows by dividing both expansions by \(|X|\), using
\(|\mathcal C(e^X)|\ge |X|-C|X|^3\), and shrinking the chart.

Write \(\mathbf C_{\rm c}(V)\) for the analogous direct and
principal coarse curvature rows, without noise. Their path budget
gives
\begin{equation}
 \|\mathbf C_{\rm c}(V)\|_\infty\le C_{\rm c}h_{\rm c}
             \quad\text{on coarse good fields}.
 \label{eq:calibrated-coarse-row-price}
\end{equation}
Set \(\mathsf A=D_y\mathbf F(0,0)\) and
\(\mathsf B=D_x\mathbf F(0,0)\). In these actual tree
coordinates \(H_0=\mathsf A^{\mathsf T}\mathsf A\).

\begin{theorem}[Calibration uniformly over Gaussian histories]
\label{thm:calibrated-row-data}
For fixed adjacent smoothing indices and positive profile prices,
there are a free-coordinate lift \(\mathcal H\) and a map
\(\Phi_0\) from coarse curvature rows to fine and noise rows such that
\begin{equation}
 \mathsf B+\mathsf A\mathcal H
       =\Phi_0D\mathbf C_{\rm c}(1),\qquad
                         \mathsf A^{\mathsf T}\Phi_0=0.
 \label{eq:calibrated-two-identities}
\end{equation}
There are \(A_\Phi,\sigma_\Phi>0\), chosen before the odd side
\(L\ge3\), for which
\begin{equation}
 |(\Phi_0)_{\rm f}|_{\sigma_\Phi}\le A_\Phi L^{-2},
 \qquad |(\Phi_0)_{\rm n}|_{\sigma_\Phi}\le A_\Phi.
 \label{eq:calibrated-preblock-prices}
\end{equation}
These bounds hold before periodisation, uniformly in the finite
fine history and every compatible period. The lift
\(\mathcal H\) has an exponential coefficient bound allowed to
depend on \(L\).
\end{theorem}
\begin{proof}
Appendix~\ref{app:calibration} constructs real columns \(E,T\),
for the fine and coarse Gaussian metrics, satisfying
\begin{equation}
 K_fE=Q^{\mathsf T}K_c,\quad QE=1-a^{-1}K_c,\quad
 T^{\mathsf T}d_f=d_cQ,\quad
 T\widetilde H_cd_c=\widetilde H_fd_fE,
 \label{eq:calibrated-installed-columns}
\end{equation}
and \(|T|_\sigma\le CL^{-2}\), with the constant chosen before
\(L\). Section~\ref{sec:calibration-installation} calculates
the counting transpose and the installed factors on one Fourier
mode, including fine differences across residue boundaries.
Here \(Q\) in this linear calculation is the scalar
line-and-block average, and \(d_f,d_c\) are the two curls.
It is not yet the derivative of the framed predictor.

We compute that distinction explicitly. For a fine tangent field
\(u\), put
\[
 t_z=-\frac1{4!s}\sum_{z',\pi}
                          \int_{\gamma_{z,z'}^\pi}u,
 \qquad (Pu)_x=t_0^x-\sum_z w(z)t_z^x.
\]
The integral is the signed sum along the coordinate-ordered path;
the frame patch and its cardinality \(s\) are those already fixed.
Differentiating the rooted regular frame gives \(t_z-t_0\).
Differentiating its two endpoints in the averaged connector gives
the actual predictor derivative
\[
                         Q_0=Q+d_{0,c}P.
\]
For a fine site array \(\omega\), path telescoping gives
\(Qd_{0,f}\omega=d_{0,c}A_0\omega\) and
\(Pd_{0,f}\omega=\omega|_{\rm roots}-A_0\omega\), where
\(A_0\omega=\sum_z w(z)\omega(Lx+z)\).
Consequently
\begin{equation}
 Q_0d_{0,f}\omega=d_{0,c}(\omega|_{\rm roots}),
                         \qquad d_cQ_0=d_cQ.
 \label{eq:calibrated-frame-correction}
\end{equation}

Let \(\omega(u)(z)\) be the signed sum of \(u\) on the fixed
tree path from \(z\) to its coarse root, minus \((Pu)\) there.
Then \(\widehat E=E+d_{0,f}\omega(E)\) vanishes on tree links,
\(d_f\widehat E=d_fE\), and
\(Q_0\widehat E=QE\). The last assertion uses
\(\omega(E)|_{\rm roots}=-PE\).
Also \(K_f\widehat E=Q_0^{\mathsf T}K_c\), because curls annihilate
gradients and \(d_{0,c}^{\mathsf T}K_c=0\).
Let \(\mathcal R_{\rm rep}\) be the derivative of the replica
map, and \(\iota\) include free coordinates in all fine links.
Define \(\mathcal H\) by
\(\mathcal R_{\rm rep}+\iota\mathcal H=\widehat E\).
This is possible and unique since both fields vanish on trees.

Now set
\begin{equation}
 \Phi_0=\begin{pmatrix}
    S_f\widetilde H_f^{-1}TS_c^{\mathsf T}\\[2pt]
    a^{-1/2}d_c^{\mathsf T}S_c^{\mathsf T}
 \end{pmatrix}.
 \label{eq:calibrated-explicit-row-lift}
\end{equation}
Indeed
\[
 \mathsf A=\begin{pmatrix}S_fd_f\iota\\-\sqrt a Q_0\iota\end{pmatrix},
 \quad
 \mathsf B=\begin{pmatrix}S_fd_f\mathcal R_{\rm rep}\\
              \sqrt a(1-Q_0\mathcal R_{\rm rep})\end{pmatrix},
 \quad D\mathbf C_{\rm c}(1)=S_cd_c.
\]
The first two identities in \eqref{eq:calibrated-installed-columns}
give matching in \eqref{eq:calibrated-two-identities}.
For orthogonality multiply out:
\[
 \mathsf A^{\mathsf T}\Phi_0
   =\iota^{\mathsf T}(d_f^{\mathsf T}T-Q_0^{\mathsf T}d_c^{\mathsf T})S_c^{\mathsf T}=0,
\]
using the third identity and \((Q_0-Q)^{\mathsf T}d_c^{\mathsf T}=0\).

All coarse factors of \eqref{eq:calibrated-explicit-row-lift}
have pre-block exponential prices. A displacement \(z\) in
the fine factor \(S_f\widetilde H_f^{-1}\) shifts the output
residue of \(T\) and its coarse base by at most
\(1+|z|_1/L\). Choose \(\sigma_\Phi\) below half the fine
coefficient width. Its convolution cost is then bounded by
the absolutely summable fine coefficient budget times
\(e^{\sigma_\Phi(1+|z|_1)}\).
This preserves \(L^{-2}\), without summing over output residues.
The noise factor is a coarse convolution of fixed cell types.
This proves \eqref{eq:calibrated-preblock-prices}.
The finite tree-path conversion gives the stated bound for
\(\mathcal H\).
\end{proof}

The construction and its quantifiers correspond to
\cite[Sections~20--22, pp.~100--114]{Main}.
For a physical background, replace each cover coefficient of
\(\Phi_0\) by that coefficient times adjoint transport along
its fixed coarse coordinate path, before periodising. Denote
the resulting row map by \(\Phi(V)\). Adjoint transport is
isometric, so the two absolute prices are unchanged.
On \(\|x\|_\infty\le2w\), a path of displacement \(u\)
changes its adjoint transport by at most
\(C_L(1+|u|)w\). The exponential reserve therefore gives
\begin{equation}
 |\Phi(V(x))-\Phi_0|_\theta\le C_Lw.
 \label{eq:calibrated-covariant-perturbation}
\end{equation}
Here and below rectangular weighted Schur norms include both
row and column sums, with the appropriate attached homes.
Columns can cost the fixed number of fine residues, depending
on \(L\); they cost no total volume.

Multiplying matching by \(\mathsf A^{\mathsf T}\) also gives
\(\mathcal H=-H_0^{-1}\mathsf A^{\mathsf T}\mathsf B\).
Thus the calibrated lift is the minimiser of the actual flat
tree-coordinate quadratic form. More precisely,
\[
 \tfrac12\|\mathsf Bx+\mathsf Ay\|_2^2
 =\tfrac12\|\Phi_0D\mathbf C_{\rm c}(1)x\|_2^2
          +\tfrac12\|\mathsf A(y-\mathcal Hx)\|_2^2.
\]
This identity explains why both matching and orthogonality,
rather than only a bounded lift, are needed in the next step.

\section{The central comparison at a general good background}
\label{sec:general-central-comparison}

We now use calibration to construct one fine configuration with
a strict curvature reserve. We then compare it directly with the
stationary centre. Throughout this section a general good background
means a coarse good field \(V(x_0)\) anywhere in the auxiliary
chart \(\|x_0\|_\infty\le w\). No restriction
\(\|x_0\|_\infty=O(h)\) is imposed, and no global logarithm on
an arbitrary physical torus is asserted.

Let the coarse nominal coupling be \(t\), and retain the prescribed
thresholds
\begin{equation}
 \widehat\beta=t+b_0\log L,\quad
 h_{\rm c}=t^{-1/2}P(t),\quad
 h=\widehat\beta^{-1/2}P(\widehat\beta),\quad
 w=\widehat\beta^{-7/16}.
 \label{eq:general-central-thresholds}
\end{equation}
At fixed \(L\), logarithmic differentiation of
\(s^{-1/2}P(s)\) gives
\[
 -\frac1{2s}+\frac{2\log\log s}{s\log s}.
\]
Integrating over the fixed interval of length \(b_0\log L\)
shows \(h_{\rm c}/h\to1\). In particular, eventually
\begin{equation}
 h_{\rm c}\le2h,\qquad h/w\to0,\qquad
       w^2/h=\widehat\beta^{-3/8}/P(\widehat\beta)\to0.
 \label{eq:general-central-scale-reserves}
\end{equation}

\subsection{An affine centre with smaller fine curvature}

For a good \(V=V(x)\), put
\[
 \mathbf f(V)=\Phi(V)\mathbf C_{\rm c}(V),\qquad
 \mathbf A(V)=D_y\mathbf F(x,0),\qquad
 \mathbf K(V)=\mathbf A(V)^{\mathsf T}\mathbf A(V).
\]
These are the guide row vector, the actual row derivative at the
replica, and its precision. They are distinct from the stationary
precision \(H(x)\). With \(A=C_{\rm c}A_\Phi\),
\begin{equation}
 \|\mathbf f_{\rm f}(V)\|_\infty\le AL^{-2}h_{\rm c},
 \qquad \|\mathbf f_{\rm n}(V)\|_\infty\le Ah_{\rm c}.
 \label{eq:general-central-guide-prices}
\end{equation}
The price \(A\) is chosen before \(L\).
Define
\begin{equation}
 m_{\mathrm{aff}}(V)=-\mathbf K(V)^{-1}\mathbf A(V)^{\mathsf T}
                               \{\mathbf F(x,0)-\mathbf f(V)\}.
 \label{eq:general-central-affine-centre}
\end{equation}

\begin{lemma}[The affine residual in the actual coordinates]
\label{lem:general-central-affine-residual}
For sufficiently small \(w\), uniformly in history and period
at the fixed profile and block,
\begin{gather}
 \mathbf K(V)\ge\eta_0/2,\qquad
 |\mathbf K(V)^{-1}|_{\theta'}\le C_L,\qquad
 \|\mathbf F(x,0)\|_\infty+\|m_{\mathrm{aff}}(V)\|_\infty\le C_Lh_{\rm c},
 \label{eq:general-central-affine-size}\\
 \|\mathbf F(x,0)-\mathbf f(V)+\mathbf A(V)m_{\mathrm{aff}}(V)\|_\infty
                                                   \le C_Lw^2.
 \label{eq:general-central-affine-residual}
\end{gather}
\end{lemma}
\begin{proof}
The local row derivative budgets imply
\(|\mathbf A(V)-\mathsf A|_\theta\le C_Lw\).
Hence \(|\mathbf K(V)-H_0|_\theta\le C_Lw\), proving the
gap and inverse bound by Lemma~\ref{lem:uniform-inverse}.
At the replica every fine plaquette is either trivial or one
coarse plaquette: crossing just one coarse boundary gives two
cancelling factors, while crossing both gives its four boundary
factors. Its fine rows are therefore \(O(h_{\rm c})\).
In the predictor patch, the path comparison and frame-score
estimate used in Lemma~\ref{lem:actual-analytic-branch} compare
all replica connectors with the root-to-root transport \(V_b\)
with error \(C_Lh_{\rm c}\). Their nearest-group projection has
the same bound. This proves the noise-row estimate, and then
the bound for \(m_{\mathrm{aff}}\) follows from its formula.

Use the trial direction \(v=\mathcal Hx\), whose maximum norm
is at most \(C_Lw\). Taylor expansion in the rooted row budgets
and \eqref{eq:calibrated-covariant-perturbation} give
\[
 \mathbf F(x,0)=\mathsf Bx+O_\infty(w^2),\quad
 \mathbf A(V)=\mathsf A+O_\theta(w),\quad
 \mathbf f(V)=\Phi_0D\mathbf C_{\rm c}(1)x+O_\infty(w^2).
\]
Thus matching gives
\(\mathbf F(x,0)-\mathbf f(V)+\mathbf A(V)v=O_\infty(w^2)\).
The operator
\[
 \Pi(V)=1-\mathbf A(V)\mathbf K(V)^{-1}\mathbf A(V)^{\mathsf T}
\]
has bounded maximum-norm operator price and annihilates
\(\mathbf A(V)\). Applying it to this last identity proves
\eqref{eq:general-central-affine-residual} by the definition of
\(m_{\mathrm{aff}}\). No Euclidean row norm proportional to volume is used.
\end{proof}

Let \(\lambda_*,a_*>0\) be the smallest direct fine and noise
row-square prices in the chosen symmetry-copy convention.
They differ from \(\lambda,a\) only by fixed finite weights.

\begin{proposition}[Curvature reserve and comparison with stationarity]
\label{prop:general-central-good-centre}
After fixing the pre-block row prices, choose an allowed odd \(L\)
and then a fixed \(W>1\) with
\[
       \frac{A}{\sqrt{\lambda_*}L^2}\ll1,\qquad
       \frac{A}{\sqrt{a_*}W}\ll1.
\]
For every good auxiliary background, eventually
\begin{gather}
 \sup_p d_G(U(x,m_{\mathrm{aff}}(V))_p)\le h/8,\qquad
 \sup_b d_G(V_bQ(U(x,m_{\mathrm{aff}}(V)))_b^{-1})\le Wh/2,
 \label{eq:general-central-affine-reserve}\\
 \|\varphi_y(x,m_{\mathrm{aff}}(V))\|_\infty\le C_Lw^2,\qquad
                         \|m_{\mathrm{aff}}(V)-m(x)\|_\infty\le C_Lw^2.
 \label{eq:general-central-stationary-residual}
\end{gather}
\end{proposition}
\begin{proof}
The affine residual and the second row derivatives give
\[
 \mathbf F(x,m_{\mathrm{aff}})=\mathbf f(V)+O_\infty(w^2+h_{\rm c}^2)
                         =\mathbf f(V)+O_\infty(w^2).
\]
First, the replica plaquette and mismatch bounds, together with
\(\|m_{\mathrm{aff}}\|_\infty\le C_Lh_{\rm c}\), put both group elements
in the fixed small chord domain of
\eqref{eq:calibrated-curvature-chord}. We may therefore use
their direct rows to obtain
\begin{align*}
 \sup_p d_G(U(x,m_{\mathrm{aff}})_p)
      &\le2\lambda_*^{-1/2}(AL^{-2}h_{\rm c}+C_Lw^2),\\
 \sup_b d_G(V_bQ(U(x,m_{\mathrm{aff}}))_b^{-1})
      &\le2a_*^{-1/2}(Ah_{\rm c}+C_Lw^2).
\end{align*}
Since \(h_{\rm c}\le2h\), the displayed bounds are at most
\[
 \left\{\frac{4A}{\sqrt{\lambda_*}L^2}
                  +\frac{2C_Lw^2}{\sqrt{\lambda_*}h}\right\}h,
 \qquad
 \left\{\frac{4A}{\sqrt{a_*}}
                  +\frac{2C_Lw^2}{\sqrt{a_*}h}\right\}h,
\]
respectively. To obtain the reserves \(h/8\) and \(Wh/2\), it is
enough to arrange
\[
 \frac{4A}{\sqrt{\lambda_*}L^2}\le\frac1{16},\qquad
 \frac{4A}{\sqrt{a_*}}\le\frac W4,\qquad
 \frac{2C_Lw^2}{\sqrt{\lambda_*}h}\le\frac1{16},\qquad
 \frac{2C_Lw^2}{\sqrt{a_*}h}\le\frac W4.
\]
The order of choices is now explicit: first fix \(A,\lambda_*,a_*\)
from the pre-block estimates, choose \(L\) for the first request
and \(W\) for the second, and only then fix the coupling threshold.
The two smallness conditions in the proposition are taken to
give those numerical bounds. At these fixed \(L,W\),
\(C_Lw^2/h\to0\), so one entry threshold gives both remaining
requests. This proves \eqref{eq:general-central-affine-reserve}.

At this field all curvatures and mismatches are \(O(h)\).
Differentiating the cubic and quartic path terms yields
\(\|\mathcal R_y(x,m_{\mathrm{aff}})\|_\infty\le C_Lh^2\); the direct
quartic defects even give \(O(h^3)\). The absolute derivative
sum converges by the native word reserve. Moreover
\(\mathbf F_y(x,m_{\mathrm{aff}})=\mathbf A(V)+O_\theta(h_{\rm c})\), and
flat orthogonality gives
\[
 \|\mathbf A(V)^{\mathsf T}\mathbf f(V)\|_\infty
   =\|\mathbf A(V)^{\mathsf T}\Phi(V)\mathbf C_{\rm c}(V)\|_\infty
                                                    \le C_Lwh_{\rm c}.
\]
Insert these estimates into
\(\varphi_y=\mathbf F_y^{\mathsf T}\mathbf F+\mathcal R_y\).
The residual is \(O(w^2+wh_{\rm c}+h_{\rm c}^2+h^2)=O(w^2)\).

Both \(m_{\mathrm{aff}}\) and \(m(x)\) lie in the stationary uniqueness
chart. Their gradient difference equals
\[
 \left[\int_0^1\varphi_{yy}(x,m(x)+s(m_{\mathrm{aff}}-m(x)))\,\dd s\right]
                                                   (m_{\mathrm{aff}}-m(x)).
\]
The bracketed matrix has the common positive gap and weighted
Schur bound, hence a volume-uniform maximum-norm inverse.
Since \(\varphi_y(x,m(x))=0\), this proves the last bound.
\end{proof}

This is the affine calculation in
\cite[Lemma~17.2 and Proposition~17.3, pp.~92--93]{Main}.
Its implication for the central ball can be taken directly,
without comparing two different Gaussian normalisations.

\subsection{Removing the physical fine indicator on its proved domain}

\begin{theorem}[General-background central comparison]
\label{thm:general-central-comparison}
Take the calibrated history profiles and choose constants in this
order: their pre-block prices, then \(L,W\), then a sufficiently
small fixed \(d_0>0\), then the coupling threshold.
Assume the coupling band, subpower loop allowance and auxiliary
size bounds of Theorem~\ref{thm:tail-auxiliary}.
For every coarse good \(V(x_0)\) with
\(\|x_0\|_\infty\le w\), every field
\[
 U(x,m(x)+gz),\qquad |z|_2\le D=d_0P(\widehat\beta),
       \quad\|x-x_0\|_\infty\le\widehat\beta^{-1},
\]
is fine good, eventually. At these nearby \(x\),
\begin{equation}
 0\le(\mathcal P_\beta F_{{\rm sf},\beta})(V(x))
              -N_g^{\mathrm{phys}}(x)Z_D(g,x)\le N_g^{\mathrm{phys}}(x)e^{-cD^2}.
 \label{eq:general-central-exact-comparison}
\end{equation}
All constants are independent of history depth and the auxiliary
period at the fixed profile. The assertion includes backgrounds
\(x\) that need not themselves be coarse good.
\end{theorem}
\begin{proof}
At \(x_0\), compare the stationary centre with the affine one.
A plaquette has four factors, so
\[
 \sup_p d_G(U(x_0,m(x_0))_p)
       \le h/8+4C_Lw^2\le h/4
\]
eventually. Comparing a general displayed field with this centre
by the same unitary telescoping gives the stronger explicit bound
\begin{equation}
 \sup_p d_G(U(x,m(x)+gz)_p)
  \le h/8+C_L(w^2+\widehat\beta^{-1})+4gD.
 \label{eq:general-central-direct-inclusion}
\end{equation}
Here \(\|m(x)-m(x_0)\|_\infty
\lesssim\widehat\beta^{-1}\); the fixed-profile constant
\(C_L\) is enlarged to include this Lipschitz bound and
the four plaquette factors.
Now \(gD/h\le\sqrt2d_0\),
\(w^2/h\to0\), and \(\widehat\beta^{-1}/h\to0\).
Choose \(4\sqrt2d_0\le1/8\) and increase the threshold.
The right side is at most \(h/2\), hence certainly at most
\(h\). The noise mismatch has a corresponding fixed-stencil
reserve, after decreasing \(d_0\) in terms of the predictor
Lipschitz constant and \(W\). All branches and active chart
cutoffs thus have their regular values on this ball.

The nearby coarse coordinates satisfy \(\|x\|_\infty\le2w\)
eventually. The proof of the auxiliary tail theorem applies on
this fixed enlargement: every occurrence of \(w\) is multiplied
by at most a fixed constant, and its powers and dimension prices
are unchanged. In \eqref{eq:tail-exact-masked-central} the fine
indicator is now identically one. The difference in
\eqref{eq:general-central-exact-comparison} is exactly its positive
physical tail. Apply that theorem.
\end{proof}

The conclusion is the comparison of
\cite[Theorem~19.4 and Corollary~19.5, p.~99]{Main}.
The direct estimate \eqref{eq:general-central-direct-inclusion}
replaces the separate intact-column lower-bound route there.
On the stated auxiliary sizes, the proof of
Lemma~\ref{lem:connected-central-comparison} also gives
\(Z_D\ge1/2\) eventually, uniformly on the enlarged chart.
Consequently
\[
 0\le\log\frac{(\mathcal P_\beta F_{{\rm sf},\beta})(V(x))}
                         {N_g^{\mathrm{phys}}(x)Z_D(g,x)}\le2e^{-cD^2}.
\]
Any restricted positive column whose retained tests equal one on
the ball has the same lower bound \(N_g^{\mathrm{phys}} Z_D\ge N_g^{\mathrm{phys}}/2\).
This last statement is only for such a verified column; it does
not delete branch factors or incoming multipliers from a different
physical sector.

There are two distinct uniformity statements. The construction of
\(\Phi_0\) and its fine gain are uniform over all finite histories
before choosing \(L\). The comparison of total central integrals
uses the auxiliary size bound. On an unrestricted volume the
reference expansion is a local family statement, while coordinate
failures and original incoming factors still require exact
physical reassembly. We turn to that representation next.


\section{Finite reads for the exact physical integral}
\label{sec:physical-selected-Gaussian}

The square root \(W\) is appropriate for the analytic connected
logarithm. A bounded Borel incoming function, however, cannot be
Taylor expanded, and composing it with the dense map \(W\xi\)
does not preserve its finite reads. We therefore keep physical
fluctuation coordinates and change their joint Gaussian density.
The correction to that density will remain in the integrand.
This is the physical fluctuation organisation of
\cite[Sections~27--29, pp.~143--166]{Main}; all matrix identities
needed here are given below.

\subsection{Sectors, holes and operated rows}

A sector \(\sigma\) records the outcomes of the local tests used
to partition the original integral. Its Borel weight satisfies
\(0\le\omega_\sigma(V,U)\le1\) and
\(\sum_\sigma\omega_\sigma(V,U)=1\). Products of a test and
its complement provide such a partition. Every original factor
is kept, including when we enlarge the frozen neighbourhood.
Let \(\mathcal H_\sigma\) contain the coarse bad flags and the
other local marks of the sector. Write \(H_r(X)\) for the
periodic coarse \(r\)-neighbourhood of \(X\), and \(\pi(e)\)
for the coarse parents of the endpoints of a fine link. Set
\[
 I_\sigma=\{e\in I:\pi(e)\cap H_{10B}(\mathcal H_\sigma)=\varnothing\},
 \qquad J_\sigma=I\setminus I_\sigma.
\]
The free links in \(J_\sigma\) keep their original Haar values
\(Y_e\). Active links have coordinates \(a_e\in\g\), with
\begin{equation}
 \mathcal U_\sigma(V,Y,a)_e=
 \begin{cases}
  e^{a_e}R_e(V),&e\in I_\sigma,\\
  Y_e,&e\in J_\sigma,\\
  1,&e\in T.
 \end{cases}
 \label{eq:physical-selected-installation}
\end{equation}
The chart partition of Section~\ref{sec:sectors} is refined if
necessary so that every active link is in the fixed injective
exponential chart. Links with a chart-complement factor are
frozen. All factors of that original partition remain in
\(\omega_\sigma\). In particular this convention covers the
case in which every link is frozen.
For the physical packet partition of
Section~\ref{sec:joint-hole-packets}, the actual fine and noise
tests already imply this active chart condition. There the frozen
set is determined by physical seeds; chart choices on its links
do not introduce additional geometric holes.

Choose an orthonormal basis of \(\g\), and attach the coarse
parent to each real coordinate in
\(\mathcal I=I\times\{1,\ldots,\dim_\R\g\}\).
Matrices below use these counting coordinates, with inactive
coordinates retained as an auxiliary identity block.
Put \(R_B=\lfloor B/100000\rfloor\) and let \(R_{\rm pred}\)
be the fixed predictor reach. Increase the profile threshold so
that, for \(r_B=R_B+R_{\rm pred}+2\),
\(4r_B+10<B/4\).

Take the direct curvature rows, principal rows cut at original
native length \(R_B\), and the actual flat framed-predictor noise
rows. To define their linear maps on a background \(V\),
transport each input colour to the row anchor along its declared
coordinate path and apply its flat row coefficients. Reversed
edges have their orientation sign and adjoint transport. Denote
these finite maps by \(d_r(V)\). Only rows anchored outside
\(H_{8B}(\mathcal H_\sigma)\) are operated. Define
\[
 D_{re}=1_{\{r\text{ operated}\}}1_{\{e\in I_\sigma\}}(d_r(V))_e.
\]
The potential stencil belongs to the read set even if a numerical
coefficient vanishes. The noise row here uses \(Q_0\), including
\eqref{eq:calibrated-frame-correction}.

For definiteness the affine source is also chosen locally.
On frozen links use a smooth bounded extension \(s_e(V,Y)\)
of \(\log(Y_eR_e(V)^{-1})\). More explicitly, put \(u=Bh\)
and multiply this logarithm by a smooth function of its squared
norm divided by \(u^2\), equal to one up to one and zero from
four onwards. Extend by zero off the fixed injective chart;
eventually \(2u\) lies strictly inside it. Thus the extension
equals the true logarithm up to radius \(u\), has norm at most
\(2u\), and is smooth in the coarse field even at its cutoff.
It is radial and hence gauge covariant. Put \(s=0\) on active
and tree links.
On an operated row let \(\alpha_r\) be its actual cut row at
the replica, minus its cut calibrated coarse guide, plus its
full transported flat row applied to \(s\). All paths in this
source, including the guide, have the same native cut; original
coarse tests are retained. In the physical packet representation
the guide is precisely the cut transported thinned vector
defined in Section~\ref{sec:packet-row-comparison}; on healthy
reads it is the calibrated coarse curvature guide. On other rows put \(\alpha_r=0\).
This auxiliary extension never changes the actual link \(Y_e\)
in \eqref{eq:physical-selected-installation}. Now set
\begin{equation}
 b=D^{\mathsf T}\alpha,\qquad \mathsf M=D^{\mathsf T}D+P_{\rm in},
 \label{eq:physical-selected-precision}
\end{equation}
where \(P_{\rm in}\) projects onto inactive coordinates.
The precision depends on \(V\) and sector geometry, not on
\(Y\) or branch tags. The source may depend on both.

\begin{lemma}[A gap unaffected by holes]
\label{lem:physical-selected-gap}
For a fixed block and path profile there are \(\eta,C,\theta>0\)
such that eventually, on periods \(c\ge100(1+B)\),
\[
                 \eta1\le\mathsf M\le C1,\qquad
                                  |\mathsf M|_\theta\le C.
\]
These constants are independent of holes, history depth and period.
\end{lemma}
\begin{proof}
Absolute row and column sums of the cut row matrix are bounded
by the original exponential path budgets. Products give the
weighted Schur upper bound and hence the operator upper bound.
For the lower bound use the local calculation in
\eqref{eq:link-reconstruction-bound}. For any free link \(e\)
it bounds \(|v_e|\), for a tree-zero array, by a fixed finite
sum of direct curls and actual flat predictor rows in a fixed
patch. Its radius and coefficients depend only on the block.
For active \(e\), all of this patch is outside the \(8B\)
neighbourhood of holes when \(B\) is large. Its required rows
have not been deleted. The principal row cut is irrelevant to
this argument, which uses only direct and noise rows.

In a local axial gauge the healthy coarse field on that patch
has link logarithms at most \(C_Lh_{\rm c}\). The local gauge
comparison of Section~\ref{sec:local-physical-bridge} and the
finite transported row formula therefore give, for arrays
supported on active coordinates,
\[
 |v_e|\le C_*\max_{r\text{ near }e}|(Dv)_r|
                 +C_*'h_{\rm c}\max_{f\text{ near }e}|v_f|.
\]
Square, replace each maximum squared by the sum of squares,
and sum over \(e\). A row or coordinate occurs in a bounded
number \(\nu_r,\nu_e\) of these fixed patches, so
\[
 \|v\|_2^2\le2C_*^2\nu_r\|Dv\|_2^2
                     +2(C_*')^2\nu_e h_{\rm c}^2\|v\|_2^2.
\]
Absorb the last term when it is at most \(\|v\|_2^2/2\).
For a general array, the active and inactive blocks are
orthogonal and its quadratic form is
\(\|Dv_{\rm ac}\|_2^2+\|v_{\rm in}\|_2^2\).
This proves the common gap, including the empty active set.
\end{proof}

\subsection{A finite-range covariance from principal inverse columns}

For each coordinate \(j\) choose the fixed window \(\Omega_j\)
consisting of all colours on links whose endpoint parents meet
the coarse \(\lfloor B/4\rfloor\)-ball about \(j\).
It contains a ball of that radius and is contained in the ball
of radius \(B/4+2\). Define
\begin{equation}
 \begin{gathered}
 J_{ij}=1_{\{i\in\Omega_j\}}(\mathsf M_{\Omega_j,\Omega_j}^{-1})_{ij},
 \qquad E=\mathsf MJ-1,\\
 C_{\rm sel}=(J+J^{\mathsf T})/2,\qquad
 S=\mathsf MC_{\rm sel}-1,\qquad \mu=-J^{\mathsf T}b.
 \end{gathered}
 \label{eq:physical-selected-windows}
\end{equation}
The principal matrix retains all operated rows:
\(\mathsf M_{\Omega,\Omega}=D_{*,\Omega}^{\mathsf T}D_{*,\Omega}
+(P_{\rm in})_{\Omega,\Omega}\). Restricting row anchors to
\(\Omega\) would give a different matrix.

\begin{proposition}[One positive finite-range Gaussian]
\label{prop:physical-selected-law}
At a smaller fixed width,
\begin{gather}
 |J|_{\theta'}+|C_{\rm sel}|_{\theta'}\le C,\qquad
 |E|_{\theta'}+|J-\mathsf M^{-1}|_{\theta'}
     +|C_{\rm sel}-\mathsf M^{-1}|_{\theta'}+|S|_{\theta'}
                                                 \le Ce^{-cB},
 \label{eq:physical-selected-small-errors}\\
 C_{\rm sel}\simeq1,\qquad C_{\rm sel}^{-1}\simeq1.
 \label{eq:physical-selected-precision-reserve}
\end{gather}
The covariance is zero between coordinates at distance
greater than \(B/4+2\). It defines one normalised joint
Gaussian, denoted by \(\mu_{C_{\rm sel}}\).
\end{proposition}
\begin{proof}
Principal matrices retain the same gap and weighted row bound.
The weighted inverse proof of Lemma~\ref{lem:uniform-inverse}
therefore gives the common envelope \(Ce^{-\nu d(i,j)}\)
for their inverses and the full inverse \(G=\mathsf M^{-1}\).
Inside \(\Omega=\Omega_j\), the exact comparison is
\[
 J_{ij}-G_{ij}=
       (\mathsf M_{\Omega,\Omega}^{-1}
                         \mathsf M_{\Omega,\Omega^c}G_{\Omega^c,j})_i.
\]
Every path in this product visits a coordinate outside
\(\Omega_j\), at distance at least \(\lfloor B/4\rfloor\)
from \(j\). Its total length also bounds \(d(i,j)\).
Split the available decay width between these two bounds and
the summation over intermediate coordinates. Four-dimensional
shell sums give
\(|J_{ij}-G_{ij}|\le Ce^{-cB}e^{-\nu' d(i,j)}\).
Outside the window the same estimate follows from \(J_{ij}=0\)
and the full inverse envelope. This pointwise estimate sums in
either index, even though the window varies with \(j\).
It proves the weighted error without a number-of-windows factor.

For \(i\in\Omega_j\), \((\mathsf MJ)_{ij}=\delta_{ij}\)
exactly. For the norms use
\(E=\mathsf M(J-G)\), symmetrise, and use
\(S=\mathsf M(C_{\rm sel}-G)\).
Since \(G\ge C^{-1}1\), the operator error less than
\((2C)^{-1}\) gives the positive covariance bound.
The upper bound gives its precision reserve. Its finite range
follows directly from the two selected windows.
\end{proof}

The same bounds hold uniformly for
\(\mathsf M_t=(1-t)1+t\mathsf M\), \(0\le t\le1\).
On a fixed healthy germ, derivatives keep the sector geometry
and windows fixed. The identity
\(D(\mathsf M_\Omega^{-1})=-\mathsf M_\Omega^{-1}
(D\mathsf M_\Omega)\mathsf M_\Omega^{-1}\)
and its iterated product rule give every fixed-order derivative.
The inverse comparison above still crosses the window boundary
in every differentiated error term. Summing its connected chains
at a smaller width preserves \(C_qe^{-cB}\).
These are derivatives of the unmasked matrix formulas, not of
sector indicators.

An inverse column reads the entries of its entire principal
matrix, including row stencils crossing the window. In
\(\mu_j=-\sum_{i\in\Omega_j}J_{ij}b_i\), each source entry
has finite reads as well. Thus \(\mu_j\) reads only a fixed
\(C B\)-neighbourhood of \(j\), in \(V,Y\), tags and sector
queries. A bounded Borel function of finitely many physical
links, evaluated at \(a=\mu+gz\), has finite reads in those
same variables and the physical Gaussian coordinates \(z_e\).
There is no dense square root acting on them.

The original exponentially relaxed retained-row test also has
finite effective reach. Written with the present notation it is
\[
 k_{t,\varrho}(x;U)=
 1_{\{d_G(U_p)\le t e^{d_\infty(x,o(p))/\varrho}\text{ for all }p\}}.
\]
Here \(\varrho>0\) is its fixed relaxation length, and \(o(p)\)
is the plaquette base point. Since \(d_G(U_p)\le2\), queries
at distance at least \(\varrho\log(2/t)\) are automatic.
The boundary-link reads lie in radius
\(\lceil\varrho\log(2/t)\rceil+2\). For \(t=h,h_{\rm c}\)
this is \(O(1+\log\widehat\beta)\), hence less than \(B/100\)
eventually at fixed profile. This uses the original inequalities,
including equality at the threshold.

All constructions are Borel: path operations are continuous,
the tests are Borel, and positive matrix inversion is continuous.
They also respect gauge transformations. If \(O\) acts by the
adjoint at each coordinate's parent, the row transport convention
gives \(D^u=O_{\rm row}DO^{\mathsf T}\), \(\alpha^u=O_{\rm row}\alpha\).
Windows include all colours, so
\(\mathsf M^u=O\mathsf MO^{\mathsf T}\), \(J^u=OJO^{\mathsf T}\),
\(C_{\rm sel}^u=OC_{\rm sel}O^{\mathsf T}\), and \(\mu^u=O\mu\).
The same Gaussian and frozen Haar measures are consequently
carried into their transformed measures; contracted component
integrals are gauge invariant.

\subsection{The exact change of density and its normalisation}

Let \(n_{\rm ac}=|I_\sigma|\), \(n_{\rm fr}=|J_\sigma|\),
\(n_0=|I|\), and put
\[
 s_g=\zeta_G(\sqrt{2\pi}g)^{\dim_\R\g},\qquad
 r=\mathsf M\mu+b.
\]
Use the original action \(A_t\) of \eqref{eq:general-action},
with all its prescribed clipping and masks. Define, without
estimating or changing that action,
\begin{align}
 \mathcal D_\sigma(V,Y,a)
       &=A_t(V,\mathcal U_\sigma(V,Y,a))
                                      -\tfrac12\|Da+\alpha\|_2^2,
 \label{eq:physical-selected-exact-phase}\\
 \mathcal Q_\sigma(z)
       &=-g^{-1}\langle r,z\rangle
               +\tfrac12\langle z,(C_{\rm sel}^{-1}-\mathsf M)z\rangle.
 \label{eq:physical-selected-density-correction}
\end{align}
All uncut tails and all differences from the auxiliary affine
rows remain in \(\mathcal D_\sigma\).
For a chart-supported Borel multiplier \(B_\sigma(a)\),
the exact identity is
\begin{align}
 &\zeta_G^{n_{\rm ac}}\int B_\sigma(a)
                 e^{-A_t(V,\mathcal U_\sigma)/g^2}\,\dd a\notag\\
 &\quad=s_g^{n_{\rm ac}}(\det C_{\rm sel})^{1/2}
      e^{-\|D\mu+\alpha\|_2^2/(2g^2)}
       \mathbb E_{C_{\rm sel}}\!\left[
       B_\sigma(\mu+gz)e^{\mathcal Q_\sigma(z)
                  -\mathcal D_\sigma(V,Y,\mu+gz)/g^2}\right].
 \label{eq:physical-selected-exact-identity}
\end{align}
Indeed expanding \(\|D(\mu+gz)+\alpha\|^2\) gives its constant
term, \(2g\langle r,z\rangle\), and
\(g^2\langle z,\mathsf Mz\rangle\) on the active subspace.
Multiplying the normalised Gaussian density cancels its
\(C_{\rm sel}^{-1}\) term and gives exactly this exponent.
The active Jacobian is
\(\zeta_G^{n_{\rm ac}}(2\pi)^{n_{\rm ac}\dim\g/2}
g^{n_{\rm ac}\dim\g}=s_g^{n_{\rm ac}}\).
Inactive coordinates form an identity block in \(\mathsf M,J\)
and \(C_{\rm sel}\), with \(b=\mu=r=0\) there. Their auxiliary
Gaussian integral is one; their determinant is one. Thus they
introduce no additional power of \(g\).
The identity holds for nonnegative Borel functions by change of
variables, and for absolutely integrable complex functions by
applying it to absolute values and real and imaginary parts.

In the original sector take
\[
 B_\sigma(a)=\omega_\sigma(V,\mathcal U_\sigma)
       \lambda_t(\mathcal U_\sigma)\mathcal B(\mathcal U_\sigma)
       e^{-I_f(\mathcal U_\sigma)}\prod_{e\in I_\sigma}j_G(a_e),
\]
zero off its active chart. The full push is exactly
\begin{align}
 (\mathcal P_\beta F)(V)
  ={}&e^{vf^4}z_{\rm n}(\beta)^{-b_c}
   \sum_{\sigma,t}s_g^{n_{\rm ac}}\int_{G^{J_\sigma}}
       (\det C_{\rm sel})^{1/2}e^{-\|D\mu+\alpha\|^2/(2g^2)}
                                                        \notag\\
 &\hspace{9mm}\cdot\mathbb E_{C_{\rm sel}}
       [B_\sigma(\mu+gz)e^{\mathcal Q_\sigma
                    -\mathcal D_\sigma(V,Y,\mu+gz)/g^2}]
                                      \,\dd h_G^{J_\sigma}(Y).
 \label{eq:physical-selected-full-push}
\end{align}
In particular every incoming factor and every original test is
evaluated on the single physical field
\(\mathcal U_\sigma(V,Y,\mu+gz)\).

For comparison with a common flat factor, let \(\mathsf M_0\)
be the all-active flat precision \(H_0\) of the same profile.
The exact identity
\(\det C_{\rm sel}=\det(1+S)/\det\mathsf M\) shows that
factoring \(s_g^{n_0}(\det\mathsf M_0)^{-1/2}\) leaves
\begin{equation}
 s_g^{-n_{\rm fr}}
  \exp\!\left\{\tfrac12\log\det\mathsf M_0
       -\tfrac12\log\det\mathsf M+\tfrac12\log\det(1+S)\right\}.
 \label{eq:physical-selected-frozen-ratio}
\end{equation}
Since \(|S|_\theta<1\), the last logarithm is the real
trace series \(\sum_{n\ge1}(-1)^{n+1}\tr S^n/n\).
Its traces are closed finite-stencil chains with their original
addresses. Likewise
\[
 C_{\rm sel}^{-1}-\mathsf M
                  =\sum_{n\ge1}(-S)^n\mathsf M
\]
gives finite-chain records for the quadratic Gaussian correction.
Its absolute matrix sum is \(O(e^{-cB})\); integrating products
of its Gaussian letters still requires the joint precision
reserve used below. For the other determinant one may use exactly
\[
 \log\det\mathsf M=\int_0^1
                 \tr[\mathsf M_t^{-1}(\mathsf M-1)]\,\dd t,
 \qquad
 \mathsf M_t^{-1}=J_t\sum_{n\ge0}(-E_t)^n.
\]
The uniform selected-window bound justifies this series.
It gives local chain records for the determinant; subtraction of
their flat values isolates its extensive scalar. This is an
identity, prior to any small physical packet estimate.
In particular the frozen ratio costs
\(n_{\rm fr}|\log s_g|\); it cannot be put into a universal
constant. By \eqref{eq:tree-count} the common flat scalar per
coarse site is exactly
\[
       vL^4-4\log z_{\rm n}(\beta)
        +(3L^4+1)\log s_g-\frac{\log\det\mathsf M_0}{2c^4}.
\]

On an all-active common chart the relation with the stationary
Gaussian is also exact. With \(\delta=g^{-1}(\mu-m)\),
\[
 \frac{\dd\mu_{H^{-1}}(z+\delta)}
      {\dd\mu_{C_{\rm sel}}(z)}
  =\sqrt{\det H\det C_{\rm sel}}
    \exp\{-\tfrac12(z+\delta)^{\mathsf T}H(z+\delta)
                         +\tfrac12z^{\mathsf T}C_{\rm sel}^{-1}z\}.
\]
Both sides evaluate a physical multiplier at the same link values
\(m+g(z+\delta)=\mu+gz\). A bounded white-coordinate
reference agrees only on its stated no-failure event; all its
complementary events stay in the physical expression.


\section{Exact masked reassembly}
\label{sec:exact-masked-reassembly}

We specialise the multiplier in \eqref{eq:general-input} to its
actual regular and forced activities, then give the exact
algebraic reassembly. Its convergence criterion is stated
separately from the identities. In particular, an unpaid physical
hole weight will not be treated as a probability bound.

\subsection{The original masks and selection rules}

Let \(\mathcal B_h(U)\) be the set of all vertices incident
to at least one plaquette with \(d_G(U_p)>h\), and put
\(\chi_X(U)=1_{\{X\cap\mathcal B_h(U)=\varnothing\}}\).
These are plaquette tests, rather than individual link bounds.
Equality belongs to the good event. To read flags in \(X\)
one includes all links of plaquettes incident to \(X\); denote
their vertex enlargement by \(X^+\).

A regular activity \(E_i\) has footprint \(X_i\), anchor
\(o_i\in X_i\), and the good-set germ and norm used above.
A forced activity \(K_\alpha\) has original support
\(Y_\alpha\ne\varnothing\), nonempty inventory
\(I_\alpha\subset Y_\alpha\), and its prescribed additional
Borel test. A nonzero value requires
\(Y_\alpha\cap\mathcal B_h(U)=I_\alpha\) and that additional
test. It reads the original links and flags, hence is local on
\(Y_\alpha^+\). Define
\begin{align}
 E^\#(U)&=\sum_i\chi_{X_i}(U)E_i(U),\notag\\
 Z_K(U)&=\sum_{\substack{\Gamma\text{ finite},\ Y_\alpha
                         \text{ pairwise disjoint}\\
                  \bigcup_{\alpha\in\Gamma}I_\alpha=\mathcal B_h(U)}}
                                      \prod_{\alpha\in\Gamma}K_\alpha(U).
 \label{eq:reassembly-input-gas}
\end{align}
Thus \(\mathcal B=e^{-E^\#}Z_K\) in the full physical formula.
The inventories specify exactly which fine flags the forced
selection accounts for. On an all-good fine field only the empty
selection is possible, and \(Z_K=1\). A good coarse field alone
does not give that conclusion.

Write \(e_i\) for a bound on \(|E_i|\) on its good set,
and \(k_\alpha=\|K_\alpha\|_\infty\). Assume their original
weighted family bounds, so their sums are finite on each finite
torus. If \(e_i\le1\), its masked Mayer factor satisfies
\begin{equation}
 m_i=e^{-\chi_{X_i}E_i}-1,\qquad |m_i|\le e^{e_i}-1\le2e_i,
 \qquad e^{-E^\#}=\sum_{A\text{ finite}}\prod_{i\in A}m_i.
 \label{eq:reassembly-incoming-Mayer}
\end{equation}
The last series converges uniformly absolutely, since
\(\prod_i(1+2e_i)\le e^{2\sum_i e_i}\).
Similarly \(\sum_\Gamma\prod k_\alpha\le e^{\sum k_\alpha}\)
bounds the forced series before any exclusion or eligibility
restriction. No hard mask is differentiated.

Exact coverage can itself be written as local factors. For an
originally disjoint selection \(\Gamma\), let
\(Y_\Gamma=\bigcup_{\alpha\in\Gamma}Y_\alpha\) and put
\[
 q_x(U;\Gamma)=
 \begin{cases}1,&x\in Y_\Gamma,\\
               \chi_{\{x\}}(U),&x\notin Y_\Gamma.
 \end{cases}
\]
Then exactly
\begin{equation}
 Z_K(U)=\sum_{\Gamma\text{ originally disjoint}}
       \prod_{\alpha\in\Gamma}K_\alpha(U)\prod_x q_x(U;\Gamma).
 \label{eq:reassembly-local-prohibition}
\end{equation}
Every nonzero selected \(K_\alpha\) imposes its own inventory;
the remaining \(q_x\)'s prohibit all flags outside their union.
These two conditions are exactly the coverage condition in
\eqref{eq:reassembly-input-gas}. We may expand
\(q_x=1+(q_x-1)\), retaining its selection query even if the
numerical factor happens to equal one. A query and every selected
original support containing its site belong to the same component.

Original supports govern \(K\)-versus-\(K\) exclusion.
Enlarged reads govern integration components. Compatible original
labels whose enlarged reads overlap are merged into a component;
they are not removed from the original sum.

\subsection{Full reads and the marginal law of a component}

A letter is one addressed factor in the expansion of the exact
physical integrand. Before expanding its broad chart factors,
retain one such factor on the union of all selected physical
Borel reads, as specified by the exact rule in
Section~\ref{sec:regular-output-derivatives}. A chart-failure
letter on that retained union is prohibited; its selection
query is part of the full footprint. Its full footprint includes the physical
links and flags it reads, the Gaussian coordinates, frozen links
and shared tags, all sector and selection queries, all windows
defining its centre and covariance, and one additional Gaussian
range neighbourhood. For a long native path keep its entire path
record and these enlargements; do not replace its native length
by periodic distance.
Choose the last enlargement so that disjoint full-footprint
components have zero cross covariance in \(C_{\rm sel}\).
A fixed additional \(B\)-neighbourhood suffices.

\begin{proposition}[Factorisation with the same physical law]
\label{prop:reassembly-component-factorisation}
Fix \(V\) and sector geometry. Integrate letters against the
single joint Gaussian \(\mu_{C_{\rm sel}}\), the original
product frozen Haar law, and the product prior of the shared
local tags. If the components \(A_1,\ldots,A_k\) of a finite
word have disjoint full footprints, and each component product
is absolutely integrable under its marginal, then
\begin{equation}
          \int\prod_{\alpha\in A}f_\alpha
                    =\prod_{\nu=1}^k\int\prod_{\alpha\in A_\nu}f_\alpha.
 \label{eq:reassembly-component-factorisation}
\end{equation}
Each integral on the right uses the marginal of that same law.
\end{proposition}
\begin{proof}
On the union of the read Gaussian coordinates the covariance is
block diagonal. Its characteristic function
\(e^{-t^{\mathsf T}C_{\rm sel}t/2}\) factors, proving independence of the
component marginals. Their frozen and tag reads are disjoint,
so those product laws factor as well. The covariance does not
depend on frozen values or tags; the centre's dependence on
them is already included in each full read. Absolute component
integrability and Tonelli's theorem give a finite product of
absolute integrals. Fubini's theorem then proves the identity for
signed or complex factors. Absolute integrability of the whole
word alone would not rule out a zero component times an infinite
component integral, which is why the hypothesis is componentwise.
\end{proof}

Local branch likelihoods stay inside the letters. More explicitly,
their original factorisation uses one frame variable \(t_x\) per
coarse site and one prediction tag \(t_e\) per coarse link:
\[
       \lambda_t(U)=\prod_x p_x(U,t_x)
                         \prod_e p_e(U,t_x,t_{x+e_i},t_e).
\]
One can sum the tags directly or use their uniform product prior
and multiply each factor by its tag cardinality. A shared frame
tag remains the same variable in all incident factors.
Sector consistency is also local: if two sectors agree on a
component's full reads, the principal matrices, sources, tests
and covariance marginal used there agree. Thus distant sector
choices do not redefine that component integral. Their local
choice multiplicities must still be included in its price.

There is a useful geometric price for the incoming enlargement.
Let \(B_f\le2B\) be the incoming covering radius. Choose a fixed
minimal connected \(B_f\)-ball cover of a nonempty original
footprint \(X\), project its union to coarse parents, and enlarge
by the required fixed \(CB\)-neighbourhood. Call the result
\(\widehat X\). There is a geometric integer \(c_1\) with
\begin{equation}
                         n_B(\widehat X)\le c_1n_{B_f}(X).
 \label{eq:reassembly-cover-enlargement}
\end{equation}
Indeed each original covering ball projects into radius at most
\(2B+2\); its further enlargement is covered by a fixed number
of \(B\)-balls, and the resulting cover is connected.
If \(\kappa_{\rm in}\ge c_1(\kappa_{\rm out}+a)+1\),
\(a>0\), the enlarged incidence at exponent
\(\kappa_{\rm out}+a\) is at most \(C_B\) times the original
anchored regular value norm or forced incidence norm, respectively.
To check the incidence cost put \(k=n_{B_f}(X)\). If a coarse
site lies in \(\widehat X\), the original anchor, or a selected
site of a forced support, lies in a fine region with at most
\(C_LB^4(1+k)^4\) sites. The extra old exponent supplies
\(e^{-k}\); sum \(C_LB^4(1+k)^4e^{-k}\) over \(k\ge1\).
This accounts also for labels with the same support.
For regular Mayer factors include the factor two in
\eqref{eq:reassembly-incoming-Mayer}. These are value incomes,
not a derivative contraction theorem.

\subsection{The convergence input for physical components}

The matrix correction, actual phase difference, affine constant
square and frozen ratio in
\eqref{eq:physical-selected-full-push}--\eqref{eq:physical-selected-frozen-ratio}
are part of the physical weights. Defining a component amplitude
means integrating their selected letters, the original masks and
incoming factors with their shared Gaussian, Haar and tag marginal.
Component factorisation is used only after its absolute integral
has been bounded. The following elementary estimate isolates the
Gaussian part of that obligation.

If a Gaussian precision \(P\) and a nonnegative diagonal charge
\(d\) satisfy \(P-\operatorname{diag}d\ge r_0 1\), then
\begin{equation}
 \mathbb E_{P^{-1}}e^{\frac12\sum_i d_i z_i^2}
   =\left(\frac{\det P}{\det(P-\operatorname{diag}d)}\right)^{1/2}
                  \le\exp\!\left\{\frac1{2r_0}\sum_i d_i\right\}.
 \label{eq:reassembly-joint-Gaussian-tilt}
\end{equation}
To prove it, complete the Gaussian integral and integrate
\(\frac12\operatorname{Tr}[(P-t\operatorname{diag}d)^{-1}
\operatorname{diag}d]\) from zero to one. Its inverse is at
most \(r_0^{-1}1\). The same estimate holds for covariance
marginals, whose precision has the common lower bound in
\eqref{eq:physical-selected-precision-reserve}.
Thus a selected-word bound
\[
 \left|\prod_{\alpha\in A}f_\alpha\right|
       \le\prod_{\alpha\in A}a_\alpha(Y,t)
                            e^{\frac12\sum_i d_i z_i^2}
\]
is integrated in one use of the common precision reserve.
The resulting frozen majorant still requires its actual Haar
integral. Bounded incoming factors contribute \(2e_i,k_\alpha\)
with zero Gaussian charge. Removing components decreases the
nonnegative charges, so such a joint bound proves componentwise
integrability independently of factorisation.

For clarity the corresponding elementary failure estimate is also
local. If \(|f_\alpha|\le A_\alpha\) and it vanishes whenever
\(|z_i|\le H\) for every \(i\in T_\alpha\), assume a coordinate
belongs to at most \(N_B\) of these read sets and \(P\ge\lambda1\).
Choose \(v=\lambda/(4N_B)\) and use
\[
 1_{\{\max_{T_\alpha}|z_i|>H\}}
       \le\sum_{i\in T_\alpha}e^{v(z_i^2-H^2)/2}.
\]
Expanding witness choices and applying the single joint tilt with
\(r_0=\lambda/2\) gives
\begin{equation}
 \mathbb E\prod_{\alpha\in A}|f_\alpha|
  \le\prod_{\alpha\in A}
       [A_\alpha|T_\alpha|e^{1/(4N_B)}e^{-\lambda H^2/(8N_B)}].
 \label{eq:reassembly-failure-witnesses}
\end{equation}
The largest total diagonal charge is \(N_Bv=\lambda/4\).
For the finite lists of coordinate tests,
\(|T_\alpha|,N_B\le C_B\) and
\(H=P(\widehat\beta)/C_B\) give incidence
\(C_Be^{-P(\widehat\beta)^2/C_B}\), smaller than every inverse
power of \(\widehat\beta\). A physical failure has this price
only after its stated vanishing on that Gaussian box is proved
for the selected centre of its sector. The present exact
reassembly retains such tests and does not assume that assertion.

After integrating actual letters, a stronger-exponent incidence
price \(\epsilon\) gives connected component income
\(C_B\epsilon\) by Lemma~\ref{lem:connected-positive-tree}.
The same prices justify the countable sums and exchanges of
integrals. In a finite truncation the reorganisation is finite
algebra; its countable limit requires this absolute majorant.
The full joint hole-packet price, including
\(s_g^{-n_{\rm fr}}\), the physical phase difference and branch
weights, is proved in Section~\ref{sec:joint-hole-packets}.
Its physical energy estimate supplies the input that a Gaussian
failure bound alone or a hole count would not supply.

\subsection{The mixed gas and both exact restorations}

Let \(\mathcal D=\mathcal B_{h_{\rm c}}(V)\) be the actual
coarse flag set, with the same incident-vertex convention.
An integrated component carries an inventory of coarse flags
from its mandatory coarse-flag letters. Components with empty
inventory are unforced species \(u_i(V)\), with full footprints
\(X_i\) and occurrence mask \(X_i\cap\mathcal D=\varnothing\).
Nonempty inventories give forced species \(d_j(V)\), footprints
\(Y_j\), and inventories \(I_j\ne\varnothing\); their eligibility
includes \(Y_j\cap\mathcal D=I_j\).
An incoming fine forced label or a Gaussian failure need not
carry a coarse flag and hence need not be coarse forced.

When the exact physical word expansion and these component prices
are available, its mixed gas is
\begin{equation}
 \Xi(V)=\sum_{\Gamma\text{ forced admissible}}
       \prod_{j\in\Gamma}d_j(V)\,
                 Z_{\rm u}(\Lambda\setminus Y_\Gamma;V).
 \label{eq:reassembly-mixed-gas}
\end{equation}
Forced admissibility means disjoint full footprints and exact
inventory coverage of \(\mathcal D\). The unforced gas contains
only mutually disjoint species avoiding both \(Y_\Gamma\) and
\(\mathcal D\). This is the reassembly input, with its model
price still explicit, rather than an unconditional estimate on
all physical sectors.

Suppose the unforced gas satisfies the logarithm criterion of
Lemma~\ref{lem:connected-log-bound}. Index its log clusters by
\(i\) and write their signed functions as \(\ell_i(V)\), retaining
the repeated species and Ursell factorial of
\eqref{eq:connected-Ursell-log}. Their footprints \(X_i\) are
unions of all cluster reads. With forbidden set \(Y\),
\[
 \log Z_{\rm u}(\Lambda\setminus Y;V)
    =\sum_{\substack{i:X_i\cap Y=\varnothing\\
                       X_i\cap\mathcal D=\varnothing}}\ell_i(V).
\]
This is the same convergent logarithm with forbidden species set
to zero. Put
\[
 c_i=\ell_i(1),\qquad E_i^{\rm new}=-\ell_i+c_i,\qquad
 v_{\rm log}=\sum_i c_i,
\]
where \(v_{\rm log}\) is extensive and its per-site increment
is \(v_{\rm log}/c^4\). Only the flat scalar is subtracted here;
quadratic extraction belongs to the later coupling reset.
Define
\begin{equation}
 t_i(V)=\begin{cases}
          -\ell_i(V),&X_i\cap\mathcal D=\varnothing,\\
          -c_i,&X_i\cap\mathcal D\ne\varnothing.
        \end{cases}
 \label{eq:reassembly-two-restorations}
\end{equation}
In this paragraph the superscript \(\#\) on \(E^{\rm new}\)
uses the coarse mask \(1_{\{X_i\cap\mathcal D=\varnothing\}}\).

\begin{proposition}[Exact exclusion and vacuum restoration]
\label{prop:reassembly-exact-restoration}
Under absolute convergence of the mixed gas and its logarithm,
\begin{equation}
 \Xi(V)=e^{v_{\rm log}-(E^{\rm new})^\#(V)}
   \sum_{\Gamma\text{ forced admissible}}\prod_{j\in\Gamma}d_j(V)
            \exp\!\left\{\sum_{i:X_i\cap Y_\Gamma\ne\varnothing}t_i(V)\right\}.
 \label{eq:reassembly-exact-restoration}
\end{equation}
\end{proposition}
\begin{proof}
The full healthy logarithm obeys exactly
\[
 \sum_{i:X_i\cap\mathcal D=\varnothing}\ell_i
  =v_{\rm log}-(E^{\rm new})^\#
                               -\sum_{i:X_i\cap\mathcal D\ne\varnothing}c_i.
\]
For a fixed admissible \(\Gamma\), remove from it every healthy
cluster meeting \(Y_\Gamma\). Every bad \(X_i\) already meets
\(Y_\Gamma\): choose a flag in \(X_i\) and use exact inventory
coverage. Its scalar \(-c_i\) is therefore precisely a second
branch term of \eqref{eq:reassembly-two-restorations}; the removed
healthy clusters are precisely the first branch terms. Exponentiate
and sum over \(\Gamma\). Absolute convergence justifies both
operations and proves the identity.
\end{proof}

The \(-\ell_i\) branch restores exclusion from forced footprints.
The \(-c_i\) branch compensates the unmasked flat vacuum scalar
on bad footprints. Both are required for the same original
density. The identity initially concerns local value functions.
Section~\ref{sec:regular-output-derivatives} proves regular
germs for the packet contribution by differentiating its physical
integration density. It does not differentiate a hard coarse
indicator; the other unforced contributions retain their own
derivative-estimate obligations.

\subsection{Merging the restoration factors}

Let \(a_j\ge\|d_j\|_\infty\), and choose \(b_i\) bounding
both \(|\ell_i(V)|\) on its healthy domain and \(|c_i|\).
Assume their incidence sums satisfy
\[
 \sup_x\sum_{j:x\in Y_j}e^{(\kappa+1)n_B(Y_j)}a_j\le a,
 \qquad
 \sup_x\sum_{i:x\in X_i}e^{(\kappa+1)n_B(X_i)}b_i\le b\le1.
\]
These prices include the integrated physical normalisations.
Set \(r_i=e^{t_i}-1\), so \(|r_i|\le2b_i\).
For a fixed forced selection only addresses \(i\) meeting its
original forced union may occur in the Mayer expansion.

A new label \(C\) is a distinct connected set of forced labels
and such restoration addresses. It contains a forced label,
its forced labels are mutually compatible, and every restoration
meets at least one of those forced labels. Define
\[
 Y_C=\bigcup_{j\in C\text{ forced}}Y_j
                    \cup\bigcup_{i\in C\text{ restoration}}X_i,
 \qquad I_C=\bigcup_{j\in C\text{ forced}}I_j.
\]
Its activity is the product of its \(d_j\)'s and \(r_i\)'s,
with all internal tests and
\(1_{\{Y_C\cap\mathcal D=I_C\}}\). It has nonempty inventory.

\begin{theorem}[Exact masked merger and its value norm criterion]
\label{thm:reassembly-exact-merger}
Choose \(\eta\ge0\), put \(\gamma=e^\eta a+2b\), and assume
\(D_B\gamma\le1\), where \(D_B=(2B+1)^4\).
The new forced gas equals the sum in
\eqref{eq:reassembly-exact-restoration} exactly, and
\begin{equation}
 \sup_x\sum_{C:x\in Y_C}e^{\kappa n_B(Y_C)}\|K_C^{\rm new}\|_\infty
                              \le e^{-\eta}(e^\gamma-1).
 \label{eq:reassembly-merger-price}
\end{equation}
Thus \(\Xi=e^{v_{\rm log}-(E^{\rm new})^\#}Z_{K^{\rm new}}\)
with the original coarse masks and exact flag coverage.
\end{theorem}
\begin{proof}
For finite alphabets expand the allowed restoration factors and
split each term into overlap components. Every component contains
an original forced label. A coarse flag in its footprint belongs
to an inventory from the same component: the original forced
footprint covering that flag intersects it and hence joins it.
Thus its inventory is exactly \(Y_C\cap\mathcal D\).
Conversely, flattening an admissible new family recovers its
compatible original forced labels and its allowed restoration
selection. Each restoration meets one of these original forced
labels. It cannot occur in two different components, since its
nonempty footprint would intersect both. The two operations are
inverse and preserve each signed product and every internal test.

For the norm, give a forced label price \(e^\eta a_j\) and
a restoration price \(2b_i\). Their combined incidence at
exponent \(\kappa+1\) is at most \(\gamma\).
Every new component contains a forced label, so its actual product
is at most \(e^{-\eta}\) times this priced product.
Also \(n_B(Y_C)\le\sum_{\alpha\in C}n_B(X_\alpha)\).
Lemma~\ref{lem:connected-positive-tree}, applied to the prices
including \(e^{\kappa n_B}\), bounds the rooted connected-set
sum by \(\gamma\le e^\gamma-1\). This proves
\eqref{eq:reassembly-merger-price}. The same absolute majorant
permits passage from finite to countable alphabets and proves
the claimed exact identity there.
\end{proof}

This is the two-restoration merger of
\cite[Proposition~29.3 and Theorem~29.4, pp.~164--165]{Main}.
The proof has kept fine original-support exclusion distinct from
full-read component formation throughout.

The completed units now have precise boundaries. Calibration
proves the general-good auxiliary central comparison, and
\eqref{eq:physical-selected-full-push} is an exact physical
representation with finite reads for bounded Borel incoming
activities. The masked restoration and merger are exact under
their explicit absolute-integral and incidence criteria.
The next two sections supply the actual joint hole-packet
payment, including the frozen ratio, affine square, physical
phase differences, branch weights and shared incoming factors.
Section~\ref{sec:regular-output-derivatives} then proves the
regular output germs for those packet components.
Sections~\ref{sec:central-local-data}--\ref{sec:central-strength-estimates}
supply the packet-free central estimates;
Sections~\ref{sec:affine-local-reset}--\ref{sec:sharp-regular-transport}
prove the local quadratic subtraction and sharp incoming transport.
Principal-action matching and assembly of the full coupling reset
remain distinct proof units. The algebraic merger does not replace
any of these estimates.


\section{Completing the physical row comparison}
\label{sec:packet-row-comparison}

The preceding merger requires prices for the actual physical
packets. A frozen link costs a missing Gaussian normalisation,
and the exact phase contains a negative affine square. We shall
keep this square with the physical energy that pays it.
Two inputs are needed before doing so: circulation on every
positive-likelihood branch, and a completion of the calibrated
row lift on all coarse curvature rows. Calibration on exact
curls alone does not control a masked or thinned row vector.

Proposition~\ref{prop:packet-macro-circulation}, proved in
Appendix~\ref{app:packet-branches}, gives
\[
 d_G(V_p)\le CL^{-2}\sum_{q\subset\mathcal M_p}d_G(U_q)
                     +\sum_{b\in\partial p}d_G(V_bQ_t(U)_b^{-1}).
\]
The constant and macro-box overlap count are chosen before
\(L\), and the estimate holds at rough fields for every actual
tag of positive likelihood. This is stronger than a
regular-branch derivative estimate.

\subsection{A complementary row without changing the Gaussian history}

Split the previously constructed flat lift as
\(\Phi_0=(\Phi_F,\Phi_I)^{\mathsf T}\), for fine and noise rows.
We now choose the still free nonnegative parameters \(l_j\)
of \eqref{eq:calibration-positive-extension}. This choice
preserves the potential kinetic matrices and all preceding
Gaussian identities. We use one Laurent division estimate for
both factorisations, then choose the positive extension before
increasing the block sides.

\begin{lemma}[A positive completion on the whole row space]
\label{lem:packet-complement}
The parameters \(l_j\) can be chosen before the block sides,
and the lower bounds on those sides then increased, so that
an exponentially local operator \(\Phi_\perp\) satisfies
\begin{equation}
 \Phi_F^{\mathsf T}\Phi_F+\Phi_I^{\mathsf T}\Phi_I+\Phi_\perp^{\mathsf T}\Phi_\perp=1,
                       \qquad\Phi_\perp S_cd_c=0.
 \label{eq:packet-row-completion}
\end{equation}
All coefficient prices are uniform over legal histories at the
fixed adjacent profiles. The fine and noise prices remain
\(CL^{-2}\) and \(C\), respectively.
\end{lemma}
\begin{proof}
\emph{Uniform bounds before choosing the positive extension.}
Write \(d_r=d_{r,c}\) for the coarse coboundary from
\(r\)-forms to \((r+1)\)-forms; thus \(d_1=d_c\).
We first use the unmodified metrics \(H_f,H_c\), their
inverses \(B_f,B_c\), and the column \(T\) of
Appendix~\ref{app:calibration}. Fix the adjacent profiles and
\(a>0\). The bounds in that appendix give a common width
\(\sigma_0>0\), chosen before the block side, history and
period, such that
\begin{equation}
 |H_c|_{\sigma_0}+|B_c|_{\sigma_0}\le C,\qquad
 \max_r|T_r|_{\sigma_0}\le CL^{-2},\qquad
 \max_r|(\nabla_{f,i}T)_r|_{\sigma_0}\le CL^{-3}.
 \label{eq:packet-complement-input-bounds}
\end{equation}
Here \(r\) is a fine residue in one block. The coefficient
norm is the one in Section~\ref{sec:calibrated-rows}: it sums
absolute coefficients with the weight
\(e^{\sigma_0|u|_1}\) over coarse displacements \(u\) and
input form types, and takes a maximum over output types and
residues. In particular it does not sum the \(L^4\) output
residues. The unmodified \(B_f\) has a uniform exponential
row bound in native fine units. Reducing \(\sigma_0\) once
accommodates this bound and all the finitely many difference
directions. All constants obtained up to the choice of
\(l_f,l_c\) below are independent of these two parameters.

We spell out the residue sum needed for the defect. In block
coordinates write \(B_{f;rr'}(u)\) for the covariance from
input residue \(r'\) to output residue \(r\). Its native fine
displacement is \(Lu+r-r'\), up to fixed cell-type offsets.
Since the residues are in one block,
\[
 |u|_1\le L^{-1}|Lu+r-r'|_1+4.
\]
For fixed \(r\), the pairs \((u,r')\) enumerate the fine input
sites once. The native covariance row bound therefore implies
\[
                 \max_r\sum_{r'}|B_{f;rr'}|_{\sigma_0}\le C.
\]
The fixed number of form types is absorbed in \(C\).
Transposing a fixed-residue matrix \(T_r\) reverses its
displacement and exchanges only fixed form types. Thus
\begin{equation}
 \begin{split}
 |T^{\mathsf T}B_fT|_{\sigma_0}
 &\le\sum_{r,r'}|T_r^{\mathsf T}|_{\sigma_0}
                   |B_{f;rr'}|_{\sigma_0}|T_{r'}|_{\sigma_0}\\
 &\le CL^{-4}\sum_r\sum_{r'}|B_{f;rr'}|_{\sigma_0}
 \le CL^{-4}L^4=C.
 \end{split}
 \label{eq:packet-defect-residue-sum}
\end{equation}
The inner sum is a covariance row sum. There is only one
remaining sum over \(L^4\) residues, and no sum counting the
number of blocks.

Set
\[
             Z=H_c(B_c-T^{\mathsf T}B_fT-a^{-1}d_1d_1^{\mathsf T})H_c.
\]
The coboundaries and their transposes have finitely many unit
coefficients. Choose \(C_{\rm d}\ge1\) bounding all these
coarse operators at width \(\sigma_0\). Submultiplicativity
and \eqref{eq:packet-defect-residue-sum} give
\begin{equation}
                 Z=Z^{\mathsf T},\qquad |Z|_{\sigma_0}\le C_Z.
 \label{eq:packet-unmodified-defect-bound}
\end{equation}
To verify its exact annihilation, use
\(TH_cd_1=H_fd_{1,f}E\),
\(T^{\mathsf T}d_{1,f}=d_1Q\), and
\(QE=1-a^{-1}K_c\), where \(K_c=d_1^{\mathsf T}H_cd_1\).
These are the installed column identities of Appendix~C;
the positive extensions agree with the unmodified metrics
on curls. Multiplication gives
\begin{equation}
 \begin{split}
 Zd_1
 &=H_c\bigl(d_1-T^{\mathsf T}B_fTH_cd_1-a^{-1}d_1K_c\bigr)\\
 &=H_cd_1(1-QE-a^{-1}K_c)=0.
 \end{split}
 \label{eq:packet-defect-annihilation}
\end{equation}
Symmetry also gives \(d_1^{\mathsf T}Z=0\).

\smallskip\noindent
\emph{One Laurent division estimate.}
Put \(z_i=e^{ip_i}\) and \(q_i=z_i-1\). For a scalar Laurent
series \(f(z)=\sum_{n\in\Z^4}f_nz^n\), let
\(\mathcal E_i f=f|_{z_i=1}\), and define
\(\mathcal D_i f=(f-\mathcal E_i f)/(z_i-1)\) by its
coefficient series. The elementary identities are
\[
 \frac{z_i^n-1}{z_i-1}=
 \begin{cases}
   \displaystyle\sum_{b=0}^{n-1}z_i^b,&n\ge1,\\[2pt]
   0,&n=0,\\[2pt]
   \displaystyle-\sum_{b=n}^{-1}z_i^b,&n\le-1.
 \end{cases}
\]
Each interval sum has weighted coefficient sum at most
\(|n|e^{\tau'|n|}\). Thus, for \(0<\tau'<\tau\),
\begin{equation}
 \begin{split}
 |\mathcal E_i f|_\tau&\le |f|_\tau,\\
 |\mathcal D_i f|_{\tau'}
 &\le\sup_{m\ge1}\bigl(m e^{-(\tau-\tau')m}\bigr)|f|_\tau
 \le\frac{|f|_\tau}{e(\tau-\tau')}.
 \end{split}
 \label{eq:packet-laurent-width-loss}
\end{equation}
Here the scalar norm is
\(|f|_\tau=\sum_n|f_n|e^{\tau|n|_1}\).
The evaluation combines coefficients by summation and sets
one exponent to zero, which proves its bound. Absolute
convergence justifies the divided difference even at
\(z_i=1\). No singular quotient is assigned there.

We give the exterior lift explicitly. A row \(K\) on
\(k\)-forms, \(1\le k\le4\), with \(Kd_{k-1}=0\), is
represented by its dual exterior form \(\omega\).
Contraction by \(q\), denoted by \(\iota_q\), removes one
basis covector and multiplies by the corresponding \(q_i\):
\[
 \iota_q(e^{i_1}\wedge\cdots\wedge e^{i_k})
 =\sum_{\nu=1}^k(-1)^{\nu-1}q_{i_\nu}
       e^{i_1}\wedge\cdots\wedge
       \widehat{e^{i_\nu}}\wedge\cdots\wedge e^{i_k}.
\]
The hat means omission, and \(e^1,\ldots,e^4\) are the
coordinate covectors. The row equation is
\(\iota_q\omega=0\).
Split along the first remaining coordinate \(i\) as
\(\omega=e^i\wedge a_i+b_i\), with neither \(a_i\) nor
\(b_i\) containing \(e^i\). Writing \(q'\) for the remaining
components gives
\[
 \iota_{q'}a_i=0,\qquad
 \iota_{q'}b_i=-q_i a_i,\qquad
 \iota_{q'}\mathcal E_i b_i=0.
\]
It follows, first off \(q_i=0\) and then by the absolutely
convergent series above, that
\[
 \iota_q(e^i\wedge\mathcal D_i b_i)
       =b_i-\mathcal E_i b_i+e^i\wedge a_i
       =\omega-\mathcal E_i b_i.
\]
Repeat with residual \(\mathcal E_i b_i\) in the next
coordinate. After four coordinates a positive-degree residual
with no coordinate left is zero. At top degree the cycle is
already zero: its coefficient is annihilated by the nonzero
analytic factors \(q_i\). Summing these identities constructs
a right lift \(K_+\) satisfying \(K_+d_k=K\).

Every summand contains one \(\mathcal D_i\), applied to a
component of the original row after coordinate evaluations.
Those evaluations preserve the input width \(\tau\).
Consequently the whole exterior lift, rather than each
coordinate separately, obeys
\begin{equation}
 |K_+|_{\tau'}
 \le C_0\bigl(1+(\tau-\tau')^{-1}\bigr)|K|_\tau.
 \label{eq:packet-exterior-lift-bound}
\end{equation}
The constant \(C_0\) counts only the fixed exterior form
types; the construction acts independently on every output
row. For the coarse operators used below, transposition gives
a left lift with the same form of bound, enlarging \(C_0\)
for the fixed types if needed.
For an operator on the cover, transposition reverses
displacements; on symbols it replaces \(q_i\) by
\(z_i^{-1}-1=-z_i^{-1}q_i\). Hence it preserves exponential
summability and all the stated identities. The argument is
division of analytic symbols on the cover, followed by
periodisation. It does not require invertibility on harmonic
cochains of a finite torus.

\smallskip\noindent
\emph{The four width losses in the two-sided factorisation.}
Set \(\sigma_b=2^{-b}\sigma_0\), \(b=0,\ldots,4\), and take
\[
                 C_{\rm div}=C_0(1+16/\sigma_0)\ge1,
\]
enlarging \(C_0\) once if necessary. This bounds every lift
from \(\sigma_{b-1}\) to \(\sigma_b\).
First lift \(Z\) on the left to get \(d_2^{\mathsf T}W=Z\).
Since \(d_2^{\mathsf T}(Wd_1)=Zd_1=0\), a second left lift gives
\[
                         d_3^{\mathsf T}Y=Wd_1.
\]
Then \(d_3^{\mathsf T}Yd_0=0\). At every nonzero real momentum
\(d_3^{\mathsf T}\), from four-forms to three-forms, is an injective
column: at least one \(q_i\) is nonzero. Hence \(Yd_0=0\)
there, and continuity extends this identity to zero momentum.
All its Fourier coefficients therefore vanish. A right lift
now gives \(Xd_1=Y\). Put \(W'=W-d_3^{\mathsf T}X\); then
\[
                         W'd_1=0,\qquad d_2^{\mathsf T}W'=Z.
\]
A fourth lift gives \(W'=F_0d_2\). With
\(F=(F_0+F_0^{\mathsf T})/2\), symmetry of \(Z\) yields
\(Z=d_2^{\mathsf T}Fd_2\).

For clarity the bounds at these successive widths are
\begin{equation}
 \begin{aligned}
 |W|_{\sigma_1}&\le C_{\rm div}C_Z,\\
 |Y|_{\sigma_2}&\le C_{\rm d}C_{\rm div}^{\,2}C_Z,\\
 |X|_{\sigma_3}&\le C_{\rm d}C_{\rm div}^{\,3}C_Z,\\
 |W'|_{\sigma_3}
 &\le(C_{\rm div}+C_{\rm d}^{\,2}C_{\rm div}^{\,3})C_Z,\\
 |F|_{\sigma_4}
 &\le C C_{\rm div}
            (C_{\rm div}+C_{\rm d}^{\,2}C_{\rm div}^{\,3})C_Z
   =:C_F.
 \end{aligned}
 \label{eq:packet-four-lift-prices}
\end{equation}
The fixed factor \(C\) in the last line covers symmetrisation.
We may enlarge \(C_F\) to be at least one. It is independent
of \(L,l_f,l_c\), history depth and period. In particular,
four width losses suffice for every history; none is
iterated once per RG step.

\smallskip\noindent
\emph{The mixed lift and both Schur sums.}
The row on coarse two-forms
\[
                         K=d_{2,f}B_fTH_c
\]
annihilates \(d_1\), because
\[
                Kd_1=d_{2,f}B_fH_fd_{1,f}E
                         =d_{2,f}d_{1,f}E=0.
\]
We also need its absolute coefficient gain. The unmodified
\(B_f\) is diagonal in two-form types. For \(i<j<k\),
\[
 (d_{2,f}B_fT)_{ijk}
 =B_{f,jk}\nabla_{f,i}T_{jk}
  -B_{f,ik}\nabla_{f,j}T_{ik}
  +B_{f,ij}\nabla_{f,k}T_{ij}.
\]
Each scalar convolution \(B_{f,ij}\) commutes with a fine
forward difference. We place that difference on \(T\) and
use \eqref{eq:packet-complement-input-bounds}, including its
block-boundary crossings. Convolution by \(B_f\) has the
uniform block row sum used in
\eqref{eq:packet-defect-residue-sum}; \(H_c\) has a fixed
coarse coefficient bound. Therefore
\[
                         \max_r|K_r|_{\sigma_0}\le CL^{-3}.
\]
Apply the same right lift, directly from \(\sigma_0\) to
\(\sigma_4\). It acts separately on each fine output row,
so it gives \(A d_2=K\) and
\begin{equation}
 d_{2,f}B_fT=A\,d_2B_c,\qquad
                         \max_r|A_r|_{\sigma_4}\le CL^{-3}.
 \label{eq:packet-mixed-lift-price}
\end{equation}
There is no transposition or sum over the fine output
residues in this construction.

Write \(A_{r\alpha,\gamma}(u)\) for its coefficients, where
\(\alpha,\gamma\) are the fine and coarse three-form types.
On the cover its full absolute row and column sums obey
\begin{equation}
 \begin{aligned}
 \sup_{r,\alpha}\sum_{u,\gamma}
                         |A_{r\alpha,\gamma}(u)|
       &\le CL^{-3},\\
 \sup_\gamma\sum_{r,\alpha,u}
                         |A_{r\alpha,\gamma}(u)|
       &\le CL^4L^{-3}=CL.
 \end{aligned}
 \label{eq:packet-mixed-schur-sums}
\end{equation}
Translation invariance identifies these with the sums over
all input or output sites for any fixed row or column.
Only the column sum counts the \(L^4\) residues. The
geometric-mean Schur estimate, in counting \(\ell^2\), is
\[
 \|A\|_{2\to2}
 \le\sqrt{\|A\|_{\infty\to\infty}\|A\|_{1\to1}}
 \le\sqrt{CL^{-3}\,CL}=C/L.
\]
Indeed Cauchy--Schwarz with weights given by the absolute
matrix entries bounds \(\|Av\|_2^2\) by the maximum row
sum times the maximum column sum times \(\|v\|_2^2\).
We also need an exponential coefficient bound for the square,
not just its operator norm. Since
\(A^{\mathsf T}A=\sum_r A_r^{\mathsf T}A_r\), submultiplicativity and the
fixed type count give
\[
 |A^{\mathsf T}A|_{\sigma_4}
 \le C\sum_r|A_r|_{\sigma_4}^{\,2}
 \le CL^4(L^{-3})^2=CL^{-2}.
\]
Choose \(C_A\) large enough that both conclusions read
\begin{equation}
          \|A\|_{2\to2}\le C_A/L,\qquad
                    |A^{\mathsf T}A|_{\sigma_4}\le C_A^2/L^2.
 \label{eq:packet-mixed-square-price}
\end{equation}
Absolute periodisation can only decrease these bounds.
They hold simultaneously on every compatible finite period.

\smallskip\noindent
\emph{Exact rebasing of the defect.}
Now use the allowed positive extensions
\[
 \widetilde B_f
    =B_f+l_f(d_{2,f}B_f)^{\mathsf T}(d_{2,f}B_f),\qquad
 \widetilde B_c
    =B_c+l_c(d_2B_c)^{\mathsf T}(d_2B_c),
 \qquad \widetilde H_c=\widetilde B_c^{-1}.
\]
At the literal Wilson input the effective \(l_f\) is zero,
as required by \eqref{eq:calibration-positive-extension}.
For a nonempty fine history it is its prescribed profile
parameter. The unmodified factorisation of \(Z\) and
\eqref{eq:packet-mixed-lift-price} imply exactly
\begin{equation}
 \begin{split}
 &\widetilde B_c-T^{\mathsf T}\widetilde B_fT-a^{-1}d_1d_1^{\mathsf T}\\
 &\hspace{12pt}
       =B_cd_2^{\mathsf T}(F+l_c1-l_fA^{\mathsf T}A)d_2B_c.
 \end{split}
 \label{eq:packet-covariance-defect-factor}
\end{equation}
Define \(D_0=(1+l_cd_2B_cd_2^{\mathsf T})^{-1}\).
To check the rebasing identity, multiply instead of invoking
an inverse formula:
\[
 \begin{split}
 (1+l_cd_2B_cd_2^{\mathsf T})d_2B_c\widetilde H_c
 &=d_2(B_c+l_cB_cd_2^{\mathsf T}d_2B_c)\widetilde H_c\\
 &=d_2.
 \end{split}
\]
Thus \(d_2B_c\widetilde H_c=D_0d_2\). Applying this and
its transpose to \eqref{eq:packet-covariance-defect-factor}
gives
\begin{equation}
 \begin{aligned}
 &\widetilde H_c(\widetilde B_c-T^{\mathsf T}\widetilde B_fT
                    -a^{-1}d_cd_c^{\mathsf T})\widetilde H_c=d_2^{\mathsf T}Rd_2,\\
 &\hspace{20pt}R=D_0^{\mathsf T}(F+l_c1-l_fA^{\mathsf T}A)D_0.
 \end{aligned}
 \label{eq:packet-positive-defect}
\end{equation}

\smallskip\noindent
\emph{A positive factor from a convergent power series.}
The constants \(C_F,C_A\) were obtained from unmodified
metrics. We can therefore choose, in this order,
\begin{equation}
                 l_c\ge2(C_F+1),\qquad
                             L^2\ge l_fC_A^2.
 \label{eq:packet-positive-parameter-choice}
\end{equation}
There is no dependence of \(C_F,C_A\) on this choice of
\(l_c\). Since \(F=F^{\mathsf T}\), its absolute row bound also
bounds its columns and hence \(\|F\|_{2\to2}\le C_F\).
On every real momentum and every finite period,
\[
 F+l_c1-l_fA^{\mathsf T}A
 \ge(l_c-C_F-l_fC_A^2/L^2)1
 \ge(l_c/2)1\ge1.
\]
There is a convenient coefficient proof of a square factor.
Define
\begin{equation}
 B_*=\sqrt{l_c}\sum_{n=0}^{\infty}\binom{1/2}{n}
                \left(\frac{F-l_fA^{\mathsf T}A}{l_c}\right)^n.
 \label{eq:packet-binomial-factor}
\end{equation}
This is the power series of \((1+z)^{1/2}\) at \(z=0\).
By \eqref{eq:packet-mixed-square-price} and
\eqref{eq:packet-positive-parameter-choice},
\[
 \left|\frac{F-l_fA^{\mathsf T}A}{l_c}\right|_{\sigma_4}
       \le\frac{C_F+1}{l_c}\le\frac12.
\]
The coefficients \(\binom{1/2}{n}\) have absolute value at
most one. Thus the series converges absolutely in the same
coefficient algebra, with
\[
                          |B_*|_{\sigma_4}\le2\sqrt{l_c}.
\]
Each power is symmetric, and multiplying two absolutely
convergent series gives
\begin{equation}
            B_*^{\mathsf T}=B_*,\qquad B_*^{\mathsf T}B_*=F+l_c1-l_fA^{\mathsf T}A.
 \label{eq:packet-binomial-square}
\end{equation}
This uses the scalar binomial identity with powers of one
operator. It requires no square-root construction for \(R\).

For completeness \(D_0\) has uniform exponential coefficients
after \(l_c\) is fixed. Indeed put \(W_c=d_2B_cd_2^{\mathsf T}\).
It has a fixed bound \(|W_c|_{\sigma_4}\le C_W\) and is
nonnegative at real momenta. Hence
\(\|(1+l_cW_c(p))^{-1}\|_{2\to2}\le1\) there.
For \(p=x+i\eta\), \(x\) real and
\(|\eta|_\infty\le\sigma_4/2\), its coefficient series gives
\[
 \begin{split}
 \|W_c(x+i\eta)-W_c(x)\|_{2\to2}
 &\le C|\eta|_\infty
       \sum_u |u|_1e^{(\sigma_4/2)|u|_1}|(W_c)_u|\\
 &\le \frac{2CC_W}{e\sigma_4}|\eta|_\infty
       \le C_W'|\eta|_\infty .
 \end{split}
\]
Here the coefficient matrix norm is the absolute row norm;
the fixed type factor \(C\) converts it to operator norm.
Choose \(C_W'\ge1\) with the displayed bound, and set
\[
          \delta_D=\min\{\sigma_4/4,(4l_cC_W')^{-1}\}.
\]
On \(|\operatorname{Im}p|_\infty<\delta_D\), perturbation
from the real inverse has norm at most \(1/4\).
The Neumann series therefore gives a holomorphic \(D_0\)
with operator norm at most two. Shifting each Fourier
contour as in Appendix~C yields
\[
                     |(D_0)_u|\le C e^{-\delta_D|u|_1}.
\]
Summing the four geometric series proves
\(|D_0|_{\delta_D/2}\le C\), uniformly in histories and
periods. This argument fixes a width after \(l_c\), but
still before the side \(L\). Take a common positive width
\(\sigma_*\) below \(\delta_D/2,\sigma_4\) and the widths
of \(\widetilde B_c,S_c\) from Appendix~C.

\smallskip\noindent
\emph{The complementary row and the order of all choices.}
Put \(P_S=S_c\widetilde B_cS_c^{\mathsf T}\).
Since \(S_c^{\mathsf T}S_c=\widetilde H_c\), direct multiplication
gives
\[
             P_S^{\mathsf T}=P_S,\qquad P_S^2=P_S,\qquad P_SS_c=S_c.
\]
Using also \(S_f^{\mathsf T}S_f=\widetilde H_f\), the original row
lift \eqref{eq:calibrated-explicit-row-lift} satisfies
\[
 \Phi_0^{\mathsf T}\Phi_0
     =S_c(T^{\mathsf T}\widetilde B_fT+a^{-1}d_1d_1^{\mathsf T})S_c^{\mathsf T}.
\]
Multiplying \eqref{eq:packet-positive-defect} on the left
by \(S_c\widetilde B_c\) and on the right by its transpose,
and using \eqref{eq:packet-binomial-square}, gives
\begin{equation}
 \begin{split}
 1-\Phi_0^{\mathsf T}\Phi_0
 &=1-P_S\\
 &\quad+S_c\widetilde B_cd_2^{\mathsf T}D_0^{\mathsf T}B_*^{\mathsf T}B_*D_0
                          d_2\widetilde B_cS_c^{\mathsf T}.
 \end{split}
 \label{eq:packet-full-row-defect-factor}
\end{equation}
Define
\begin{equation}
 \Phi_\perp=
 \begin{pmatrix}
       1-P_S\\ B_*D_0d_2\widetilde B_cS_c^{\mathsf T}
 \end{pmatrix}.
 \label{eq:packet-explicit-complement}
\end{equation}
Because \((1-P_S)^2=1-P_S\), its Gram operator is exactly
the right side of \eqref{eq:packet-full-row-defect-factor}.
Moreover its first block vanishes on \(S_cd_1\), and its
second block gives
\[
 B_*D_0d_2\widetilde B_cS_c^{\mathsf T}S_cd_1
                      =B_*D_0d_2d_1=0.
\]
This proves both identities in
\eqref{eq:packet-row-completion}. Every factor in
\eqref{eq:packet-explicit-complement} has an exponential
coefficient bound at the common width \(\sigma_*\).
The complementary row is therefore exponentially local on
the whole row space, including vectors which are not exact
curvatures.

Finally the choices can be made for all legal steps at once.
If the coarse index is \(t\), its fine index is
\(j\in\{t,t+1\}\). Include the literal Wilson input when
taking the corresponding uniform bounds. For each \(t\),
choose
\[
 l_t\ge2\left(1+
       \max_{\text{legal }(j,t)} C_F(j,t)\right).
\]
This is a maximum over the finitely many adjacent profile
types, not over history length. After fixing these \(l_j\),
choose \(\lambda_j\) below half the common lower spectral
bound of \(\widetilde H_v\) and construct \(S_v\) as in
\eqref{eq:calibration-positive-extension}. All resulting
coefficient budgets and common widths are now fixed.
For each legal adjacent pair increase the admissible odd
block sides to satisfy \(L^2\ge l_fC_A(j,t)^2\), together
with the previously imposed lower bounds. A same-index
repetition uses \(l_f=l_c=l_j\) and the same fixed
inequality. Thus no choice is recursive in history depth.

The fine factor \(S_f\widetilde B_f\) still has a native
fine coefficient budget fixed before \(L\). Its convolution
with \(T\), with a maximum over output residues as in
Section~\ref{sec:calibrated-rows}, preserves the \(CL^{-2}\)
price in \(\Phi_F\); the coarse factors in \(\Phi_I\)
have price \(C\). Their constants may depend on the now
fixed adjacent profiles. The potential kinetic matrices
are unchanged because \(\widetilde H_vd_1=H_vd_1\).
The completion changes no Gaussian history, frozen-link
factor or normalisation in the joint integral.
\end{proof}

This proves the completion of \cite[Lemma~30.4, pp.~170--173]{Main},
with a power-series square factor in place of the square root
of the rebased defect.
It is needed because the actual coarse row vector below is
thinned and masked, and need not be an exact curvature.

\subsection{The clipped rows that remain in the physical action}

For a scale threshold \(h\), the prescribed direct price is
\[
 c_h(s)=\begin{cases}\lambda s^2/2,&s\le h,\\
                         \lambda_{\rm l}hs,&s>h,
          \end{cases}
 \qquad 0<\lambda_{\rm l}<\lambda/16.
\]
Its value at equality is the first branch; continuity at the
threshold is not required. The principal row square retains
\(k_{h,\varrho}(o;U)\), the original exponentially relaxed test
defined above. Use the prescribed path-compatible choice
\(\varrho^{-1}<\sigma/4\), where \(\sigma\) is a common
available native width for the adjacent principal rows.
It is chosen after those widths and before \(L,W\), as in
\cite[Lemma~13.10, pp.~52--53]{Main}.
Cubic and quartic terms retain their original
fine masks. The physical loop term likewise retains its masks,
as distinguished from \(I_f^{\rm un}\) in
Section~\ref{sec:smooth-principal}. Each masked degree-\(k\)
path term is bounded everywhere by its coefficient price times
\(h^k\). These statements describe the given action, rather
than introducing a new clipping rule in the integral.
At the literal Wilson boundary keep its actual uncut quadratic
direct price; it dominates all the fine lower bounds used below.

Let \(F_f(U;p)\) be the two-layer cut fine vector consisting of
\(\sqrt{\lambda_f}\mathcal C(U_p)\) and the principal row at
that anchor, cut at native length \(R_B\). Let \(c_f(U;p)\)
be its actual clipped direct price plus its tested row square.
Define \(F_c(V;p),c_c(V;p)\) at the coarse scale in the same
way. The finite symmetry copies and square-root weights are
included in these row indices. All uncut row tails are kept as
separate physical terms with price \(Ce^{-c_LR_B}\) per anchor.

Thin the coarse row vector as follows. Its principal coordinate
is retained only when its original row test holds; its direct
coordinate is
\[
       v_p^{\rm dir}=\sqrt{\lambda_c}
                    \min(1,h_c/d_G(V_p))\mathcal C(V_p),
\]
with multiplier one at chord zero. A retained principal row has
\[
 |\ell(V;o)|\le h_c\sum_j|r_j|e^{|\gamma_j|/\varrho}
                                                     \le C_{\rm p}h_c,
\]
by this path-compatible choice; a removed row is zero.
Thus \(\|v\|_\infty\le Ch\) when \(h_c\le2h\), including
rough coarse fields. No good-set bound is applied to an
unretained row here.
Transport and cut all three row maps at the same native length,
and put
\[
 G_F=\Phi_F(V)v,\qquad G_I=\Phi_I(V)v,\qquad
                         Z_b=\sqrt a\mathcal C(V_bQ_t(U)_b^{-1}).
\]
Their absolute prices give \(|G_F|\le CL^{-2}h\),
\(|G_I|\le Ch\), with constants chosen before \(L\).
Use these same guide rows in the source \(\alpha\) of
\eqref{eq:physical-selected-precision}. On healthy reads they
coincide with the calibrated coarse guide used previously.
This choice changes no exact change-of-density identity.

Define six physical row contributions:
\begin{align}
 w_F(p)&=c_f(U;p)+\tfrac12|G_F(p)|^2-\langle G_F(p),F_f(U;p)\rangle,
                                                        \notag\\
 w_I(b)&=\tfrac a2q_b^2+\tfrac12|G_I(b)|^2-\langle G_I(b),Z_b\rangle,
       \qquad q_b=d_G(V_bQ_t(U)_b^{-1}),\notag\\
 w_C(p)&=\langle v_p,F_c(V;p)\rangle-\tfrac12|v_p|^2-c_c(V;p),\notag\\
 w_L(p)&=\langle v_p,[\Phi_F^{\mathsf T}F_f+\Phi_I^{\mathsf T}Z-F_c]_p\rangle,
       \qquad w_\perp(x)=\tfrac12|(\Phi_\perp v)_x|^2,\notag\\
 w_D(p)&=\tfrac12\langle v_p,
       [1-\Phi_F^{\mathsf T}\Phi_F-\Phi_I^{\mathsf T}\Phi_I-\Phi_\perp^{\mathsf T}\Phi_\perp]v|_p\rangle.
 \label{eq:packet-physical-rows}
\end{align}
Products and transposes here are the actual cut transported
operators. Only the full flat operators satisfy
\eqref{eq:packet-row-completion}; its transported defect has
been written explicitly in \(w_D\).
Direct cancellation of the two mixed pairs and then of all
extra squares gives the exact identity
\begin{equation}
 \sum(w_F+w_I+w_C+w_L+w_\perp+w_D)
                      =\sum c_f+\tfrac a2\sum q_b^2-\sum c_c.
 \label{eq:packet-exact-row-identity}
\end{equation}
For example the remaining extra squares are one half of
\(\|\Phi_Fv\|^2+\|\Phi_Iv\|^2+\|\Phi_\perp v\|^2-\|v\|^2\)
plus the opposite quadratic form from \(w_D\).
No approximate isometry has been used in this identity.
It is the row accounting of
\cite[Lemma~30.5, p.~172]{Main}. We keep \(C,L,D\) together
at their common coarse anchor when allocating rows to packets.


\section{The joint hole-packet estimate}
\label{sec:joint-hole-packets}

We use the same physical integral and selected Gaussian as in
\eqref{eq:physical-selected-full-push}. Here \(\widehat\beta\)
is the fine nominal coupling,
\(h=\widehat\beta^{-1/2}P(\widehat\beta)\), and
\(\widehat\beta/2\le\beta=g^{-2}\le2\widehat\beta\).
The adjacent coarse threshold satisfies \(h_c/h\to1\).
Thus \(P^2/2\le\beta h^2\le2P^2\).
When both scales occur, take the larger of their adjacent
principal and loop coefficient budgets.
Every \(C_B\) below is a fixed polynomial \(C(1+B)^p\)
at the current profile; it never contains a volume factor.

\subsection{Physical seeds and the completed packet energy}

Choose a smooth \(\chi:[0,\infty)\to[0,1]\), equal to one
on \([0,1]\) and zero on \([2,\infty)\). Resolve unity using
\(\chi(4s_p/h)\) and its complement for each fine chord
\(s_p=d_G(U_p)\), and \(\chi(q_b/(Wh))\) and its complement
for each actual noise chord. Their radial expressions are smooth
where they vary and constant near zero. Exact coarse flags
remain their original Borel tests.
A seed record is a selected complementary fine test, a selected
complementary noise test, or an actual coarse bad site. Records
with the same location are still distinct. On a nonzero sector,
\begin{equation}
 \begin{array}{c|cc}
       &\text{selected}&\text{unselected}\\ \hline
 \text{fine}&s_p\ge h/4&s_p\le h/2\\
 \text{noise}&q_b\ge Wh&q_b\le2Wh.
 \end{array}
 \label{eq:packet-seed-tests}
\end{equation}
Weak bounds here include the boundaries of the original tests.
Only positive branch likelihoods require an estimate.

Project fine vertices to their coarse parents and noise bonds
to their endpoints. Join records at distance at most \(500B\).
A nonempty connected component of this graph is a geometric
packet \(p\); let \(m(p)\) be its number of records, and
\(p^{[r]}\) the coarse \(r\)-neighbourhood of their locations.
Every selected record in \(p^{[450B]}\) belongs to this packet.
Different packets have disjoint \(200B\)-neighbourhoods.

For this physical partition choose \(\mathcal H_\sigma\) in
Section~\ref{sec:physical-selected-Gaussian} to be the union of
seed locations, with the same \(10B\) frozen neighbourhood.
Outside it all fixed predictor patches have the default tests
in \eqref{eq:packet-seed-tests}. The block-tree path comparison
then gives
\begin{equation}
 d_G(U_eR_e(V)^{-1})\le C_{L,W}h.
 \label{eq:packet-true-coordinate-entry}
\end{equation}
For internal edges close the edge with its two block-tree paths;
the resulting loop is filled by a fixed number of fine
plaquettes. For crossing edges the same loop compares with the
predicted coarse bond, then its actual mismatch adds at most
\(2Wh\). The small plaquettes put both frame and predictor in
their regular branches by the earlier chart-entry calculation.
All these paths stay in the default-test neighbourhood.
Thus active links lie in the fixed chart eventually.
Chart complements outside the frozen set have zero contribution;
inside it one may keep the original chart factors or sum their
two choices exactly to one. No independent unpaid chart seeds
are needed. Gaussian-coordinate failures do not enlarge the
frozen set.

Assign each row of \eqref{eq:packet-physical-rows} to the packet
whose \(200B\)-neighbourhood contains its anchor. Keep
\(C,L,D\) together. Only fine and noise rows have an affine
comparison; for all other rows it is zero. Write these
comparison rows as \(D_ra+\alpha_r\), using the very same
source and frozen coordinates as in the exact Gaussian formula.
The completed energy is
\begin{equation}
 \mathcal R_p=\sum_{r\text{ assigned to }p}
       \{w_r-\tfrac12|D_ra+\alpha_r|^2
                            +\tfrac12|D_r\mu+\alpha_r|^2\}.
 \label{eq:packet-completed-physical-energy}
\end{equation}
It includes the negative affine square and the assigned part
of its completed-square constant. Equation
\eqref{eq:packet-exact-row-identity} identifies this with the
assigned exact phase after extracting the coarse clipped
quadratic action. Higher principal terms and loop terms keep
their original masks and are estimated separately below.

The frozen free edges belonging to \(p\) form a set \(J_p\).
Their deficit relative to the common Gaussian factor satisfies
\begin{equation}
 |J_p|\le C_L(1+B)^4m(p),\qquad
 \log s_g^{-|J_p|}\le C_B(1+\log\widehat\beta)m(p).
 \label{eq:packet-frozen-bill}
\end{equation}
There are \(C_L\) fine links per coarse cell, and each seed
freezes only its \(10B\)-neighbourhood. The second bound follows
from the exact \(s_g=\zeta_G(\sqrt{2\pi}g)^{\dim\g}\).
Enlarging a read because an incoming activity is long does not
freeze any additional edge.

\subsection{The physical energy supplied by the seeds}

Let \(\mathcal A_p\) be the sum of rows with anchors within
\(2B\) of a seed of \(p\), with the same coarse grouping.
There are at most \(C_Bm(p)\) such anchors. We now prove its
positive energy before estimating Gaussian integrals.

\begin{lemma}[Local seed energy]
\label{lem:packet-seed-energy}
Choose the completed row profiles first, then a sufficiently
large odd \(L\), then a sufficiently large fixed \(W\).
There are \(\kappa>0\) and a coupling threshold, uniform in
history depth and period at those profiles, such that
\begin{equation}
                              \mathcal A_p\ge\kappa h^2m(p)
 \label{eq:packet-positive-seed-energy}
\end{equation}
on the actual retained physical tests.
\end{lemma}
\begin{proof}
We separate large inputs, where only absolute coefficient sums
are used, from small inputs, where the full flat cancellations
apply. Constants multiplying \(L^{-2}\) in the large-input
bill are chosen before \(L,W\).

If an original fine row test fails, a witness has
\(s_q>h e^{d(o,q)/\varrho}\). The pointwise indicator bound
\(1_{\{\text{test fails at }o\}}\le
 \sum_{q:s_q>h}(s_q/h)e^{-d(o,q)/\varrho}\)
and the row and column coefficient sums give
\begin{equation}
 h\#\{\text{failed rows}\}
       +\sum_{\text{failed rows}}|\ell_{R_B}(U;o)|
                                      \le C\sum_{q:s_q>h}s_q.
 \label{eq:packet-dropped-rows}
\end{equation}
For the second sum split input chords at \(h\): the small
ones cost \(Ch\) per failed row, and the large ones use the
same-level column bound. The exponential cones have a finite
four-dimensional sum. Locally every witness remains in the
owned halo, because the original relaxed test has effective
reach \(O(\log\widehat\beta)\).

Complete squares in retained fine rows. The direct layer and
\(|G_F|\le CL^{-2}h\) give a fixed price
\(d_fh^2\) for each selected fine chord \(h/4\le s\le h\),
with a fixed \(d_f>0\) after choosing \(L\) and entry
(one may take \(d_f=\lambda_f/128\) for the unsplit direct layer).
For \(s>h\) use its actual coefficient
\(\lambda_{{\rm l},f}hs\); its cross term and
\eqref{eq:packet-dropped-rows} cost only \(CL^{-2}hs\).
Choosing that coefficient at most \(\lambda_{{\rm l},f}/4\)
leaves
\begin{align}
 \sum w_F&\ge\tfrac12\lambda_{{\rm l},f}h
                     \sum_{\text{fine seeds}:s>h}s
         +d_fh^2\#\{\text{fine seeds}:s\le h\}-C_Bh^4m(p),\notag\\
 \sum w_I&\ge\tfrac a8\sum_{\text{noise seeds}}q_b^2.
 \label{eq:packet-fine-noise-payment}
\end{align}
For noise, \(|Z_b|\le\sqrt a q_b\) makes an unselected
completed square nonnegative; \(q_b\ge Wh\) pays the cross
term on selected bonds when \(W\) is large.

The principal part of \(w_C\) is exactly zero, both when
retained and when removed. For its direct part choose a small
fixed \(\delta_0\). The curvature expansion and direct
thinning give
\begin{equation}
 w_C\ge\begin{cases}
    -Ch_c^4,&s\le h_c,\\
    (\lambda_c/4)h_cs,&h_c<s\le\delta_0,\\
    -\lambda_{{\rm l},c}h_cs,&s>\delta_0.
          \end{cases}
 \label{eq:packet-coarse-payment}
\end{equation}
Indeed on the middle interval its two inner-product terms are
\(\lambda_c(h_c/s-h_c^2/(2s^2))|\mathcal C(V_p)|^2\).
Use \(|\mathcal C(V_p)|^2\ge s^2-Cs^4\), shrink \(\delta_0\),
and use \(\lambda_{{\rm l},c}<\lambda_c/16\). For larger
chords those two terms are still nonnegative.

Set \(T=\widehat\beta^{-49/100}\). An exceptional record is
a fine or noise seed whose chord exceeds \(T\); let their
number be \(m_{\rm exc}\). A coarse row is hit if it is within
\(4B/5\) of such a record. The hit set \(Q_{\rm hit}\) has at most
\(C(1+B)^4m_{\rm exc}\) coarse row types. All its reads lie
inside the owned halo; enlarging it by their reach changes this
constant, not its power in the seed count.

Here is the large-input bookkeeping. On row Hilbert spaces use
the absolute row and column sums of block operator norms. For
one \(A\) chosen before \(L,W\), the actual cut transported
operators obey
\begin{equation}
 \begin{array}{c|cc}
       &\text{row sum}&\text{column sum}\\ \hline
 \Phi_F&AL^{-2}&AL^2\\
 \Phi_I&A&A\\
 \Phi_\perp&A&A.
 \end{array}
 \label{eq:packet-row-column-table}
\end{equation}
The fine column has \(L^4\) residues times its per-residue
\(CL^{-2}\) bound. Transport is isometric, cuts decrease
absolute sums, and periodisation cannot increase them.
There are same-level nonnegative kernels \(k^f,k^c\), each
with row and column sums at most \(A\), majorising \(F_f,F_c\)
by the corresponding plaquette chords.

Put \(e_q=s_q1_{\{\text{fine seed},\ s_q>h\}}\) and
\(n_b=q_b1_{\{\text{noise seed}\}}\). In all relevant reads
\(s_q\le h+e_q\), \(q_b\le2Wh+n_b\).
The macro circulation bound gives
\[
 s_q^c\le r_{L,W}h+H_q,\quad
 H_q=CL^{-2}\sum_{u\subset\mathcal M_q}e_u
                       +\sum_{b\in\partial q}n_b,
 \quad r_{L,W}=A(1+L^2+W),
\]
\[
 \sum_qH_q\le A\mathcal B_p,\qquad
                  \mathcal B_p=L^{-2}\sum_qe_q+\sum_bn_b.
\]
Only the fixed macro overlap and boundary incidence are used in
the last inequality. In the absence of exceptional inputs the
same macro bound gives \(s_q^c\le C_LT\); hence all coarse
plaquettes above \(\delta_0\) are hit eventually.

The uniform small input and the excess input require different
sums. For example,
\begin{align*}
 \left|\sum_{p\in Q_{\rm hit}}\langle v_p,(\Phi_F^{\mathsf T}F_f)_p\rangle\right|
 &\le A^3L^2h^2|Q_{\rm hit}|
    +Ah\sum_qe_q\sum_r k^f_{rq}\sum_{p\in Q_{\rm hit}}\|\Phi_F(r,p)\|\\
 &\le A^3L^2h^2|Q_{\rm hit}|+A^3L^{-2}h\sum_qe_q.
\end{align*}
The first term uses a column sum. In the second, fix \(r\)
first and use the fine row sum \(AL^{-2}\), followed by the
same-level column sum of \(k^f\). No residue count enters its
large-input coefficient. The noise and coarse terms similarly
give
\begin{align*}
 \left|\sum_{Q_{\rm hit}}\langle v,\Phi_I^{\mathsf T}Z\rangle\right|
       &\le2A^2\sqrt a Wh^2|Q_{\rm hit}|+A^2\sqrt a h\sum_bn_b,\\
 \left|\sum_{Q_{\rm hit}}\langle v,F_c\rangle\right|
       &\le A^2r_{L,W}h^2|Q_{\rm hit}|+A^3h\mathcal B_p.
\end{align*}
For the completeness defect, its absolute row norm is at most
\(1+3A^2\), since the fine product costs
\((AL^2)(AL^{-2})\). Thus its scalar contribution on \(Q\)
is at most \(\frac12 A^2(1+3A^2)h^2|Q_{\rm hit}|\).
Finally the negative coarse price above \(\delta_0\) is at most
\[
 2\lambda_{{\rm l},c}r_{L,W}h^2|Q_{\rm hit}|
                              +2\lambda_{{\rm l},c}A h\mathcal B_p.
\]
Combining these explicit bounds yields
\begin{equation}
 \sum_{Q_{\rm hit}}(w_L+w_D)-\lambda_{{\rm l},c}h_c
                       \sum_{s_q^c>\delta_0}s_q^c
 \ge-\frac{C_7}{L^2}h\sum_qe_q-C_7h\sum_bn_b
                    -C_{L,W}(1+B)^4h^2m_{\rm exc},
 \label{eq:packet-hit-budget}
\end{equation}
where \(C_7=2A^3+A^2\sqrt a+2\lambda_{{\rm l},c}A\)
is chosen before \(L,W\). Only the last constant contains
the positive powers of \(L,W\).

For an unhit row all fine and noise inputs in its enlarged
\(4B/5\)-box have chord at most \(T\).
The block-tree comparison first gives the actual relative
fine coordinates of size \(C_LT\). The macro bound gives
coarse plaquette chords \(C_LT\). Applying the axial layer
lemma on a coarse box of radius \(B/2\), then lifting that
gauge constantly on blocks, gives one simultaneous fine and
coarse logarithmic chart on the smaller \(2B/5\)-box, with
\[
                          \varepsilon_B=C_{L,W}(1+B)^{10}T.
\]
For the coarse gauge bound, its side times its total
\(O(B^4)\) plaquette cost is enough. The extra polynomial
reserve covers the fixed enlargements. The fine-to-replica
comparison is invariant under this lifted gauge. All logarithms
are the injective ones, since \(\varepsilon_B\to0\).
Extend these logarithmic arrays by zero outside the box. All
cut reads are unchanged, while their uncut tails have their
original exponential price.

The flat cancellations hold for arbitrary extended arrays;
\(Q\) in these identities is the scalar linear average of
Appendix~\ref{app:calibration}:
\begin{gather*}
 F_f^{(1)}=S_fd_fu_f,\quad F_c^{(1)}=S_cd_cu_c,\quad
 Z^{(1)}=\sqrt a(u_c-Q_0u_f),\\
 \Phi_{F,0}^{\mathsf T}F_f^{(1)}+\Phi_{I,0}^{\mathsf T}Z^{(1)}-F_c^{(1)}=0,
 \qquad 1-\Phi_{F,0}^{\mathsf T}\Phi_{F,0}
             -\Phi_{I,0}^{\mathsf T}\Phi_{I,0}-\Phi_{\perp,0}^{\mathsf T}\Phi_{\perp,0}=0.
\end{gather*}
The first follows from
\(\Phi_{F,0}^{\mathsf T}S_fd_f=S_cd_cQ\),
\(\sqrt a\Phi_{I,0}^{\mathsf T}=S_cd_c\), and \(d_cQ_0=d_cQ\).
The second is \eqref{eq:packet-row-completion}; it applies
to the actual bounded thinned \(v\), without assuming that
\(v\) is a curl or differentiating its masks.

For clarity, Taylor's formula in this common chart gives row
remainder \(C_L(1+R_B)^s\varepsilon_B^2
+C_Le^{-c_LR_B}\varepsilon_B\), and the row and column
perturbation of each transported lift has price
\(C_L\{(1+R_B)^s\varepsilon_B+e^{-c_LR_B}\}\).
Two derivatives of a native word cost at most a fixed multiple
of \((1+l)^2\) times its coefficient. The exponential budgets
sum this polynomial; discarded words gain \(e^{-c_LR_B}\).
The actual predictor uses its full framed derivative \(Q_0\).
All post-ratio factors \(L^2\) are allowed in \(C_L\) here.
Subtract the two displayed flat identities, using
\(K^{\mathsf T}K-K_0^{\mathsf T}K_0=(K^{\mathsf T}-K_0^{\mathsf T})K+K_0^{\mathsf T}(K-K_0)\).
Since \(e^{-c_LR_B}\le T\) eventually, this yields
\[
 |[\Phi_F^{\mathsf T}F_f+\Phi_I^{\mathsf T}Z-F_c]_p|\le C_BT^2,
 \qquad
 |[(1-\Phi_F^{\mathsf T}\Phi_F-\Phi_I^{\mathsf T}\Phi_I-\Phi_\perp^{\mathsf T}\Phi_\perp)v]_p|
                                                \le C_Bh(T+e^{-c_LR_B}).
\]
There are \(C_Bm(p)\) owned anchors. With \(|v|\le Ah\)
and \(h\le T\), the total unhit loss is therefore
\begin{equation}
          \sum_{\rm unhit}(w_L+w_D)
             \ge-C_B\{hT^2+h^2e^{-c_LR_B}\}m(p)= -o(h^2)m(p).
 \label{eq:packet-unhit-budget}
\end{equation}

We first record the numerical requests for the absorption:
\[
       C_7/L^2\le\lambda_{{\rm l},f}/4,\qquad W\ge16C_7/a,
\]
together with the previous noise requirement. The two excess terms in
\eqref{eq:packet-hit-budget} are paid by
\eqref{eq:packet-fine-noise-payment}. Their remaining prices
are at least \(\lambda_{{\rm l},f}hs/4\) for large fine
seeds, \(d_fh^2\) for the other fine seeds, and \(aq_b^2/16\)
for noise seeds. Gain rows \(w_\perp\) stay nonnegative.
Use a common energy reserve satisfying
\[
 0<\mu_0\le\min\{\lambda_{{\rm l},f}/16,d_f/4,
                                  a/64,\lambda_c/64,1\}.
\]
Writing the record counts as \(m_F,m_I,m_C\), the first two
payments give at least
\(2\mu_0h^2(m_F+m_I)+2\mu_0hT m_{\rm exc}\).
Let \(K\) count owned coarse plaquettes with
\(h_c<s_q^c\le\delta_0\). Each contributes at least
\(4\mu_0h^2\), using \(h_c\ge h/2\).
Every coarse bad site meets a bad plaquette. A plaquette above
\(\delta_0\) has an exceptional input by the macro bound;
its fixed overlap and boundary incidence allow an assignment
to such an input with at most \(A\) assignments per input.
Since a plaquette has four vertices,
\[
                              m_C\le4K+4A m_{\rm exc}.
\]
The exceptional reserve pays its assigned costs provided
\[
    \{4A\mu_0+C_{L,W}(1+B)^4\}h^2\le\mu_0hT.
\]
This request follows from \((1+B)^4h\ll T\), with \(L,W,\mu_0\) fixed.
The exceptional reserve thus pays both the hit-set bill and
the coarse flags assigned to exceptional inputs. The remaining
energy is at least \(\mu_0h^2(m_F+m_I+m_C)\).
The quartic errors and \eqref{eq:packet-unhit-budget} each
use at most a quarter of this after increasing entry: indeed
\[
 C_Bh^2\to0,\qquad
 C_BT^2/h=C_B\widehat\beta^{-12/25}/P\to0,\qquad
                              C_Be^{-c_LR_B}\to0.
\]
Collecting the requests, the choices have the form
\[
 \begin{gathered}
 C_7/L^2\ll\lambda_{{\rm l},f},\qquad C_7/a\ll W,\qquad
 0<\mu_0\ll\min\{\lambda_{{\rm l},f},d_f,a,\lambda_c,1\},\\
 (1+B)^4h\ll T,\qquad
 C_B\{h^2+T^2/h+e^{-c_LR_B}\}\ll\mu_0.
 \end{gathered}
\]
First fix the completed profiles and pre-block prices, then choose
\(L\), then \(W\) satisfying all its noise requirements, and then
\(\mu_0\). The thresholds in these first three comparisons give
the numerical bounds used above. After these choices, all
constants and polynomial degrees in the final two requests are
fixed. The prescribed scales give \((1+B)^4h/T\to0\) and the
three displayed error limits, so a single entry threshold meets
them. Its fixed-data dependence includes the overlap price \(A\);
it is independent of history depth and period.
This proves \eqref{eq:packet-positive-seed-energy} with
\(\kappa=\mu_0/2\). All sums used absolute cover prices
before any volume sum. Finite symmetry averaging preserves
the inequalities and only changes fixed type constants.
\end{proof}

This develops the seed estimate of
\cite[Lemma~30.6, pp.~173--181]{Main}. The two summation orders
in \eqref{eq:packet-row-column-table} preserve the pre-block
fine gain; constants depending on \(L,W\) are used only for
the subsequent coupling threshold.


\subsection{The annulus and the negative affine squares}

The seed estimate concerns actual nonlinear rows. We now compare
them with the affine rows subtracted in the very same Gaussian
representation; a comparison at a different frozen field would
not suffice.

\begin{lemma}[The completed energy retains the seed charge]
\label{lem:packet-annulus-comparison}
On the physical sector tests,
\begin{equation}
 \mathcal R_p\ge\mathcal A_p-e_\beta m(p),\qquad
 e_\beta=C_B(h^3+h^2e^{-c_LR_B}),\qquad
                      (1+B)^n\beta e_\beta\longrightarrow0
 \label{eq:packet-annulus-comparison}
\end{equation}
for every fixed \(n\). All physical coordinates and protected
frozen logarithms used in this comparison are the actual ones.
\end{lemma}
\begin{proof}
An operated row is anchored outside the \(8B\)-neighbourhood
of every seed. Its full cut reads, including the relaxed tests,
are farther than \(2B\) from every seed. They have fine chords
at most \(h/2\), noise chords at most \(2Wh\), and healthy
coarse plaquettes. The same tree-strip calculation as in
\eqref{eq:packet-true-coordinate-entry} gives actual relative
link logarithms bounded by \(C_{L,W}h\). Since the protected
radius is \(u=Bh\), its extension equals the true frozen
logarithm there eventually. No truncation of a frozen physical
link has occurred.

In a local coarse gauge, the rows at the replica and their
guides are \(C_Bh\). Expand the actual row at the replica in
these same active and frozen coordinates. Its Taylor remainder
is \(C_Bh^2\); the replica derivative differs from the
transported flat derivative by \(C_Bh\), whose action on the
coordinates costs another \(C_Bh^2\). Thus its completed row
differs from \(D_ra+\alpha_r\) by \(C_Bh^2\), while both
have size \(C_Bh\). Their squared difference is bounded by
\(C_Bh^3\). On each operated fine or noise row, therefore,
\[
                  w_r-\tfrac12|D_ra+\alpha_r|^2\ge-C_Bh^3.
\]
The direct chord-versus-curvature difference is nonnegative.
All row tests on this stencil have their original value one.

Outside the core, unoperated fine and noise rows are the same
nonnegative completed physical squares. The coarse group
\(C,L,D\) has loss at most \(C_B(h^3+h^2e^{-c_LR_B})\),
by the healthy local chart, its flat cancellations, and the
first line of \eqref{eq:packet-coarse-payment}. Gain rows
are nonnegative. Core rows have no affine subtraction, since
they are within \(2B\) of a seed. There are \(C_Bm(p)\)
assigned anchors. Finally the completed-square constants
\(\frac12|D_r\mu+\alpha_r|^2\) are retained in the exact
identity and may now be discarded for a lower bound.
This proves the first two assertions.
For the last use \(\beta h^3\le2P^2h\) and
\(\beta h^2e^{-c_LR_B}\le2P^2e^{-c_LR_B}\). Both vanish
faster than every inverse power of \(1+B\).
\end{proof}

This is \cite[Lemma~30.7, pp.~181--182]{Main}. In particular,
the negative affine square is paid before taking an absolute
Gaussian integral, rather than bounded as an unrelated
exponential.

\subsection{The local integration bills, with all factors retained}

Keep the common vacuum factor
\(e^{vf^4}z_{\rm n}(\beta)^{-b_c}s_g^{n_0}
(\det\mathsf M_0)^{-1/2}\) outside the packet expansion.
The coarse principal factor extracted by
\eqref{eq:packet-exact-row-identity} also keeps its specified
value; the later coupling reset is a different operation.
A packet multiplier, relative to these factors, has the form
\begin{equation}
 \Theta_p=1_p L_p J_p^{\rm Haar}s_g^{-|J_p|}
                e^{-\beta\mathcal R_p+\Delta_p^{\rm loc}}\mathsf M_p.
 \label{eq:packet-retained-multiplier}
\end{equation}
Here \(1_p\) contains its original sector and flag tests,
\(L_p\) its actual scaled local likelihoods,
\(J_p^{\rm Haar}\) its active Jacobians and broad chart factors,
and \(\mathsf M_p\) its assigned incoming and microscopic
factors. Every factor uses
\(U=\mathcal U_\sigma(V,Y,\mu+gz)\).
The integration law is the one product frozen Haar law, the
one product uniform tag prior, and the single selected Gaussian
\(\mu_{C_{\rm sel}}\). The likelihoods have not been moved
into a resampled tag law.

We specify determinant locality here to keep the packet covering
bound valid. In the determinant flow following
\eqref{eq:physical-selected-frozen-ratio}, the term
\(J_t(\mathsf M-1)\) has a fixed \(CB\) read window.
Assign its diagonal contributions near the packet, paired with
their flat contributions, to \(\Delta_p^{\rm loc}\).
All positive-defect chains in that flow, and all chains in
\(\log(1+S)\), retain their separate addresses and native
covering weights. Their incidence price is \(C_Be^{-cB}\),
by \eqref{eq:physical-selected-small-errors}, the Schur product
bound and the geometric series. They are additional local
letters with zero Gaussian charge. They are not bundled into
a geometric packet of bounded radius. This separation is an
exact expansion of the original determinant, including its flat
scalar, rather than an estimate replacing that determinant.

The remaining elementary bills are as follows:
\begin{equation}
 \begin{split}
 \log^+(L_pJ_p^{\rm Haar})+|\Delta_p^{\rm loc}|&\le C_Bm(p),\\
 \text{regular and microscopic bill}
    &\le C_B\{\delta+(1+\Lambda)h^2
                       +\beta C_{\rm p}h^3+\beta e^{-c_LR_B}\}m(p).
 \end{split}
 \label{eq:packet-local-bills}
\end{equation}
For the first inequality, a scaled likelihood is at most its
finite tag cardinality, the Haar density is bounded on the
fixed broad chart, and there are \(C_Bm(p)\) local factors.
A zero likelihood simply gives a zero integrand. A diagonal
determinant-flow term is bounded by the two uniform Schur
prices, integrated over an interval of length one; only
\(C_Bm(p)\) local diagonals are assigned. No determinant
volume outside the packet is charged here.

For the second inequality, use the regular value incidence
bound on its mask and the full-read enlargement estimate
\eqref{eq:reassembly-cover-enlargement}. Masked quadratic loops
cost \(\Lambda h^2\), and masked higher principal terms cost
\(\beta C_{\rm p}h^3\) and the smaller quartic price.
The original native row tails cost \(\beta e^{-c_LR_B}\).
Each neighbourhood has \(C_Bm(p)\) possible incidence sites.
An assigned long factor still retains its own full footprint
and coefficient price when covering norms are taken. The bill
is its summed value bound, not a replacement of that footprint.

Incoming forced occurrences have a particularly direct bound.
Every nonzero selected \(K_\alpha\) has a bad vertex in its
nonempty inventory. An incident fine plaquette has chord
greater than \(h\), so the corresponding fine sector test
must select that seed. For packets with total record count
\(m\), there are at most \(4m\) such fine vertices.
Assign the occurrence once to its first eligible vertex.
If \(w\) is the incoming forced incidence norm, then
\begin{equation}
 \sum_{\Gamma\text{ attached}}\prod_{\alpha\in\Gamma}k_\alpha
 \le\prod_{\alpha\text{ eligible}}(1+k_\alpha)
                                                        \le e^{4wm}.
 \label{eq:packet-incoming-occurrences}
\end{equation}
Original exclusions and all inventory tests remain in the signed
integral; dropping them here only increases an absolute majorant.
An occurrence reading two packets is still one occurrence and
joins those packets in a full component. It is not copied.
If one retains covering weights in this calculation, replace
its price by the corresponding full-read weighted price and
use the incoming exponent reserve first.

Equations \eqref{eq:packet-frozen-bill}--\eqref{eq:packet-incoming-occurrences}
show that one bill
\(H(1+B)^n(1+\log\widehat\beta)m\), for fixed sufficiently
large \(H,n\), pays all these factors on a prescribed bounded
profilewise ball of incoming norms. The small terms in
\eqref{eq:packet-local-bills} tend to zero after every fixed
polynomial loss in \(B\). These are the bills of
\cite[Lemmas~30.1--30.2, pp.~168--170]{Main}, with the long
determinant chains kept as individual letters.

\subsection{One joint Gaussian--Haar estimate}

\begin{theorem}[Actual joint packet payment]
\label{thm:packet-joint-payment}
Choose the completed profiles, then \(L\), then \(W\), and
finally the coupling threshold as above. There is \(c_0>0\),
uniform in period and legal history depth at those profiles,
such that for every fixed \(H,n\), eventually
\begin{equation}
 \beta\mathcal R_p\ge c_0P^2m(p)
                +H(1+B)^n(1+\log\widehat\beta)m(p).
 \label{eq:packet-paid-energy}
\end{equation}
No additional physical energy hypothesis is required.
For any finite set of packets, fixed \(r_0>0,u_0\ge0\), and
additional letters satisfying
\[
 |F_\gamma|\le A_\gamma
          e^{\frac12\sum_i d_{\gamma i}z_i^2},\quad
 d_{\gamma i}\ge0,\quad \sum_\gamma d_{\gamma i}\le u_0,
 \qquad C_{\rm sel}^{-1}\ge(r_0+u_0)1,
\]
their joint absolute integral, including the original Haar
variables, tags and the attached incoming selection sum, obeys
\begin{equation}
 \int\left|\prod_p\Theta_p\prod_\gamma F_\gamma\right|
 \le e^{-c_0P^2\sum_p m(p)}\prod_\gamma A_\gamma
        \exp\!\left\{\frac1{2r_0}\sum_{\gamma,i}d_{\gamma i}\right\}.
 \label{eq:packet-joint-integral}
\end{equation}
Here the integral notation includes that incoming sum, and
\(H,n\) are taken large enough for its local bills.
For a fixed incoming selection one can instead retain the
single product of its original occurrence prices on the right.
\end{theorem}
\begin{proof}
The seed and annulus estimates give
\(\beta\mathcal R_p\ge\kappa P^2m(p)/4\) eventually.
For each fixed \(n\),
\[
            (1+B)^n(1+\log\widehat\beta)/P^2\longrightarrow0,
\]
because \(B=O((\log\widehat\beta)^2)\) and
\(\log P=(\log\log\widehat\beta)^2\).
Taking, for example, \(c_0=\kappa/16\) leaves enough
reserve for any specified \(H,n\), proving
\eqref{eq:packet-paid-energy}.

Apply the energy inequality and all local bills pointwise on
the original retained tests. They give a bound uniform in the
actual frozen values and shared tags. Tonelli's theorem permits
the incoming sums on this nonnegative majorant. Their Haar and
tag laws retain probability mass one.
For the remaining Gaussian use
\eqref{eq:reassembly-joint-Gaussian-tilt} with the single
diagonal charge \(d_i=\sum_\gamma d_{\gamma i}\).
Its exact determinant ratio is bounded by the exponential in
\eqref{eq:packet-joint-integral}. This is one use of the
precision reserve for the whole product, without any
independence assertion between touching packets.
\end{proof}

This proves the packet estimate of
\cite[Theorem~30.8, p.~182]{Main}. A smaller fixed \(c_0\)
may be used to leave energy for further fixed polynomial
insertions. In particular a Gaussian monomial of total degree
\(q\), under the one tilted law, has absolute moment at most
\(C_qr_0^{-q/2}\); apply the one-dimensional Gaussian moment
bound and H\"older to its positive multiplicities. For a
polynomial multiply this by its absolute coefficient sum.
This moment bound is independent of the number of unused
Gaussian coordinates. A bound on the coefficient sum must
still use the full read footprint, as the following example
explains.

\subsection{A shared long incoming occurrence and covering payment}

Take two geometric packets \(p_1,p_2\), joined by one incoming
occurrence whose support reads both. Set \(m=m(p_1)+m(p_2)\)
and remove only that occurrence from \(\Theta_{p_j}\), writing
the remaining factors as \(\Theta_{p_j}^0\). Before any
factorisation their literal integral is
\[
 \mathcal I_\alpha=\int\dd\nu(t)\dd h_G^{J_\sigma}(Y)
        \dd\mu_{C_{\rm sel}}(z)\,
       \chi_{\rm rest}\Theta_{p_1}^0\Theta_{p_2}^0
                      K_\alpha(\mathcal U_\sigma)\prod_\gamma F_\gamma.
\]
The inventory indicator of \(K_\alpha\) and all original
fine-flag prohibitions are retained in this expression.
The frozen deficit is exactly
\(s_g^{-(|J_{p_1}|+|J_{p_2}|)}\); the bridge adds no frozen
links. After the joint physical payments, its absolute integrand
is bounded by
\(k_\alpha e^{-c_0P^2m}\prod_\gamma A_\gamma
 e^{\frac12\sum_i d_i z_i^2}\).
Writing \(P_0=C_{\rm sel}^{-1}\), the remaining integration is
exactly
\[
 \left(\frac{\det P_0}{\det(P_0-\operatorname{diag}d)}\right)^{1/2},
 \quad
 \log\left(\frac{\det P_0}{\det(P_0-\operatorname{diag}d)}\right)^{1/2}
       \le\frac1{2r_0}\sum_i d_i.
\]
The trace proof of this inequality requires no commutation of
the precision with the diagonal charge. Thus this example has
one \(k_\alpha\), one Gaussian tilt and the same common vacuum
normalisation as the full physical integral.

Let \(X\) be the component's full read, and
\(\widetilde Y_\alpha\) the enlarged bridge support.
Covering subadditivity gives
\[
 e^{\kappa_1 n_B(X)}k_\alpha
 \le e^{\kappa_1 n_B(X_{p_1})+\kappa_1 n_B(X_{p_2})}
                         e^{\kappa_1 n_B(\widetilde Y_\alpha)}k_\alpha
\]
when these are its reads; other letters retain their own
covering factors. A geometric packet cover costs
\(n_B(X_p)\le Cm(p)\), by its connected \(500B\)-seed
graph and fixed read halo. The bridge is paid by its incoming
covering reserve from \eqref{eq:reassembly-cover-enlargement}.
If a fixed-order insertion also has a coefficient price
\(C(1+B)^p\widehat\beta^a k^{p_q}m^q\), with
\(k=n_B(X)\), first use a reserved \(\eta>0\) to pay
\[
       k^{p_q}e^{\kappa_1k}\le
                  C_{p_q,\eta}e^{(\kappa_1+\eta)k}.
\]
Only afterwards may the remaining
\(C_{p_q,\eta}(1+B)^p\widehat\beta^am^q\) be absorbed by
unused \(P^2m\) energy. Its logarithm is at most
\(C[1+\log\widehat\beta+\log(1+B)]m+qm\).
There is no estimate of the bridge length in terms of \(m\).
This calculation gives the payment order used in
Section~\ref{sec:regular-output-derivatives}. Its density-score
calculation proves the required full-cover derivative coefficient
bound there, before consuming the packet energy reserve.

Finally, at a root site the number of geometric seed packets
of size \(m\) is at most \((C_B)^m\), after a polynomial
root-to-first-seed count. Encode a spanning exploration tree,
its depth-first parenthesis word, the bounded local record
types, and a displacement within \(500B\) at each step.
The root count is polynomial in \(B,m\) and is absorbed by
the same exponential bound. Therefore
\begin{equation}
 \sup_o\sum_{p:o\in X_p}e^{\kappa_1 n_B(X_p)}
                      e^{-c_0P^2m(p)}\le C_Be^{-c_1P^2}.
 \label{eq:packet-covering-income}
\end{equation}
Long determinant chains and incoming bridges keep their own
addressed prices when forming larger components. The joint
bound \eqref{eq:packet-joint-integral} supplies their product
majorants under the same precision budget. Once the remaining
letter incomes satisfy the connected criterion of
Section~\ref{sec:exact-masked-reassembly}, these estimates
justify its component integrability and its countable packet
sums. Disjoint full components factor by
Proposition~\ref{prop:reassembly-component-factorisation};
geometric packets joined by a read do not factor separately.

All constants were chosen before the UV cutoff or physical
volume. The local chart and window boxes fit in the period
reserve. Large packet neighbourhoods may wrap; union counts
and connected covering counts remain upper bounds on the torus.
Native cuts always count the original walk length. The slow
profile clock pays the finitely many additional entry thresholds
at each profile and uses the same Wilson family at all volumes.
This completes the physical joint packet payment. We next
derive the ordered regular germs of packet components under
this same integration law. Their scalar and quadratic jets
remain labelled for the later exact coupling reset.


\section{Regular germs after physical integration}
\label{sec:regular-output-derivatives}

An unforced coarse component can still contain fine hole packets
and bounded Borel incoming forced activities. Its integrated
amplitude must have the ordered derivatives of
\eqref{eq:local-physical-activity-norm}. Differentiating an
incoming factor after substitution of the selected centre would
require derivatives it need not possess. We instead hold its
physical fine links fixed and differentiate the smooth integration
density. This section proves the packet contribution to the
regular output family. The leading transport of a component
without a packet is treated separately in
Sections~\ref{sec:central-local-data}--\ref{sec:sharp-regular-transport}.
This is the scope of \cite[Section~31, pp.~187--194]{Main}.
The amplitudes below are relative to the prescribed coarse
principal factor and common vacuum normalisation; both retain
their exact places in the physical density.

\subsection{Retaining the physical charts before expanding a word}

Fix the word, sector geometry, all potential read sets, the
coupling and all thresholds before taking derivatives. At an
anchor \(o\), the allowed directions are
\(|A_e|_{\g}\le h_c(1+d(e,o))^s\). Derivatives are the ordered
left derivatives already defined, without commuting their order.
On a healthy coarse footprint the regular extension uses coarse
good indicators equal to one, the quadratic coarse clipping
formula, and the unthinned coarse curvature vector. These
unmasked formulas extend across equality in the original coarse
tests. The coarse mask is multiplied back only after taking the
derivatives. Fine masks, branch tests and incoming inventories
remain functions of the physical integration variable.
An unforced component has no mandatory coarse bad-site
inventory. A word imposing a coarse bad flag belongs to the
forced sector, rather than to these regular germs.

For the packet partition we may use a broad active chart of
width \(H=hB^{17}\), without changing a small-field threshold.
Indeed \eqref{eq:packet-true-coordinate-entry} gives
\(|a_e|\le C_{L,W}h<H\) on every nonzero original sector for
active links. Insert \(b_e(a_e)=\chi(|a_e|/H)\), with
\(\chi=1\) on \([0,1]\) and zero on \([2,\infty)\), before
expanding the sector. Its product is exactly one there.
Eventually \(2H\) lies in the fixed injective exponential
chart, where the original fixed chart factor is one. Its
product with \(b_e\) is therefore exactly \(b_e\).
This insertion is an identity on the original physical
integrand, not a change of \(h\), \(h_c\), or of a frozen set.

For a fixed selection \(\Gamma\) of physical Borel factors let
\(T(\Gamma)\) contain all active physical links they read,
including every physical read of a packet. Before expanding the
broad factors use the exact identity
\begin{equation}
 \left(\prod_{\alpha\in\Gamma}F_\alpha\right)\prod_e b_e
  =\left(\prod_{\alpha\in\Gamma}F_\alpha\right)
       \prod_{e\in T(\Gamma)}b_e
       \sum_{A\subset T(\Gamma)^c}\prod_{e\in A}(b_e-1).
 \label{eq:output-retained-chart-rule}
\end{equation}
Thus a chart-failure letter is forbidden on a retained physical
read. Overlapping reads retain one broad factor on their union,
not several copies. The selection query belongs to the same
full component as the factors answering it. Disjoint full
components have disjoint retained products, so the earlier
factorisation remains valid. For absolute majorants the retained
factors have cap one; dropping only the exclusion restriction
increases those majorants. No chart has been discarded in the
signed identity.

\subsection{The same Gaussian marginal in physical variables}

Let \(T\) be the retained active physical links in a component,
including its incoming reads. In a matrix index, identify a
link set with all its colour coordinates; put
\(D_T=(\dim_\R\g)|T|\). At fixed frozen links \(Y\) and shared
tags, use the original marginal
\(C_T=(C_{\rm sel})_{T,T}\), with mean \(\mu_T\).
This is a principal covariance, not the inverse of a restricted
precision. It is not a conditional covariance.
On the retained chart define
\[
 a_e(V,U)=\log(U_eR_e(V)^{-1}),\qquad
                  r_T=a_T-\mu_T,\qquad P_T=C_T^{-1}.
\]
The density with respect to probability Haar measure on \(G^T\)
is exactly
\begin{equation}
 k_T(V,U;Y,t)=
 \frac{e^{-r_T^{\mathsf T}P_Tr_T/(2g^2)}}
 {(\sqrt{2\pi}g)^{D_T}(\det C_T)^{1/2}
                     \prod_{e\in T}\zeta_Gj_G(a_e)}.
 \label{eq:output-physical-density}
\end{equation}
The affine Gaussian marginal has the numerator and Lebesgue
normalisation shown here, while
\(\dd h_G(U_e)=\zeta_Gj_G(a_e)\dd a_e\); right translation
by \(R_e(V)^{-1}\) preserves Haar measure. These two identities
prove the formula, including every constant. A numerator
Jacobian cancels its Haar denominator only if that numerator
factor is present in the selected word.

If the component also reads Gaussian coordinates \(R\) outside
\(T\), retain them. With \(S=T\cup R\), \(D_S=D_T+|R|\),
\(C_S=(C_{\rm sel})_{S,S}\), and
\(w=((a_T-\mu_T)/g,z_R)\), the mixed density is
\begin{equation}
 k_{T,R}=
 \frac{e^{-w^{\mathsf T}C_S^{-1}w/2}}
 {(2\pi)^{D_S/2}g^{D_T}(\det C_S)^{1/2}
                     \prod_{e\in T}\zeta_Gj_G(a_e)}.
 \label{eq:output-mixed-density}
\end{equation}
Its measure is probability Haar on \(G^T\) times Lebesgue
measure on \(z_R\). All frozen links retain their original
probability Haar measures. In particular the two blocks of
\(w\) still have their original cross covariances.

A component integral now has the form
\begin{equation}
 \mathcal I(V)=\int F(U,Y,t)G(V,U,Y,t,z_R)k_{T,R}
               \,\dd h_G^T(U)\dd z_R\dd h_G(Y)\dd\nu(t).
 \label{eq:output-fixed-physical-integral}
\end{equation}
The bounded Borel factor \(F\) contains the fine hard tests,
fine actions and incoming factors; its physical variables are
fixed. The smooth factor \(G\) contains the remaining smooth
exponents, retained broad cutoffs and smooth noise partitions.
When \(R\) is empty use \(k_T\) and omit \(\dd z_R\).
Unused frozen links and tags integrate to one. All density
traces below are over \(S\), rather than over the entire torus.

\begin{lemma}[Differentiation at fixed physical links]
\label{lem:output-Borel-differentiation}
At a fixed cutoff and volume, the unmasked integral
\eqref{eq:output-fixed-physical-integral} has a \(C^q\) germ
if its selected matrices retain their positive spectral bounds
on a neighbourhood, and its mixed Gaussian factors have the
strict common precision reserve of the preceding packet bound.
Ordered derivatives act on \(Gk_{T,R}\); they do not act on
\(F\). For the mixed case the Gaussian factors are those of
the priced physical words, whose fixed-order derivatives add
polynomial insertions as verified below.
\end{lemma}
\begin{proof}
The retained broad cutoffs have compact support inside the
injective chart. Extend their product with the chart density
by zero outside it. On a smaller background neighbourhood the
chart maps and inverse matrices are smooth with a fixed
positive gap. For \(R=\varnothing\), all physical variables
range over compact groups and the tag set is finite, giving an
integrable majorant for the finitely many derivatives at this
fixed cutoff.
For the mixed case the remaining Gaussian factors have, after
their single combined quadratic charge, a strictly positive
precision. Continuity leaves part of this reserve on a smaller
neighbourhood. The differentiated factors below add only
fixed-degree polynomials; these have an integrable Gaussian
majorant there. Apply the fundamental theorem of calculus to
difference quotients and dominated convergence, then repeat
in the specified order. Countable label families are first
truncated and subsequently restored using the uniform
majorants proved below. This argument proves existence of
derivatives; compactness by itself proves no cutoff-uniform
estimate.
\end{proof}

\subsection{Selected data with polynomial covering losses}

Let \(X\) be the full component footprint, including all
windows, frozen and tag queries, and put \(k=n_B(X)\ge1\).
A connected cover gives
\(D_S\le C_Bk\) and \(1+d(e,o)\le C_Bk\) on its reads.
Thus allowed directions have supremum at most
\(C_Bh_ck^s\) there. Only a polynomial loss in \(k\) is
needed for a packet. Operator bounds on each fixed principal
window suffice for this purpose.

\begin{lemma}[Local bounds for marginal matrices and means]
\label{lem:output-selected-data}
For every fixed \(q,s\), some finite \(p_q\) gives, for
\(0\le j\le q\),
\begin{gather}
 \|D_{A_1}\cdots D_{A_j}C_S\|_{2\to2}
       +\|D_{A_1}\cdots D_{A_j}C_S^{-1}\|_{2\to2}
                                               \le C_Bk^{p_q},
 \label{eq:output-marginal-matrix-jets}\\
 \|D_{A_1}\cdots D_{A_j}\mu_T\|_\infty
       +\|D_{A_1}\cdots D_{A_j}a_T\|_\infty
                                               \le C_Bh k^{p_q}
 \label{eq:output-mean-chart-jets}
\end{gather}
on the retained broad support and healthy coarse germs.
Constants are independent of the shape and cardinality of
\(S\), the period and the terminal cutoff; their size dependence
is the displayed polynomial in the full covering size.
\end{lemma}
\begin{proof}
Restrict directions to the declared coarse reads. An ordered
derivative of a transported native word chooses positions
along that word. Its order-\(j\) absolute price is bounded
by \(C_j(1+l)^j\) times its coefficient and the direction
suprema. The native exponential reserve sums this polynomial.
Fixed sector geometry is not differentiated. The row and
column Schur bounds therefore control all derivatives of
\(\mathsf M\), and hence of any principal matrix, in operator
norm without a volume sum.

For an invertible matrix \(A\), write \(D_i=D_{A_i}\).
The ordered inverse formula is
\begin{equation}
 D_1\cdots D_j A^{-1}
  =\sum_{b=1}^j(-1)^b
       \sum_{(B_1,\ldots,B_b)}
       A^{-1}(D_{B_1}A)A^{-1}\cdots(D_{B_b}A)A^{-1}.
 \label{eq:output-ordered-inverse-formula}
\end{equation}
The inner sum is over ordered partitions of the derivative
positions into nonempty blocks; inside a block retain their
original order. This follows inductively from
\(D(A^{-1})=-A^{-1}(DA)A^{-1}\), with no commutation of
derivatives. Every principal inverse has the same operator
gap bound. Thus each differentiated selected entry of \(J\)
has a dimension-independent cap. Each row and column of
\(J\) has only \(C_B\) possible nonzero entries, because its
window range is \(B/4+2\). The Schur test gives the required
operator derivative bounds for \(J\) and \(C_{\rm sel}\).
Restriction gives them for \(C_S\). The covariance spectral
gap is inherited by every marginal; applying
\eqref{eq:output-ordered-inverse-formula} again gives the
bound for \(C_S^{-1}\). This uses no number-of-windows factor.

For the source use the smooth protected logarithm fixed in
Section~\ref{sec:physical-selected-Gaussian}. At radius
\(u=Bh\), an order-\(j\) derivative of its radial truncation
costs \(C_ju^{1-j}\). The variation supplies \(h_c^j\), so
the scaled bound is \(C_Bh k^{p_q}\), including its transition
region. The replica rows and the unmasked coarse guide have
the same bound: on their healthy local reads their values are
\(O(h_c)\), and their scaled positive-order derivatives have
the native word budgets. The regular predictor source is
analytic there. It follows that \(b=D^{\mathsf T}\alpha\) has this
maximum-norm derivative bound uniformly in frozen \(Y\) and
tags. Now use \(\mu=-J^{\mathsf T}b\) and the finite-range column sums
just proved.
On the retained chart \(|a_e|\le2H=C_Bh\). The inverse
exponential chart and the replica formula give the same bounds
for its scaled derivatives at fixed physical \(U_e\).
This proves \eqref{eq:output-mean-chart-jets}.
\end{proof}

The smaller Gaussian correction prices also survive
differentiation. Differentiate the principal inverse comparison
in the proof of Proposition~\ref{prop:physical-selected-law}:
every term still crosses the selected window boundary. The
weighted inverse envelopes and the native derivative budgets
therefore give \(C_qe^{-cB}\), at a reduced width, in either
matrix index. The identities
\[
 r=-(2S-E)b,\qquad C_{\rm sel}^{-1}-\mathsf M=-S C_{\rm sel}^{-1}
\]
then preserve the small residual coefficients and their
derivatives. The same proof applies to \(\mathsf M_t\)
uniformly in the determinant interpolation parameter.
Derivatives of a chain add a fixed polynomial in its native
length, paid by its spare exponential width. Hence a
differentiated Gaussian correction or determinant letter keeps
its coefficient price, up to the fixed-order polynomial
allowances, rather than becoming an uncontrolled extensive
error. The operator-norm calculation above is sufficient for
the marginal scores; these boundary estimates are used only to
retain the correction-letter incomes.


\subsection{The first and second ordered density variations}

The score is the logarithmic derivative of the integration
density, not a derivative of a physical Borel multiplier.
For \(\ell=\log k_T\), its first variation at fixed
\(U,Y,t\) is
\begin{align}
 D_A\ell={}&-\tfrac12\operatorname{Tr}(P_TD_AC_T)
       +g^{-2}\langle P_Tr_T,D_A\mu_T-D_Aa_T\rangle\notag\\
 &+\tfrac1{2g^2}\langle P_Tr_T,(D_AC_T)P_Tr_T\rangle
                         -\sum_{e\in T}D_A\log j_G(a_e).
 \label{eq:output-first-score}
\end{align}
The coupling, dimension and Haar constant are fixed, so their
logarithms contribute no derivative. The mean variation and
the variation of the physical logarithm are both present.

For the next calculation only write \(C=C_T\), \(K=C^{-1}\),
\(r=r_T\), and \(J=\sum_{e\in T}\log j_G(a_e)\).
A subscript \(AB\) means \(D_AD_B\) in that order; for
example \(r_{AB}=a_{AB}-\mu_{AB}\). Matrix inversion gives
\[
 K_A=-KC_AK,\qquad
 K_{AB}=KC_AKC_BK+KC_BKC_AK-KC_{AB}K,
\]
and
\[
 D_AD_B\log\det C=\operatorname{Tr}(KC_{AB}-KC_AKC_B).
\]
Symmetry of \(K\) and its first derivatives, without exchanging
the two directions, gives the complete quadratic variation
\[
 \begin{split}
 D_AD_B(r^{\mathsf T}Kr)={}&2r_{AB}^{\mathsf T}Kr+2r_B^{\mathsf T}Kr_A
             +2r_B^{\mathsf T}K_Ar+2r_A^{\mathsf T}K_Br+r^{\mathsf T}K_{AB}r.
 \end{split}
\]
Substitution yields
\begin{equation}
 \begin{split}
 \ell_{AB}={}&\tfrac12\operatorname{Tr}(KC_AKC_B-KC_{AB})
       -g^{-2}r_{AB}^{\mathsf T}Kr-g^{-2}r_B^{\mathsf T}Kr_A\\
 &+g^{-2}r_B^{\mathsf T}KC_AKr+g^{-2}r_A^{\mathsf T}KC_BKr\\
 &+\tfrac1{2g^2}r^{\mathsf T}
       (KC_{AB}K-KC_AKC_BK-KC_BKC_AK)r-J_{AB}.
 \end{split}
 \label{eq:output-second-score}
\end{equation}
The two different covariance products and the ordered mean
and chart variations are all required. The Haar term is
\[
 J_{AB}=\sum_{e\in T}
    \{\dd^2\log j_G(a_e)[a_{e,A},a_{e,B}]
                       +\dd\log j_G(a_e)[a_{e,AB}]\},
\]
where \(\dd,\dd^2\) are ordinary Lie-algebra differentials.
An original numerator \(j_G(a_e)\), if present, cancels its
denominator and their product-rule derivatives; an absent
numerator cannot be cancelled.

For the density itself and for its smooth multiplier we obtain
\begin{gather}
 D_Ak_T=k_T\ell_A,\qquad
 D_AD_Bk_T=k_T(\ell_{AB}+\ell_A\ell_B),\notag\\
 D_AD_B(Gk_T)=k_T\{G_{AB}+G_A\ell_B+G_B\ell_A
                                      +G(\ell_{AB}+\ell_A\ell_B)\}.
 \label{eq:output-second-density-product}
\end{gather}
In particular the product of the first scores is a density
term and must not be omitted after estimating the second log
score.

To include the cutoff cross terms explicitly, write a smooth
word as \(G=e^\Psi\prod_{\lambda\in\mathcal Q}b_\lambda\),
where \(\mathcal Q\) is its fixed finite list of smooth factors.
It may contain broad cutoffs, smooth noise partitions and Haar
numerators. Then
\begin{equation}
 \begin{split}
 G_{AB}=e^\Psi\bigg\{&
  (\Psi_{AB}+\Psi_A\Psi_B)\prod_\lambda b_\lambda\\
 &+\Psi_A\sum_\lambda b_{\lambda,B}\prod_{\tau\ne\lambda}b_\tau
  +\Psi_B\sum_\lambda b_{\lambda,A}\prod_{\tau\ne\lambda}b_\tau\\
 &+\sum_\lambda b_{\lambda,AB}\prod_{\tau\ne\lambda}b_\tau
  +\sum_{\lambda\ne\tau}b_{\lambda,B}b_{\tau,A}
                          \prod_{\xi\ne\lambda,\tau}b_\xi\bigg\}.
 \end{split}
 \label{eq:output-smooth-word-second-variation}
\end{equation}
These formulas use derivatives of each cutoff on its own
support. They never divide by a cutoff that may vanish.
The first variation of one left vector field along the other
flow is included in \(a_{AB},\mu_{AB},C_{AB}\) and the smooth
factor derivatives. They are not symmetric Taylor coefficients.
The fixed physical Borel factor \(F\) multiplies all these
terms without being differentiated. This is the calculation
of \cite[Section~31.3, pp.~189--191]{Main}, here with the same
marginal and physical Haar denominator throughout.

\subsection{Uniform score estimates and the remaining Gaussian coordinates}

\begin{lemma}[Scores cost a full-cover polynomial]
\label{lem:output-score-bounds}
For fixed \(q,s\), at every healthy coarse background and on
the retained broad support, some finite \(p_q\) gives
\begin{align}
 |D_{A_1}\cdots D_{A_j}\log k_T|
              &\le C_B(1+P^2)k^{p_q},&&1\le j\le q,\notag\\
 |D_{A_1}\cdots D_{A_j}k_T|
              &\le k_T[C_B(1+P^2)]^q k^{p_q},&&0\le j\le q.
 \label{eq:output-pure-score-price}
\end{align}
For the mixed density the second bound is replaced by
\begin{equation}
 |D_{A_1}\cdots D_{A_j}k_{T,R}|
       \le k_{T,R} C_Bk^{p_q}(1+P^2+Q_R)^q,
                      \qquad Q_R=\sum_{i\in R}z_i^2.
 \label{eq:output-mixed-score-price}
\end{equation}
Here the unindexed \(P=P(\widehat\beta)\) remains the
subpower scale, distinct from the precision \(P_T\).
There is no exponential cost in \(|T|\) or \(|R|\).
\end{lemma}
\begin{proof}
Use Lemma~\ref{lem:output-selected-data}. A trace costs at
most \(D_T\) times the operator norm of its matrix product.
A quadratic term costs at most
\(D_T\|u\|_\infty\|M\|_{2\to2}\|v\|_\infty\).
Thus in \eqref{eq:output-first-score} and
\eqref{eq:output-second-score} the traces cost a polynomial
in \(k\); every residual quadratic costs
\(C_Bg^{-2}h^2k^{p_q}\le C_BP^2k^{p_q}\).
The Haar sum has only \(|T|\le C_Bk\) terms and bounded
ordinary derivatives on the fixed chart.

At higher fixed order, derivatives of the log determinant
are traces of finitely many inverse and covariance derivative
products. The residual quadratic still has exactly two vector
positions, with ordered derivatives in either position. The
inverse partition formula bounds every matrix product, and
the same dimension count bounds every trace and quadratic.
This proves the log-score estimate. The density identity
\begin{equation}
 k_T^{-1}D_1\cdots D_jk_T
       =\sum_{\pi\in\operatorname{Part}(j)}
                           \prod_{B\in\pi}D_B\log k_T
 \label{eq:output-density-partitions}
\end{equation}
then proves the second estimate. Within each block keep the
original derivative order. The number of partitions depends
only on \(q\).

For \(k_{T,R}\) make the same calculation with
\(w=((a_T-\mu_T)/g,z_R)\). Only its \(T\)-part has
derivatives at fixed physical variables. Its norm squared is
at most \(C_B P^2 k\); the remaining norm squared is \(Q_R\).
All mixed products are bounded by \(2|uv|\le u^2+v^2\).
The matrix bounds for \(C_S^{-1}\) give one quadratic
polynomial \(C_Bk^{p_q}(1+P^2+Q_R)\) for each log score.
The partition formula gives \eqref{eq:output-mixed-score-price}.
No independence between the two blocks is used.
\end{proof}

To integrate the extra polynomial, use exactly the tilted law
already charged in \eqref{eq:packet-joint-integral}. Its
precision is the original marginal precision minus the sum of
all diagonal charges, and is at least \(r_0 1\). Linear
exponential letters have already been bounded by this common
quadratic majorant; this tilted law is centred. For every
integer \(j\ge1\), H\"older and the one-dimensional Gaussian
moment give
\begin{equation}
 \mathbb E_{\rm tilted}Q_R^j
                    \le(2j-1)!!\,r_0^{-j}|R|^j.
 \label{eq:output-mixed-score-moments}
\end{equation}
Indeed expand \(Q_R^j\) as a sum over \(j\) coordinate
positions and use H\"older with exponent \(j\) for each
square; each coordinate has variance at most \(r_0^{-1}\).
For example
\[
 \mathbb E_{\rm tilted}(1+P^2+Q_R)^2
   \le(1+P^2)^2+2(1+P^2)r_0^{-1}|R|+3r_0^{-2}|R|^2.
\]
Since \(|R|\le C_Bk\), the general fixed-order moment costs
only \(C_B(1+P^2)^q k^{p_q}\). The determinant ratio is
the one already paid in the original joint tilt. No second
precision reserve or Gaussian copy is introduced. After the
pointwise estimate on retained broad supports one may convert
back to \(z_T\), bound those supports by one, and apply this
same joint Gaussian moment bound.

The germ is local also at a coarse threshold equality. Use
the fixed unmasked formulas where all declared principal
matrices and covariance marginals retain half their healthy
spectral bounds and where the replica source is regular.
This open set contains every healthy background. If an
everywhere defined Borel extension is needed, multiply by
smooth spectral and chart protectors constant one near the
healthy data and extend by zero. They depend only on the
declared matrices and local reads, and are invariant under
their orthogonal gauge conjugations. Their derivatives on
the healthy set vanish, so none of the displayed norms is
changed. This constructs genuine neighbourhood germs without
differentiating a coarse flag.


\subsection{Derivatives of the actual smooth factors}

At fixed physical \(U,Y,t\), the fine actions, fine masks,
incoming activities and branch likelihoods do not depend on
\(V\). The actual \(Q_t(U)\) is held fixed, even where it is
only Borel in \(U\). The noise action is nevertheless smooth
in its coarse argument. For example,
\[
 |D_A\{\tfrac a2d_G(V_bQ_t(U)_b^{-1})^2\}|
                                                   \le C_a|A_b|.
\]
Each further ordered derivative inserts another bounded group
matrix and another direction. Summing the finite component
reads gives a bound \(C_B\beta h_c k^{p_q}\) for each
positive-order derivative of its noise exponent. This bound
also suffices for the physical mixed row terms: their rough
fine factors are bounded at fixed \(U\), while a differentiated
coarse guide supplies a scaled \(h_c\). The smooth affine
squares cost only \(C_BP^2k^{p_q}\). Masked fine loop and
higher terms stay in the Borel multiplier; the corresponding
coarse germs have their existing path-budget bounds.

Noise partition derivatives cost
\(C_q(h_c/(Wh))^j\) times their direction weights. Their
supports retain the selected inequality \(q_b\ge Wh\) or
the unselected inequality \(q_b\le2Wh\) from
\eqref{eq:packet-seed-tests}. For a derivative on a broad
factor the ratio is \(h_c/H\); for the protected logarithm
it is \(h_c/u\). These are bounded by fixed polynomials
in \(B\) at every fixed order. A differentiated cutoff is
bounded by its derivative cap and support, with its
undifferentiated partners still present. In particular a
factor which vanishes is never used as a denominator.

For the actual Gaussian correction letters, the residual and
precision identities following Lemma~\ref{lem:output-selected-data}
preserve their small coefficient budgets. Differentiating their
linear or quadratic exponent in the mixed coordinates adds a
fixed-degree polynomial, with the same coefficient price and
a polynomial in \(P,k\). Differentiating a Mayer factor
\(e^v-1\) at positive order leaves \(e^v\) times products of
derivatives of \(v\). The exponential cap of the original word
is unchanged. At most \(q\) differentiated positions add
polynomial insertions, which are integrated under the already
charged tilted law. Determinant chains have no Gaussian
quadratic charge. Native-length factors arising from their
derivative positions use the spare native width. Thus the
following proposition concerns the constructed matrix letters,
not arbitrary smooth factors with only a value bound.

\begin{proposition}[A packet pays the derivatives after integration]
\label{prop:output-packet-derivatives}
Fix \(q,s\). Consider an unforced physical component of
Sections~\ref{sec:exact-masked-reassembly}--\ref{sec:joint-hole-packets}
with \(m\ge1\) packet seed records and full footprint \(X\),
\(k=n_B(X)\). Keep the physical tests required by the joint
packet estimate and its common precision budget. On the
healthy coarse footprint its integrated unmasked amplitude
has a \(C^q\) germ. In the joint word majorant its derivative
norm costs at most
\begin{equation}
             [C_B(1+P^2+\beta h_c)]^q k^{p_q}
 \label{eq:output-raw-packet-jet-price}
\end{equation}
in addition to the original letter prices, with a fixed
inflation per letter allowed. Consequently its packet factor
in that norm can be bounded by
\begin{equation}
                         C_Bk^{p_q}e^{-c_2P^2m}
 \label{eq:output-paid-packet-jet-price}
\end{equation}
for a fixed \(c_2>0\), without a new Gaussian precision
charge. The full footprint polynomial has not been replaced
by a seed-count polynomial.
\end{proposition}
\begin{proof}
Retain broad factors on every physical Borel read by
\eqref{eq:output-retained-chart-rule}. Convert these variables
to physical Haar measure using \eqref{eq:output-physical-density}
or \eqref{eq:output-mixed-density}. All physical reads of the
packet phase and its retained tests are included in \(T\).
Lemma~\ref{lem:output-Borel-differentiation} applies to their
fixed unmasked coarse formulas. Derivatives act on the smooth
density, affine data, determinant letters, smooth noise action
and smooth partitions; they do not act on incoming Borel
values or fine hard masks.

The score and smooth-factor estimates give
\eqref{eq:output-raw-packet-jet-price}. At second order all
five terms in \eqref{eq:output-second-density-product} and
all terms in \eqref{eq:output-smooth-word-second-variation}
are included. At order \(q\) there are only \(q\)
differentiated positions, but there may be arbitrarily many
letter species on one footprint. If the word has \(r\)
letters, their placement factor is at most a fixed multiple
of \(r^q\). For any reserved \(\eta>0\),
\[
                r^q\le C_{q,\eta}e^{\eta r}.
\]
Inflate each letter price by \(e^\eta\) before applying the
connected criterion. No inequality \(r\le C_Bk\) is assumed.
Polynomial native-length costs are paid by their separate
native reserve, rather than by the number of distinct links.

Each differentiated partition still has the weak physical
inequalities required in the core and annulus estimates.
The remaining retained tests are unchanged. Thus
Theorem~\ref{thm:packet-joint-payment} applies term by term to
the original physical exponent, with some \(P^2m\) reserve
left. The fine Borel weight stays multiplied by that smooth
exponent throughout the energy estimate. Polynomial insertions
in the unconverted Gaussian
coordinates use \eqref{eq:output-mixed-score-moments} under
that same tilted marginal. Its \(|R|^q\) cost is included
in \(k^{p_q}\). The determinant ratio and the sum of quadratic
charges are exactly the previous ones. This verifies also
the Gaussian majorants used in the differentiability lemma
for these mixed words.

Finally, at fixed profile and order,
\[
 \log[C_B(1+P^2+\beta h_c)]^q
                \le C_{q,L}(1+\log\widehat\beta+\log(1+B))
                                                        =o(P^2).
\]
Since \(m\ge1\), any fixed fraction of its unused packet
energy absorbs this coefficient. This proves
\eqref{eq:output-paid-packet-jet-price}.
For summation retain the full-cover polynomial and pay it
first by
\[
 k^{p_q}e^{\kappa k}
              \le C_{p_q,\eta}e^{(\kappa+\eta)k},\qquad
 C_{p_q,\eta}=\sup_{x\ge1}x^{p_q}e^{-\eta x}<\infty.
\]
In particular a long incoming bridge is paid by its own
covering reserve. Its length is not bounded by \(m\), and
is not paid by \(P^2m\). The remaining geometric packet
sum is \eqref{eq:packet-covering-income} at the larger
covering exponent. These uniform majorants justify the
countable component sums and their ordered derivatives.
\end{proof}

This is the rare-component conclusion of
\cite[Proposition~31.3, pp.~193--194]{Main}. The argument spends
a packet's physical energy. Applying its rough score price
to a component without a packet would lose the leading
small-field order and would not prove the sharp regular
transport estimate.

\subsection{The packet part of the regular logarithm}

We now pass from physical words to integrated component species,
then to the ordinary hard-core logarithm. Use the factorial
jet seminorm \eqref{eq:connected-algebra-norm}; its product
constant is one. The whole-product estimate
\eqref{eq:connected-one-reserve} changes the anchor once for
the entire cluster and pays its full-cover polynomial with
one reserve. Thus the logarithm does not acquire one
\(2^q\) factor per vertex.

At the physical word level,
\eqref{eq:output-paid-packet-jet-price} and
\eqref{eq:packet-covering-income} give marked packet income
\(C_Be^{-c_2P^2}\), after paying its cover and native reserves.
Combining these marked prices with the remaining inflated
letter prices uses Lemma~\ref{lem:connected-positive-tree}.
When that connected criterion holds, the integrated species
which contain a packet therefore have derivative income
\(f_q\le C_Be^{-cP^2}\) at the stronger exponent. A packet
shared by several incoming reads is still counted through
its one physical component integral.

For completeness the marking step in the logarithm can be
proved directly from its already established positive
majorant. Suppose the unmarked species have derivative
incidence income \(b_q\), the marked species have income
\(f_q\), both with the covering reserve of
Lemma~\ref{lem:connected-log-bound}, and
\begin{equation}
                         D_Bb_q\le1/4,\qquad D_Bf_q\le1/4.
 \label{eq:output-marked-log-criterion}
\end{equation}
If \(f_q=0\) there is no marked contribution. Otherwise
multiply each marked price by
\(t=(4D_Bf_q)^{-1}\ge1\). The modified total income is
\(b_q+tf_q\le1/(2D_B)\), so the ordinary logarithm majorant
applies. Every ordered Ursell tuple containing a marked
species gains at least one factor \(t\). Dividing its absolute
majorant by \(t\) bounds the marked logarithmic income by
\begin{equation}
                   t^{-1}C_B(b_q+tf_q)\le2C_Bf_q.
 \label{eq:output-marked-log-price}
\end{equation}
Repeated species retain their usual Ursell factorials. This
is a marking of the logarithm's nonnegative majorant, not
a second distinct-label component expansion.

\begin{corollary}[The integrated packet contribution is a regular family]
\label{cor:output-packet-regular-family}
Fix \(q,s,\kappa\) and the stated covering and native reserves.
Suppose the other integrated unforced species have the required
ordered derivative incomes in
\eqref{eq:output-marked-log-criterion}. The logarithmic clusters
containing a physical packet have unmasked \(C^q\) germs on
their healthy coarse footprints, including threshold equality.
For their signed action functions \(f_\alpha=-\ell_\alpha\),
with fixed anchors \(o_\alpha\),
\begin{equation}
 \sup_o\sum_{\alpha:o_\alpha=o}e^{\kappa n_B(X_\alpha)}
       \|f_\alpha\|_{q,h_c,s;o_\alpha,X_\alpha}
                                                    \le C_Be^{-c_3P^2}.
 \label{eq:output-packet-regular-family}
\end{equation}
Flat subtraction gives \(E_\alpha^{\rm hole}=f_\alpha-f_\alpha(1)\)
with the same order of bound and \(E_\alpha^{\rm hole}(1)=0\).
Its extensive vacuum contribution is kept explicitly; its
absolute value per coarse site is at most \(C_Be^{-c_3P^2}\).
\end{corollary}
\begin{proof}
Attach the component norm prices just proved, including the
fixed letter inflation and the full-cover polynomial payment.
The marking argument gives
\eqref{eq:output-marked-log-price}. Use the factorial jet
product rule and the whole-cluster anchor estimate inside
the existing stronger covering reserve. The rooted logarithm
proof then gives the stated anchored norm, with a cutoff- and
volume-independent majorant for all ordered derivatives.
It permits the differentiated series and the countable
species limits on the regular germs. Subtracting a flat value
costs at most twice the raw norm and no derivative price.
Summing their values at a fixed anchor bounds the vacuum
scalar per site; it does not make that total scalar uniformly
bounded in the volume.

Collect the finite colour contractions, symmetry copies and
shared tag sums with the same laws as in the physical
representation. Read sets include complete colour spaces.
The selected matrices transform by orthogonal gauge
conjugation, the broad and protected cutoffs are radial, and
physical Haar measure is invariant. Consequently the
integrated species and their logarithmic germs are gauge
invariant. The local spectral extension described above
preserves this covariance and agrees on overlapping healthy
germs. No hard coarse flag has been differentiated.
\end{proof}

This supplies the derivative-norm part of
\cite[Corollary~31.4, p.~194]{Main}. We keep the packet
family's scalar and quadratic jets as labelled terms.
Their generic affine-Hessian extraction and exact masked
compensation are proved in Section~\ref{sec:affine-local-reset};
their principal coefficients are retained for the subsequent
coupling reset. No quadratic term has been discarded here.

The packet components now have both their joint physical
integral price and their regular germs. Borel incoming forced
activities required no derivatives. The other unforced derivative incomes required above are
supplied by the packet-free central estimates and sharp incoming
transport in the following sections. Their leading orders are
proved independently of rare packet energy. Principal-action
matching and assembly of the quadratic coupling reset remain
subsequent one-step units.


\section{Local central data without a physical packet}
\label{sec:central-local-data}

A central component has its full read set disjoint from the
physical seeds of Section~\ref{sec:joint-hole-packets}. It
therefore reads no frozen link. This statement concerns the full
reads, including the rows defining its selected windows, rather
than just the support of its displayed function. Its integration
law is still the original marginal of \(C_{\rm sel}\); intersecting
insertions keep their common variables and shared tags.

The packet estimate spent an amount of energy proportional to
\(P^2\). There is no such income here. We first prove two more
precise facts about central records: their selected centre has
small fine curvature, and their dependence on the block side can
be bounded before that side is chosen. These facts will serve
different purposes. Bounds denoted by \(C_B=C(1+B)^p\) may depend
on the fixed ratio \(L\). In contrast, constants denoted by
\(C_j\), and their exponents, are chosen before \(L\), at the
fixed adjacent profiles and specified derivative orders.

\subsection{A local affine completion and its first jet}

Retain the notation \(\mathbf F\), \(\mathbf C_c\),
\(\mathsf A,\mathsf B,\mathcal H\) of
Section~\ref{sec:calibrated-rows}. In a healthy central window
the source is the cut replica row minus the cut calibrated guide;
no protected frozen logarithm occurs. The two cancellations
\begin{equation}
 \mathsf B+\mathsf A\mathcal H
       =\Phi_0D\mathbf C_c(1),\qquad \mathsf A^{\mathsf T}\Phi_0=0
 \label{eq:central-local-flat-cancellation}
\end{equation}
are identities for the combined row vector. We do not assume
that every term in one of its path expansions has a vanishing
first jet.

\begin{lemma}[Local selected-centre comparison]
\label{lem:central-local-centre}
Fix a central physical test and the windows and row reads defining
its selected mean. On a sufficiently large healthy, nonwrapping
box containing those reads, choose a local axial gauge. Write
\(x\) for its coarse link logarithms. On a smaller box containing
the test, the following bounds, including all prescribed ordered
scaled derivatives, hold:
\begin{align}
 \mu&=\mathcal Hx+O_\infty\bigl(C_B(h_c^2+e^{-c_LB})\bigr),
                                                       \notag\\
 D\mu+\alpha
       &=O_\infty\bigl(C_B(h_c^2+e^{-c_LB})\bigr),
 \label{eq:central-local-mean-jet}\\
 \mathbf F(V,\mu)
       &=\Phi(V)\mathbf C_c(V)
                         +O_\infty\bigl(C_B(h_c^2+e^{-c_LB})\bigr).
 \label{eq:central-local-calibrated-row}
\end{align}
Here \(\mathcal Hx\) is evaluated on a small local completion
as in the proof; changing that completion outside the healthy
box changes its central value by the displayed exponential error.
The fine part of the leading row on the right of
\eqref{eq:central-local-calibrated-row} has bound
\(A L^{-2}h_c\), with \(A\) chosen before \(L\). Its noise
part has bound \(Ah_c\). All other constants in this lemma may
depend on the fixed block and profiles, but not on the cutoff,
period, sector, or finite history length.
\end{lemma}
\begin{proof}
The local gauge construction gives \(\|x\|_\infty\le C_Bh_c\).
Extend these logarithms by zero outside a larger box, without
changing the central selected windows or any of their row reads.
Call the resulting small global link array \(x^{\rm loc}\).
Its boundary curvature can have size \(C_Bh_c\), rather than
\(h_c\). We use the unmasked row formulas on this auxiliary
array. This is a comparison of formulas on a small chart; it
does not assert that the completion belongs to the physical
all-good sector, or replace a physical integral.

Let \(\delta_0=C_Bh_c\) bound this chart. On the auxiliary
no-hole completion let \(D_0,\alpha_0\) be the cut affine rows
and source, and put
\[
 M_0=D_0^{\mathsf T}D_0,\qquad q_0=-M_0^{-1}D_0^{\mathsf T}\alpha_0,
                  \qquad v=\mathcal Hx^{\rm loc}.
\]
The subscript zero in this paragraph denotes the auxiliary
completion, not evaluation at the identity. Its precision has
the flat gap after a perturbation of size
\(C_L\delta_0+Ce^{-c_LB}\), so its inverse has bounded row
and column sums by Lemma~\ref{lem:uniform-inverse}.
Flat matching, Taylor's formula for the row records, and their
native cut give
\[
 \begin{split}
 D_0&=\mathsf A+O(C_L\delta_0+Ce^{-c_LB}),\\
 \alpha_0&=(\mathsf B-\Phi_0D\mathbf C_c(1))x^{\rm loc}
               +O_\infty(C_L\delta_0^2+Ce^{-c_LB}\delta_0).
 \end{split}
\]
Consequently
\[
 \alpha_0+D_0v
                 =O_\infty(C_L\delta_0^2+Ce^{-c_LB}\delta_0),
 \qquad
 q_0-v=-M_0^{-1}D_0^{\mathsf T}(\alpha_0+D_0v).
\]
The last identity gives the same bound for \(q_0-v\), and
\[
 \alpha_0+D_0q_0
 =\{1-D_0M_0^{-1}D_0^{\mathsf T}\}(\alpha_0+D_0v)
\]
gives it for the affine residual. These are maximum-norm
operator bounds; no volume-dependent Euclidean row norm occurs.

Compare the selected mean on this completion with \(q_0\).
For each selected column the inverse difference is the boundary
formula proved in Proposition~\ref{prop:physical-selected-law}.
Its chain crosses the window boundary, at distance a fixed
multiple of \(B\) from its central coordinate. The source is
bounded, so the difference is \(C_Be^{-c_LB}\). The actual
sector has the same principal matrix and source entries on
these central window reads. Hence its selected mean there is
the same as that of the completion. The native tails of the
guide comparison have the same exponential bound.

Finally the actual energetic row has expansion
\(\mathbf F(V,a)=\mathbf F(V,0)+D_0a+O_\infty(C_L\delta_0^2)
\) for \(|a|\le C_L\delta_0\), including the cut-tail error.
This proves \eqref{eq:central-local-calibrated-row} and the
two mean bounds after absorbing fixed powers of \(B\).
The guide's fine and noise bounds are exactly
\eqref{eq:calibrated-preblock-prices} and
\eqref{eq:calibrated-coarse-row-price}.

For derivatives, keep the base gauge fixed. Its logarithmic
coordinates and each of their scaled derivatives through the
prescribed order are bounded by \(C_Bh_c\). The combined
uncut error has zero value and first jet at the identity.
Taylor's integral formula and the ordered chain rule therefore
bound every such derivative by \(C_Bh_c^2\). Differentiate
the inverse identities, using
\(D_A M^{-1}=-M^{-1}(D_AM)M^{-1}\). In a differentiated
boundary error, every inverse chain still crosses the same
boundary; inserted directions cost only a fixed polynomial
in the chain length. A smaller positive exponential width
pays that polynomial. Thus the cut and window errors retain
\(C_Be^{-c_LB}\), proving the derivative assertions.
\end{proof}

This argument uses the same actual affine rows as the selected
law. It also proves the local intact-box statement without
requiring a globally good logarithm of the physical field.
Choose \(L\) and the noise reserve \(W\) as in
Proposition~\ref{prop:general-central-good-centre}, and then
raise the entry coupling. Since
\[
 C_B(h_c^2+e^{-c_LB})/h\longrightarrow0,
\]
the mean has fine chord less than \(h/4\) and noise mismatch
less than \(3Wh/4\) on each central test. If
\(\max_e|z_e|\le\varepsilon_Lh/g\) on that test's fine reads,
the installed field \(a=\mu+gz\) retains a strict margin in
both tests, for a sufficiently small fixed
\(\varepsilon_L>0\). Its original regular frame and predictor
branches are then retained. The broad chart factors at width
\(hB^{17}\) are one. The Lipschitz prices here are local
stencil prices, allowed to depend on the fixed block.

The centre comparison is the local content of
\cite[Lemma~25.4, pp.~137--138, and Lemma~28.5,
pp.~156--158]{Main}; the proof above compares
the affine centres directly. It does not use the global
central-integral comparison as a derivative estimate.

\subsection{Record bounds whose constants precede the ratio}

The rough window-cardinality estimates of
Lemma~\ref{lem:output-selected-data} suffice for a rare packet.
They would place powers of \(B\) in a sharp transport
coefficient. We instead use decay before counting a read set.
A polynomial spatial moment of a kernel \(K\) means a row
or column sum of \(\|K_{ab}\|(1+d_c(a,b))^m\), for a specified
fixed nonnegative integer \(m\). Distances here are in coarse
units, also for the parents of fine coordinates.

\begin{lemma}[Polynomial bounds for central records]
\label{lem:central-polynomial-ratio}
At fixed adjacent profiles and fixed derivative orders, there
are \(C_j,M_j<\infty\), chosen before \(L\), such that, after
choosing a ratio-dependent entry coupling, the central operated
matrices, selected inverse columns, covariances and their
principal marginal inverses have all the required fixed
polynomial spatial moments bounded by \(C_jL^{M_j}\).
Every positive-order coarse derivative insertion supplies its
factor \(h_c\), together with its prescribed spatial weight.
The central mean has value bound
\begin{equation}
                         |\mu_e|\le C_jL^{M_j}h.
 \label{eq:central-polynomial-mean}
\end{equation}
These bounds are independent of \(B\), window size, history
length, cutoff and period. Traces and numbers of read
coordinates are counted separately.
\end{lemma}
\begin{proof}
Central sources are regular replica rows minus their
calibrated coarse guides. Path comparison in the fixed block
patch bounds them by \(C_jL^m h\). Their positive-order
derivatives have the same polynomial ratio bounds with the
corresponding direction factors. All threshold decisions have
their unmasked regular germ on this healthy domain.

A block tree has at most \(L^4\) vertices. Its paths, their
elementary square interchanges, and the overlap numbers in
the reconstruction estimate are bounded by fixed powers of
\(L\). The proof of Lemma~\ref{lem:physical-selected-gap}
therefore gives a lower precision bound \(c_jL^{-m}\), with
\(c_j,m\) fixed before \(L\). Direct and noise row metrics
have their prescribed fixed positive lower bounds. The
curvature perturbation is smaller than half this gap after
the entry choice; the inactive identity block causes no loss.

Absolute row budgets for these kernels and their fixed-order
derivatives have the form
\(C_jL^m e^{-\sigma_jd_c(a,b)}\), with \(\sigma_j>0\)
chosen before \(L\). The native cut decreases these absolute
budgets. Conjugating by \(e^{\theta d_c(\cdot,a)}\), with
\(\theta=c_jL^{-m'}\) sufficiently small, perturbs a principal
precision by less than half its gap. The weighted inverse
argument gives
\[
       \|(M^{-1})_{ab}\|
               \le C_jL^{m'}e^{-c_jL^{-m'}d_c(a,b)}.
\]
There are at most \(C_GL^4\) coordinates above one coarse
site. Summing any fixed polynomial spatial moment costs
another fixed power of \(L\), independent of the window.
The selected covariance differs from the full inverse by
the boundary error of Proposition~\ref{prop:physical-selected-law}.
At its entry threshold it retains half the positive spectral
margins. A principal marginal covariance inherits these
margins and the kernel bounds; applying the same weighted
inverse argument on that arbitrary subset bounds its inverse.

An ordered derivative of an inverse is a finite sum of chains
of inverses and differentiated matrices. The fixed number
of direction weights can be moved along each chain by the
triangle inequality and the just proved polynomial moments.
For the regular frame and predictor, their finite-dimensional
implicit equations have Hessian gaps bounded below by inverse
powers of \(L\). Their path derivatives and fixed-order
implicit differentiation consequently have polynomial ratio
prices as well. The same is true for the replica, tree-slice,
and exponential coordinate changes. Enlarge \(M_j\) to cover
these finitely many operations. Finally \(\mu=-J^{\mathsf T}b\) and the
central source bound give \eqref{eq:central-polynomial-mean}.
This proof counts neither the number of windows nor the
physical volume.
\end{proof}

This is \cite[Lemma~35.4, p.~216]{Main}. The distinction
between this pre-ratio bound and a post-ratio \(C_B\) estimate
will be retained in the transport argument.

\section{The central logarithm and its two strengths}
\label{sec:central-strength-estimates}

We now estimate the warning-free central logarithm with its
leading small-field order. A warning is one of the retained
failures of the broad chart, default row, or physical tests in
the exact column expansion. Warning letters keep their original
prices and are not rescaled below. Their marked estimates use
the common Gaussian reserve already present in the physical
construction. This section treats the subfamily containing no
such letter. Ordinary coefficient matching is a separate task.

\subsection{Compression of the actual read cover}

In the connected sums, an incoming insertion can meet another
factor anywhere in its full coarse read. We therefore need a
bound on the sum of weighted sizes of all incoming insertions
reading a specified coarse site. The incoming bound sums the
same sizes by their fine anchors.
Write \(e_i\) for the fine good-set norm of the incoming
regular activity \(E_i\), and assume
\begin{equation}
 \delta=\sup_x\sum_{i:o_i=x}e^{\kappa_f n_{B_f}(X_i)}e_i
                       \le\widehat\beta_f^{-11/10}.
 \label{eq:central-incoming-family}
\end{equation}

Let \(d_i=n_{B_f}(X_i)\), let \(\bar o_i\) be the coarse
parent of the fine anchor \(o_i\), and denote the full coarse
read by \(Y_i\). As in Section~\ref{sec:exact-masked-reassembly}
take \(B_f\le2B\) and a fixed minimal connected cover of
\(X_i\) by \(d_i\) fine balls of radius \(B_f\).
The read \(Y_i\) includes \(\bar o_i\), the mean and covariance
windows, row paths, tags and physical tests. It lies in a
fixed \(rB\)-neighbourhood of the coarse parents of the
cover. Here \(r\) measures the required read enlargement.

The bound \eqref{eq:reassembly-cover-enlargement}
allows a fixed number of coarse balls for each fine ball.
Its use in the weighted sum can require a larger fine
covering exponent. To allow the same exponent at successive
scales, first traverse the connected fine cover. Its
length shrinks by \(L\) in coarse units, so sampling
the whole walk gives a smaller coarse cover.

\begin{lemma}[Cover compression and value incidence]
\label{lem:sharp-cover-compression}
There are \(K_r,C_r<\infty\), fixed before \(B,L\), such that
\begin{equation}
 n_B(Y_i)\le K_r(1+C_rd_i/L).
 \label{eq:sharp-cover-compression}
\end{equation}
Consequently, for any fixed output exponent \(\kappa_*\),
choosing \(L\ge2C_rK_r\kappa_*/\kappa_f\) gives
\begin{equation}
 \kappa_* n_B(Y_i)\le\kappa_*K_r+\tfrac12\kappa_f d_i.
 \label{eq:sharp-cover-exponent}
\end{equation}
At that ratio the incoming value incidence satisfies
\begin{equation}
 \sup_x\sum_{i:x\in Y_i}e^{\kappa_*n_B(Y_i)}e_i
                         \le C_{j,L}(1+B)^4\delta.
 \label{eq:sharp-compressed-incidence}
\end{equation}
These statements also hold for wrapping footprints and walks.
\end{lemma}
\begin{proof}
Join the centres of fine balls that overlap or are
adjacent. The graph is connected and its edges have
periodic maximum length at most \(2B_f+1\).
Start at the fine anchor, go to a ball centre and
traverse a spanning tree in both directions.
Using successive shortest lifts, this gives a
continuous fine walk of length at most \(CB_fd_i\).

Divide its coordinates by \(L\). The resulting walk
has length at most \(CB_fd_i/L\) in coarse units.
A coarse parent differs from the scaled fine site
by a bounded lattice distance. Thus the coarse
parents of every ball are within \(B_f/L+C\) of
the scaled walk. Sample every \(B/4\) units of length,
including the endpoints, and round to the nearest coarse
sites. Since \(B_f\le2B\), this takes at most
\(C(1+d_i/L)\) samples. Consecutive rounded centres
are at distance at most \(B/4+1\le B\), so their
radius-\(B\) balls have connected union.
Every point of \(Y_i\) is within \(C_r B\) of some
sample, including the read enlargement. Cover each
such neighbourhood by a connected grid of radius-\(B\)
balls that includes the ball at the sample. A grid
of spacing \(B\) needs only a number of balls depending
on \(r\). The resulting connected cover proves
\eqref{eq:sharp-cover-compression}. If the walk wraps,
project the balls back to the torus. The compatible
fine and coarse periods make the parent map consistent
with the lifts, and projection preserves connectivity.
Multiplying by \(\kappa_*\) and applying the stated
lower bound on \(L\) gives \eqref{eq:sharp-cover-exponent}.

If \(x\in Y_i\), the same walk shows that the coarse
parent of \(o_i\) is within \(C_rB(1+d_i/L)\) of \(x\).
For \(d_i=d\), the number of possible fine anchors
is therefore at most \(C_rL^4(1+B)^4(1+d)^4\).
At each such fine anchor, \eqref{eq:central-incoming-family}
bounds the sum of \(e_i\) in this class by
\(e^{-\kappa_f d}\delta\). Multiply by the bound
\(e^{\kappa_*K_r+\kappa_f d/2}\) for the coarse weight
in \eqref{eq:sharp-cover-exponent} and sum
\((1+d)^4e^{-\kappa_f d/2}\) over \(d\ge1\).
This proves \eqref{eq:sharp-compressed-incidence}.
\end{proof}

These are the covering and value estimates of
\cite[Lemma~35.1 and (35.17), pp.~218--219 and 227]{Main}.
The locators were checked on 20 November 2026.
The same positive fine and coarse exponents can now
be used at repeated profiles.

\subsection{The norms and the two small incomes}

For a local function with anchor \(o\) and healthy footprint
\(X\), retain the norm
\(\|f\|_{q,h_c,s;o,X}\) of
\eqref{eq:local-physical-activity-norm}. A family before flat
and affine-Hessian subtraction will be measured by
\begin{equation}
 \|f\|_{\mathrm{raw},\kappa}
   =\sup_o\sum_{\alpha:o_\alpha=o}
        e^{\kappa n_B(X_\alpha)}
              \|f_\alpha\|_{q,h_c,s;o_\alpha,X_\alpha}.
 \label{eq:central-raw-family-norm}
\end{equation}
The word ``raw'' only records that the two normalisations
have not yet been imposed. Intermediate products use the
factorial jet seminorm \eqref{eq:connected-algebra-norm};
we convert to the prescribed maximum norm once for a completed
family. Hard masks are not differentiated.

The incoming part of the estimate uses the value incidence
\eqref{eq:sharp-compressed-incidence} at the required, fixed
larger output exponent. The enlargement estimate of
Section~\ref{sec:exact-masked-reassembly} also gives this type
of bound when its input exponent condition holds.

Use the actual \(g=\beta^{-1/2}\) in its original nominal
coupling band, \(h=\widehat\beta_f^{-1/2}P\), the original
\(B\), and the existing principal-coefficient bound
\(\Lambda\). Thus \(h/g\) is bounded above and below by fixed
positive multiples of \(P\). In this section \(C_B=C(1+B)^p\) is allowed
to depend on the fixed block, profiles and specified native
and covering reserves. It is independent of period, cutoff,
sector and finite history. Choose an envelope
\begin{align}
 A_\beta&=C(1+B)^{p_0}(1+P)^{p_1}(1+\Lambda)^{p_2}\ge1,
                                                        \notag\\
 \epsilon_B&=C_Be^{-c_LB},\qquad
 b=A_\beta\{gP^3+(1+\Lambda)h^2+\epsilon_B\},
                \qquad d=A_\beta\delta .
 \label{eq:central-strength-incomes}
\end{align}
The constants and fixed powers are enlarged finitely many
times in the proof, and their final values are fixed before
the strength radii are defined. These are chosen upper
envelopes, rather than the possibly smaller actual sums.
The prescribed slow-clock bounds give
\begin{equation}
 A_\beta=\widehat\beta_f^{o(1)},\qquad
 b=\widehat\beta_f^{-1/2+o(1)},\qquad b\ge gP^3\ge g.
 \label{eq:central-strength-rates}
\end{equation}

\subsection{Protecting the smooth difference after cancellation}

Let \(H=hB^{17}\). Choose a smooth radial map
\(a\mapsto a^{\rm pr}\) on each Lie-algebra link coordinate,
equal to \(a\) on the support of its broad chart factor,
bounded by \(CH\), and obeying
\begin{equation}
           \|\dd^r a^{\rm pr}\|\le C_rH^{1-r}\quad(r\ge1).
 \label{eq:central-protector-jets}
\end{equation}
A smooth radial cutoff times \(a\), with transition outside
the original support, has these properties. It commutes with
orthogonal adjoint transport.

Let \(\mathcal R_{\rm sm}(V,a)\) be the combined smooth row
energy after subtracting the installed affine squares. Extend
it as \(\mathcal R_{\rm sm}(V,a^{\rm pr})\). The affine
reference \(\frac12\|Da+\alpha\|^2\) still uses the original
\(a\), and the selected law is unchanged. In particular we
do not estimate two large protected squares separately.
This replacement is made in a central column before its broad
and default tests are expanded. On their joint support the
fine clipping is quadratic and the row predicates are one;
the smooth difference equals the original difference there.
Multiplication by those same tests is therefore an exact
identity. Their subsequent expansion retains all warnings.

Every original masked fine path term is instead declared a
Borel physical read, with its original argument. Incoming
functions are never evaluated at protected arguments. On the
union of these physical reads retain the broad factors by
\eqref{eq:output-retained-chart-rule}. A chart-failure letter
on that retained union is prohibited by that exact rule.

We will use the following elementary scaled Taylor fact.
If a smooth function has zero jets below degree \(m\), all
its arguments have size \(C_Bh\), and every positive-order
scaled argument jet has size \(C_Bh\), then every composed
jet through fixed order \(q\) is \(C_Bh^m\). For a chain-rule
term with \(r\le m\) derivatives of the function, Taylor's
bound \(C|y|^{m-r}\) is multiplied by \(r\) argument factors;
for \(r>m\) their product is \(C_Bh^r\le C_Bh^m\).
The auxiliary row formulas are smooth through all required
orders on this chart. Derivatives of the protector contribute
\(H^{1-r}\) together with \(r\) small variations, giving
the same conclusion with fixed powers of \(B\). Equivariance
allows this argument on the bounded protected coordinate
itself when a local gauge depends on the background.

\begin{lemma}[The bounded ordinary logarithmic prices]
\label{lem:central-bounded-logs}
The protected smooth row differences, paired determinant
differences and original bounded masked fine terms have
addressed logarithms \(v_\alpha\) and nonnegative prices
\(t_\alpha\) such that, in mixed coordinates with every Borel
physical read held fixed,
\begin{equation}
 \|v_\alpha\|_{\rm alg}\le t_\alpha,\qquad
 \sup_x\sum_{\alpha:x\in X_\alpha}
    e^{(\kappa+a)n_B(X_\alpha)}t_\alpha
       \le C_B\{gP^3+(1+\Lambda)h^2+\epsilon_B\}.
 \label{eq:central-bounded-log-income}
\end{equation}
Here \(a>0\) denotes the fixed unused covering reserve, and
the bounds are uniform in the remaining Gaussian coordinates
and shared tags. The unbounded corrections from the selected
Gaussian are excluded from this lemma.
\end{lemma}
\begin{proof}
Keep the six-row grouping \eqref{eq:packet-physical-rows} and the
flat completion \eqref{eq:packet-row-completion}. For fine
and noise rows write the actual residual as \(r+e\), where
\(r=Da+\alpha\) is its installed affine row. The protected
chart expansion gives, in every scaled jet under consideration,
\[
 |r|\le C_Bh,\qquad |e|\le C_Bh^2,\qquad
 \tfrac12|r+e|^2-\tfrac12|r|^2
                         =\langle r,e\rangle+\tfrac12|e|^2.
\]
The last expression has price \(C_Bh^3\); the direct
chord-versus-curvature defect starts at degree four.

For the other rows let \(v_1=D\mathbf C_c(1)x\) and
\(F_1=\mathsf Aa+\mathsf Bx\). Flat calibration and completion
give
\[
 \Phi_0^{\mathsf T}\mathsf A=0,\qquad \Phi_0^{\mathsf T}F_1=v_1,
                                  \qquad \Phi_{\perp,0}v_1=0.
\]
Indeed \(\Phi_0^{\mathsf T}\Phi_0+\Phi_{\perp,0}^{\mathsf T}\Phi_{\perp,0}=1\),
and \(\Phi_{\perp,0}v_1=0\) on exact coarse curls by
\eqref{eq:packet-row-completion}. Thus the bracket in \(w_L\)
has zero first jet, so \(w_L\) starts at degree three;
\(w_\perp\) starts at degree four. The coefficient bracket
in \(w_D\) vanishes at the identity, so \(w_D\) starts at
degree three. The coarse direct defect in \(w_C\) has degree
four. Native cuts leave an exponentially small defect with
the same fixed-order derivative price at a reduced width.
Hence the combined exponent costs
\[
 C_B\beta h^3+C_Be^{-c_LB}
                       \le C_BgP^3+C_Be^{-c_LB}.
\]
Any fixed power of the coupling multiplying a cut error is
absorbed by reducing \(c_L\), since \(B\) grows quadratically
in the logarithm of that coupling.

Lemma~\ref{lem:central-local-centre} gives
\(|D\mu+\alpha|\le C_B(h_c^2+e^{-c_LB})\), also in scaled
jets. Its retained constant-square exponent consequently
costs \(C_B\beta h_c^4+C_Be^{-c_LB}\), within the same
envelope. Its sign and value are retained in the exact density.

On the original fine mask, every curvature slot has norm
at most \(h\); all intervening transports are isometries.
The masked degree-three and degree-four terms have value
prices \(C_B\beta h^3\) and \(C_B\beta h^4\); quadratic
loop terms have price \(C_B(1+\Lambda)h^2\). In held physical
variables these entire masked expressions are fixed, so no
derivative of their fine masks or of their functions is used.
The protected Haar logarithm has price \(C_Bh^2\), since
\(j_G(0)=1\) and \(\dd j_G(0)=0\).

The zeroth-window determinant is paired with its flat
reference before estimating it. Its trace identity and
\(M-M_{\rm flat}=O(C_Bh_c)\) in a local gauge give price
\(C_Bh_c\). Ordered inverse differentiation preserves this
bound with fixed covering and native reserves, and
\(h_c\le2gP\) is within \(C_BgP^3\). Positive-defect
determinant walks retain their exponential smallness by
the boundary-chain argument in
Lemma~\ref{lem:output-selected-data}.
The extensive flat determinant remains a common scalar
normalisation and is not counted as a small letter.

Absolute path budgets and their unused native reserves sum
these local prices, including differentiated path lengths.
The connected cover estimate then proves
\eqref{eq:central-bounded-log-income} at the stated reserve.
Every cancellation above preceded taking absolute values.
\end{proof}

\subsection{A single Gaussian reserve for a whole word}

The residual and precision corrections are
\[
 r=-(2S-E)b,\qquad
 C_{\rm sel}^{-1}-\mathsf M=-S C_{\rm sel}^{-1}
                    =\sum_{n\ge1}(-S)^n\mathsf M.
\]
In the first formula \(E=\mathsf MJ-1\) is the selected-column
defect of \eqref{eq:physical-selected-windows}, not an incoming
activity. The precision series keeps every intermediate matrix
index as part of its address.

Group each finite-window matrix entry before taking its absolute
cap. The differentiated entry bounds have exponential decay
and an exponentially small boundary factor. For a typical defect
entry their cap is \(C_Be^{-c_LB}e^{-\nu_Ld(i,j)}\). Its square
root has summable row and column budgets at any width below
\(\nu_L/2\), still exponentially small in \(B\). The square-root
entry budgets of \(\mathsf M\) are bounded. Thus a walk with
\(n\) defect entries is summed by the \(n\)-th power of a small
square-root budget, followed by this bounded final matrix budget.
Endpoint occupation introduces at most a linear factor in
\(n\), and differentiated chains introduce fixed powers;
their geometric sums converge. Applying this to the
matrix-walk series gives correction addresses \(\alpha\) with
\begin{equation}
 |v_\alpha|\le c_\alpha Q_\alpha,\qquad
 Q_\alpha=1+|z_i|^2+|z_j|^2
       \ \hbox{or}\ 1+|z_i|^2,
 \label{eq:central-Gaussian-letter-cap}
\end{equation}
whose endpoint occupation and weighted incidence sums of
\(\sqrt{c_\alpha}\) are at most \(\epsilon_B\), including
fixed-order polynomial insertions. Repeated endpoints are
counted with multiplicity. Native bounded path terms from
Lemma~\ref{lem:central-bounded-logs} keep their linear prices
\(t_\alpha\). An absolute sum of path coefficients alone
would not justify summing their individual square roots;
that inference is not used here.

For a complex scalar \(u\), the elementary inequalities
\(|e^w-1|\le |w|e^{|w|}\) and \(x\le e^x\), \(x\ge0\), give
\begin{equation}
 |e^{uv_\alpha}-1|
 \le |u|\sqrt{c_\alpha}
       \exp\{(\sqrt{c_\alpha}+|u|c_\alpha)Q_\alpha\}.
 \label{eq:central-Gaussian-complex-Mayer}
\end{equation}
Fixed-order derivatives add only their prescribed polynomial
insertions and a fixed price inflation. The incoming factor
\(e^{-v\chi_{X_i}E_i}-1\) has value price \(2|v|e_i\)
when \(|v|e_i\) is small; it adds no quadratic Gaussian charge.

Let \(C_{\rm sel}^{-1}\ge\lambda_G1\).
Complete the fixed polynomial inflations in \(D_B\) first,
then choose \(0<c\ll1\) for the common precision and
connected-logarithm bounds. Its smallness threshold depends
only on the fixed profile and the specified reserves.
Define the strength radii by
\begin{equation}
 |u|\le R=\frac{c}{D_Bb},\qquad
 |v|\le S_*=\frac{c}{D_Bd},\qquad R,S_*\ge2,
 \label{eq:central-strength-discs}
\end{equation}
These radii are defined from the incomes, not chosen
independently. If \(\delta=0\), omit the second variable.
Increase entry so that \(D_Bb\ll c\) and, when \(d>0\),
\(D_Bd\ll c\). This gives \(R,S_*\ge2\) whenever defined.
Since \(b\ge g\), \(R\le Cg^{-1}\), and the same entry
choice makes \((1+R)\epsilon_B\ll\lambda_G\).
We write \(S_*\) to distinguish the scalar radius from
the matrix defect \(S\).

\begin{lemma}[The common precision and derivative majorant]
\label{lem:central-common-reserve}
Throughout a neighbourhood of the closed discs
\eqref{eq:central-strength-discs}, every finite distinct-address
central word is absolutely integrable with all its prescribed
ordered coarse derivatives. Its derivative integral has a
nonnegative majorant of the form
\begin{equation}
 A_\beta(1+n_B(X(\Gamma)))^{m_q}
                               \prod_{\alpha\in\Gamma}p_\alpha(u,v),
 \label{eq:central-word-derivative-price}
\end{equation}
where the inflated primitive prices have incidence income at
most \(C(|u|b+|v|d)\) at the specified larger covering
exponent. All charges are taken from the same Gaussian
precision reserve, independent of the word.
\end{lemma}
\begin{proof}
Sum \eqref{eq:central-Gaussian-complex-Mayer} over the
distinct correction addresses in a word. Their quadratic
charge at any coordinate is at most
\(C_q(1+|u|)\epsilon_B\). At the entry threshold it is
less than \(\lambda_G/16\). Reserve at most another
\(\lambda_G/16\) for a local differentiation neighbourhood
and at most \(\lambda_G/8\) for any of the already priced
warning witnesses. Even doubling these allowances leaves
at least half the original precision. No charge is assigned
separately to each appearance of the same physical coordinate.

Move every Borel read back to physical Haar variables using
\eqref{eq:output-fixed-physical-integral}, retaining its broad
factor. The mixed covariance is the principal marginal of
the same selected Gaussian. Lemma~\ref{lem:output-score-bounds}
and the ordered density partition formula bound all scores
by a fixed polynomial in \(B,P\), the full-cover size, and
\(1+\sum |z_e|^2\). No derivative lands on a Borel incoming
multiplier or a hard fine mask. The smooth logarithms have
the prices of Lemma~\ref{lem:central-bounded-logs}; small
Gaussian corrections keep their coefficient caps after this
change of variables. In particular \((a_T-\mu_T)/g\) and its
fixed-order variations cost powers of \(P,B\), absorbed in
\(A_\beta\), rather than powers of the coupling.

For clarity, domination under differentiation does not
introduce a volume factor. Fix a finite read set \(I\) and
a base background \(V_0\). By continuity of its positive
covariance and chart maps, shrink its neighbourhood until
the density ratio is bounded by
\[
 2\exp\left\{\frac{\lambda_G}{32(1+|I|)}
                         \left(1+\sum_{e\in I}|z_e|^2\right)\right\}.
\]
The displayed coefficient uses only a fixed part of the
remaining precision. Young's inequality handles the small
mean shift. Such a neighbourhood may depend on
\(I\); the derivative bound at its base point does not.
It supplies a common integrable function for the finite
read integral, also at equality in a coarse good-field test.

Polynomial moments are taken only after all word charges
have been combined. For a nonnegative diagonal charge
matrix \(A\) on the read marginal with
\(2C^{1/2}AC^{1/2}\le\frac12 1\),
\[
 \mathbb E_Ce^{z^{\mathsf T}Az}
   =\det(1-2C^{1/2}AC^{1/2})^{-1/2}
   \le \exp\{C\operatorname{Tr}(CA)\}.
\]
The trace is bounded by a constant times the sum of the
endpoint charges, not the number of unused coordinates.
Assign its exponential to the individual correction
letters; their charges are uniformly bounded, so this is
a fixed inflation per selected letter. Under the remaining
positive-precision Gaussian, a polynomial of fixed degree
in \(\sum_{I}|z_e|^2\) costs a fixed power of \(1+|I|\).
Since \(|I|\le C_B n_B(X(\Gamma))\), this is precisely
the cover polynomial in
\eqref{eq:central-word-derivative-price}.

There can be arbitrarily many addressed letters on one
footprint. If their number is \(r\), the ordered product
rule contributes at most \(C_qr^q\). For a fixed small
\(\eta>0\), \(r^q\le C_{q,\eta}e^{\eta r}\), which is
another fixed inflation per letter. We do not bound \(r\)
by the covering number. The factor \(A_\beta\) for a
nonempty word is bounded by assigning \(A_\beta\ge1\)
to each of its letters. The primitive prices may thus be
taken as fixed inflations of
\[
 |u|t_\alpha,\qquad |u|\sqrt{c_\alpha},\qquad |v|e_i.
\]
Pay the common-anchor and cover polynomial once from the
unused covering reserve, using
\(\sup_{k\ge1}(1+k)^{m_q}e^{-\eta k}<\infty\).
The bounded-letter income, correction endpoint sums and
incoming value incidence give the asserted primitive
income. This proves the majorant and integrability.
\end{proof}

\subsection{The local analytic family and coefficient extraction}

Replace every ordinary logarithmic correction by
\(e^{uv_\alpha}\), and every incoming occurrence by
\(e^{-v\chi_{X_i}E_i}\). All laws, means, shared tag
probabilities, broad-factor retention rules and normalising
constants are fixed while \(u,v\) vary. In particular a
strength is only a scalar coefficient-counting parameter.
Let \(F(u,v)\) be the warning-free connected logarithmic
family, with these parameters. At \((0,0)\) there is no
selected letter; the reference marginal and the original
tag prior are normalised, so \(F(0,0)=0\).

\begin{theorem}[The central joint-strength bound]
\label{thm:central-joint-strength}
The family \(F(u,v)\) is analytic on a neighbourhood of the
closed bidisc \eqref{eq:central-strength-discs}, with an
absolutely summable local-family majorant satisfying
\begin{align}
 \sup_{|u|\le R}\|F(u,0)\|_{\mathrm{raw},\kappa}
                                      &\le C_BRb,\notag\\
 \sup_{|u|\le R,\,|v|\le S_*}
             \|F(u,v)\|_{\mathrm{raw},\kappa}
                                      &\le C_B(Rb+S_*d).
 \label{eq:central-joint-strength-bound}
\end{align}
The bounds hold after every prescribed ordered derivative,
before summing the family labels, uniformly in cutoff,
physical period and admitted finite history.
\end{theorem}
\begin{proof}
First perform the distinct-address component expansion
with Lemma~\ref{lem:central-common-reserve}. The positive
rooted-tree estimate
Lemma~\ref{lem:connected-positive-tree} gives integrated
component prices at the remaining stronger exponent.
Its smallness parameter is bounded by
\(D_B(|u|b+|v|d)\le2c\).
Next apply Lemma~\ref{lem:connected-log-bound} to these
integrated species, with their ordinary Ursell factorials
and possible repetitions. Repetition at this second step
does not repeat Gaussian charges in a physical word: each
species integral has already been bounded. The factorial
jet product rule and the whole-product anchor estimate
\eqref{eq:connected-one-reserve} use only one common-anchor
reserve for a cluster. They yield
\eqref{eq:central-joint-strength-bound}.

Finite-read integrals are analytic in the scalar strengths
by the same common precision bound. The strict smallness
margins hold on a slightly larger bidisc. Uniform absolute
majorants justify the sums and ordered derivatives there,
proving analyticity and the stated good-set germs. This is
a bound on local families, rather than a bound on the total
finite-volume integral by a volume-independent number.
\end{proof}

\begin{proposition}[Central warning income]
\label{prop:central-warning-income}
At the physical strengths \(u=v=1\), central logarithmic
clusters containing a retained local warning, and no physical
packet, have raw derivative-family income
\begin{equation}
                C_B\exp\{-P^2/C_B\}
                         \le C_{j,L,J_0}\widehat\beta_f^{-J_0}
                 \quad\hbox{for every fixed }J_0>0.
 \label{eq:central-warning-income}
\end{equation}
The original warning factors and physical masks are retained.
\end{proposition}
\begin{proof}
The intact-box conclusion following
Lemma~\ref{lem:central-local-centre} proves that every local
broad-chart, default-row and physical-failure warning is
zero when all Gaussian coordinates in its test read obey
\(|z_e|\le\varepsilon_Lh/g\). Shrink this width by a fixed
factor so that the same strict margin holds throughout a
local coarse germ. Its derivatives, when a smooth warning
is differentiated, have support outside that smaller box.
Hard physical tests are held fixed in Haar variables.

These warnings form a finite local list. A coordinate
belongs to at most \(N_B\le C_B\) test reads, each containing
at most \(C_B\) coordinates. Give all selected warning
witnesses together at most \(\lambda_G/8\) of the common
precision, by using the witness exponent
\(v_0=\lambda_G/(8N_B)\). The proof of
\eqref{eq:reassembly-failure-witnesses}, now with this
smaller reserve, bounds their joint integral by the product
of their individual prices
\(C_B\exp\{-P^2/C_B\}\). Gaussian normalisation costs are
assigned to these witnesses as in that proof. The other
ordinary and incoming letters use the remaining reserve
of Lemma~\ref{lem:central-common-reserve}.

After changing Borel reads to physical variables, fixed-order
scores and smooth warning derivatives cost only a fixed
polynomial in \(B,P\), the full-cover size and Gaussian
coordinates. Take the corresponding fixed-degree moments
under this one tilted law. The full-cover polynomial is
paid by the unused covering exponent before summation,
and the remaining polynomial is absorbed by reducing the
warning exponent. This does not bound a long incoming
bridge by the number of warnings.

The integrated component prices now have the stated marked
income. Apply the marked logarithm argument
\eqref{eq:output-marked-log-price} to obtain the same order
for clusters containing a warning. Countable sums and
ordered derivatives have this common local majorant.
Finally \(\log P=(\log\log\widehat\beta_f)^2\), whereas
\(C_B\) is a fixed power of \(1+B\). Thus
\(P^2/(C_B\log\widehat\beta_f)\to\infty\), proving the
last inequality in \eqref{eq:central-warning-income}.
\end{proof}

This uses the finite local witness estimate of
\cite[Lemma~28.6, p.~158]{Main}, now with the selected-centre
vanishing condition verified above and the full ordered
density scores retained.

The ingredients of the joint-strength theorem are
\cite[Lemmas~34.1--34.4 and Theorem~34.5, pp.~205--211]{Main}.
The strength argument below is
\cite[Propositions~32.1--32.2, pp.~197--198]{Main}. Its
purpose is to estimate higher connected coefficients once,
without repeating a separate diagram expansion for each.

\begin{proposition}[Two ordinary coefficients and one incoming insertion]
\label{prop:central-strength-coefficients}
With the incomes \eqref{eq:central-strength-incomes},
\begin{align}
 F(1,0)&=\partial_uF(0,0)+\tfrac12\partial_u^2F(0,0)
                                             +R_{\ge3},\notag\\
 \|R_{\ge3}\|_{\mathrm{raw},\kappa}&\le C_Bb^3,
 \label{eq:central-third-remainder}\\
 F(1,1)-F(1,0)&=\partial_vF(0,0)+R_{\rm cmp},\qquad
 \|R_{\rm cmp}\|_{\mathrm{raw},\kappa}
                                             \le C_B(bd+d^2).
 \label{eq:central-incoming-comparison}
\end{align}
The exact single insertion is
\begin{equation}
 \partial_vF(0,0)=-\sum_i A_i(V),\qquad
 A_i(V)=\operatorname{Av}\int
       \left(\prod_{e\in T_i}b_e\right)
                \chi_{X_i}(U)E_i(U)\,\dd\gamma_{C_i}(z).
 \label{eq:central-exact-incoming-insertion}
\end{equation}
Here \(T_i\) is its full physical fine read set of active
free coordinates, \(C_i=(C_{\rm sel})_{T_i,T_i}\), and
\(U=\mathcal U_\sigma(V,\mu+gz)\) is the original installed
field. The symbol \(\operatorname{Av}\) is the original
normalised shared tag average, with any prescribed discrete
averages unchanged. Tree links retain their installed values.
\end{proposition}
\begin{proof}
Cauchy's formula applies to each labelled cluster and each
ordered derivative. Its nonnegative family majorant permits
the contour integrals before all sums and suprema. The
coefficient of \(u^n\) in \(F(u,0)\) has norm at most
\(C_BRbR^{-n}\), so
\[
 \sum_{n\ge3}C_BRbR^{-n}
        \le2C_BbR^{-2}\le C_Bb^3.
\]
For the second difference remove
\(F(u,0)+v\partial_vF(0,0)\). The remaining terms have
either positive degree in both variables, or zero
\(u\)-degree and \(v\)-degree at least two. If
\(M=C_B(Rb+S_*d)\), their sum at \((1,1)\) is bounded by
\[
 \frac{M}{(R-1)(S_*-1)}
             +M\frac{S_*^{-2}}{1-S_*^{-1}}
                              \le C_B(bd+d^2),
\]
because \(Rb=S_*d=c/D_B\). Fixed powers of \(D_B\)
are included in \(C_B\). The derivative at \((0,0)\)
selects exactly one original incoming occurrence and no
ordinary interaction letter. Its retained-chart rule gives
\eqref{eq:central-exact-incoming-insertion}, including its
minus sign. No differentiation of an incoming function
was used in this scalar argument.
\end{proof}

When the ordinary series is the logarithm of a fixed protected
Gaussian integral, its first two coefficients, summed as total
actions at finite volume, are \(\mathbb E\mathcal V\) and
\(\frac12\{\mathbb E\mathcal V^2-(\mathbb E\mathcal V)^2\}\),
where \(\mathcal V\) is its ordinary logarithmic potential.
In the physical expansion the coefficients are defined by
the actual addressed rules, including the retained broad
factors. Their comparison with the formal Gaussian coefficients,
and their classical and one-loop cancellations, have not been
proved by Cauchy's formula. Section~\ref{sec:matching-taylor-data} identifies these retained
coefficients and computes the stationary Taylor data needed for
their matching. The remainders just established have orders
\begin{equation}
 \|R_{\ge3}\|_{\mathrm{raw},\kappa}
                  \le\widehat\beta_f^{-3/2+o(1)},\qquad
 \|R_{\rm cmp}\|_{\mathrm{raw},\kappa}
                  \le\widehat\beta_f^{-1/2+o(1)}\delta.
 \label{eq:central-remainder-orders}
\end{equation}
These estimates supply the previously deferred central
unforced derivative incomes. The leading single incoming
average still requires its sharper pre-ratio contraction.

\section{The affine-Hessian subtraction and its masks}
\label{sec:affine-local-reset}

The leading incoming contribution in
\eqref{eq:central-exact-incoming-insertion} needs a
gain after summing over the \(L^4\) fine anchors
with a given coarse parent. A constant value \(E_i(1)\)
carries no power of \(L\), so first extract it.
The flat gauge identity below removes the linear
Taylor term. The next term to examine is the Hessian
of \(E_i\) at the identity.

Consider the power count on a small nonwrapping
footprint. The free lift decomposes into a constant
part, an affine part and a smaller remainder. After
a local gauge change removes the constant part, the
affine part has size \(O(L^{-2})\) in the fine direction
norm and the remainder has size \(O(L^{-3})\).
The lift and the gauge change are calculated in
Section~\ref{sec:sharp-regular-transport}. At the
level of powers, bounding the Hessian on two affine
directions by the incoming norm gives \(O(e_iL^{-4})\).
Summing by the coarse anchor spends that factor:
\[
                       L^4(L^{-2})^2=1.
\]
If the Hessian vanishes on pairs of affine directions,
a remaining Hessian term contains at least one
remainder and has size \(O(e_iL^{-5})\).
It then gives \(L^4L^{-5}=L^{-1}\) after summation.
This calculation tells us which quadratic response to
extract. We now specify the normalisation and prove
the local subtraction, retaining the extracted scalar
and quadratic terms in the exact identity.

For an anchor \(o\) on a torus of side \(\ell\), let
\(\xi(x)\) be the centred integer lift of \(x-o\). A footprint
is nonwrapping about \(o\) if
\(10d_\infty(x,o)<\ell\) on that footprint. The period
reserve also keeps its endpoint halo in the same local chart.
An affine link field has the form
\begin{equation}
 A_i(x)=a_i+\sum_{j=1}^4\xi_j(x)b_{ij},\qquad
                         a_i,b_{ij}\in\g.
 \label{eq:reset-affine-field}
\end{equation}
For a gauge invariant local function \(f\), write
\[
 \widetilde f(a)=f((e^{a_e})_e),\qquad
 H_f(A,C)=\dd^2\widetilde f(0)[A,C].
\]
This is its exponential-chart Hessian, a symmetric bilinear
form; it is distinct from an ordered composition of left
derivatives away from the identity. The required regular
normalisations are
\begin{equation}
 E_i(1)=0,\qquad H_{E_i}(A,C)=0
       \quad\hbox{for affine }A,C
       \hbox{ on a nonwrapping }X_i.
 \label{eq:reset-incoming-normalisations}
\end{equation}
They are induction data, not consequences of the size norm.
At the initial Wilson step there is no incoming correction.
Their propagation uses the subtraction proved below.

\subsection{What gauge invariance removes at the identity}

\begin{lemma}[The first two flat gauge identities]
\label{lem:reset-flat-gauge-identities}
For the compact simple Lie algebra fixed in Section~\ref{sec:finite},
\(\dd\widetilde f(0)=0\). If
\((d_0\omega)_i(x)=\omega(x+e_i)-\omega(x)\) for a site field
\(\omega\), then \(H_f(d_0\omega,A)=0\). On a nonwrapping
footprint, the Hessian on affine fields depends only on their
constant linearised curls
\begin{equation}
                (d_1A)_{ij}=b_{ji}-b_{ij},\qquad i<j.
 \label{eq:reset-affine-curl}
\end{equation}
\end{lemma}
\begin{proof}
Constant gauge transformations act by the same adjoint map
in every link slot. Each component of the linear functional
\(\dd\widetilde f(0)\) is therefore adjoint invariant and
vanishes on every commutator. The commutator span is a
nonzero ideal and hence equals \(\g\); the linear functional
is zero. Next apply the site gauge \(e^{t\omega(x)}\) to
\(e^{sA_e}\) and differentiate gauge invariance once in
\(s\) and once in \(t\) at zero. The mixed derivative of
the coordinate change is multiplied by the just proved
zero first derivative. The remaining identity is
\(H_f(d_0\omega,A)=0\), with an irrelevant sign from the
gauge convention. Locality restricts this identity to the
links actually read by \(f\).

If an affine field has zero curl, \(b_{ij}=b_{ji}\). On the
local read patch and its endpoint sites put
\[
 \omega(x)=\sum_i\xi_i(x)(a_i-\tfrac12b_{ii})
                  +\tfrac12\sum_{i,j}\xi_i(x)\xi_j(x)b_{ij}.
\]
Direct differencing gives \(d_0\omega=A\) on every read
link. Extend \(\omega\) arbitrarily to the remaining torus
sites. Nonwrapping ensures that these local differences
do not cross a seam. The gauge identity and locality
annihilate this field in either Hessian slot. Subtracting
two affine fields with equal curl proves the last assertion.
\end{proof}

\subsection{A clover with the same Hessian}

For a six-component antisymmetric curvature array
\(\omega=(\omega_{ij})_{i<j}\in\g^6\), choose the affine
template
\[
             A^\omega_j(x)=\sum_{i<j}\xi_i(x)\omega_{ij}.
\]
Its curl is \(\omega\). It has size at most
\(C(1+d_\infty(x,o))\max_{i<j}|\omega_{ij}|\), with a
fixed dimension constant.

Recall the smooth covariant curvature
\(\mathcal C(W)=\Pi_\g((W-W^{-1})/2)\). For
\(\varepsilon,\varepsilon'\in\{1,-1\}\), let
\(\gamma_{\varepsilon i,\varepsilon'j}\) be the four-link
loop based at \(o\) with steps
\(\varepsilon e_i,\varepsilon'e_j,-\varepsilon e_i,
-\varepsilon'e_j\). Backward links use their inverses.
Define the clover curvature by
\begin{equation}
 \Phi_{ij}(V;o)=\frac14\sum_{\varepsilon,\varepsilon'=\pm1}
       \varepsilon\varepsilon'\,
                    \mathcal C(V_{\gamma_{\varepsilon i,\varepsilon'j}}).
 \label{eq:reset-clover}
\end{equation}
Thus ``clover'' means the oriented average of the four
elementary \(ij\)-plaquettes incident to the anchor, all
based there. It reads \(B_\infty(o,2)\), is conjugated at
\(o\) by a gauge transformation, and has size at most \(h\)
on a healthy anchor. Differentiating its four links at the
identity gives
\(\dd\Phi_{ij}(1)A=(d_1A)_{ij}\) for an affine \(A\).
The signs in \eqref{eq:reset-clover} remove the orientation
signs in these four derivatives.

Put \(X'=X\cup B_\infty(o,2)\). If \(X'\) is nonwrapping,
define
\begin{equation}
            Q_f(V)=\tfrac12 H_f(A^{\Phi(V;o)},A^{\Phi(V;o)});
 \label{eq:reset-quadratic-counterterm}
\end{equation}
otherwise put \(Q_f=0\). This formula uses only the fixed
flat Hessian coefficients and local clovers, not link
logarithms of \(V\). Adjoint invariance of the Hessian
proves gauge invariance of \(Q_f\). Lemma~\ref{lem:reset-flat-gauge-identities}
shows that \(Q_f\) has the same Hessian as \(f\) on two
affine fields whenever \(X'\) is nonwrapping. This is the
counterterm of \cite[(13.6), p.~48]{Main}, written directly
as a bilinear form so no colour coefficients need be introduced.

\begin{lemma}[Uniform local subtraction]
\label{lem:reset-local-norm}
Let \(q\ge2\), \(s\ge2\), \(h\le1\), and
\(\|f\|_{q,h,s;o,X}\le u\). The residual
\(\bar f=f-f(1)-Q_f\), local on \(X'\), satisfies
\begin{equation}
 \|\bar f\|_{q,h,s;o,X'}\le C_r u,
                         \qquad C_r=C_r(G,q,s).
 \label{eq:reset-uniform-norm}
\end{equation}
It has zero flat value and zero affine Hessian on a
nonwrapping \(X'\). There is also the exact masked identity
\begin{equation}
 \chi_X f=f(1)+\chi_{\{o\}}Q_f-u_f+\chi_{X'}\bar f,
 \label{eq:reset-exact-mask}
\end{equation}
where \(u_f=0\) on a healthy \(X'\), and otherwise
\(u_f=f(1)+\chi_{\{o\}}Q_f-\chi_Xf\).
This Borel local function has value bound \(C_r'u\), with
\(C_r'\) independent of the footprint, threshold, block,
period and cutoff.
\end{lemma}
\begin{proof}
For \(\max|\omega_{ij}|,\max|\nu_{ij}|\le1\), the scaled
templates \(hA^\omega/C\), \(hA^\nu/C\) are allowed
directions. At the identity their chart Hessian equals
the ordered second derivative, since the correction is
multiplied by \(\dd\widetilde f(0)=0\). Thus
\[
                    |H_f(A^\omega,A^\nu)|\le C u h^{-2}.
\]
Equivalently one can use polarisation of the diagonal
second derivative; either estimate has a fixed dimension
constant. Every clover and each of its positive-order
scaled derivatives through order \(q\) has size
\(C_{G,q,s}h\) on a healthy anchor. At order zero use its
four healthy plaquettes. At positive orders differentiate
the four-link words; their read directions have size at
most \(h3^s\), and \(h^r\le h\). Hence every derivative
of the bilinear clover product has size \(Ch^2\).
The apparent coefficient price \(h^{-2}\) is exactly
cancelled. It follows that
\(\|Q_f\|_{q,h,s;o,\{o\}}\le C_{G,q,s}u\), with its
declared clover halo. On the healthy set of \(X'\),
\(f\) retains its original norm and \(|f(1)|\le u\).
The triangle inequality proves \eqref{eq:reset-uniform-norm}.
The normalisations follow from the equality of Hessians.

On \(\chi_{X'}=1\), both smaller masks equal one and
\eqref{eq:reset-exact-mask} is the definition of \(\bar f\).
On its complement it is the definition of \(u_f\).
The bounds \(|\chi_Xf|\le u\), \(|f(1)|\le u\), and
\(|\chi_{\{o\}}Q_f|\le Cu\) prove the asserted value
price. No value of \(f\) away from its good set is needed
in this masked identity.
\end{proof}

These are the local operations of
\cite[Lemmas~13.1--13.2 and 13.5--13.6, pp.~46--49]{Main}. Scalars,
clover counterterms and masked discrepancies are all retained
as labelled terms. In a family, \(n_B(X')\le n_B(X)+1\)
when \(B\ge2\) and \(o\in X\). The extraction therefore
costs at most \(C_re^\kappa\) in its anchored norm, with
no new factor depending on \(B\) or \(L\). This generic
subtraction will also apply to the packet and central
remainder germs; summing their principal coefficients and
matching the full principal action remain separate work.

\section{Sharp transport of the incoming regular family}
\label{sec:sharp-regular-transport}

The leading term in
\eqref{eq:central-exact-incoming-insertion} requires a constant
chosen before the ratio. We prove the contraction for a principal
average on the same selected Gaussian, then retain the exact
cutoff comparison as a separate small family. Throughout,
\begin{equation}
 q_f\ge\max(q,3),\qquad s_f\ge s_c\ge2,
 \qquad e_i=\|E_i\|_{q_f,h,s_f;o_i,X_i}.
 \label{eq:sharp-input-profile}
\end{equation}
The incoming functions obey
\eqref{eq:reset-incoming-normalisations}, and \(B_f\le2B\).
Equality \(q_f=q\ge3\) is allowed. In particular this estimate
does not lose one derivative at each blocking step.

First fix the profiles, derivative orders, positive covering
exponents and pre-ratio direct and noise prices. Constants
\(C_j,M_j\) are fixed from these data before choosing \(L\).
Next choose the ratio, then the auxiliary inner width below,
and finally the entry coupling. The inner width and entry
threshold may depend on \(L\); the leading contraction
constants are fixed before it. The coefficients
\(V_j,\widetilde V_j\) are therefore kept explicit; an
\(L\)-dependent implicit constant is allowed only in the
subordinate post-ratio estimates. All choices use the same
Wilson family, physical masks and selected integration law.

Use \(d_i,\bar o_i,Y_i\) from
Lemma~\ref{lem:sharp-cover-compression} and its
covering weight bound \eqref{eq:sharp-cover-exponent}.
The two size regimes below are \(d_i^3<L\) and
\(d_i^3\ge L\).

\subsection{The constant and affine parts of the free lift}

Use \(E\) for the free interpolation column of
Appendix~\ref{app:calibration}. It is an operator, not the
incoming family. Its first two fine differences satisfy
\begin{equation}
 |\nabla_f^rE|_\sigma\le C_jL^{-1-r},\qquad 0\le r\le2,
 \label{eq:sharp-free-column-differences}
\end{equation}
by \eqref{eq:calibration-column-differences}; these bounds
include residue changes across block boundaries. For a fine
field on the original read patch put
\[
             |u|_i=\sup_e
               \frac{|u_e|}{h(1+d_f(e,o_i))^{s_f}}.
\]
This abbreviation only describes the existing incoming
direction bound. It is not an additional activity norm.

\begin{lemma}[Three different orders in a lifted direction]
\label{lem:sharp-three-parts}
If \(|v(y)|\le Ch(1+d_c(y,\bar o_i))^{s_c}\), then on a
nonwrapping fine patch
\(Ev=c+b+\rho\), where \(c\) is constant and
\(b_\mu(x)=\sum_\nu\xi_\nu(x)b_{\mu\nu}\), with
\begin{equation}
 |c_\mu|\le C_jhL^{-1},\qquad
 |b_{\mu\nu}|\le C_jhL^{-2},\qquad
                         |\rho|_i\le C_jL^{-3}.
 \label{eq:sharp-three-parts}
\end{equation}
The constants are independent of the read radius, \(B,L\),
history and period.
\end{lemma}
\begin{proof}
Polynomial moments of
\eqref{eq:sharp-free-column-differences} give
\[
 |\nabla_f^rEv(x)|\le C_jhL^{-1-r}
                    (1+d_f(x,o_i)/L)^{s_c}.
\]
The ratio of two polynomial weights is bounded by a fixed
polynomial in their separation; its moment against the
exponential kernel is summable. The fine anchor's residue
within its coarse block changes only a fixed factor.

Take \(c_\mu=(Ev)_\mu(o_i)\) and
\(b_{\mu\nu}=\nabla_{f,\nu}(Ev)_\mu(o_i)\).
For \(r=d_f(x,o_i)\le L\), telescope first differences
along coordinate paths. At most \(C(1+r)^2\) second
differences occur, giving
\[
              |Ev(x)-c-b(x)|\le C_jhL^{-3}(1+r)^2.
\]
For \(r\ge L\), use the value bound and the two subtracted
terms instead:
\[
 |\rho(x)|\le C_jh\{L^{-1}(1+r/L)^{s_c}
                                      +L^{-1}+L^{-2}r\}.
\]
After division by \(h(1+r)^{s_f}\), each term is bounded
by \(C_jL^{-3}\), since \(s_f\ge s_c\ge2\). The same
division handles \(r\le L\). No read-radius factor remains.
\end{proof}

\subsection{Why the affine normalisation gives a fifth power}

Estimating the constant part as an arbitrary fine direction
would retain only \(L^{-1}\). A local gauge transformation
removes it and leaves an affine term of order \(L^{-2}\).
The incoming affine Hessian cancels the square of that term.
The following calculation makes these two operations explicit.

\begin{lemma}[The local fifth-order gain]
\label{lem:sharp-fifth-order}
Let \(f=E_i\) on a nonwrapping footprint. Suppose that in
local exponential coordinates
\(a_\mu(x)=c_\mu+b_\mu(x)+\rho_\mu(x)\), where \(b\) is
linear in \(\xi\). For \(0<\tau\le1\), assume the constant
coefficients, linear coefficients, and \(|\rho|_i\), together
with their parameter jets through order \(q\), are bounded
respectively by \(Ch\tau,Ch\tau^2,C\tau^3\).
Suppose the queried fine plaquettes have chord at most
\(\eta h\), for sufficiently small fixed \(\eta>0\).
If the fine read radius \(R_i\) satisfies
\[
        h(1+R_i)^{2s_f+2q+10}\le\varepsilon(C,q,G),
\]
then every parameter jet through order \(q\) of
\(f(\exp a)\) has size at most \(C_je_i\tau^5\).
Only derivatives of \(f\) through \(\max(q,3)\) are used.
\end{lemma}
\begin{proof}
Put \(S_x=\sum_\nu\xi_\nu(x)c_\nu\) and use the site
gauge \(g_x=e^{S_x}\) on the read patch. It extends to a
torus gauge by a fixed spatial cutoff outside a larger
nonwrapping patch. On each read edge the transformed link is
\(e^{S_x}e^{a_\mu(x)}e^{-S_{x+e_\mu}}\).
The stated smallness condition puts all factors in a fixed
logarithmic chart. Expanding the logarithm of this product
of exponentials, with its convergent remainder, gives
\begin{equation}
 \widetilde a_\mu(x)
      =b_\mu(x)+\tfrac12[S_x,c_\mu]+\widetilde\rho_\mu(x),
                  \qquad |\widetilde\rho|_i\le C_j\tau^3.
 \label{eq:sharp-constant-gauge}
\end{equation}
This is the Baker--Campbell--Hausdorff expansion through
degree two. Indeed its linear constant terms cancel, and
their degree-two commutators sum to \(\frac12[S_x,c_\mu]\).
The first omitted constant-only terms are bounded by
\(Ch^3\tau^3(1+|\xi|)^2\), and the first mixed terms
with \(b\) by \(Ch^2\tau^3(1+|\xi|)^2\). After division
by \(h(1+|\xi|)^{s_f}\), both are at most \(C\tau^3\).
Terms containing \(\rho\) have its weighted bound times
bounded analytic coefficients. Higher commutators are
summed using \(\sup|S_x|\ll1\). Fixed-order chain and
product rules prove the same bounds for all parameter
jets. The first two terms of
\eqref{eq:sharp-constant-gauge} form an affine field
of weighted size \(C_j\tau^2\).

We must check the radial segments used for Taylor's formula.
For four small logarithms on a queried plaquette,
\[
        (\exp u)_p=1+(d_1u)_p+O(\|u\|_\infty^2).
\]
The endpoint \(\widetilde a\) has the same plaquette
chord as \(a\). This identity bounds
\(|(d_1\widetilde a)_p|\) by
\(C_G(\eta h+\|\widetilde a\|_\infty^2)\).
Apply it once more to \(t\widetilde a\), \(0\le t\le1\).
Choose \(\eta\) first; the stated smallness condition makes
\(\|\widetilde a\|_\infty^2/h\) small enough that every
such plaquette has chord less than \(h/2\).
Thus the whole segment is inside the incoming good domain.
Its four-edge mask halo was included in the physical reads.

Write \(F(u)=f((e^{u_e})_e)\). On this segment,
\[
 |\dd^mF(u)[v_1,\ldots,v_m]|
                         \le C_q e_i\prod_\nu|v_\nu|_i
                  \quad(0\le m\le\max(q,3)).
\]
The differential of the exponential converts a chart
direction to a left direction by an average of adjoint
isometries. Higher conversions multiply pointwise
directions; their extra factors are bounded by powers of
\(h(1+R_i)^{s_f}\le1\). The number of ordered terms depends
only on the derivative order, so there is no edge-count
factor. Lemma~\ref{lem:reset-flat-gauge-identities} gives
\(\dd F(0)=0\), and the incoming normalisations give
\(F(0)=0\) and \(H_0=\dd^2F(0)=0\) on two affine fields.

Write \(\widetilde a=a_2+a_3\), with \(a_2\) affine of
size \(C_j\tau^2\), and \(a_3\) of size \(C_j\tau^3\),
also in all parameter jets. Use the three separate identities
\begin{align*}
 F(u)&=\tfrac12H_0(u,u)
       +\tfrac12\int_0^1(1-t)^2\dd^3F(tu)[u,u,u]\,\dd t,\\
 \dd F(u)v&=H_0(u,v)
       +\int_0^1(1-t)\dd^3F(tu)[u,u,v]\,\dd t,\\
 \dd^2F(u)[v,w]&=H_0(v,w)
       +\int_0^1\dd^3F(tu)[u,v,w]\,\dd t.
\end{align*}
In each Hessian term the pair of affine parts cancels.
Every remaining pair has size \(C_j\tau^5\); each cubic
term has size \(C_j\tau^6\). For a parameter derivative
of \(F(\widetilde a)\), an ordered chain-rule term with
one or two argument factors uses the last two identities.
A term with \(m\ge3\) factors is bounded directly by
\(C_je_i\tau^{2m}\). This proves the asserted bound
through order \(q\). We have not differentiated a third-order
Taylor remainder through \(q\) additional orders.
Finally, gauge invariance gives the exact equality
\(f(\exp a)=f(\exp\widetilde a)\).
\end{proof}

These two calculations are
\cite[Lemmas~35.2--35.3, pp.~213--215]{Main}. Both the
affine-Hessian condition and the healthy radial segments
are essential to the fifth power.

\subsection{An inner protector on the original marginal}

Let \(T_i\) contain the active free coordinates of the
original fine insertion and all four edges of every
plaquette queried by its mask. Integrate out the remaining
Gaussian coordinates first, obtaining precisely
\(C_i=(C_{\rm sel})_{T_i,T_i}\). This is a marginal
covariance, not the inverse of a principal precision.
Tree links receive no Gaussian coordinate. The number of
real coordinates is at most \(C_j(1+B)^4(1+d_i)\).
This count uses fine \(B_f\)-balls on the original support;
it introduces no additional \(L^4\) factor.

Choose a fixed smooth \(\theta:[0,\infty)\to[0,1]\), one
on \([0,1/4]\) and zero on \([1,\infty)\). Fix \(M_j\)
from Lemma~\ref{lem:central-polynomial-ratio}, including
the finitely many enlargements needed below. After choosing
\(L\), choose \(\eta_L\) once from the fixed adjacent profiles
and \(L\), with
\begin{equation}
 0<\eta_L\ll\eta L^{-M_j-6},\qquad H_L=\eta_Lh/g,
 \qquad \Gamma_i(z)=\prod_{e\in T_i}\theta(|z_e|^2/H_L^2).
 \label{eq:sharp-inner-protector}
\end{equation}
Here \(\eta>0\) is the small-field constant already fixed
in Lemma~\ref{lem:sharp-fifth-order}. The principal estimate
uses only this upper bound on \(\eta_L\).
Prices involving \(\eta_L^{-1}\) belong to the post-ratio
cutoff comparison.
All colours on a read link are included. The factors are
radial, and for fixed discrete reads \(\Gamma_i\) is
independent of \(V\) at fixed \(z\).

The three relevant contributions must satisfy
\[
 C_jL^{-2}\ll\eta,\qquad
 C_jL^{M_j}\eta_L\ll\eta,\qquad
 C_{j,L}(1+B)^m(h_c^2+e^{-c_LB})\ll\eta h.
\]
Choose \(L\) for the first, \(\eta_L\) for the second,
and then increase entry for the third. On
\(\Gamma_i\ne0\), every queried fine plaquette has chord
less than the \(\eta h\) required in
Lemma~\ref{lem:sharp-fifth-order}, and the original broad
factors equal one. To check this quantitatively, the guide's
fine row has size \(C_jL^{-2}h_c\). The selected-centre error
in Lemma~\ref{lem:central-local-centre} has size
\(C_{j,L}(1+B)^m(h_c^2+e^{-c_LB})\). The added fluctuation
changes a fixed fine test by at most
\(C_jL^{M_j}\eta_Lh\). First make the two leading
coefficients smaller than \(\eta/4\), then make the centre
error smaller than \(\eta h/4\). This is a local argument
at each tested plaquette and works also for wrapping
\(X_i\). The central value bound
\eqref{eq:central-polynomial-mean}, together with
\(B^{17}\to\infty\), puts these coordinates in the region
where the broad factors are one.

Define the principal average
\begin{equation}
 P_i(V)=\operatorname{Av}\int\Gamma_i(z)
       E_i(\mathcal U_\sigma(V,\mu+gz))\,\dd\gamma_{C_i}(z).
 \label{eq:sharp-principal-average}
\end{equation}
The average is the same normalised shared tag average as
in \eqref{eq:central-exact-incoming-insertion}. All original
fine arguments and the actual marginal determinant are
retained. The inner protector is an auxiliary device for
this principal term; it has not narrowed a physical cutoff.

\subsection{Small and large incoming covers}

\begin{lemma}[Two address bounds for the principal average]
\label{lem:sharp-principal-two-cases}
There are \(C_j,m_j\), chosen before \(L\), such that above
its entry threshold
\begin{align}
 \|P_i\|_{q,h_c,s_c;\bar o_i,Y_i}
                      &\le C_jL^{-5}e_i,&&d_i^3<L,\notag\\
 \|P_i\|_{q,h_c,s_c;\bar o_i,Y_i}
                      &\le C_jL^{m_j}(1+d_i)^{m_j}e_i,
                                                     &&d_i^3\ge L.
 \label{eq:sharp-principal-two-cases}
\end{align}
Each function has a gauge invariant \(C^q\) germ on its
healthy coarse domain, including threshold equality.
\end{lemma}
\begin{proof}
Suppose first \(d_i^3<L\). A connected cover bounds the
fine read radius by \(C(1+B)d_i\), and its coarse projection
by \(CBd_i/L+C\). For sufficiently large \(L\), the latter
lies in a fixed small fraction of a radius-\(B\) box. The
declared full reads contain the healthy local box and
windows needed for the following comparison. The period
reserve makes the fine patch nonwrapping, since its radius
is at most \(C(1+B)L^{1/3}\), while its side is \(L\)
times the coarse period.

Choose the coarse axial gauge at the base background and
keep it fixed while differentiating. Its base logarithms
obey the sharper pointwise bound
\begin{equation}
                   |x_e|\lesssim_G h_c(1+d_c(e,\bar o_i)).
 \label{eq:sharp-pointwise-axial-bound}
\end{equation}
For this bound, compare the coordinate path to the initial
endpoint, followed by the one edge, with the coordinate
path to the final endpoint. Move that one edge past the
later coordinate segments. At most \(C(1+d_c(e,\bar o_i))\)
elementary plaquettes are interchanged; all earlier path
segments cancel. The product telescoping estimate bounds
its chord by that number times \(h_c\), and the small-chart
logarithm has the same bound. The implicit constant depends
only on the fixed group and dimension, independently of
the block ratio. This refines, for these
particular one-edge loops, the sufficient quadratic-radius
bound in \eqref{eq:local-gauge-value}.

The first logarithmic variation, with this fixed gauge,
has the prescribed coarse direction weight. Higher chart
variations contain at least two direction factors; on the
local box their extra normalised size is
\(C_{j,L}(1+B)^m h\). Thus the base field and first
variations satisfy the hypothesis of
Lemma~\ref{lem:sharp-three-parts}, with a pre-ratio constant;
the higher chart variations can be included in a small
error after the entry choice.

We give the remaining linear comparison explicitly.
Lemma~\ref{lem:central-local-centre} identifies the first
mean jet with \(\mathcal Hx\), with error
\(C_{j,L}(1+B)^m(h_c^2+e^{-c_LB})\) in the required jets.
The replica derivative \(\mathcal R_{\rm rep}\) and the
tree-slice inclusion \(\iota\) satisfy
\[
 \mathcal R_{\rm rep}+\iota\mathcal H
             =\widehat E=E+d_{0,f}\omega(E),
\]
as proved after \eqref{eq:calibrated-frame-correction}.
A fine gauge change inside the gauge invariant function
\(E_i\) removes this gradient. It changes neither the
selected marginal nor \(\Gamma_i\). Taylor expansion of
the finite coordinate products and logarithms adds only
terms with at least two small coordinate or direction
factors. By Lemma~\ref{lem:central-polynomial-ratio}, the
protected fluctuation has incoming weighted size at most
\(C_jL^{M_j}\eta_L\), also in its coarse jets. Consequently
the fine logarithm and each prescribed ordered coarse jet
are a free interpolation jet plus that fluctuation and
an error with incoming size at most
\begin{equation}
 C_{j,L}(1+B)^m(1+d_i)^m
                       \{h+e^{-c_LB}/h\}.
 \label{eq:sharp-local-linear-error}
\end{equation}
The division by \(h\) leaves a factor \(h\) from each
quadratic chart error; the window error gives the second
term. Fixed-order chain rules and the inverse boundary
identity justify these estimates also for the ordered jets.
There is no covariance square root acting on \(z\): the
installed fluctuation is exactly \(gz\).

Lemma~\ref{lem:sharp-three-parts} decomposes the free jets
with sizes \(L^{-1},L^{-2},L^{-3}\). The fluctuation fits
in the last part by the choice of \(\eta_L\). Since
\(d_i<L^{1/3}\), the error
\eqref{eq:sharp-local-linear-error} is at most \(L^{-3}\)
at a ratio-dependent entry coupling, uniformly over this
class. The read-radius smallness condition of
Lemma~\ref{lem:sharp-fifth-order} holds there as well.
That lemma, with \(\tau=L^{-1}\), bounds the composition
and all its jets on the protector by \(C_jL^{-5}e_i\).

We must also differentiate the marginal density. Keep
\(z\) fixed and leave \(\mu\) inside the fine field.
For the centred Lebesgue density \(\varphi_{C_i}\),
\begin{equation}
 D_A\log\varphi_{C_i}(z)
  =-\tfrac12\operatorname{Tr}(C_i^{-1}D_AC_i)
   +\tfrac12 z^{\mathsf T}C_i^{-1}(D_AC_i)C_i^{-1}z.
 \label{eq:sharp-centred-score}
\end{equation}
There is no \(\mu/g\) term in this score. Every
positive-order covariance derivative supplies at least
one factor \(h_c\). Lemma~\ref{lem:central-polynomial-ratio},
the coordinate count and \(|z_e|\le H_L\le CP\) bound
every positive-order density score on the protector by
\begin{equation}
 C_jL^{m_j}h(1+B)^{m_j}(1+P)^{m_j}(1+d_i)^{m_j}.
 \label{eq:sharp-positive-score-price}
\end{equation}
The ordered density partition formula includes all score
products. For \(d_i<L^{1/3}\) their displayed bound is
at most one after the entry choice; \(h\) dominates every
fixed power of \(B,P\). The zeroth density has integral
one and \(\Gamma_i\) is independent of \(V\). The product
rule therefore preserves \(C_jL^{-5}e_i\), with a constant
chosen before \(L\).

If \(d_i^3\ge L\), no affine normalisation or global fine
chart is used. All queried fine plaquettes are still
healthy on the protector. Differentiate the original
incoming function using its fine regular norm. The first
pulled-back left direction has incoming size
\(C_jL^{m_j}\), by the polynomial-moment bounds and
\(s_f\ge s_c\). For an \(r\)-fold direction coefficient,
its direction-weight contribution divided by one allowed
fine direction is bounded by
\[
 C_r\frac{h_c^r(1+\rho/L)^{rs_c}}
                 {h(1+\rho)^{s_f}}
       \le C_rh^{r-1}(1+\rho)^{(r-1)s_c},
       \qquad \rho\le C(1+B)(1+d_i).
\]
For \(r\ge2\), every resulting power of \(B\) is paired
with a positive power of \(h\), and is at most one at
the entry threshold after the finite fixed orders are
chosen. The remaining loss is a fixed polynomial in
\(1+d_i\), with a pre-ratio polynomial price in \(L\).
The same reasoning applies to
\eqref{eq:sharp-positive-score-price}. Ordered chain
and product rules now give the second line of
\eqref{eq:sharp-principal-two-cases}, including wrapping
supports.

The strict fine healthy margin and positive covariance
persist in a coarse neighbourhood for each fixed finite
read set. Compact support of \(\Gamma_i\) justifies
differentiation under the marginal integral and proves
the germs, also at equality in a coarse threshold. The
radial protector, full colour reads, covariant records
and the unchanged tag probability give gauge invariance.
\end{proof}

This is \cite[Lemma~35.5, pp.~217--219]{Main}. In the
small-cover proof, the entry choice makes the positive-order
scores bounded by one; it does not absorb an uncontrolled
post-ratio constant into the leading \(L^{-5}\) coefficient.

\subsection{Restoring the exact incoming insertion}

On \(\Gamma_i\ne0\), both \(\chi_{X_i}\) and all retained
\(b_e\)'s are one. The original average
\eqref{eq:central-exact-incoming-insertion} therefore has
the exact decomposition
\begin{equation}
 \begin{split}
 A_i&=P_i+R_i^{\rm cut},\\
 R_i^{\rm cut}(V)&=\operatorname{Av}\int
    \left(\prod_{e\in T_i}b_e\right)(1-\Gamma_i(z))
             \chi_{X_i}(U)E_i(U)\,\dd\gamma_{C_i}(z).
 \end{split}
 \label{eq:sharp-exact-cutoff-remainder}
\end{equation}
All these identities are on the healthy coarse domain of
\(Y_i\), before restoring their original coarse occurrence
masks in the exact reassembly.

\begin{lemma}[The rare cutoff comparison]
\label{lem:sharp-cutoff-remainder}
For every fixed \(J_0>0\), after a ratio-dependent entry
choice,
\begin{equation}
 \|R_i^{\rm cut}\|_{q,h_c,s_c;\bar o_i,Y_i}
        \le C_{j,L,J_0}\widehat\beta_f^{-J_0}
                                      (1+d_i)^{m_j}e_i.
 \label{eq:sharp-cutoff-address-price}
\end{equation}
The exponent \(m_j\) can be fixed independently of
\(J_0,L,B\); the constants are uniform in period and cutoff.
\end{lemma}
\begin{proof}
Convert all Borel physical reads back to Haar variables
using Section~\ref{sec:regular-output-derivatives}. Now
\(\chi_{X_i}(U)E_i(U)\) is fixed and bounded by \(e_i\).
Derivatives land on the exact density, retained broad
factors and smooth protector. Their scores cost a fixed
polynomial in \(B,P,1+d_i\), with a ratio-dependent
constant. Protector derivatives cost powers of
\(h_c/(gH_L)\le C/\eta_L\), allowed in that constant.

Every derivative term retains the event
\(\max_{e\in T_i}|z_e|\ge H_L/2\). For an undifferentiated
protector this follows from \(1-\Gamma_i\ne0\); for a
differentiated factor it follows from its transition
support. The broad factors bound all read coordinates by
\(C_{j,L}(1+B)^mP\), so the score polynomials have a
pointwise bound of the asserted fixed degree. Each
coordinate variance is at most \(C_jL^{M_j}\). A
one-coordinate Gaussian tail and a union bound give
\[
 \Pr\{\max_{e\in T_i}|z_e|\ge H_L/2\}
        \le C_j(1+B)^4(1+d_i)e^{-c_{j,L}P^2}.
\]
No independence assumption is used. All colour dimensions
are fixed. Combining this with the score polynomials
gives a fixed polynomial times this exponential. Since
\(P^2/\log\widehat\beta_f\to\infty\) and
\(\log B,\log P=o(P^2)\), that expression is at most
\(C_{j,L,J_0}\widehat\beta_f^{-J_0}(1+d_i)^{m_j}e_i\)
for every fixed \(J_0\). This proves the estimate with
its exact marginal determinant and Haar normalisation;
no extensive density constant was introduced.
\end{proof}

\subsection{The anchored gain and exact normalised output}

\begin{theorem}[Sharp principal regular transport]
\label{thm:sharp-principal-transport}
There are \(V_j<\infty\) and a geometric lower ratio bound,
chosen before \(L\), such that for each larger allowed odd
ratio, and then sufficiently large entry coupling,
\begin{align}
 \|P\|_{\mathrm{raw},\kappa_c}&\le\frac{V_j}{L}\delta,
 \label{eq:sharp-anchored-contraction}\\
 \|R^{\rm cut}\|_{\mathrm{raw},\kappa_c}
                &\le C_{j,L,J_0}\widehat\beta_f^{-J_0}\delta
                         \quad\hbox{for every fixed }J_0>0.
 \label{eq:sharp-cutoff-family-price}
\end{align}
The families are local and gauge invariant, and retain
their lattice covariance and original shared tag averages.
The bounds are uniform in finite history, cutoff and period.
\end{theorem}
\begin{proof}
If \(d_i^3\ge L\), then \(L^{m_j}\le d_i^{3m_j}\) and
\(L^5\le d_i^{15}\). The two cases in
\eqref{eq:sharp-principal-two-cases} consequently give,
for a fixed exponent \(p_j\),
\[
 \|P_i\|_{q,h_c,s_c;\bar o_i,Y_i}
                      \le C_jL^{-5}(1+d_i)^{p_j}e_i
                                      \quad\hbox{for every }i.
\]
At a fixed coarse anchor, use
\eqref{eq:sharp-cover-exponent} with \(\kappa_*\ge\kappa_c\)
and pay the remaining polynomial by
\(\sup_{d\ge1}(1+d)^{p_j}e^{-\kappa_fd/2}<\infty\).
Exactly \(L^4\) possible fine anchors project to a given
coarse anchor. At each one the incoming anchored sum is
bounded by \(\delta\). Thus \(L^{-5}L^4=L^{-1}\) proves
\eqref{eq:sharp-anchored-contraction}, with its constant
fixed before \(L\). We do not first convert this sharp
bound into incidence or charge every site of \(Y_i\).
The same calculation applied to
\eqref{eq:sharp-cutoff-address-price} proves
\eqref{eq:sharp-cutoff-family-price}; its factor \(L^4\)
is allowed in the post-ratio constant.

All read enlargements were the original ones. Radial
protector factors, covariant marginal matrices and the
same normalised tag laws preserve gauge invariance and
lattice covariance. Tag cardinality adds no factor,
since its weights form a probability average. Uniformity
follows from the record, column and period bounds used
above.
\end{proof}

\begin{corollary}[Normalised transport with all mask terms retained]
\label{cor:sharp-normalised-transport}
For \(q\ge2\), put \(Y_i'=Y_i\cup B_\infty(\bar o_i,2)\), and use
Section~\ref{sec:affine-local-reset} to define
\(\bar P_i=P_i-P_i(1)-Q_{P_i}\). There is
\(\widetilde V_j\), chosen before \(L\), such that
\begin{equation}
             \|\bar P\|_{\mathrm{raw},\kappa_c}
                            \le\frac{\widetilde V_j}{L}\delta.
 \label{eq:sharp-normalised-contraction}
\end{equation}
This family has zero flat value and zero affine Hessian
on every nonwrapping enlarged footprint. Exactly,
\begin{equation}
 \chi_{Y_i}P_i=P_i(1)+\chi_{\{\bar o_i\}}Q_{P_i}
                        -u_i^{\rm reset}+\chi_{Y_i'}\bar P_i.
 \label{eq:sharp-normalisation-mask}
\end{equation}
The Borel function \(u_i^{\rm reset}\) vanishes on a
healthy \(Y_i'\) and has value bound
\(C_jL^{-5}(1+d_i)^{p_j}e_i\).
\end{corollary}
\begin{proof}
Lemma~\ref{lem:reset-local-norm} costs only
\(C_r(G,q,s_c)\), independent of \(Y_i,B,h_c,L\).
The covering number increases by at most one. Hence
\(\widetilde V_j=C_re^{\kappa_c}V_j\) suffices. Its
exact masked identity is
\eqref{eq:sharp-normalisation-mask}, with
\(u_i^{\rm reset}=0\) on a healthy \(Y_i'\) and otherwise
\(P_i(1)+\chi_{\{\bar o_i\}}Q_{P_i}-\chi_{Y_i}P_i\).
The same address bound proves its value price. The
quadratic coefficient price \(h_c^{-2}\) is retained;
the healthy clover product contributes \(h_c^2\), as
proved in the local subtraction lemma. Scalars,
quadratic terms and these masked discrepancies are
kept for the exact scalar and principal-action reset.
\end{proof}

The theorem and corollary are
\cite[Theorem~35.7 and Corollary~35.8, pp.~219--221]{Main}.
Combining the exact cutoff identity with
\eqref{eq:central-incoming-comparison} gives the full
warning-free incoming comparison
\begin{equation}
 \begin{split}
 F(1,1)-F(1,0)&=-\sum_iP_i+R_{\rm inh},\\
 \|R_{\rm inh}\|_{\mathrm{raw},\kappa}
   &\le C_B(bd+d^2)
               +C_{j,L,J_0}\widehat\beta_f^{-J_0}\delta
     \le\widehat\beta_f^{-1/2+o(1)}\delta .
 \end{split}
 \label{eq:sharp-full-incoming-comparison}
\end{equation}
The sign agrees with \(e^{-\chi_{X_i}E_i}\). For the last
bound choose one fixed \(J_0>1/2\). Applying the same
bounded local subtraction to this comparison family
preserves its small norm while retaining its own scalar,
quadratic and masked terms.

Collect the choices in their order of dependence.
First fix \(V_j,\widetilde V_j\), all pre-ratio prices
and the geometric lower bounds for the finitely many
adjacent profile pairs. Choose an allowed odd \(L\) with
\[
                         \frac{\max(V_j,\widetilde V_j)}L\ll1,
\]
where the smallness threshold fits the prescribed incoming
fraction. Next fix the inner width
\eqref{eq:sharp-inner-protector}. Only then select entry.
The chart and score conditions have the form
\[
 C_{j,L}(1+B)^{m_j}(1+P)^{m_j}
                  \{h+e^{-c_LB}/h\}\ll1,
\]
after the finitely many required powers and post-ratio
constants have been collected. Fixed powers of \(L\)
in these entry conditions are allowed in \(C_{j,L}\);
they do not alter the already fixed leading contraction
constant. For a prescribed \(J_0>1/2\), the rare-cutoff
coefficient also satisfies
\(C_{j,L,J_0}\widehat\beta_f^{-J_0}\ll1\) after increasing
the same threshold. The definitions of \(B,h,P\) verify
these comparisons simultaneously. These choices use the
same slow clock and period reserve, independently of
the requested volume subsequence.

This completes the packet-free higher central estimates
and the sharp incoming regular transport. Section~\ref{sec:matching-taylor-data} now identifies the first
two retained ordinary coefficients and derives the exact classical
and one-loop Taylor data. Their curvature matching must precede
the auxiliary-action bound and its finite-period localisation. The scalar and quadratic
ledgers must then be assembled with the exact masked
reassembly. Neither the estimates in this section nor
compilation of their text establishes closure of the
complete nonlinear RG step or a continuum construction.

\section{Taylor data for matching the principal action}
\label{sec:matching-taylor-data}

The incoming regular family now contracts with its constant
fixed before the block ratio. The warning-free central
logarithm has a small remainder after its first two ordinary
strength coefficients are removed. Those coefficients are
still present in the exact action. To decide what belongs
in the next principal action, we need their precise retained
factors and the low field derivatives of the stationary
prefactor.

There are two different differentiations here. An ordinary
strength \(u\) multiplies interaction letters while the
physical \(g=\beta^{-1/2}\), cutoffs and covariance are fixed.
A field derivative varies the coarse exponential coordinates
\(x\). We first keep the former coefficients exactly, then
compute the latter derivatives from the stationary equation.
Their comparison and curvature matching are subsequent steps;
the two operations are not identified term by term.

Fix the adjacent profiles, their prescribed path budgets and
an allowed odd ratio \(L\). Unless a pre-ratio bound is
explicitly stated, constants in this section may depend on
these fixed data. The exponential kernel bounds will be
uniform in finite history, cutoff and compatible period.
We continue to use the actual unmasked phase
\eqref{eq:stationary-model-action}, its stationary graph
\(m\), and its positive free-coordinate Hessian
\(H(x)\). All statements below concern their common
regular chart.

\subsection{The two coefficients with every chart factor retained}

Let \(A\) be a finite set of ordinary letter labels.
Write \(v_\alpha\) for the logarithmic interaction assigned
to label \(\alpha\), and \(T_\alpha\) for the physical
Gaussian coordinates at which that occurrence retains a
broad chart factor. Put
\[
                   b_T=\prod_{e\in T}b_e,\qquad b_\varnothing=1.
\]
The \(b_e\)'s are the original broad factors of
Section~\ref{sec:physical-selected-Gaussian}. A smooth letter
with no such factor has \(T_\alpha=\varnothing\).
Expectations below use the original common Gaussian law
and the same normalised shared tag averages. Overlapping
letters keep the same variables.

Before taking a logarithm, the finite ordinary expansion is
\begin{equation}
 \begin{split}
 Z_A(u)&=1+\sum_{\varnothing\ne J\subset A}
       \mathbb E\left[
         b_{\bigcup_{\alpha\in J}T_\alpha}
                     \prod_{\alpha\in J}(e^{uv_\alpha}-1)\right],\\
 F_A(u)&=\log Z_A(u),\qquad F_A(0)=0 .
 \end{split}
 \label{eq:matching-retained-partition}
\end{equation}
The union in the subscript is important: each physical chart
factor occurs once, even when several labels read it.
Labels remain distinct if their functions happen to agree.

\begin{lemma}[The first two ordinary strength coefficients]
\label{lem:matching-retained-coefficients}
For any ordering of \(A\), define
\begin{align}
 a_1&=\sum_{\alpha\in A}\mathbb E[b_{T_\alpha}v_\alpha],
                                                        \notag\\
 a_2&=\sum_{\alpha\in A}\mathbb E[b_{T_\alpha}v_\alpha^2]
       +2\sum_{\alpha<\gamma}
         \mathbb E[b_{T_\alpha\cup T_\gamma}
                                      v_\alpha v_\gamma].
 \label{eq:matching-retained-moments}
\end{align}
Then
\begin{equation}
               F_A'(0)=a_1,\qquad F_A''(0)=a_2-a_1^2.
 \label{eq:matching-retained-cumulants}
\end{equation}
These formulas extend to the ordinary family of
Section~\ref{sec:central-strength-estimates}, with its
prescribed ordered physical derivatives and anchored
family norm.
\end{lemma}
\begin{proof}
Since every \(e^{uv_\alpha}-1\) vanishes at zero, a first
derivative can leave only a singleton \(J\).
A second derivative can leave either one label differentiated
twice or two different labels each differentiated once.
In the latter case the two derivative orders give the
factor two. Thus \(Z_A(0)=1\), \(Z_A'(0)=a_1\) and
\(Z_A''(0)=a_2\). Direct differentiation gives
\[
 (\log Z_A)'=\frac{Z_A'}{Z_A},\qquad
 (\log Z_A)''=\frac{Z_A''}{Z_A}
                              -\frac{(Z_A')^2}{Z_A^2},
\]
which proves \eqref{eq:matching-retained-cumulants}.
The ordinary second Taylor coefficient is
\((a_2-a_1^2)/2\), including its factorial.

To pass to the full family, use the absolute labelled
majorant already proved in
Section~\ref{sec:central-strength-estimates}.
At each fixed period the anchored sum also gives a finite
total sum after summing its finitely many anchors.
For a circle \(|u|=r\) strictly inside the proved strength
radius, the coefficient formula
\[
 [u^k]F_A(u)=\frac1{2\pi i}
                  \int_{|u|=r}\frac{F_A(u)}{u^{k+1}}\,\dd u
\]
is dominated by that same connected-family majorant.
The bound holds after every prescribed ordered physical
derivative. Absolute convergence therefore permits the
contour integral, family sum and fixed derivative in either
order. Exhausting the labels proves the assertion, with the
volume-uniform anchored estimate. Physical Borel masks
remain fixed in the Haar representation used there.
\end{proof}

For example, two letters with the same read \(T\) contribute
\(\mathbb E[b_Tv_\alpha v_\gamma]\), not an expression with
\(b_T^2\). Smooth broad factors need not be idempotent.
Only when the retained factors are identically one on the
integration domain do \eqref{eq:matching-retained-moments}
reduce to the usual two moments of
\(\sum_\alpha v_\alpha\).
Thus their later replacement by full Gaussian formulas
requires a comparison estimate.

The ordinary remainder already has norm
\(\widehat\beta_f^{-3/2+o(1)}\), by
\eqref{eq:central-remainder-orders}; the present formulas
identify precisely the two coefficients it leaves.
They are the retained coefficients of
\cite[Lemma~36.1, p.~228]{Main}.
No chart factor or shared tag probability has been
discarded in identifying them.

\subsection{The classical action and its exact quadratic term}

Retain the existing notation
\[
 \Psi(x)=\varphi(x,m(x)),\qquad
 H(x)=\varphi_{yy}(x,m(x)),\qquad
                            H_0=H(0).
\]
The function \(\Psi\) is the stationary classical action.
It does not contain the inherited loop term or the Haar
density. In particular the common \(\beta\) multiplying
the leading phase cancels from its stationary equation.
The fine inverse \(H_0^{-1}\) is the one in the stationary
construction, with the counting-coordinate and exponential
bounds already proved; no new covariance is chosen.

We first check the quadratic term because it is the part
that the coarse row action must already reproduce.

\begin{lemma}[The calibrated quadratic classical action]
\label{lem:matching-classical-quadratic}
The flat value and first derivative of \(\Psi\) vanish.
For a coarse direction \(v\),
\begin{equation}
          D^2\Psi(0)[v,v]
                  =\langle d_cv,H_cd_cv\rangle .
 \label{eq:matching-classical-quadratic}
\end{equation}
Here \(H_c\) is the unmodified coarse curvature metric.
\end{lemma}
\begin{proof}
The flat phase calculation gives
\(\varphi(0,0)=0\), \(D\varphi(0,0)=0\), and
\(m(0)=0\). This proves the first assertion.
Use the actual row derivatives \(\mathsf A,\mathsf B\)
of Section~\ref{sec:calibrated-rows}. At the flat field,
\[
 D^2\varphi(0,0)[(v,y),(v,y)]
                           =|\mathsf Bv+\mathsf Ay|_2^2,
                     \qquad H_0=\mathsf A^{\mathsf T}\mathsf A .
\]
The higher curvature atoms have no quadratic term, so
these are the Hessian and mixed Hessian of the full phase.
Differentiating the stationary equation once gives
\[
                    H_0Dm(0)v=-\mathsf A^{\mathsf T}\mathsf Bv.
\]
The calibrated identities
\eqref{eq:calibrated-two-identities} imply
\(\mathsf A^{\mathsf T}(\mathsf B+\mathsf A\mathcal H)=0\).
Uniqueness of the solution therefore gives
\(Dm(0)=\mathcal H\). Since the fine gradient is zero
at the stationary point,
\[
 D^2\Psi(0)[v,v]
      =|(\mathsf B+\mathsf A\mathcal H)v|_2^2
      =|\Phi_0S_cd_cv|_2^2.
\]
Lemma~\ref{lem:packet-complement} supplies the exact
full-row completion and \(\Phi_\perp S_cd_c=0\). Applying
its Gram identity to this particular vector gives
\[
 |\Phi_0S_cd_cv|_2^2=|S_cd_cv|_2^2
      =\langle d_cv,\widetilde H_cd_cv\rangle
      =\langle d_cv,H_cd_cv\rangle .
\]
The last equality uses
\(\widetilde H_cd_c=H_cd_c\).
Only the isometry on coarse curls was needed.
\end{proof}

Thus the quadratic agreement follows from the calibrated
Gaussian history. The new matching data begin at field
degrees three and four.

\subsection{Stationarity computes the cubic and quartic terms}

For coarse directions \(v_i\) set
\[
 m_i=Dm(0)v_i,\qquad
 m_{ij}=D^2m(0)[v_i,v_j],\qquad
                               w_i=(v_i,m_i).
\]
The pair \(w_i\) is a direction in the joint \((x,y)\)
coordinate space, tangent to the stationary graph.
Define the fine linear functional
\begin{equation}
 B_{ij}[u]=D^3\varphi(0,0)[w_i,w_j,(0,u)].
 \label{eq:matching-mixed-cubic}
\end{equation}
We identify a fine linear functional with its vector in the
fixed real inner product. Accordingly
\(B_{ij}H_0^{-1}B_{kl}\) means
\(\langle B_{ij},H_0^{-1}B_{kl}\rangle\).

\begin{lemma}[The cubic and quartic stationary jets]
\label{lem:matching-classical-jets}
The first two differentiated stationary identities are
\begin{equation}
 \begin{split}
 D^2\varphi(0,0)[w_i,(0,u)]&=0,\\
                         H_0m_{ij}&=-B_{ij}.
 \end{split}
 \label{eq:matching-stationary-identities}
\end{equation}
The classical derivatives are
\begin{align}
 D^3\Psi(0)[v_1,v_2,v_3]
    &=D^3\varphi(0,0)[w_1,w_2,w_3],\notag\\
 D^4\Psi(0)[v_1,v_2,v_3,v_4]
    &=D^4\varphi(0,0)[w_1,w_2,w_3,w_4]\notag\\
    &\quad-\sum_{(ij|kl)}B_{ij}H_0^{-1}B_{kl}.
 \label{eq:matching-classical-jets}
\end{align}
The sum has the three unordered pairings
\((12|34),(13|24),(14|23)\).
\end{lemma}
\begin{proof}
For a fixed fine direction \(u\), the function
\(D_y\varphi(x,m(x))[u]\) is identically zero.
Its first derivative is the first line of
\eqref{eq:matching-stationary-identities}.
Differentiating it once more in directions \(v_i,v_j\)
gives
\[
 0=D^3\varphi(0,0)[w_i,w_j,(0,u)]
             +D^2\varphi(0,0)[(0,m_{ij}),(0,u)].
\]
The second term is \(\langle H_0m_{ij},u\rangle\).
This proves the second identity and fixes its minus sign.

To see the symmetry factors without an expanded multivariable
chain rule, set every coarse direction equal to \(v\).
Put \(z(t)=(tv,m(tv))\), so that
\[
 z'(0)=w=(v,Dm(0)v),\quad
 z''(0)=(0,r),\quad z'''(0)=(0,s),
\]
where \(r=D^2m(0)[v,v]\) and
\(s=D^3m(0)[v,v,v]\).
The third derivative of \(\varphi(z(t))\) is
\[
 D^3\varphi[w,w,w]
       +3D^2\varphi[w,(0,r)]
       +D\varphi[(0,s)]
\]
at zero. The last two terms vanish by stationarity and
the first identity in
\eqref{eq:matching-stationary-identities}.
This gives the diagonal cubic identity.

At fourth order the complete chain rule is
\[
 \begin{split}
 \frac{\dd^4}{\dd t^4}\varphi(z(t))\big|_{t=0}
  ={}&D^4\varphi[w,w,w,w]
        +6D^3\varphi[w,w,(0,r)]\\
    &+3D^2\varphi[(0,r),(0,r)]
        +4D^2\varphi[w,(0,s)]
        +D\varphi[z^{(4)}(0)].
 \end{split}
\]
Every phase derivative in this display is evaluated at
\((0,0)\). The fourth and fifth terms vanish for the
same reasons. If \(B_{vv}\) denotes
\eqref{eq:matching-mixed-cubic} with both directions
equal to \(v\), the remaining correction is
\[
 \begin{split}
 6B_{vv}[r]+3\langle r,H_0r\rangle
 &=6B_{vv}[-H_0^{-1}B_{vv}]
            +3\langle H_0^{-1}B_{vv},B_{vv}\rangle\\
 &=-3B_{vv}H_0^{-1}B_{vv}.
 \end{split}
\]
Polarisation of the symmetric cubic and quartic forms
gives \eqref{eq:matching-classical-jets}. In the quartic
identity the three pairings each become the same
\(B_{vv}H_0^{-1}B_{vv}\) on the diagonal, which accounts
for their coefficient one.
\end{proof}

For later matching define the homogeneous Taylor term of
any analytic \(f\) by
\(T_kf(x)=D^kf(0)[x,\ldots,x]/k!\).
The preceding calculation reads explicitly
\begin{equation}
 \begin{split}
 \Psi(tv)
  ={}&\frac{t^2}{2}\langle d_cv,H_cd_cv\rangle
       +\frac{t^3}{6}D^3\varphi(0,0)[w,w,w]\\
    &+t^4\left\{\frac1{24}D^4\varphi(0,0)[w,w,w,w]
                     -\frac18B_{vv}H_0^{-1}B_{vv}\right\}
       +O(t^5).
 \end{split}
 \label{eq:matching-classical-diagonal}
\end{equation}
This is a Taylor expansion at fixed finite volume.
The coefficients have the rooted bounds below; no
volume-uniform total scalar remainder is asserted here.
The calculation is
\cite[Lemma~36.2 and (36.5), p.~229]{Main}, expressed
using the existing \(H_0\). No third or fourth derivative
of \(m\) has to be estimated to obtain these jets.

\subsection{The Haar jet and the full one-loop action}

The next field-dependent term in the stationary prefactor
has no \(\beta\) in front of it. Define
\begin{equation}
 \begin{split}
 \mathcal I(x,y)&=I_f^{\rm un}(U(x,y)),\\
 T(x)&=\mathcal I(x,m(x))
           -\sum_{e\in I}\log j_G(m_e(x))
                         +\frac12\log\det H(x).
 \end{split}
 \label{eq:matching-one-loop-action}
\end{equation}
Thus \(T=\psi_*+\tfrac12\log\det H\).
We call this combination the one-loop action: it includes
the inherited loop term, the Haar coordinate density and
the determinant of the same stationary Hessian.
The determinant is taken on the actual free fine
coordinate space of real dimension \(d=|I|\dim_\R\g\).

The full prefactor \eqref{eq:tail-full-prefactor} gives
the exact scalar ledger
\begin{equation}
 \begin{split}
 -\log N_g^{\mathrm{phys}}(x)
   ={}&\beta\Psi(x)+T(x)+b_c\log z_{\mathrm n}(\beta)
           -|I|\log\zeta_G\\
     &\hspace{17pt}-\frac d2\log(2\pi)-d\log g .
 \end{split}
 \label{eq:matching-prefactor-ledger}
\end{equation}
Consequently
\[
 -\log\frac{N_g^{\mathrm{phys}}(x)}{N_g^{\mathrm{phys}}(0)}
                    =\beta\Psi(x)+T(x)-T(0),
 \qquad T(0)=\tfrac12\log\det H_0 .
\]
The field-independent factors have been separated, with
their original signs, and remain in the vacuum
normalisation.

We compute the Haar contribution before differentiating
\(T\). This also fixes its sign in the quadratic read.

\begin{lemma}[The quadratic Haar jet]
\label{lem:matching-haar-jet}
For \(X,Y\in\g\),
\begin{equation}
 D\log j_G(0)=0,\qquad
 D^2\log j_G(0)[X,Y]
                  =\frac1{12}\operatorname{Tr}_{\g}
                                     (\ad X\,\ad Y).
 \label{eq:matching-haar-jet}
\end{equation}
Here \(\ad X(Y)=[X,Y]\), and
\(\operatorname{Tr}_{\g}\) is the trace of a real linear
map on \(\g\) in the fixed orthonormal basis.
\end{lemma}
\begin{proof}
Appendix~\ref{app:haar} gives
\(j_G(X)=\det A_X\) near zero, where
\(A_X=\int_0^1e^{-s\ad X}\,\dd s\).
Expanding the finite-dimensional operator exponential,
\[
 A_X=1-\tfrac12\ad X+\tfrac16(\ad X)^2+O(|X|^3).
\]
The logarithm series and
\(\log\det A_X=\operatorname{Tr}_{\g}\log A_X\)
therefore give
\[
 \log j_G(X)
   =-\tfrac12\operatorname{Tr}_{\g}(\ad X)
       +\left(\tfrac16-\tfrac18\right)
                       \operatorname{Tr}_{\g}((\ad X)^2)
       +O(|X|^3).
\]
The invariant real inner product makes \(\ad X\)
skew-adjoint, so its trace is zero. The quadratic
coefficient is \(1/24\). Substitute \(sX+tY\);
the \(st\) coefficient is
\[
 \frac1{24}\operatorname{Tr}_{\g}
               (\ad X\,\ad Y+\ad Y\,\ad X)
       =\frac1{12}\operatorname{Tr}_{\g}(\ad X\,\ad Y).
\]
This proves both identities.
In particular the contribution of \(-\log j_G\) is
nonnegative on a repeated real direction, since
\(\operatorname{Tr}_{\g}((\ad X)^2)\le0\).
\end{proof}

\begin{lemma}[The complete one-loop quadratic read]
\label{lem:matching-one-loop-jet}
Let \(R_0=H_0^{-1}\), and write
\(H_i=DH(0)[v_i]\) and
\(H_{ij}=D^2H(0)[v_i,v_j]\).
For fine directions \(u_1,u_2\), these operators are
determined by
\begin{align}
 \langle u_1,H_i u_2\rangle
  &=D^3\varphi(0,0)[w_i,(0,u_1),(0,u_2)],\notag\\
 \langle u_1,H_{ij}u_2\rangle
  &=D^4\varphi(0,0)[w_i,w_j,(0,u_1),(0,u_2)]\notag\\
  &\quad+D^3\varphi(0,0)
                    [(0,m_{ij}),(0,u_1),(0,u_2)].
 \label{eq:matching-hessian-variations}
\end{align}
The complete second derivative of \(T\) is
\begin{align}
 D^2T(0)[v_i,v_j]
   ={}&D^2\mathcal I(0,0)[w_i,w_j]\notag\\
      &-\frac1{12}\sum_{e\in I}
           \operatorname{Tr}_{\g}(\ad m_{i,e}\,\ad m_{j,e})
                                                        \notag\\
      &+\frac12\operatorname{Tr}(R_0H_{ij})
          -\frac12\operatorname{Tr}
                            (R_0H_iR_0H_j).
 \label{eq:matching-complete-loop-read}
\end{align}
The traces in the last line are on the free fine
coordinate space. The quadratic Taylor term is
\(T_2T(x)=D^2T(0)[x,x]/2\).
\end{lemma}
\begin{proof}
For fixed \(u_1,u_2\), differentiate
\[
 \langle u_1,H(x)u_2\rangle
   =D_y^2\varphi(x,m(x))[u_1,u_2].
\]
The first derivative inserts \(w_i\).
The second inserts \(w_j\) in the phase derivative and
also differentiates the tangent, inserting
\((0,m_{ij})\). This gives both lines of
\eqref{eq:matching-hessian-variations}.
By \eqref{eq:matching-stationary-identities},
\(m_{ij}=-R_0B_{ij}\), so only phase derivatives
through order four enter.

Every inherited loop atom has two curvature slots,
both zero at the flat field. Hence \(D\mathcal I(0,0)=0\).
The chain rule along the graph consequently gives
\[
 D_{ij}\{\mathcal I(x,m(x))\}\big|_{x=0}
                         =D^2\mathcal I(0,0)[w_i,w_j].
\]
There is no term with \(D\mathcal I\) and
\(m_{ij}\). Likewise \(D\log j_G(0)=0\);
the Haar contribution is exactly the second line of
\eqref{eq:matching-complete-loop-read}, by
\eqref{eq:matching-haar-jet}.

For the determinant, expand
\(\det(1+tK)=1+t\operatorname{Tr}K+O(t^2)\)
to obtain
\[
 D_i\log\det H
                 =\operatorname{Tr}(H^{-1}D_iH).
\]
Differentiating \(H^{-1}H=1\) gives
\[
 D_jH^{-1}
            =-H^{-1}(D_jH)H^{-1}.
\]
One more derivative therefore yields
\[
 D_{ij}\log\det H
   =\operatorname{Tr}\{
       H^{-1}D_{ij}H
       -H^{-1}(D_jH)
                  H^{-1}(D_iH)\}.
\]
At zero, cyclicity of the trace gives the last line
of \eqref{eq:matching-complete-loop-read}, including
its factor \(1/2\) and its minus sign. This calculation
is on the finite physical coordinate space, so no
trace exchange has been assumed.
\end{proof}

The four contributions in
\eqref{eq:matching-complete-loop-read} are the inherited
loop, the Haar term, the Hessian second variation, and
the product of two Hessian first variations.
The inherited term differentiates the actual \(U(x,y)\),
including its replica and frame dependence.
The third includes both a fourth phase derivative and
the second stationary response. The fourth uses two
third phase derivatives. In particular the second
variation is not determined by the two-third-derivative
trace alone. These are the retained one-loop identities
of \cite[Lemma~36.7 and (36.20), pp.~244--246]{Main},
with the Haar jet evaluated explicitly.

\subsection{Gauge identities and locality of these Taylor data}

The one-loop combination, rather than any one of its
coordinate terms, is the object to match.
We first verify that the coordinate calculation has
retained its gauge meaning.

\begin{lemma}[Coordinate independence and the leading gauge identity]
\label{lem:matching-gauge-jets}
The one-loop action \(T\) is unchanged by a smooth change
of free fine coordinates, when the same Haar density
is used. In the common regular chart it is gauge
invariant as a function of the coarse connection.
In particular
\[
 DT(0)=0,\qquad
             D^2T(0)[d_{0,c}\omega,v]=0.
\]
More generally, let \(k\ge1\) be an integer and \(f\) an
invariant analytic functional whose field Taylor terms of degrees less
than \(k\) vanish. Its leading \(k\)-linear form
annihilates \(d_{0,c}\omega\) in every leg.
\end{lemma}
\begin{proof}
Let \(y=y(x,\widetilde y)\) be a nonsingular change of
fine coordinates and put \(J=D_{\widetilde y}y\) at the
stationary point. The fine gradient of the phase is
zero there. Thus terms with second derivatives of this
coordinate change disappear from the Hessian, giving
\[
                    \widetilde{H}=J^{\mathsf T}HJ.
\]
Write \(j=\prod_{e\in I}j_G(y_e)\), with the fixed
factor \(\zeta_G^{|I|}\) kept in
\eqref{eq:matching-prefactor-ledger}. Its transformed
coordinate density is
\(\widetilde j=j|\det J|\). Consequently
\[
 \frac12\log\det\widetilde{H}-\log\widetilde j
                   =\frac12\log\det H-\log j.
\]
The inherited loop value at the stationary configuration
is a scalar and is unchanged as well.

A coarse gauge change induces a change of the fine
configuration followed by the root-fixed transformation
returning it to the chosen tree slice.
Proposition~\ref{prop:tree-integration} gives unit
gauge volume for this operation. The phase and inherited
loop are invariant, and the same Haar measure is
retained. Uniqueness of the stationary graph in the
common chart identifies the transformed stationary point.
The preceding coordinate identity proves coarse gauge
invariance of \(T\). Its first derivative vanishes and
its Hessian annihilates linear gauge increments by the
first two flat gauge identities proved in
Lemma~\ref{lem:reset-flat-gauge-identities}.

For the last assertion, the infinitesimal gauge variation
of a coarse logarithm is
\(d_{0,c}\omega+O(x\omega)\), up to the fixed overall
gauge sign. Differentiate invariance in that direction.
Since the first possible term in the functional has
degree \(k\), the homogeneous degree-\((k-1)\) identity
is
\[
                   D^kf(0)[d_{0,c}\omega,x,\ldots,x]=0.
\]
All contributions from \(O(x\omega)\) have higher
degree. Polarising the remaining \(k-1\) directions
and using symmetry proves the assertion in every leg.
\end{proof}

There is no new spatial norm in these calculations.
The phase jets, stationary responses and \(H_0^{-1}\)
carry the rooted exponential bounds of
Sections~\ref{sec:rooted-stationary}--%
\ref{sec:local-stationary-prefactor}.
The contractions above preserve those bounds by
Lemma~\ref{lem:local-rooted-joining}.
For instance \(B_{ij}H_0^{-1}B_{kl}\) joins the two
rooted cubic vertices by one inverse edge. Fixing
one external coarse root, sum the first vertex,
then the inverse edge by its weighted row bound,
then the second vertex. A smaller remaining width
pays the connecting paths, so no free volume sum
is introduced.

For the loop traces the same point deserves an
explicit check. In
\[
 \operatorname{Tr}(R_0H_iR_0H_j)
       =\sum_{a,b,c,d}(R_0)_{ab}(H_i)_{bc}
                         (R_0)_{cd}(H_j)_{da},
\]
root at the external coarse direction in
\(H_i\). The vertex paths and inverse edges
form a connected closed chain reaching the other
external direction. Their total connecting length
dominates the distance between these two roots.
Successive absolute row and rooted vertex sums,
with a fixed smaller exponential weight, bound
the displayed kernel. Closing the trace does
not release the external anchor for a second
sum over coarse sites. There are at most
\(4L^4\dim_\R\g\) fine coordinate residues per
coarse block, which is allowed in a fixed-\(L\)
constant.

For example, take the external directions to be the
coarse link-colour basis vectors. The unweighted
row sum of this two-leg trace kernel is bounded by
\[
 \begin{split}
 &\sum_j\left|\operatorname{Tr}
                      (R_0H_iR_0H_j)\right|\\
 &\quad\le\sum_{b,c}|(H_i)_{bc}|
       \sum_a|(R_0)_{ab}|\sum_d|(R_0)_{cd}|
                         \sum_j|(H_j)_{da}|
 \lesssim1 .
 \end{split}
\]
The first vertex is summed with its external root fixed;
the last vertex is rooted at a fine input. The inverse
uses its column and row bounds, respectively. These are
exactly the rooted bounds already available; the implicit
constant may depend on the fixed profiles and \(L\).
Keeping the connecting-length weights proves the same
estimate at a smaller positive exponential width.

Thus the classical third and fourth kernels and
the one-loop quadratic kernel have finite
exponential coefficient bounds per external
anchor, uniform in cutoff, period and admissible
history. The inherited one-loop part is linear
in its budget \(\Lambda\); the other one-loop
terms have fixed-profile bounds. These are
fixed-ratio bounds. Establishing a path budget
chosen before \(L\) requires the additional
smooth-lift and vertex estimates in the next
matching step. All finite-torus closed traces
have been retained here, including those whose
lifts wind around the period.

The first matching problem is to express the cubic
difference between \(\Psi\) and the coarse row action
by the prescribed coarse curvature slots. At degree
three \eqref{eq:matching-classical-diagonal} only needs
the stationary tangent. This allows a direct estimate
of the cubic kernel on the smooth free lift before
choosing \(L\).

\section{Matching the cubic stationary term by curvatures}
\label{sec:cubic-curvature-matching}

The coarse action already has its quadratic part. Write
\(S_{\rm row}\) for the unmasked coarse Wilson and row-square
terms of \eqref{eq:unmasked-principal-action}, using the
calibrated row factor \(S_c\) of \eqref{eq:calibrated-row-metric}
and the prescribed symmetry average. Its flat value and first
derivative vanish. Since \(\widetilde H_cd_c=H_cd_c\),
\[
 T_2\bigl(S_{\rm row}(V(x))\bigr)
 =\tfrac12\|S_cd_cx\|_2^2
 =T_2\Psi(x)
\]
by \eqref{eq:matching-classical-quadratic}. The first term
to match is therefore \(T_3(\Psi-S_{\rm row}\circ V)\).

Why should this term be representable by the paths of the
principal action? A slot \(Z_{\gamma,\alpha}\) in
\eqref{eq:principal-slot} vanishes at the identity. Its
derivative in a coarse link direction \(v\) is
\((d_cv)_\alpha\) at the terminal site of \(\gamma\):
the derivative of the transport multiplies the zero
curvature, and the four plaquette edges give their signed
sum. Thus the cubic part of a product of three such slots
has one linear curl in each of its three arguments. We will
look for exactly this representation of the cubic
difference. Its coefficients must have a positive
exponential width and a bound chosen before \(L\), so that
the new term can enter the same principal-action class.

\subsection{Using the smooth lift to compute the cubic jet}

For this calculation it suffices to find a fine trial field
with the stationary tangent. Indeed, on any one finite
torus, if \(y(0)=m(0)=0\) and \(Dy(0)=Dm(0)\), then
\[
 \begin{split}
 \varphi(x,y(x))-\Psi(x)
   &=\int_0^1(1-t)D_y^2\varphi\bigl(x,m(x)+t(y(x)-m(x))\bigr)\\
   &\hspace{26mm}\bigl[y(x)-m(x),y(x)-m(x)\bigr]\,\dd t.
 \end{split}
\]
Stationarity removed the first-order term in \(y-m\).
Since \(y(x)-m(x)=O(\|x\|_\infty^2)\), these two functions
have the same Taylor terms through degree three. This lets
us use the smooth free lift \(E\) from Appendix~C instead
of estimating the cubic derivative in the tree coordinates.

Start with \(U^{\rm sm}_e(x)=\exp((Ex)_e)\). The actual
rooted predictor has differential \(Q_0=Q+d_{0,c}P\)
as computed before \eqref{eq:calibrated-frame-correction},
where \(Q\) on the right is the straight linear average.
To cancel the extra gradient, put
\(r_o(x)=\exp((PEx)_o)\) at the coarse root \(Lo\).
Use a fine gauge with these root values and then return to
the tree slice by a root-fixed gauge. At first order its
full fine field is \(Ex+d_{0,f}\omega(Ex)=\widehat E x\)
from the proof of Theorem~\ref{thm:calibrated-row-data}.
Indeed its gauge increment has root values \(-PEx\) in
the convention \(u\mapsto u+d_{0,f}\omega\), and its
tree-link increments are zero. These conditions determine
the tree correction. The resulting fine chart \(y(x)\)
therefore has \(Dy(0)=\mathcal H=Dm(0)\).
By gauge invariance of the fine action and
\eqref{eq:actual-branch-covariance}, the value of this trial
phase is exactly
\begin{equation}
 S_f^{\rm un}(U^{\rm sm}(x))
   +a\sum_{b=(o,o')\in E_c^+}
           q_G\bigl(V(x)_b
                [r_o(x)Q(U^{\rm sm}(x))_b r_{o'}(x)^{-1}]^{-1}\bigr).
 \label{eq:cubic-smooth-trial}
\end{equation}
The root rotation in this expression changes the predictor
used in the trial evaluation. The coarse input is still
\(V(x)\).

\smallskip\noindent
\emph{The fine-action contribution.}
To bound a cubic coefficient, insert coarse plane waves
\(e^{ip_i\cdot o}v_i\), with three momenta whose sum is zero
modulo \(2\pi\). Here \(v_i\) is a unit vector of coarse link
types and colours. Denote the three lifted inputs by \(u_i\).
We use the real Taylor tensors extended complex multilinearly.
Choose a positive momentum width \(\eta\) small enough for
\eqref{eq:calibration-column-differences} and the fixed
coarse Gaussian row bounds. Thus, for independent momenta
on a strip of width \(\eta\), all three momenta lie within
the common bounds, including their sum. At a fine residue
over the base coarse block the amplitudes of \(u_i\) and of
their first and second fine differences are bounded by
\(C_jL^{-1},C_jL^{-2},C_jL^{-3}\), respectively.
At a point at fine path distance \(n\) from the block,
multiply these bounds by
\(\exp\{C\eta(1+n/L)\}\): its coarse base point has distance
at most \(C(1+n/L)\). This explains how a path in fine units
uses the strip measured in coarse units.

Consider a slot based at a fine anchor in this block, with
terminal displacement \(z\) and path length \(n=|\gamma|\).
Before multiplying its derivatives with other slots, compare
its first and second derivatives to those of the slot with
the same orientation at the anchor and empty path. The
first derivative is the curl at the terminal site. The
difference of the two curls is a sum along the path of
first differences of the curl. These are second
differences of \(u_i\), so the difference is bounded by
\(C_j n L^{-3}e^{C\eta(1+n/L)}\).
For the second derivative, let \(A_a,B_a\), \(1\le a\le4\),
be the two signed input values along the four ordered
edges of \(p\). Differentiating the forward and inverse
words in \(\mathcal C(U_p)\) cancels the terms with two
derivatives on the same edge. The mixed derivative is
\[
 \tfrac12\operatorname{pr}_{\g}
     \sum_{a<b}\bigl([A_a,B_b]+[B_a,A_b]\bigr).
\]
Moving the base point back to the anchor, replace one
input in each bracket at a time. Each change
is bounded by \(C_j n L^{-2}e^{C\eta(1+n/L)}\) by
telescoping, and the other input gives \(C_jL^{-1}\)
with the corresponding phase bound. Differentiating the
transport gives the additional two terms
\[
 \left[\sum_{e\in\gamma} u_i(e),(d_fu_k)_\alpha(o+z)\right],
 \qquad
 \left[\sum_{e\in\gamma} u_k(e),(d_fu_i)_\alpha(o+z)\right],
\]
where path sums count orientation. Each has price bounded by
a product of \(C_jnL^{-1}\) and \(C_jL^{-2}\) with the same
type of phase bound. In particular the second-derivative
difference also costs \(C_j(n+1)L^{-3}e^{C\eta(1+n/L)}\),
with an enlarged fixed \(C\). Both derivatives of an
empty-path slot have price \(C_jL^{-2}\).

Now freeze each slot of a row at its anchor, and call the
resulting row \(\ell_b^0\). Write
\(r^0_{b\alpha}=\sum_{k:\alpha_{bk}=\alpha}r_{bk}\)
for its coefficients. At zero momentum the filled metric
of \eqref{eq:calibration-filled-metric} is the identity,
because \(G_{v,i}(0)=b_v(0)=1\). The extension in
\eqref{eq:calibration-positive-extension} has the same
value there, since \(d_2(0)=0\). Consequently
\[
                  \lambda_f\,1+(r^0)^{\mathsf T}r^0=1
\]
for the fine rows of a nonempty history. At the empty
history the action is the literal Wilson term. With the
direct plaquette price included, the frozen rows thus
give the Wilson cubic derivative: the difference between
\(q_G(U_p)\) and \(|\mathcal C(U_p)|^2/2\) starts at
degree four, as in \eqref{eq:calibrated-square-decomposition}.
We can do this calculation before the signed symmetry
average. On each symmetry copy the same estimates apply
to the transformed inputs.

The first and second row-derivative errors after freezing
are bounded by \(A_3L^{-3}\), for a finite bound \(A_3\)
depending on the incoming path data. To see the uniformity,
first require \(C\eta/L\le\sigma/2\), where \(\sigma\)
is the incoming width. Then the phase factor and
\(n+1\) in the preceding estimates are paid by
\eqref{eq:stationary-row-budget}. At the flat field a
row-square derivative is
\[
 D^3\bigl(|\ell_b|^2/2\bigr)[u_1,u_2,u_3]
     =\sum_{i=1}^3
       \big\langle D\ell_b[u_i],D^2\ell_b[u_j,u_k]\big\rangle,
 \qquad \{j,k\}=\{1,2,3\}\setminus\{i\}.
\]
The empty-path row terms on the right have size \(C_jL^{-4}\).
In the difference, one frozen factor of size \(C_jL^{-2}\)
and one error give \(C_jA_3L^{-5}\). Two errors give
\(A_3^2L^{-6}\). The incoming cubic atoms at a fine anchor
are products of three first slot derivatives, costing
\(C_jL^{-6}\) times their incoming budget. Quartic atoms
have vanishing cubic derivative. Summing the \(L^4\)
fine anchors in one coarse block gives a fine contribution
bounded by \(C_j+L^{-1}A_3\), after increasing \(A_3\)
by a fixed function of the already fixed incoming data.

\smallskip\noindent
\emph{The actual noise contribution.}
There are only four outgoing noise terms per coarse site,
so we need a bound on the first two derivatives of their
mismatches. For a route in the frame patch or a connector,
the length is at most \(C_jL\). Placing one derivative on
a factor gives at most \(C_jL\cdot L^{-1}\), and placing
two gives at most \(C_jL^2\cdot L^{-2}\), including two
derivatives at the same factor. The phases are bounded
on the patch by a fixed constant, because it meets a
fixed number of coarse blocks. Thus the route derivatives
through order two are bounded before \(L\).

We check that the frame derivatives have the same
property. Write, within a patch,
\[
 T_{zz'}=\frac1{4!}\sum_\pi U_{\gamma_{z,z'}^\pi}^{-1},
 \quad
 \bar h=\frac1s\sum_{z'}h_{z'},\quad
 A_z=\frac1s\sum_{z'}h_{z'}(T_{zz'}-I_N),
 \quad M_z=\bar h+A_z .
\]
Route reversal and coordinate minimisation of
\eqref{eq:actual-frame-score} give \(h_z=\pi_G(M_z)\)
near the flat field, by the same coordinate argument in
Lemma~\ref{lem:regular-frame-unique}. The equation also
holds at \(z=0\), since a rooted minimiser is an unrooted
minimiser after allowing common rotations, as in that proof.
Use \(h_0=I_N\). With \(\Pi=D\pi_G(I_N)\), subtract the
equations at \(z\) and \(0\) and differentiate:
\begin{align*}
 D_i h_z &=\Pi(D_i A_z-D_i A_0),\\
 D_{ik}h_z
   &=\Pi(D_{ik}A_z-D_{ik}A_0)
       +D^2\pi_G(I_N)[D_iM_z,D_kM_z]
       -D^2\pi_G(I_N)[D_iM_0,D_kM_0].
\end{align*}
All derivatives are at the flat field in the indicated
lifted directions. Since \(T_{zz'}=I_N\) there, the
derivative \(D_i A_z\) is the average of \(D_i T_{zz'}\).
The second derivative is the average of
\[
 D_{ik}T_{zz'}+(D_i h_{z'})(D_k T_{zz'})
                       +(D_k h_{z'})(D_i T_{zz'}).
\]
The first equation bounds \(D_i h_z\) and hence
\(D_iM_z\). The second uses only these first derivatives
and the two route derivatives. All the averages have
mass one. In particular, the bounds have no growing
factor \(s\).

The connector formula \eqref{eq:actual-connectors} bounds
its first two derivatives by the product rule, including
the inverse frame at the second endpoint. If \(A=\sum_z w(z)Y_{b,z}\)
as in the regular predictor formula, then
\[
 D_iQ_b=\Pi D_iA,\qquad
 D_{ik}Q_b=\Pi D_{ik}A+D^2\pi_G(I_N)[D_iA,D_kA]
\]
at the identity. Each derivative of \(A\) is a probability
average of the connector derivatives, so these are bounds
for the actual predictor \(Q\). There is also a bound for \(PE\)
from its two path averages in Section~\ref{sec:calibrated-rows}:
each path has length \(O_j(L)\) and \(Ex\) gives \(L^{-1}\)
per input. This bounds the derivatives of the two root
rotations in \eqref{eq:cubic-smooth-trial} through order two.
Writing the resulting mismatch as \(e^{\xi_b(x)}\) near
the flat field, \(\xi_b(0)=0\) and both its derivatives
are bounded before \(L\). The cubic derivative of the
noise price is the sum of the three pairings of
\(D_i\xi_b\) with \(D_{jk}\xi_b\) times \(a\), because
\(q_G(e^\xi)=|\xi|^2/2+O(|\xi|^4)\). Its bound per coarse
site is \(C_j\). This calculation uses only flat
derivatives of the actual frames and predictor, and then
their complex multilinear extensions.

\smallskip\noindent
\emph{The resulting cubic symbol.}
Let \(K(p_1,p_2)\) be the symmetric trilinear Fourier
kernel of \(T_3(\Psi-S_{\rm row}\circ V)\), with
\(p_0=-p_1-p_2\). Concretely, take the mixed cubic
derivative of the energy assigned to one coarse site in
\eqref{eq:cubic-smooth-trial} minus the coarse row energy
at that site. Evaluate it on the three waves with these
momenta and divide by \(3!\). Summing the translates
picks precisely the triples of total momentum zero.
This gives the kernel on the integer cover. Evaluating
the same paths and convolutions periodically gives its
periodisation, with all repeated link occurrences retained.
The expansions used above are absolutely summable by
the input width and the bounds of Appendix~C. They give
periodic holomorphic dependence on the independent
momenta. The coarse row contribution is bounded on this
strip by its fixed Gaussian path budget. In all,
\begin{equation}
             \sup_{\rm strip}|K(p_1,p_2)|
                      \le C_j+A_3/L .
 \label{eq:cubic-kernel-before-ratio}
\end{equation}
The strip and \(C_j\) use the fixed Gaussian and geometric
profiles. After fixing the incoming path data, take \(L\)
above the two requirements \(C\eta/L\le\sigma/2\) and
\(A_3/L\le1\). The bound to be used in dividing \(K\) is
then fixed before this choice of \(L\).

\subsection{Dividing by one curl in every argument}

The stationary value is invariant under small coarse gauge
changes: change the fine field with the prescribed root
values, return it to the tree slice, and use uniqueness of
the stationary point on the chart. The coarse row action
is gauge invariant as well. Their leading difference
therefore annihilates \(d_{0,c}\omega\) in every argument
by Lemma~\ref{lem:matching-gauge-jets}. For a Fourier
mode the gradient is multiplication by \(q(p)\) from
Appendix~C. Write \(\iota_i\) for contraction of the
coefficient in argument \(i\) with \(q(p_i)\), with the
exterior convention of Section~\ref{sec:packet-row-comparison}.
In particular the contraction of \(e^a\wedge e^b\) is
\(q_a e^b-q_b e^a\), dual to the \(ab\) component of
\(q(p_i)\wedge v_i\) on the input. The
Ward identities give
\[
                      \iota_0K=\iota_1K=\iota_2K=0 .
\]
Indeed, on any compatible torus substitute three Fourier
modes of total momentum zero and a gradient in the
indicated argument. The sum over translates is the number
of coarse sites, which cancels. The identities therefore
hold on the discrete momentum grids, and then everywhere
by density of the grids and holomorphic continuation.
Contraction in one argument commutes with contraction in
a different argument, since these are separate tensor legs.

For one independent momentum we already have the needed
primitive, namely a coefficient whose contraction is the
given transverse coefficient. Use the Laurent division of
\eqref{eq:packet-laurent-width-loss} and the exterior
splitting in Lemma~\ref{lem:packet-complement}. Denote
its four-coordinate operation in leg \(i\) by \(\mathcal J_i\).
It successively takes \(e^a\wedge\mathcal D_a b\) on the
remaining tensor \(e^a\wedge u+b\). Each step leaves the
coordinate evaluation of \(b\). After the four steps, the
evaluation term vanishes in positive exterior degree.
For a scalar in that leg it is the value at \(p_i=0\).
Thus the scalar operation is also a primitive when that
value vanishes. All other contractions whose momenta
are independent of \(p_i\) pass through this operation.

This suggests starting with leg \(0\) and then dividing the
two independent legs. The first operation needs care:
using \((p_0,p_2)\) as variables makes \(p_1\) dependent.
Put
\[
                     X=\mathcal J_0 K,\qquad
                     Y=\mathcal J_0(\iota_1 X).
\]
Their exterior degrees in the three arguments are
\((2,1,1)\) and \((3,0,1)\). We have
\(\iota_0X=K\), \(\iota_2X=0\), and
\(\iota_0Y=\iota_1X\), \(\iota_2Y=0\):
the second primitive exists because
\(\iota_0\iota_1X=\iota_1K=0\).
To remove the defect \(\iota_1X\) while keeping
\(\iota_0X=K\), we want to write \(Y\) as a contraction
in leg \(1\), up to a \(\iota_0\)-exact term. The scalar
division shows exactly the obstruction to doing this:
the restriction of \(Y\) to \(p_1=0\).

On that restriction \(\iota_0Y=0\) and \(p_0=-p_2\).
Take the leg-\(0\) primitive \(Z_0\) of this restriction
in the single momentum \(p_0\). It has degrees \((4,0,1)\)
and satisfies \(\iota_2Z_0=0\), since
\(\iota_0\iota_2Z_0=\iota_2(Y|_{p_1=0})=0\).
To check the implication, contraction of a top exterior
form by \(q(p_0)\) is injective for \(q(p_0)\ne0\):
its four components are the top coefficient times the
four components of \(q(p_0)\), with signs. Continuity
handles the zero. Extend \(Z_0\) as \(Z(p_1,p_2)=Z_0(p_2)\).
Then \(Y-\iota_0 Z\) vanishes at \(p_1=0\) and
is still killed by \(\iota_2\). Returning to independent
variables \((p_1,p_2)\), set
\[
 D=\mathcal J_1(Y-\iota_0Z),\qquad
 X'=X-\iota_0D .
\]
Scalar division gives \(\iota_1D=Y-\iota_0Z\) and
\(\iota_2D=0\). Hence
\[
 \iota_0X'=K,\qquad
 \iota_1X'=\iota_1X-\iota_0(Y-\iota_0Z)=0,\qquad
 \iota_2X'=0 .
\]
Now divide \(X'\) in the two independent legs:
\(\iota_1\mathcal J_1X'=X'\), and
\(\iota_2\mathcal J_2\mathcal J_1 X'=\mathcal J_1X'\).
We have obtained a two-form coefficient in each leg,
\begin{equation}
          K^{\rm curv}=\mathcal J_2\mathcal J_1 X',
          \qquad K=\iota_0\iota_1\iota_2 K^{\rm curv}.
 \label{eq:cubic-three-curls}
\end{equation}
The order of reconstruction in this formula matters:
the last two primitives were made in \((p_1,p_2)\),
where \(p_0\) is again dependent.

Here are the locality estimates for this finite
calculation. Use the coefficient norm
\[
 |F|_\tau=\sum_{z_1,z_2\in\Z^4}
          \|F_{z_1,z_2}\|e^{\tau(|z_1|_1+|z_2|_1)}
 \quad\text{for }F=\sum_{z_1,z_2}F_{z_1,z_2}
                             e^{i(p_1\cdot z_1+p_2\cdot z_2)}.
\]
Take the sum of absolute tensor components, including
the fixed colour types, in \(\|\cdot\|\) here. The
analogue with one displacement is used on \(p_1=0\).
Contour shifting at half the width of
\eqref{eq:cubic-kernel-before-ratio} bounds the
coefficients by their strip bound times
\(e^{-\eta(|z_1|_1+|z_2|_1)/2}\). Summation at
\(\tau_0=\eta/4\) is finite. By
\eqref{eq:packet-laurent-width-loss}, one \(\mathcal J_i\)
costs at most \(4/[e(\tau-\tau')]\) between widths
\(\tau>\tau'>0\): each of its four terms acts after
coordinate evaluations of norm at most one. A change
from \((p_1,p_2)\) to \((p_0,p_2)\) replaces
\((z_1,z_2)\) by \((-z_1,z_2-z_1)\), so half the width
pays this change or its inverse. Restriction, extension
and the change of sign in the single momentum have
norm at most one in the corresponding coordinates.
Multiplication by a \(q\) in the displayed formulas
costs a fixed factor bounded in terms of \(\tau_0\)
and the exterior types. There are six applications of
\(\mathcal J_i\) and at most four changes of paired
coordinates along any term in the construction.
Halving the width at each such operation, for example,
leaves the fixed width \(\tau_* = 2^{-12}\tau_0\)
with each denominator in a division paid by a width
loss at least \(\tau_*\). It follows that
\(|K^{\rm curv}|_{\tau_*}\lesssim |K|_{\tau_0}\),
with the constant fixed by the strip and tensor types.
The division includes zero momentum and can be
periodised on every compatible torus.

This is the cubic case of the correction in
\cite[Lemma~36.3, pp.~230--232]{Main}. The coefficient
bounds here reuse the Laurent estimate already proved
for the positive complement. The smooth trial
calculation gives the cubic part of
\cite[Proposition~36.6 and (36.12), pp.~240--243]{Main}
using only all-smooth inputs and two frame derivatives.
These locators were checked on 20 November 2026.

\subsection{Installing the full curvature correction}

For a coefficient of \(K^{\rm curv}\) with offsets
\((z_1,z_2)\), use three slots of
\eqref{eq:principal-slot} anchored at a common coarse
site. Use the empty path for argument \(0\), and fixed
coordinate paths of lengths \(|z_1|_1,|z_2|_1\) for
the other two. Replace the two-form inputs by the
corresponding full transported curvatures. Take real
parts of the coefficients in a real colour basis and
average the coefficient contractions over the simultaneous
adjoint action of \(G\). The linear curls commute with
these operations, and the original \(K\) is real in
position space and invariant under constant gauges:
conjugating by a constant element of \(G\) preserves
the tree chart and carries the stationary solution to
the one for the conjugated input.
Thus these projections preserve
\eqref{eq:cubic-three-curls}. The adjoint representation
is isometric, so their cost is at most a fixed
finite-dimensional factor in the component norm.
Sum over anchors and take the prescribed lattice symmetry
average. The cubic difference has these symmetries: the
regular frame score is invariant under relabelling its
paths by signed permutations, and the connector weights
are symmetric. Reversing a connector exchanges the two
endpoint frames and takes inverse, which the regular
projection preserves by uniqueness. The stationary
value then has these symmetries by returning the
transformed fine field to the tree slice. Thus the
averaged correction has the same cubic term.

\begin{proposition}[Cubic matching of the classical action]
\label{prop:cubic-curvature-matching}
There is a real invariant cubic slot correction \(P_3(V)\)
of the form allowed in \eqref{eq:stationary-higher-budget}
such that
\begin{equation}
 T_k\bigl(\Psi-S_{\rm row}\circ V-P_3\circ V\bigr)=0,
                            \qquad 0\le k\le3 .
 \label{eq:cubic-matching}
\end{equation}
The correction is periodised from the integer cover
and is independent of the actual coupling. Its output
path width and budget are fixed before \(L\), with the
lower bound on \(L\) chosen from the fixed incoming path
data. These bounds are uniform in compatible periods and
legal finite histories at the fixed adjacent profiles.
\end{proposition}
\begin{proof}
Use the coefficients and paths just constructed.
Their total word length is \(12+|z_1|_1+|z_2|_1\),
counting each occurrence before periodisation. Hence at
any fixed output width up to \(\tau_*\) their budget
is bounded by \(e^{12\tau_*}|K^{\rm curv}|_{\tau_*}\)
times a fixed type and symmetry factor. This is the
claimed bound by \eqref{eq:cubic-kernel-before-ratio}
and the width calculation. Each slot vanishes at the
identity. In degree three one must therefore take
the linear term of all three slots. Their contraction
is exactly \(K\) by \eqref{eq:cubic-three-curls}.
Quadratic agreement already holds for the row action.
This gives \eqref{eq:cubic-matching}. The phases and
flat derivatives throughout this construction came
from \(\varphi\), with its common factor \(\beta\)
removed. Its branch and path data and the Gaussian
history determine the coefficients before a period
or actual coupling is selected.
\end{proof}

\section{Matching the quartic stationary term by curvatures}
\label{sec:quartic-curvature-matching}

After the cubic correction the functional
\(\Psi-S_{\rm row}\circ V-P_3\circ V\) is invariant and its
first possible term has degree four. By
Lemma~\ref{lem:matching-gauge-jets} this term has a Ward identity
in every argument. To use these identities with a path budget
fixed before \(L\), we first need a uniform bound for its
Fourier kernel. In the quartic formula
\eqref{eq:matching-classical-jets} we must bound the
fourth derivative on smooth lifts and the cubic functional
with a fine variation to be paired through the flat Hessian.

\subsection{The stationary correction in centred frame coordinates}

A rooted frame can have a large derivative at its root when the
fine variation is controlled only in a summed norm. The connectors
use an average, so it is useful to keep the common frame rotation
as a separate variable. At a coarse site \(o\), let \(h_z^o\) be
the actual regular frame and define its weighted centre and
centred frame by
\[
 c_o(U)=\pi_G\Bigl(\sum_z w(z)h_z^o(U)\Bigr),\qquad
 h_z^{c,o}(U)=c_o(U)^{-1}h_z^o(U).
\]
The weights are those of \eqref{eq:convolution-weights}.
Replacing both endpoint frames in
\eqref{eq:actual-connectors} by the centred frames gives the
predictor \(\overline Q(U)_b=c_o(U)^{-1}Q(U)_b c_{o'}(U)\)
for \(b=(o,o')\), by equivariance of \(\pi_G\).
With \(U_e=e^{A_e}\) for a fine Lie-algebra field \(A\), and with
\(\vartheta_o\) a coarse site Lie-algebra variable, put
\begin{equation}
 \begin{split}
 \Phi(x,A,\vartheta)
  ={}&S_f^{\rm un}(U)+a\sum_{b=(o,o')}
     q_G\bigl(V(x)_b
          [e^{\vartheta_o}\overline Q(U)_b e^{-\vartheta_{o'}}]^{-1}\bigr).
 \end{split}
 \label{eq:quartic-centred-phase}
\end{equation}
Apply a fine gauge with root values \(e^{\vartheta_o}c_o(U)^{-1}\)
and then the root-fixed tree gauge. The resulting field has a
unique chart coordinate \(y\) at fixed \(x\). The fine action
is invariant, and the predictor after these transformations is
exactly \(e^{\vartheta_o}\overline Q(U)_b e^{-\vartheta_{o'}}\).
Thus \(\Phi\) equals \(\varphi(x,y)\) under this analytic map.
For a field already in the tree chart, taking
\(e^{\vartheta_o}=c_o(U)\) makes the root rotation the
identity. This gives a local right inverse of the map.

At flatness \(Dc_o[u]=-(Pu)_o\), since \(Dh_z^o[u]=t_z^o-t_0^o\)
in the calculation before \eqref{eq:calibrated-frame-correction}.
Consequently \(D\overline Q[u]=Q_0u-d_{0,c}Pu=Qu\), with \(Q\)
the straight linear average. For a fine site direction \(\omega\),
the infinitesimal gauge in the joint variables is
\[
            (u,\zeta)=(d_{0,f}\omega,A_0\omega).
\]
Indeed, under a fine gauge \(\Omega\), choose the new coarse
frame by requiring
\(e^{\vartheta'_o}c_o(U^\Omega)^{-1}\Omega(Lo)
   =e^{\vartheta_o}c_o(U)^{-1}\).
This gives exact invariance of \(\Phi\) with \(x\) fixed.
Differentiation for \(\Omega=e^{-t\omega}\) uses
\(Pd_{0,f}\omega=\omega|_{\rm roots}-A_0\omega\)
and gives the stated tangent. These gauge directions are precisely
the kernel of the differential to the tree chart. They are
contained in the kernel by gauge covariance, and their map
from site fields is injective because a zero gradient is constant
and \(A_0\) preserves constants. The difference of the joint
and tree-chart dimensions counts one tree link per nonroot
site and one frame variable per root, each with \(\dim\g\)
components. This is also the dimension of the gauge directions.

The quadratic energy of a vertical variation \((u,\zeta)\),
meaning \(x\) is fixed, is its Hessian quadratic form
\begin{equation}
     \langle d_fu,H_fd_fu\rangle+a\|Qu-d_{0,c}\zeta\|_2^2.
 \label{eq:quartic-block-energy}
\end{equation}
It follows from the calibrated fine row Gram form and the
first mismatch derivative \(v-Qu+d_{0,c}\zeta\) for a full
variation \((v,u,\zeta)\). In particular
\(w_i^{c}=(v_i,Ev_i,0)\) is a stationary tangent. Its
Hessian pairing with \((0,u,\zeta)\) is
\[
 \langle K_f Ev_i,u\rangle
 -a\langle v_i-QEv_i,Qu-d_{0,c}\zeta\rangle=0
\]
by \eqref{eq:calibrated-installed-columns} and
\(d_{0,c}^{\mathsf T}K_c=0\).

At fixed \(L\) and period, take any real linear complement of the
vertical gauge directions. The affine slice through \((A,\vartheta)
=(Ex,0)\) parallel to that complement maps locally invertibly
to the tree chart. Its stationary solution has a correction to
\((Ex,0)\) starting in degree two. Apply
Lemma~\ref{lem:matching-classical-jets} on this slice. Define
\[
 B^c_{ij}[u,\zeta]=D^3\Phi(0)[w_i^c,w_j^c,(0,u,\zeta)].
\]
Write \(H_{\rm bl}^{-1} B^c_{ij}\) for the solution of the Hessian
equation on the quotient by vertical gauge directions. The formula is
\begin{equation}
 D^4\Psi(0)[v_1,v_2,v_3,v_4]
   =D^4\Phi(0)[w_1^c,w_2^c,w_3^c,w_4^c]
                    -\sum_{(ij|kl)} B^c_{ij}H_{\rm bl}^{-1}B^c_{kl}.
 \label{eq:quartic-centred-jet}
\end{equation}
To justify the quotient in the pairing, let \(Z\) be the vertical
vector field of the exact gauge transformation above, with a fixed
fine generator. Differentiate \(D\Phi[Z]=0\) twice along \(w_i^c,w_j^c\)
at zero. The terms with \(DZ\) vanish by the stationary tangent
identity, since \(DZ\) is vertical. The term with \(D^2Z\) vanishes
since \(D\Phi(0)=0\). This leaves \(B^c_{ij}[Z(0)]=0\).
Both types of derivative in \eqref{eq:quartic-centred-jet}
are taken in these centred coordinates.

\subsection{A uniform inverse in block energy}

Use one coarse Fourier mode \(p\), with alias momenta
\(k_\ell=(p+2\pi\ell)/L\) indexed by the centred \(L^4\) residue
set of Appendix~\ref{app:calibration}. We use the unitary Fourier
coordinates \(u_\ell\) on one fine block, so the amplitude of a
single fine wave is \(L^{-2}u_\ell\). For real \(p\), use the norm
\[
 \omega_\ell^2=L^{-2}+|q(k_\ell)|^2,\qquad
 \|(u,\zeta)\|_{\rm bl}^2=\sum_\ell\omega_\ell^2|u_\ell|^2+|\zeta|^2.
\]
The quotient norm is the infimum over representatives modulo
the displayed infinitesimal gauges. We work on a section of this
quotient that satisfies \(q(k_\ell)^*u_\ell=0\) for \(\ell\ne0\)
and \(\zeta=0\). Subtract a gradient at each \(\ell\ne0\) to set
the first conditions, using \(|q(k_\ell)|^2>0\). Then use the
\(\ell=0\) gauge to remove \(\zeta\): the principal column of
\(A_0\) is \(L^{-2}a_L(k_0)^J\), with \(a_L(k_0)^J\) bounded
away from zero by the sine quotient in Appendix~C. This also
fixes the principal gauge at \(p=0\).

Let \(B_\ell=L^{-1}h_1(k_\ell)\) be the alias column of \(Q\)
in these coordinates. The factor is the length \(L\) of the
straight average times \(L^{-2}\). The filter in \(h_1\) and
each principal line average are bounded away from zero on
the real principal alias. For the nonprincipal aliases,
\(|q(k_\ell)|\gtrsim L^{-1}(1+|\ell|_\infty)\). The sine
quotient and the bounds in Appendix~C give
\[
 |B_0u_0|\gtrsim_j L^{-1}|u_0|,\qquad
 \sum_\ell\frac{\|B_\ell\|^2}{\omega_\ell^2}
     \lesssim_j 1+\sum_{\ell\ne0}
          \frac{w_j(\ell)^2}{(1+|\ell|_\infty)^2}\lesssim_j 1.
\]
Here \(w_j(\ell)=\prod_a(1+|\ell_a|)^{-j}\) is the alias
decay weight of \eqref{eq:calibration-history-filter-bound}.
On the selected section the curvature energy controls
\(\sum_{\ell\ne0}\omega_\ell^2|u_\ell|^2\), because
\(\|q(k_\ell)\wedge u_\ell\|^2=|q(k_\ell)|^2|u_\ell|^2\)
on the transverse aliases and \(H_f\simeq 1\).
As \(Qu=\sum_\ell B_\ell u_\ell\), we also have
\[
 \omega_0^2|u_0|^2
 \lesssim_j |Qu|^2+
       \left(\sum_{\ell\ne0}\frac{\|B_\ell\|^2}{\omega_\ell^2}\right)
       \sum_{\ell\ne0}\omega_\ell^2|u_\ell|^2 .
\]
Here \(\omega_0\lesssim L^{-1}\). This proves a coercive bound
\(c_j\|u\|_{\rm bl}^2\) on the section. Conversely, on any
representative, the curvature energy is at most
\(C_j\sum_\ell\omega_\ell^2|u_\ell|^2\), and Cauchy--Schwarz
with the weighted column bound controls the noise term by
\(C_j\|(u,\zeta)\|_{\rm bl}^2\). Taking the infimum gives the
upper bound in the quotient norm and shows that the selected
section has norm comparable to that quotient norm.

We also need this inverse for complex coarse momenta.
Freeze the section and weights at a real \(p\) and use the
conjugate section at \(-p\) for the test variable of the
complex bilinear form. If the coarse momentum changes by
\(\delta p\), then \(q(k_\ell)\) changes by at most
\(C|\delta p|/L\). After division by \(\omega_\ell\) this
costs \(C|\delta p|\). The \(H_f\) symbols have a uniformly
bounded holomorphic extension by \eqref{eq:calibration-filled-metric}.
The weighted column bound for \(Q\) holds on the same type
of complex cube, using the complex sine quotient and the
frozen real weights. Cauchy's estimate bounds the change of
\(Q\) as a map from block energy to the coarse link norm by
\(C_j|\delta p|\). Applying these bounds to the two factors
in each quadratic term shows that the Hessian restricted to
the frozen sections changes by at most \(C_j|\delta p|\) in
the operator norm from block energy to its dual.

These sections remain complements. For each nonprincipal alias
the pairing of its real \(q(k_\ell)^*\) with the new gradient
changes by at most \(C|\delta p|\,|q(k_\ell)|/L\), a fixed
small fraction of \(|q(k_\ell)|^2\) for \(|\delta p|\ll1\).
The principal column of \(A_0\) stays nonzero by the same filter
bound and Cauchy's estimate. Choose the momentum width so that
\(C_j|\delta p|<c_j/2\) in the Hessian bound as well. The
Neumann series gives an inverse from the dual block norm to
the block norm, bounded by a constant fixed before \(L\).
On overlaps the pairings of gauge-annihilating functionals agree
by uniqueness of the solution on the quotient. Gauge identities
extend to complex momenta by analytic continuation of the
Fourier kernels. The pairings thus give a holomorphic symbol
on the common strip, with periodicity from alias relabelling.


\subsection{The mixed cubic and smooth quartic bounds}

We evaluate the phase derivatives per coarse block. For real
momenta on a compatible torus this means inserting waves of
total momentum zero and dividing by the number of coarse blocks.
The corresponding formulas on the integer cover continue to
complex momenta. For smooth inputs \(Ev_i\) with unit coarse
amplitudes, the fine field and its first two differences have
supremum bounds \(C_jL^{-1},C_jL^{-2},C_jL^{-3}\) on a block by
\eqref{eq:calibration-column-differences}. For a vertical
input in the selected section with unit block energy, Parseval
and Cauchy--Schwarz give, on one fine block \(B\),
\begin{equation}
 \|u\|_{2,B}\lesssim L,\qquad
 \|\nabla_f u\|_{2,B}\lesssim1,\qquad
 \|u\|_{1,B}\lesssim L^3,\qquad
 \|\nabla_f u\|_{1,B}\lesssim L^2 .
 \label{eq:quartic-block-counting}
\end{equation}
The norms on \(B\) here use counting measure, so the last two
bounds multiply the first two by \((L^4)^{1/2}\).
These estimates hold with the frozen real weights on the
complex strip, since the continued Bloch factors on one block
are bounded and \(|q(k_\ell)|\lesssim\omega_\ell\).
On an enlarged patch of diameter \(O_j(L)\) they cost a fixed
factor. At displacement of fine length \(n\) they cost at most
\(e^{C\eta(1+n/L)}\), where \(\eta\) bounds the imaginary
parts of the coarse momenta.

\smallskip\noindent
\emph{Fine-action derivatives.}
For four smooth inputs, the fourth derivative of a row square
at zero is the sum of three products of second derivatives
and four products of a first and a third derivative.
For a slot of path length \(n\), the third derivative is bounded
by \(C_j(n+1)^3L^{-3}e^{C\eta(1+n/L)}\): place the three
derivatives in the ordered word, allowing repeated positions.
The first slot derivative costs \(C_jL^{-2}\) with its phase
factor. Hence every first--third product costs \(L^{-5}\)
times a path polynomial and the phase factors. In the
second--second products, the empty-path factors cost \(L^{-2}\)
each and the change of either factor from freezing the slot is
bounded by \(C_j(n+1)L^{-3}e^{C\eta(1+n/L)}\), as in the
cubic calculation. Thus the error in these products also has
price \(L^{-5}\). Both calculations can be summed against the
incoming exponential budgets by requiring \(C\eta/L\le\sigma/2\),
with \(C\) enlarged to bound the phase factors from all slots
in a term. For the frozen rows the identity
\(\lambda_f\,1+(r^0)^{\mathsf T}r^0=1\) gives the local
curvature square. Its fourth derivative has size \(C_jL^{-4}\)
per anchor. The difference
\(q_G(e^\xi)-|\mathcal C(e^\xi)|^2/2\) is \(O(|\xi|^4)\),
so at degree four it uses four first derivatives of
\(\xi=\log U_p\), each a linear curl of size \(C_jL^{-2}\).
An incoming cubic atom at degree four uses slot derivative
orders \((2,1,1)\); an incoming quartic atom uses
\((1,1,1,1)\). Their respective bounds per anchor are
\(L^{-6}\) and \(L^{-8}\) times the weighted path prices.
The count of \(L^4\) anchors therefore gives a bound
\(C_j+A_4/L\) for the smooth fourth derivative. Here and
below \(A_4\) is a finite function of the fixed profile and
incoming path data.

For the mixed cubic take two smooth inputs \(v=Ev_i\),
\(v'=Ev_j\) and the arbitrary input \(u\) in
\eqref{eq:quartic-block-counting}. Write \(Z^0\) for the slot
of the same orientation with empty path. With \(\ell^1(B)\)
on the arbitrary factor, the slot bounds, apart from the
phase factors, are
\[
 \|DZ[u]\|_{1,B}\lesssim_j L^2,\qquad
 \|D^2Z[u,v]\|_{1,B}\lesssim_j L^2+nL .
\]
Indeed the first derivative is the curl at the terminal
site. The plaquette brackets in the second derivative cost
\(L^3L^{-1}=L^2\) in this norm. The differentiated transport
gives one path sum times one curl, with costs
\(nL^3L^{-2}\) and \(nL^{-1}L^2\), according to
which input goes on the path.

In a cubic row product \(DZ[u]D^2Z[v,v']\), first freeze
the smooth factor. The just derived bound and its smooth
freezing error give a cost \(L^2L^{-3}\) times a path polynomial.
For the remaining term, \(DZ[u]\) is a translate of \(DZ^0[u]\).
Move that shift of the anchor to the frozen smooth factor in
the sum over anchors. The product has total momentum zero,
so this change of variable is exact on one full residue block.
The first difference of \(D^2Z^0[v,v']\) is \(O_j(L^{-3})\):
each plaquette bracket gives one difference \(O_j(L^{-2})\)
times one field \(O_j(L^{-1})\). Telescoping a shift of
length at most \(n\) therefore costs \(O_j(nL^{-1})\)
after multiplying by the \(\ell^1(B)\) curl bound.
Thus the calculation places the extra difference on the
smooth factor.

For the products \(DZ[v]D^2Z[u,v']\) and
\(DZ[v']D^2Z[u,v]\), again freeze the smooth factor at
cost \(L^{-3}(L^2+nL)\) times the path polynomial.
The differentiated transport of the other slot costs at most
\(nL\cdot L^{-2}\) in the resulting product. Its remaining
plaquette term is a translate of \(D^2Z^0[u,v']\) or
\(D^2Z^0[u,v]\). Move the shift to the smooth curl, whose first
difference has size \(O_j(L^{-3})\). This costs \(O_j(nL^{-1})\)
using the \(L^2\) bound for the plaquette term. These shift
estimates carry the same phase prices on the strip. In all
three terms the fully frozen contribution has size \(C_j\),
by the row identity and the empty-path bounds. An incoming
cubic atom adds at most \(L^2 L^{-2}L^{-2}=L^{-2}\) times its
weighted budget. Summing the path polynomials as above yields
a bound \(C_j+A_4/L\) for the mixed cubic as a functional on
the selected section. The calculation used a first difference
of \(u\); all second differences fall on smooth lifts.
The estimates apply on each symmetry copy by relabelling
the same link and plaquette types.

\smallskip\noindent
\emph{The actual frame and noise derivatives.}
The norm for a fine arbitrary input on a patch of diameter
\(O_j(L)\), with uniform probability measure, is at most
\(C_jL^{-1}\) in \(L^2\), by the first estimate of
\eqref{eq:quartic-block-counting}. In the coordinate-ordered
routes of \eqref{eq:actual-frame-score} use the product
uniform measure on their two endpoints. At a specified
position in a coordinate segment the starting coordinates
are drawn coordinate by coordinate from the two endpoints.
Their marginal is uniform on the patch because the endpoints
are independent and the patch is a box. The read at that
position is a fixed translate of this marginal, restricted
to routes that have the step, after splitting the two signs
of the segment. Hence an arbitrary input at a derivative
position costs \(C_jL^{-1}\) in \(L^2\) of the endpoint pair.
A smooth input costs \(C_jL^{-1}\) in supremum norm. On
the straight connector in \eqref{eq:actual-connectors} the
same assertion follows from injectivity of the translate of
the patch variable at each step.

For a derivative of order \(r\le3\), distributing its
labelled directions among the \(O_j(L)\) positions costs
at most \((C_jL)^r\) terms. Derivatives at the same position
are symmetrised products of the inserted fields by the
exponential series. Use the sup norm for all-smooth
derivatives and \(L^2\) for a derivative with one arbitrary
input. In each product at most one factor is in \(L^2\);
H\"older's inequality costs at most a fixed constant times
\(L^{-r}\) per term. Missing positions can be padded with
identity matrices, so the sum uses the same set of positions
on the full probability space. This proves bounds \(C_j\)
for the route and connector derivatives in the indicated
norms. Averaging a variable is a contraction. Convolution
with probability weights gives \(\sup_z w(z)\le L^{-4}\),
so the ratio of \(w\) to uniform probability on a patch of
\(O_j(L^4)\) points is bounded.

It remains to pass the route bounds through the implicit
frame equation. Suppress the coarse site from the indices
and let \(\langle\cdot\rangle_\nu\) be the uniform patch average,
\(\langle\cdot\rangle_w\) the average with the predictor weights.
Write \(\Pi=\operatorname{pr}_{\g}\) and
\(R(Z)=\pi_G(Z)-\Pi Z\), so \(DR(I_N)=0\) by
\eqref{eq:projection-derivative}. With \(T_{zz'}\) the
average of the inverse route words from \(z\) to \(z'\), set
\[
 a_z=\big\langle h^c_{z'}(T_{zz'}-I_N)\big\rangle_{\nu,z'},
 \qquad t_z=\langle h^c\rangle_\nu+a_z,\qquad
 s_z=\Pi a_z+R(t_z).
\]
By minimising the score in each frame component,
\(h_z^c=\pi_G(t_z)\). The rooted minimiser also minimises
over unrestricted frames, since a common left rotation
can set the root to \(1\); reversing the routes yields the
same average in the two endpoint contributions to the score.
Equivariance of the projection gives
\(\pi_G(\langle h^c\rangle_w)=I_N\), so
\(\Pi\langle h^c\rangle_w=\Pi I_N\) by the nearest-point
normal equation. Since \(h_z^c=\Pi\langle h^c\rangle_\nu+s_z\)
this yields the exact formula
\begin{equation}
             h_z^c=s_z+\Pi I_N-\langle\Pi s\rangle_w .
 \label{eq:quartic-centred-frame-recursion}
\end{equation}

For a set \(I\) of labelled directions, a subscript \(I\)
denotes their mixed derivative at zero. The product rule
gives
\[
 (a_z)_I=\sum_{A\subsetneq I}
       \big\langle h^c_A(z') (T_{zz'})_{I\setminus A}
                 \big\rangle_{\nu,z'} ,
\]
with \(h^c_\varnothing=I_N\). The highest derivative of
\(h^c\) drops out because \(T_{zz'}-I_N=0\) at zero.
For the orders needed here the other terms of the recursion
are
\[
 \begin{aligned}
 s_i&=\Pi a_i,\\
 s_{ij}&=\Pi a_{ij}+D^2R(I_N)[t_i,t_j],\\
 s_{123}&=\Pi a_{123}+D^3R(I_N)[t_1,t_2,t_3]
               +\sum_{(ij|k)}D^2R(I_N)[t_{ij},t_k].
 \end{aligned}
\]
The sum has the three choices of the singled-out index \(k\),
and \(t_I=\langle h^c_I\rangle_\nu+a_I\).
In each line only lower derivatives of \(h^c\) occur.
Use the probability product bound and then average as in
\eqref{eq:quartic-centred-frame-recursion}. Successively
through order three this gives bounds \(C_j\) for
\(h^c_I\) in sup norm with all-smooth inputs, or \(L^2\)
with one arbitrary input. The projection derivatives are
fixed group tensors. All estimates also apply to their
complex multilinear extensions.

Use these derivatives on both ends of the straight connector,
with inversion at a unitary matrix given by adjoint on the
real group. The connector bounds and the product rule give
the same probability bounds for the centred connector.
Now apply the original weighted average and then \(\pi_G\)
to obtain \(\overline Q\). For example, the second derivative
is the projection of the second derivative of the weighted
average plus \(D^2\pi_G(I_N)\) on the two first derivatives
of that average. The third derivative has, in addition to
the projected third derivative, the three terms with
\(D^2\pi_G(I_N)\) on a second and a first derivative and the
one with \(D^3\pi_G(I_N)\) on three first derivatives. All
these coarse link tensors have bounded size before \(L\).

On the chosen section the trial coarse frame and all its
variations are zero. Let \(\xi_b=\log(V_b\overline Q_b^{-1})\)
be the noise mismatch logarithm in these coordinates.
Its needed first, second and third derivatives follow from
the preceding bounds and the fixed group maps. Since
\(q_G(e^{\xi})=|\xi|^2/2+O(|\xi|^4)\), the mixed cubic
noise derivative uses only the three pairings of a first
and a second mismatch derivative. The all-smooth fourth
derivative uses the three second--second pairings, the four
first--third pairings, and the fourth derivative of
\(q_G\circ\exp\) on four first mismatch derivatives.
Each of these has size \(C_j\) per coarse link, with four
outgoing links per site. The norm estimates use the actual
predictor, because the change of frames was an exact change
of variables in \eqref{eq:quartic-centred-phase}.

The centred recursion and the weighted inverse follow the
construction in \cite[Lemmas~36.4--36.5 and
Proposition~36.6, pp.~236--243]{Main}. Only one arbitrary input
is needed for the quartic stationary correction. In the
fine-word bound the shifts of the anchor place the extra
difference on the smooth factor.


\subsection{Four curvature slots and the classical remainder}

The mixed bound on each section gives a bound on the dual
quotient norm by the comparison of section and quotient
norms. A source generated by waves \(i,j\) acts on vertical
waves at \(-p_i-p_j\). The total Hessian pairs opposite modes
and has the same factor \(|\Lambda_c|\) as the total source
when both are written per block. Dividing the Hessian
equation by this factor gives the fibre equation. Thus the
fourth derivative and the paired terms per block in
\eqref{eq:quartic-centred-jet}, on unit coarse amplitudes, have
a sum bounded by
\[
              C_j+A_4/L+C_j(C_j+A_4/L)^2 .
\]
The coarse row action and \(P_3\) have fixed coarse path
budgets. A word of length \(n\) costs at most a fixed
polynomial in \(n+1\) through order four, and the phase
factor costs \(e^{C\eta(n+1)}\).
Use a width with \(C\eta\) at most half the available coarse
path widths. The fourth derivative of the full \(P_3\) is
included in this bound.

Choose the common momentum width using the free and
geometric bounds and the fixed coarse path widths. It is
chosen small enough that all four momenta and all paired sums
in the quartic jet lie in the established strips. With the
incoming width and budget fixed, impose the classical floor
\(C\eta/L\le\sigma/2\) and \(L\ge\max(A_3,A_4)\), where \(A_3\)
is from \eqref{eq:cubic-kernel-before-ratio}. These
conditions permit arbitrarily large odd \(L\). They bound
the symmetric four-linear kernel
\[
       K_4\ \text{of}\quad
       T_4(\Psi-S_{\rm row}\circ V-P_3\circ V)
\]
on that strip by a constant fixed before \(L\), with
normalisation \(1/4!\). The formulas give its periodic
holomorphic extension on the integer cover and its
periodisations on the compatible tori. This follows
also by comparison at the discrete momenta with
\eqref{eq:quartic-centred-jet} in the tree chart. At fixed
\(L\) the rooted Taylor series already gives a periodised
kernel, and agreement on the dense grids identifies
its symbol with the one bounded here.

Write \(p_0=-p_1-p_2-p_3\).
The Ward identities give \(\iota_iK_4=0\) for all four
legs, by the same Fourier test as in the cubic case.
The Laurent primitives \(\mathcal J_i\) used there give
a short four-leg calculation. Work first in
\((p_0,p_2,p_3)\) and set
\(X=\mathcal J_0K_4\), \(Y=\mathcal J_0\iota_1X\).
Then \(\iota_0X=K_4\) and \(\iota_0Y=\iota_1X\), and
both \(X,Y\) are killed by \(\iota_2,\iota_3\).
On \(p_1=0\), the restriction of \(Y\) is killed
also by \(\iota_0\). Take its leg-\(0\) primitive \(Z_0\)
in the independent variables \((p_0,p_3)\). This is
a four-form in leg \(0\). For \(i=2,3\),
\(\iota_0\iota_iZ_0=\iota_i(Y|_{p_1=0})=0\).
Injectivity of contraction on the top degree away from
\(q(p_0)=0\) and continuity at zero give \(\iota_iZ_0=0\).

In \((p_1,p_2,p_3)\) extend \(Z_0\) independent of \(p_1\)
to a coefficient \(Z\), keeping \(\iota_2Z=\iota_3Z=0\).
Set \(D=\mathcal J_1(Y-\iota_0Z)\). The argument of
\(\mathcal J_1\) is scalar in leg \(1\) and vanishes at
\(p_1=0\), so \(\iota_1D=Y-\iota_0Z\). Now \(X-\iota_0D\)
is killed by \(\iota_1,\iota_2,\iota_3\) and contracts
to \(K_4\) in leg \(0\). Dividing in the three independent
legs gives the two-form coefficient in each leg,
\[
 K_4^{\rm curv}=\mathcal J_3\mathcal J_2\mathcal J_1(X-\iota_0D),
 \qquad K_4=\iota_0\iota_1\iota_2\iota_3 K_4^{\rm curv}.
\]
For the widths use the coefficient norm in the cubic
argument with a third displacement \(z_3\) included.
Contour shifts give a finite norm at any fixed width
up to one quarter of the kernel strip. A change to
\((p_0,p_2,p_3)\) replaces \((z_1,z_2,z_3)\) by
\((-z_1,z_2-z_1,z_3-z_1)\) and thus costs at most a
factor three in the width. The inverse change and
the changes on \(p_1=0\) have the same bound.
Restriction and extension cost at most one in the
corresponding coordinates. Apply
\eqref{eq:packet-laurent-width-loss} at each of the
seven divisions, keeping half the current width,
and use the finite coordinate changes just described.
This gives a positive output width and a bound for
\(K_4^{\rm curv}\) at that width, both fixed before \(L\).
This is the four-leg case of the curl division in
\cite[Lemma~36.3, pp.~230--232]{Main}.

\begin{proposition}[Quartic matching of the classical action]
\label{prop:quartic-curvature-matching}
There is a real invariant quartic slot correction \(P_4(V)\)
of the form allowed in \eqref{eq:stationary-higher-budget}.
With \(S_c=S_{\rm row}+P_3+P_4\) the classical agreement is
\begin{equation}
                T_k(\Psi-S_c\circ V)=0,\qquad 0\le k\le4 .
 \label{eq:quartic-matching}
\end{equation}
The corrections periodise from the integer cover and are
independent of the actual coupling. Their output path
widths and budgets can be fixed before \(L\), with the
classical floor for \(L\) depending on the fixed incoming
path data. The bounds are uniform in the period, the
terminal cutoff and the legal finite histories at the
fixed adjacent profiles.

On the common analytic chart at a common reserved spatial
width, \(\beta(\Psi-S_c\circ V)\) has an exact expansion
\(\sum_\alpha r_\alpha\) by connected footprints of the type in
Corollary~\ref{cor:local-stationary-remainder}, with
\begin{equation}
 \sup_o\sum_{\alpha:o_\alpha=o}e^{\kappa n_B(X_\alpha)}
    \|r_\alpha\|^{\rm chart}_{q,r,s;o_\alpha,X_\alpha}
                    \lesssim_{q,s,\kappa}\beta r^5 .
 \label{eq:quartic-matching-remainder}
\end{equation}
\end{proposition}
\begin{proof}
For a coefficient with displacements \(z_1,z_2,z_3\), use
the empty-path slot at the anchor and the three slots
reached by coordinate paths to these displacements.
Their total length is \(16+\sum_{i=1}^3|z_i|_1\).
As in Proposition~\ref{prop:cubic-curvature-matching},
take real coefficients and average the contractions over
constant adjoint rotations and the prescribed signed lattice
symmetries. The quartic residual has both invariances.
The projections commute with its reconstruction from linear
curls, so the resulting full correction has leading term \(K_4\)
and the same type of exponential budget. The four slots all
vanish at the identity, which proves \eqref{eq:quartic-matching}
using the cubic agreement. The phase and both matching
calculations use the action with its common \(\beta\) removed.

For the remainder, use the fixed-radius rooted analytic
majorant for \(\Psi\) from
Corollary~\ref{cor:local-stationary-remainder}. The coarse
row terms and both corrections have a scalar majorant with
the same type of spatial labels. In fact their coefficient
sum for a word of length \(n\) is bounded by its weight
times \((n+1)^C e^{C(r_*+\kappa)(n+1)}\), by the
exponential series argument of
Lemma~\ref{lem:local-path-majorant}. Choose the fixed radius
\(r_*\) and spatial reserve within the common path widths.
The sum over labels converges at that radius. By
\eqref{eq:quartic-matching} the sum of labels of degree at
most four is zero, so keeping the higher-degree labels
gives an exact expansion of the difference. For every
remaining degree \(k\ge5\) and \(2r\le r_*\), use
\((2r)^k\le(2r/r_*)^5r_*^k\), with the extra factor two
paying the derivative polynomial as in
Corollary~\ref{cor:local-stationary-remainder}.
This gives \eqref{eq:quartic-matching-remainder}.
\end{proof}

This completes the classical part of
\cite[Proposition~36.6 and (36.8)--(36.12), pp.~240--243]{Main}.
These locators were checked on 20 November 2026.
The one-loop quadratic matching concerns \(T_2(T)\),
where \(T\) is the one-loop combination in
\eqref{eq:matching-one-loop-action}.

\appendix
\section{The exponential chart and its Haar Jacobian}
\label{app:haar}

The coordinate argument above needs a local chart and its exact
measure density. We give the background calculation in full,
starting from the stated matrix-group data. The chart construction
is the argument of \cite[Lemma~1.2, pp.~12--13]{Main}; the volume
and differential calculations below make the Jacobian explicit.

\subsection{A chart without a connectedness assumption}

If \(X\in\g\), differentiating
\((e^{tX})^*e^{tX}=I_N\) at zero gives \(X^*=-X\).
Thus \(\g\) lies in the real vector space
\(\mathcal A=\{X:X^*=-X\}\). Since \(\g\) is a finite-dimensional
linear subspace, it is closed. For every \(R\in G\),
\[
               e^{tRXR^{-1}}=Re^{tX}R^{-1}\in G,
\]
so \(RXR^{-1}\in\g\). In particular
\(e^{tX}Ye^{-tX}\in\g\) for \(X,Y\in\g\), and its derivative
at zero is \([X,Y]\). Closedness of \(\g\) gives
\([X,Y]\in\g\).

Near \(I_N\), the convergent series
\[
              \log(I_N+Z)=\sum_{j\ge1}\frac{(-1)^{j+1}}j Z^j
\]
is an analytic inverse of matrix exponential. For a unitary matrix
near \(I_N\), its logarithm is skew-Hermitian: by the convergent
matrix identities, \((\log U)^*=\log U^*=\log U^{-1}=-\log U\).
Hence exponential identifies neighbourhoods in \(\mathcal A\)
and \(U(N)\).

Let \(\g^\perp\) be the orthogonal complement in \(\mathcal A\)
for the normalised Frobenius inner product
\eqref{eq:physical-norm}. In the preceding logarithm chart,
the map
\[
       (X,Y)\longmapsto e^Xe^Y,
       \qquad X\in\g,\quad Y\in\g^\perp,
\]
has derivative \((X,Y)\mapsto X+Y\) at \((0,0)\), which is
invertible. The finite-dimensional analytic inverse function
theorem therefore gives unique small coordinates \((X,Y)\) for
every unitary matrix sufficiently near \(I_N\).

We claim that group elements sufficiently near \(I_N\) have
\(Y=0\). If this failed, there would be \(R_j\in G\) with
\(R_j\to I_N\) and coordinates \((X_j,Y_j)\to(0,0)\),
\(Y_j\ne0\). Because \(e^{X_j}\in G\),
\[
                          e^{Y_j}=e^{-X_j}R_j\in G.
\]
After a subsequence,
\(Y_j/|Y_j|_{\g}\to Y\in\g^\perp\), \(|Y|_{\g}=1\);
here the same normalised Frobenius norm is used on \(\mathcal A\).
For each real \(t\), choose integers \(k_j\) such that
\(k_j|Y_j|_{\g}\to t\). Then
\[
                e^{k_jY_j}=(e^{Y_j})^{k_j}\in G,
                \qquad e^{k_jY_j}\longrightarrow e^{tY}.
\]
Closedness of \(G\) implies \(e^{tY}\in G\) for every real \(t\).
By the definition of \(\g\), this means \(Y\in\g\), contradicting
\(Y\in\g^\perp\) and \(|Y|_{\g}=1\).

It follows that exponential restricts to an analytic chart of \(G\)
near the identity. Conversely every small \(X\in\g\) has
\(e^X\in G\) by definition. Left translating this chart gives
charts near every element of \(G\), including all its components.
Shrinking the initial ball as necessary gives the radius \(r_G\)
used above. This argument has added no connectedness hypothesis.

\subsection{From coordinate volume to the derivative of exponential}

The fixed inner product on \(\g\) defines a left-invariant
Riemannian metric by
\[
           \langle UV,UW\rangle_{T_UG}=\langle V,W\rangle_{\g},
           \qquad V,W\in\g.
\]
Let \(\dd\operatorname{vol}_G\) be its Riemannian volume measure.
It is left invariant. Compactness of \(G\) gives
\(0<\Vol(G)<\infty\), so Haar uniqueness yields
\begin{equation}
     \dd h_G=\zeta_G\,\dd\operatorname{vol}_G,
     \qquad \zeta_G=\Vol(G)^{-1}.
 \label{eq:haar-volume}
\end{equation}
The volume here is that of the entire group for this metric.

Choose an orthonormal basis \(E_1,\ldots,E_k\) of \(\g\),
where \(k=\dim_{\R}\g\), and write
\(X=\sum_{a=1}^kx^aE_a\). In any coordinate parametrisation
\(\Phi\), the Riemannian volume density is the square root of
the determinant of its Gram matrix:
\begin{equation}
 \dd\operatorname{vol}_G(\Phi(x))
   =\sqrt{\det\bigl(\langle\partial_a\Phi,
                                 \partial_b\Phi\rangle\bigr)_{a,b}}
                                                 \,\dd x^1\cdots\dd x^k.
 \label{eq:coordinate-volume}
\end{equation}
The Gram matrix records the lengths and mutual angles of the
coordinate tangent vectors. For \(\Phi(X)=e^X\), left translation
back to the identity gives its entries as
\begin{equation}
   \big\langle e^{-X}D\exp_X[E_a],
                       e^{-X}D\exp_X[E_b]\big\rangle_{\g}.
 \label{eq:exponential-Gram}
\end{equation}
It remains to compute the linear map in this expression.

\begin{lemma}[Left-trivialised exponential differential]
\label{lem:exp-differential}
For \(X,Y\in\g\),
\begin{equation}
     e^{-X}D\exp_X[Y]
        =\int_0^1e^{-tX}Ye^{tX}\,\dd t
        =\left(\int_0^1e^{-t\ad X}\,\dd t\right)Y.
 \label{eq:exponential-differential}
\end{equation}
\end{lemma}
\begin{proof}
For a real parameter \(s\), differentiate the matrix product
\(e^{(1-t)X}e^{t(X+sY)}\) in \(t\). The terms containing \(X\)
cancel, giving
\[
 \frac{\dd}{\dd t}\left(e^{(1-t)X}e^{t(X+sY)}\right)
                     =s e^{(1-t)X}Ye^{t(X+sY)}.
\]
Integration from zero to one gives
\[
     e^{X+sY}-e^X=s\int_0^1e^{(1-t)X}Ye^{t(X+sY)}\,\dd t.
\]
Divide by \(s\) and pass to zero; the integrand converges uniformly
on the compact \(t\)-interval. Hence
\[
               D\exp_X[Y]=\int_0^1e^{(1-t)X}Ye^{tX}\,\dd t.
\]
Multiplication on the left by \(e^{-X}\) proves the first identity.
For the second, \(Z(t)=e^{-tX}Ye^{tX}\) satisfies
\[
                         Z'(t)=-[X,Z(t)],\qquad Z(0)=Y.
\]
The solution of this linear equation on \(\g\) is
\(Z(t)=e^{-t\ad X}Y\), proving the claim.
\end{proof}

Set \(A_X=\int_0^1e^{-t\ad X}\,\dd t\), a real linear map
on \(\g\). Equations~\eqref{eq:exponential-Gram} and
\eqref{eq:exponential-differential} show that the Gram matrix is
\(A_X^{\mathsf T}A_X\) in the chosen orthonormal basis. Therefore
\[
                 \sqrt{\det(A_X^{\mathsf T}A_X)}=|\det A_X|.
\]
Substituting into \eqref{eq:coordinate-volume} and
\eqref{eq:haar-volume} proves
\[
          \dd h_G(e^X)=\zeta_G|\det A_X|\,\dd X.
\]
As \(A_0=1\), its determinant equals one at zero and stays
positive on a smaller chart. It is smooth, so the function
\(j_G(X)=|\det A_X|\) is smooth, positive, and satisfies
\(j_G(0)=1\). The last assertion of Lemma~\ref{lem:haar-chart}
follows from smoothness on a neighbourhood of the compact ball
\(\{|X|_{\g}\le r_G/2\}\) and its positive minimum there.
Finally right Haar invariance gives the translated formula
\(\dd h_G(e^XR)=\dd h_G(e^X)\) for fixed \(R\).

\section{Nearest group points in matrix space}
\label{app:projection}

The branch construction uses nearest points for the normalised
Frobenius norm \eqref{eq:physical-norm}. We establish their local
uniqueness and regularity directly from the analytic matrix-group
charts proved in Appendix~\ref{app:haar}.

\begin{lemma}[A uniform nearest-point neighbourhood]
\label{lem:nearest-group-point}
There is \(r_*>0\), depending only on the fixed matrix group,
such that every matrix \(A\) with
\(\operatorname{dist}(A,G)<r_*\) has a unique nearest point
\(\pi_G(A)\in G\). This map is real analytic on that neighbourhood.
It satisfies \eqref{eq:projection-derivative} and, for a sufficiently
small ball about \(I_N\),
\eqref{eq:projection-small-variation}. In addition,
\begin{equation}
                      \pi_G(RAS)=R\pi_G(A)S\qquad(R,S\in G).
 \label{eq:projection-equivariance}
\end{equation}
\end{lemma}
\begin{proof}
Regard \(\operatorname{Mat}_N(\C)\) as a real Euclidean space
with inner product \(N^{-1}\operatorname{Re}\tr(A^*B)\).
The tangent space of \(G\) at \(U\) is \(T_UG=U\g\).
A normal vector at \(U\) means a matrix orthogonal to this
tangent space. Consider the normal-coordinate map
\[
                    (U,Z)\longmapsto U+Z,
                    \qquad Z\perp T_UG.
\]
Analytic group coordinates and local orthonormal coordinates for
these normal spaces make this an ordinary analytic map between
real spaces of the same dimension. At \((U,0)\) its derivative
is \((V,Z)\mapsto V+Z\), the invertible orthogonal decomposition
into the tangent and normal spaces. The analytic inverse function
theorem gives a local analytic inverse. By compactness of \(G\),
we can choose normal vectors uniformly small so that the derivative
remains invertible in such charts about the whole zero section.

A nearest group point exists for every matrix \(A\), by compactness.
If \(U\) is nearest, first variation of squared distance along
curves in \(G\) gives
\[
                   \langle A-U,V\rangle=0\qquad(V\in T_UG).
\]
Thus \(A=U+Z\) with \(Z\perp T_UG\).
There is a uniform positive distance within which nearest points
are unique. Otherwise there would be matrices \(A_j\) with
\(\operatorname{dist}(A_j,G)\to0\) and two distinct nearest
points \(U_j,V_j\). After a subsequence, compactness gives
\(U_j\to U\), and
\[
        \|U_j-V_j\|_{\rm F}
                 \le2\operatorname{dist}(A_j,G)\longrightarrow0.
\]
Hence \(V_j\to U\) as well. Both normal representations of
\(A_j\) eventually lie in one inverse-function neighbourhood
about \((U,0)\), contradicting uniqueness of its inverse.
Shrink the distance radius further so that the preceding analytic
normal-coordinate inverses apply throughout it. They identify
the unique nearest point and prove that \(\pi_G\) is analytic.

At \(I_N\), the derivative of the normal-coordinate inverse
takes a matrix to its tangent and normal components. Its group
component therefore gives \(D\pi_G(I_N)=\operatorname{pr}_{\g}\).
Choose a convex ball about \(I_N\), contained in the nearest-point
neighbourhood, on which
\(\|D\pi_G(A)-\operatorname{pr}_{\g}\|\le1/4\).
Integrating that derivative along the segment from \(B\) to
\(A\) proves \eqref{eq:projection-small-variation}.
The operator norm here is the one induced by the real matrix
Euclidean norm fixed above.

Finally, left and right multiplication by \(R,S\in G\) are
isometries of the matrix space which map \(G\) bijectively to
itself. They preserve distance to \(G\). Applying them to the
unique nearest-point minimisation proves
\eqref{eq:projection-equivariance}.
\end{proof}

\section{The Gaussian history and the two calibrated columns}
\label{app:calibration}

This appendix proves the linear input used in
Theorem~\ref{thm:calibrated-row-data}. The calculation has two
parts. Cumulative filters control every finite history without a
constant raised to its depth. Two explicit columns then identify
the Gaussian minimiser and its dual curvature map, with the
fine price \(L^{-2}\) required by the central comparison.
All symbols act identically on colour components.

\subsection{Filters, counting measure and the filled curvature metric}

For \(p\in\C^4\) set \(q_i(p)=e^{ip_i}-1\), so the symbol of
\(d_r\) is exterior multiplication by \(q(p)\). For odd \(M\)
define
\[
 r_M(z)=M^{-1}\sum_{\nu=0}^{M-1}e^{i\nu z},\quad
 a_M(p)=\prod_i e^{-i(M-1)p_i/2}r_M(p_i),\quad
 \Delta_M(p)=M^2\sum_iq_i(p/M)q_i(-p/M).
\]
The transpose of a real kernel has analytic symbol
\(A^\dagger(p)=A(-p)^{\mathsf T}\), which is the Hermitian adjoint on
real momenta. At an alias, the dagger reverses the full alias
momentum. For a list \(z_i\), let \(\operatorname{diag}_r(z)\)
be diagonal on \(r\)-cell types, with entry \(\prod_{i\in I}z_i\)
at type \(I\). Thus \(\operatorname{diag}_2(z)\) uses products
of two entries, not their sums.

Regrouping finite sums gives exactly
\[
 r_{LM}(z)=r_L(Mz)r_M(z),\qquad
                         a_{LM}(p)=a_L(Mp)a_M(p).
\]
A history \(v\) has mass \(M_v\) and filter \(s_v\). Initially
\(M_v=1,s_v=1\). Appending a side \(L\) at fine index \(j\)
sets
\begin{equation}
 M_{v'}=LM_v,\qquad s_{v'}(p)=a_L(M_vp)^js_v(p).
 \label{eq:calibration-filter-history}
\end{equation}
The coarse index is \(t\in\{j,j-1\}\). If the fine history
has \(s_v=a_{M_v}^j e_v\), the product identity gives, in
the two cases,
\begin{equation}
 s_{v'}=
 \begin{cases}
  a_{LM_v}^{j}e_v,&t=j,\\
  a_{LM_v}^{j-1}(a_{LM_v}e_v),&t=j-1.
 \end{cases}
 \label{eq:calibration-drop-factorisation}
\end{equation}
The empty history has \(a_1=e_v=1\), so induction proves
the representation at every depth. A repetition of an
index leaves the extra factor \(e_v\) unchanged; a decrease
adds exactly one cumulative filter. Thus for a final history
\(s_v=a_{M_v}^j e_v\), with
\(e_v=\prod_\nu a_{m_\nu}\), the extra masses, ordered
from largest to smallest, satisfy
\(m_\nu/M_v\le3^{-\nu}\). In particular their sum of
ratios is at most \(3/2\), independently of the number
of index repetitions or decreases.

For this history take \(M=M_v\), and for centred aliases
\(p_\ell=(p+2\pi\ell)/M\) put
\(w_j(\ell)=\prod_i(1+|\ell_i|)^{-j}\).
We work first on enlarged cubes
\(|\Re p_i|<\pi+\delta,\ |\Im p_i|<\delta\).
The sine quotient
\[
 r_M(z)=e^{i(M-1)z/2}\frac{\sin(Mz/2)}{M\sin(z/2)}
\]
gives \(|r_M((p_i+2\pi\ell_i)/M)|\le C/(1+|\ell_i|)\).
For the principal alias on the real cube the centred quotient
is positive and bounded below. For each extra mass \(m\le M\),
its complex upper bound is \(e^{C\delta m/M}\), and its real
principal lower bound is \(e^{-Cm/M}\). For the latter bound use the centred quotient directly:
\[
 a_m(p/M)=\prod_i
 \frac{\sin(mp_i/(2M))/(mp_i/(2M))}
      {\sin(p_i/(2M))/(p_i/(2M))}
       \ge e^{-C(m/M)^2}\ge e^{-Cm/M}.
\]
All zero quotients are read by continuity. On this fixed real
range, \(\sin t/t\ge e^{-Ct^2}\): Taylor expansion gives the
bound near zero, and positivity on the remaining compact interval
gives one larger constant. The denominator quotients are at most
one. This proves the lower bound uniformly in \(m\le M\). Multiplication of the extra
factors now gives, without a constant per factor,
\[
\begin{aligned}
 |e_v(p_\ell)|&\le\exp\!\left(C\delta\sum_\nu m_\nu/M_v\right)
                  \le e^{3C\delta/2},\\
 e_v(p/M_v)&\ge\exp\!\left(-C\sum_\nu m_\nu/M_v\right)
                  \ge e^{-3C/2}\quad(p\text{ real}).
\end{aligned}
\]
Together with the \(j\)-th power of the cumulative filter,
this proves, for fixed final index,
\begin{equation}
 |s_v(p_\ell)|\le C_jw_j(\ell),\qquad
 s_v(p/M_v)\ge c_j>0\quad(p\text{ real}).
 \label{eq:calibration-history-filter-bound}
\end{equation}
The constants and cube width do not depend on the history.
Also \(\sum_{\ell\in\Z^4}w_2(\ell)<\infty\), by the product
of four convergent one-dimensional sums.

Put \(b_{v,i}=s_vr_{M_v}(p_i)\), and for scalar \(F\) set
\[
 (\mathcal Q_{v,i}F)(p)=M_v^{-2}\sum_\ell
                  b_{v,i}(p_\ell)b_{v,i}(-p_\ell)F(p_\ell).
\]
We verify the prefactor. With counting Fourier transform
\(\widehat u(k)=\sum_x e^{-ikx}u(x)\), the rescaled block
average is \((\mathcal B_vu)_i(X)=M\sum_z b_i(z)u_i(MX+z)\).
If the fine period is \(M\) times the coarse period, summing
over \(X\) gives
\[
 \widehat{\mathcal B_vu}_i(p)=M^{-3}\sum_\ell
                                      b_{v,i}(p_\ell)\widehat u_i(p_\ell).
\]
For a finite positive diagonal Gaussian covariance, opposite
aliases pair uniquely. Dividing the resulting covariance by
the coarse Fourier volume gives
\(M^{-6}M^4=M^{-2}\), exactly \(\mathcal Q_{v,i}\).
There is no further average over the alias set.
This computation does not assign an infinite variance to a zero
mode; the singular scalar expression below is used only away
from zero and its needed combinations are subsequently filled.

Let \(\mathcal P_v\) list the noise suffix histories: it is empty
initially, and a step appends that step to every old member and
adds the empty newest suffix. Define, for nonzero real \(p\),
\begin{equation}
 D_{v,i}=\mathcal Q_{v,i}\Delta_1^{-1}
                       +a^{-1}\sum_{u\in\mathcal P_v}\mathcal Q_{u,i}1.
 \label{eq:calibration-history-covariance}
\end{equation}
These diagonal coefficients are auxiliary covariances, not the
inverse of a gauge-degenerate potential kinetic operator.
Their noise part is uniformly bounded: the suffix masses are
\(1\) followed by masses increasing by at least three, and
\[
 |\mathcal Q_{u,i}1|\le C_jM_u^{-2},\qquad
                         \sum_{u\in\mathcal P_v}M_u^{-2}\le9/8.
\]
For one step set \(f(q)=a_L(q)^j\) and
\(h_r(q)=f(q)\operatorname{diag}_r(r_L(q_i))\).
Splitting an \(LM\)-alias into an \(L\)-alias and an
\(M\)-alias in \eqref{eq:calibration-history-covariance} gives
\begin{equation}
 D_{v',i}(p)=L^{-2}\sum_\ell f(q_\ell)f(-q_\ell)
       r_L(q_{\ell,i})r_L(-q_{\ell,i})D_{v,i}(q_\ell)+a^{-1},
 \qquad q_\ell=(p+2\pi\ell)/L.
 \label{eq:calibration-one-step-covariance}
\end{equation}
The filter identity is
\(b_{v',i}(k)=f(Mk)r_L(Mk_i)b_{v,i}(k)\); the two counting
factors multiply to \(L^{-2}M^{-2}\). The new empty suffix
is exactly the last \(a^{-1}\).

Separate the principal alias and define regular functions
\begin{align*}
 D_{v,i}&=L_{v,i}/\Delta_{M_v}+R_{v,i},\qquad
 L_{v,i}=b_{v,i}(p/M_v)b_{v,i}(-p/M_v),\\
 G_{v,i}&=L_{v,i}+\Delta_{M_v}R_{v,i},\qquad
 U_v=s_v(p/M_v)s_v(-p/M_v),\\
 b_v&=U_v^{-1}-\sum_i
             \frac{q_i(p)q_i(-p)R_{v,i}}{L_{v,i}G_{v,i}}.
\end{align*}
The scalar \(b_v\) has no axis index and is distinct from the
filter \(b_{v,i}\). For nonprincipal aliases the sine bound gives
\[
 |M_v^2\Delta_1((p+2\pi\ell)/M_v)|
                                      \ge c(1+|\ell|_\infty)^2.
\]
On a thin complex cube its real part remains positive: the change
from imaginary parts is at most \(C\delta(1+|\ell|_\infty)\).
Together with the summable filter bound and geometric noise sum,
this bounds \(R_{v,i}\) holomorphically and uniformly.
On the real cube \(R_{v,i}\ge0\), so
\(G_{v,i}\ge L_{v,i}\ge c_j\).
The line identity
\(q_i(p)=M_vr_{M_v}(p_i/M_v)q_i(p/M_v)\) implies the exact
cancellation
\begin{equation}
 \Delta_{M_v}(p)b_v(p)
                         =\sum_i\frac{q_i(p)q_i(-p)}{G_{v,i}(p)}.
 \label{eq:calibration-filled-identity}
\end{equation}
Indeed replace \(1/G_{v,i}\) by
\(1/L_{v,i}-\Delta_{M_v}R_{v,i}/(L_{v,i}G_{v,i})\)
and sum. Both \(\Delta_{M_v}(p)\) and \(\sum_i|q_i(p)|^2\)
are comparable to \(|p|^2\) on the real cube, uniformly in
\(M_v\). Thus \(c_j\le b_v\le C_j\).
At zero the regular definitions give \(G_{v,i}(0)=b_v(0)=1\).
Cauchy estimates on a larger cube bound first derivatives
uniformly. Shrinking its width in terms of these bounds and the
real lower bounds makes \(G_{v,i}^{-1},b_v^{-1}\) holomorphic
and uniformly bounded. No division by \(\Delta_{M_v}\) at zero
is used in this argument.

Define
\begin{equation}
 H_v=b_v^{-1}\operatorname{diag}_2(G_{v,i}^{-1}),\qquad
 B_v=H_v^{-1}=b_v\operatorname{diag}_2(G_{v,i}).
 \label{eq:calibration-filled-metric}
\end{equation}
These positive metrics glue periodically. To see this before
taking coefficients, away from zero put
\(\rho_v=\sum_iq_i(p)q_i(-p)/D_{v,i}(p)\).
Then \(D_{v,i}=G_{v,i}/\Delta_{M_v}\),
\(\rho_v=\Delta_{M_v}^2b_v\), and
\(H_v=\rho_v^{-1}\operatorname{diag}_2(D_{v,i}^{-1})\).
The latter expression is periodic by alias relabelling.
Analytic continuation of the regular formula therefore gives
periodicity across the filled zero as well.

A bounded periodic symbol holomorphic on a strip of width
\(\delta\) has coefficient bound \(Ce^{-\delta'|u|_1}\) for
any \(\delta'<\delta\), by shifting each Fourier contour
opposite to the sign of \(u_i\). Summing four geometric series
gives its exponential coefficient norm at every smaller width.
Absolute convolution is submultiplicative, and periodisation
cannot increase an absolute row sum. Applying this observation
to \eqref{eq:calibration-filled-metric} gives prices uniform in
history and period.

For nonempty histories use the prescribed nonnegative \(l_j\)
to define
\begin{equation}
 \widetilde H_v=
  [B_v+l_j(d_2B_v)^\dagger(d_2B_v)]^{-1};
 \qquad \widetilde H_v=1\text{ at the empty history}.
 \label{eq:calibration-positive-extension}
\end{equation}
For fixed \(l_j\), its positive upper and lower bounds and
its holomorphic inverse are uniform over histories and periods.
The compatible choices of the free \(l_j\) for physical packet
payment are made in Lemma~\ref{lem:packet-complement};
\eqref{eq:packet-positive-parameter-choice} uses bounds from
unmodified metrics to select these parameters before the block
sides. The row-factor choice below is made afterwards.
Multiplying \(H_vd_1\) by the inverted matrix
and using \(d_2d_1=0\) proves
\[
 \widetilde H_vd_1=H_vd_1,\qquad
                         d_1^\dagger\widetilde H_v=d_1^\dagger H_v.
\]
Choose \(\lambda_j>0\) below half the common real lower
spectral bound for \(\widetilde H_v\).
The positive square root \(R_v=(\widetilde H_v-\lambda_j1)^{1/2}\)
has a common holomorphic strip: use one contour in the right
half-plane surrounding the common spectral interval in
\[
 R_v=\frac1{2\pi i}\int_\Gamma
       z^{1/2}(z-\widetilde H_v+\lambda_j1)^{-1}\,\dd z.
\]
The real resolvent bound persists on a smaller strip by Cauchy's
estimate and the resolvent identity. The contour gives bounded
holomorphic symbols with the real-kernel symmetry, hence
exponential coefficients. This proves all metric and row-factor
assertions in \eqref{eq:calibrated-row-metric} before choosing \(L\).

\subsection{The two columns and their cancellations}
\label{sec:calibration-column-symbols}

Fix a fine history \(v\) of index \(j\ge4\), an odd side
\(L\ge3\), and its successor \(v'\), whose index is
\(t\in\{j,j-1\}\), with \(t\ge4\). Constants denoted by
\(C_j,c_j\) below may depend on this fixed adjacent pair
and \(a\), but are chosen before \(L\), history depth,
cumulative mass and period. The aliases are indexed by
\[
 \mathcal A_L=\{\ell\in\Z^4:|\ell_i|\le(L-1)/2\},
                    \qquad q_\ell=(p+2\pi\ell)/L.
\]
There are \(L^4\) aliases. We first calculate on nonzero
real coarse momenta, then remove the apparent poles before
Fourier inversion.

Retain the one-step filters \(f(q)=a_L(q)^j\) and
\(h_r(q)=f(q)\operatorname{diag}_r(r_L(q_i))\). Put
\[
 F_\ell=\operatorname{diag}_1(L^{-2}D_{v,i}(q_\ell)),\qquad
 D=\operatorname{diag}_1(D_{v',i}(p)),
\]
\[
 \mathbf q_\ell=Ld_0(q_\ell),\quad \mathbf q=d_0(p),\quad
 c_\ell=Ld_1(q_\ell),\quad c=d_1(p),\quad K=c^\dagger H_cc,
\]
where \(H_f=H_v(q_\ell)\) and \(H_c=H_{v'}(p)\). Thus
\(H_f\) is evaluated at the alias in which it occurs.
Define the scalars and row
\[
 \rho_\ell=\mathbf q_\ell^\dagger F_\ell^{-1}\mathbf q_\ell,
 \quad \rho=\mathbf q^\dagger D^{-1}\mathbf q,
 \quad \pi=\rho^{-1}\mathbf q^\dagger D^{-1},
 \quad Z=\sum_\ell f(q_\ell)f(-q_\ell)\tr F_\ell,
\]
and set \(u_\ell=f(-q_\ell)\tr F_\ell/Z\) and
\(y_\ell=f(q_\ell)f(-q_\ell)\rho/\rho_\ell\).
For real nonzero \(p\), the covariances and the principal
filter are positive, so all scalar denominators here are
nonzero. The two columns are
\begin{align}
 \mathcal E_\ell&=F_\ell h_1(q_\ell)^\dagger K
                                      +u_\ell\mathbf q_\ell\pi,
 \label{eq:calibration-E-alias}\\
 \mathcal T_\ell&=H_fc_\ell F_\ell h_1(q_\ell)^\dagger c^\dagger
             +y_\ell h_2(q_\ell)^\dagger(1-H_ccDc^\dagger).
 \label{eq:calibration-T-alias}
\end{align}
These formulas specify the columns also on gradient and
nonexact curvature directions. The individual pole terms
will not be Fourier transformed separately.

\smallskip\noindent
\emph{Four identities from one diagonal calculation.}
The diagonal two-form metric has entry
\((\rho_\ell F_{\ell,i}F_{\ell,k})^{-1}\) at type
\(\{i,k\}\). Expanding the curl matrix therefore gives
\begin{equation}
 c_\ell^\dagger H_fc_\ell
 =F_\ell^{-1}-F_\ell^{-1}\mathbf q_\ell\rho_\ell^{-1}
                                      \mathbf q_\ell^\dagger F_\ell^{-1}.
 \label{eq:calibration-Maxwell-identity}
\end{equation}
Indeed its \(ii\) entry is
\[
 \sum_{k\ne i}\frac{|\mathbf q_{\ell,k}|^2}
                         {\rho_\ell F_{\ell,i}F_{\ell,k}}
     =F_{\ell,i}^{-1}
                -\frac{|\mathbf q_{\ell,i}|^2}{\rho_\ell F_{\ell,i}^2},
\]
using \(\rho_\ell=\sum_k|\mathbf q_{\ell,k}|^2/F_{\ell,k}\).
Its \(ik\) entry, \(i\ne k\), is
\(-\mathbf q_{\ell,i}\overline{\mathbf q_{\ell,k}}/
(\rho_\ell F_{\ell,i}F_{\ell,k})\), as on the right.
The same computation on the coarse lattice gives
\begin{equation}
 K\mathbf q=0,\qquad DK=1-\mathbf q\pi,\qquad
                 KD=1-D^{-1}\mathbf q\rho^{-1}\mathbf q^\dagger.
 \label{eq:calibration-coarse-projectors}
\end{equation}

The line identity
\(Lr_L(z)(e^{iz}-1)=e^{iLz}-1\) gives
\[
 h_{r+1}(q_\ell)Ld_r(q_\ell)=d_r(p)h_r(q_\ell),
 \quad h_1(q_\ell)\mathbf q_\ell=f(q_\ell)\mathbf q,
 \quad h_2(q_\ell)c_\ell=ch_1(q_\ell).
\]
Together these imply
\begin{align}
 (c_\ell^\dagger H_fc_\ell)\mathcal E_\ell
      &=h_1(q_\ell)^\dagger K,
 &\sum_\ell h_1(q_\ell)\mathcal E_\ell&=1-a^{-1}K,\notag\\
 c_\ell^\dagger\mathcal T_\ell
      &=h_1(q_\ell)^\dagger c^\dagger,
 &\mathcal T_\ell H_cc&=H_fc_\ell\mathcal E_\ell.
 \label{eq:calibration-alias-identities}
\end{align}
Here are the multiplications. In the first identity the
gradient term is killed by \(c_\ell\mathbf q_\ell=0\).
The other term is
\[
 \{1-F_\ell^{-1}\mathbf q_\ell\rho_\ell^{-1}
                     \mathbf q_\ell^\dagger\}
                       h_1(q_\ell)^\dagger K=h_1(q_\ell)^\dagger K,
\]
because \(\mathbf q_\ell^\dagger h_1(q_\ell)^\dagger K
=f(-q_\ell)\mathbf q^\dagger K=0\).
The covariance recursion and the definition of \(u_\ell\) give
\[
 \sum_\ell h_1(q_\ell)F_\ell h_1(q_\ell)^\dagger=D-a^{-1}1,
                      \qquad \sum_\ell f(q_\ell)u_\ell=1.
\]
Thus the second identity is
\((D-a^{-1}1)K+\mathbf q\pi=1-a^{-1}K\).

For the third identity, multiply the first term of
\eqref{eq:calibration-T-alias} by \(c_\ell^\dagger\).
Its gradient projector vanishes since
\(\mathbf q_\ell^\dagger h_1(q_\ell)^\dagger c^\dagger
=f(-q_\ell)\mathbf q^\dagger c^\dagger=0\).
The second term vanishes since
\[
 c^\dagger(1-H_ccDc^\dagger)
   =(1-KD)c^\dagger
   =D^{-1}\mathbf q\rho^{-1}\mathbf q^\dagger c^\dagger=0.
\]
For the fourth identity use
\[
 (1-H_ccDc^\dagger)H_cc=H_cc(1-DK)=H_cc\mathbf q\pi=0.
\]
Its remaining term is
\(H_fc_\ell F_\ell h_1(q_\ell)^\dagger K
=H_fc_\ell\mathcal E_\ell\), since the extra gradient
in \(\mathcal E_\ell\) is killed by \(c_\ell\).
All matrix factors keep their stated order.

\smallskip\noindent
\emph{Cancel the principal poles before estimating.}
Put
\[
 s=\Delta_{LM_v}(p),\qquad q_0=p/L,\qquad
 g_i=G_{v,i}(q_0),\qquad G_i=G_{v',i}(p),
\]
\[
 \tau_i=a^{-1}+\sum_{\ell\ne0}f(q_\ell)f(-q_\ell)
                   r_L(q_{\ell,i})r_L(-q_{\ell,i})F_{\ell,i},
 \quad H=\sum_{\ell\ne0}f(q_\ell)f(-q_\ell)\tr F_\ell.
\]
The symbol \(H\) in this calculation is a scalar alias tail.
Since \(\Delta_{M_v}(q_0)=s/L^2\), the covariance recursion
is exactly
\begin{equation}
 G_i=f(q_0)f(-q_0)r_L(q_{0,i})r_L(-q_{0,i})g_i+s\tau_i.
 \label{eq:calibration-principal-cancellation}
\end{equation}
Define
\[
 A=f(q_0)f(-q_0)\sum_i g_i+sH,\qquad
 W=b_{v'}^{-1}\mathbf q^\dagger\operatorname{diag}_1(G_i^{-1}).
\]
Then \(D^{-1}=s\operatorname{diag}_1(G_i^{-1})\),
\(\rho=s^2b_{v'}\), \(Z=A/s\), and \(\pi=W/s\).
These identities are initially used only away from the pole.

To regularise \(\mathcal E_\ell\), insert
\(K=D^{-1}-D^{-1}\mathbf q\pi\) into its definition:
\[
 \mathcal E_\ell=F_\ell h_1(q_\ell)^\dagger D^{-1}
       +\{u_\ell\mathbf q_\ell
                -F_\ell h_1(q_\ell)^\dagger D^{-1}\mathbf q\}\pi.
\]
The line identity reduces the second term componentwise.
For \(\ell\ne0\) its coefficient of \(\mathbf q_{\ell,i}W\) is
\[
 \gamma_{\ell,i}=f(-q_\ell)
       \left(\frac{\tr F_\ell}{A}
        -\frac{F_{\ell,i}r_L(q_{\ell,i})r_L(-q_{\ell,i})}{G_i}\right).
\]
For \(\ell=0\), \(F_0=s^{-1}\operatorname{diag}_1(g_i)\),
so its apparently singular coefficient is
\begin{align*}
 &\frac{f(-q_0)}s
        \left(\frac{\sum_i g_i}{A}
                -\frac{g_i r_L(q_{0,i})r_L(-q_{0,i})}{G_i}\right)\\
 &\quad=\frac{f(-q_0)}{s f(q_0)f(-q_0)}
                    \{(1-sH/A)-(1-s\tau_i/G_i)\}\\
 &\quad=f(q_0)^{-1}(\tau_i/G_i-H/A)=:\gamma_{0,i}.
\end{align*}
Consequently the regular formulas are
\begin{align}
 \mathcal E_0&=\operatorname{diag}_1(g_i)h_1(q_0)^\dagger
                    \operatorname{diag}_1(G_i^{-1})
                         +\operatorname{diag}_1(\gamma_{0,i})\mathbf q_0W,
                                                        \notag\\
 \mathcal E_\ell&=sF_\ell h_1(q_\ell)^\dagger
                    \operatorname{diag}_1(G_i^{-1})
                         +\operatorname{diag}_1(\gamma_{\ell,i})\mathbf q_\ell W,
                                      \qquad\ell\ne0.
 \label{eq:calibration-filled-E-column}
\end{align}
Neither right side divides by \(s\).

For the dual column put
\[
 B=h_2(q_0)^{-1},\quad x_i=\tau_i/G_i,\quad
 Y=f(q_0)f(-q_0)b_{v'}/b_v(q_0),\quad
 G_1=\operatorname{diag}_1(G_i),\quad\tau_1=\operatorname{diag}_1(\tau_i),
\]
and let \(X_2\) have entry \(x_i+x_k\) at type \(\{i,k\}\).
In the principal alias,
\(\rho_0=s^2b_v(q_0)\), \(y_0=Y\), and
\(c_0=Bch_1(q_0)\). Also
\(h_1(q_0)\operatorname{diag}_1(g_i)h_1(q_0)^\dagger
=G_1-s\tau_1\). Hence its two original terms are
\begin{align*}
 H_fc_0F_0h_1(q_0)^\dagger c^\dagger
      &=s^{-1}H_fBcG_1c^\dagger-H_fBc\tau_1c^\dagger,\\
 y_0h_2(q_0)^\dagger(1-H_ccDc^\dagger)
      &=Yh_2(q_0)^\dagger
                       -s^{-1}Yh_2(q_0)^\dagger H_ccG_1c^\dagger.
\end{align*}
Their only pole numerator is \(H_fB-Yh_2(q_0)^\dagger H_c\).
It is diagonal on two-form types. Fixing one type \(\{i,k\}\),
its two diagonal entries are
\[
 (H_fB)_{ik}
   =\frac1{b_v(q_0)g_ig_k f(q_0)r_L(q_{0,i})r_L(q_{0,k})},
\]
\[
 (Yh_2(q_0)^\dagger H_c)_{ik}
   =\frac{f(q_0)f(-q_0)^2r_L(-q_{0,i})r_L(-q_{0,k})}
                                                   {b_v(q_0)G_iG_k}.
\]
The ratio of the second to the first is
\[
 \frac{(G_i-s\tau_i)(G_k-s\tau_k)}{G_iG_k}
                        =(1-sx_i)(1-sx_k).
\]
This calculation covers every type \(i<k\). Since
\(1-(1-sx_i)(1-sx_k)=s(x_i+x_k-sx_ix_k)\), it gives
\[
 H_fB-Yh_2(q_0)^\dagger H_c
          =sH_f\{X_2-s\operatorname{diag}_2(x_i)\}B.
\]
Cancel this common factor before estimating. The resulting
regular formulas, including the nonprincipal aliases, are
\begin{align}
 \mathcal T_0={}&Yh_2(q_0)^\dagger+
       [H_f\{X_2-s\operatorname{diag}_2(x_i)\}BcG_1
                               -H_fBc\tau_1]c^\dagger,\notag\\
 \mathcal T_\ell={}&H_fc_\ell F_\ell h_1(q_\ell)^\dagger c^\dagger
  +\frac{f(q_\ell)f(-q_\ell)b_{v'}}{\rho_\ell}
         h_2(q_\ell)^\dagger(s^2 1-sH_ccG_1c^\dagger),
                                      \qquad\ell\ne0.
 \label{eq:calibration-filled-dual}
\end{align}
In the first line \(H_f=H_v(q_0)\). No diagonal factor has
been commuted through \(c\) or \(c^\dagger\).

\smallskip\noindent
\emph{One common alias bound for every history.}
Take the common enlarged complex cube supplied by the filled
metric bounds in C.1. The fine and coarse filled metrics,
their inverses, and all principal filter inverses are bounded
there by fixed adjacent-profile constants. The quantity
\(s=\Delta_{LM_v}(p)\) is bounded uniformly because \(p\)
stays in this fixed cube. For \(\ell\ne0\), the elementary
sine bounds give
\begin{equation}
 |F_{\ell,i}|\le C_j(1+|\ell|_\infty)^{-2},\qquad
 |\rho_\ell^{-1}|\le C_j(1+|\ell|_\infty)^{-4},\qquad
 |\mathbf q_\ell|+\|c_\ell\|\le C(1+|\ell|_\infty).
 \label{eq:calibration-nonprincipal-prices}
\end{equation}
For completeness, on real momenta
\(L^2\Delta_{M_v}(q_\ell)\ge c(1+|\ell|_\infty)^2\).
An imaginary displacement of \(p\) of size \(\delta\)
changes this quantity by at most
\(C\delta(1+|\ell|_\infty)+C\delta^2\).
A fixed small \(\delta\) therefore preserves its positive
real part, uniformly in \(L,M_v,\ell\ne0\). Now use
\(F_{\ell,i}=L^{-2}G_{v,i}(q_\ell)/\Delta_{M_v}(q_\ell)\)
and
\(\rho_\ell=L^4\Delta_{M_v}(q_\ell)^2b_v(q_\ell)\)
to obtain the first two bounds. The last follows directly
from \(|e^{iz}-1|\le |z|e^{|\Im z|}\).

The filter estimate gives \(\|h_r(q_\ell)\|\le C_jw_j(\ell)\)
for \(r=0,1,2\). In particular the nonprincipal sums in
\(\tau_i,H\) are bounded by a constant times
\(\sum_{\ell\ne0}w_j(\ell)^2(1+|\ell|_\infty)^{-2}\),
which is bounded independently of the finite alias set.
On real momenta
\(A\ge f(q_0)f(-q_0)\sum_i g_i\ge c_j>0\).
Its uniform holomorphic bound and Cauchy's derivative
estimate give a common smaller cube with bounded \(A^{-1}\),
as for the filled denominators in C.1. Only finitely many
such width reductions are made, independently of history
depth. Every term of the principal formulas is now bounded
by \(C_j\).

For \(\ell\ne0\), the first term of
\eqref{eq:calibration-filled-E-column} has price
\(C_jw_j(\ell)(1+|\ell|_\infty)^{-2}\).
In its second term,
\(|\gamma_{\ell,i}|\le C_jw_j(\ell)(1+|\ell|_\infty)^{-2}\),
which is multiplied by at most \(C(1+|\ell|_\infty)\)
from \(\mathbf q_\ell\). In
\eqref{eq:calibration-filled-dual}, the first term has price
\(C_j(1+|\ell|_\infty)(1+|\ell|_\infty)^{-2}w_j(\ell)\).
The second has price
\(C_jw_j(\ell)^2(1+|\ell|_\infty)^{-4}\), using a bounded
price for its last two matrix factors. Thus
\begin{equation}
                 \|\mathcal E_\ell\|+\|\mathcal T_\ell\|
                                                   \le C_jw_j(\ell).
 \label{eq:calibration-column-alias-price}
\end{equation}
The constants and cube width were chosen before \(L\) and
do not grow with history depth. The regular formulas also
give the harmonic values
\begin{equation}
 \mathcal E_0(0)=\mathcal T_0(0)=1,\qquad
           \mathcal E_\ell(0)=\mathcal T_\ell(0)=0\quad(\ell\ne0).
 \label{eq:calibration-harmonic-columns}
\end{equation}
Indeed \(g_i(0)=G_i(0)=b_v(0)=b_{v'}(0)=1\),
\(f(0)=r_L(0)=1\), and \(c(0)=\mathbf q(0)=s(0)=0\).
The identities \eqref{eq:calibration-alias-identities}
therefore continue across zero. This assigns finite column
and kinetic values, without assigning a covariance to a pole.
These computations expand \cite[Sections~21.1--21.2,
Lemma~21.1 and Lemma~21.2, pp.~107--110]{Main}.

\subsection{Installation, differences and closure of the history}
\label{sec:calibration-installation}

The columns in C.2 are dimensionless alias matrices. To obtain
operators in the counting inner products, we calculate the
normalisation on one Fourier mode. Throughout this calculation
\(Q\) is the scalar line-and-block average of
Theorem~\ref{thm:calibrated-row-data}. The correction from
\(Q\) to the actual framed derivative \(Q_0\) is retained
there in \eqref{eq:calibrated-frame-correction}.

\smallskip\noindent
\emph{A single mode fixes the counting factors.}
Let the coarse and fine periods be \(c\) and \(Lc\), and use
the centred cube \(\mathcal A_L\) also for fine residues
\(r\). Every fine site has a unique expression \(LX+r\)
with \(X\in\Lambda_c\), modulo its period. The counting
Fourier transforms and inverses are
\begin{equation}
 \begin{split}
 \widehat u(k)&=\sum_{x\in\Lambda_{Lc}}e^{-ik\cdot x}u(x),
 \qquad u(x)=(Lc)^{-4}\sum_k e^{ik\cdot x}\widehat u(k),\\
 \widehat v(p)&=\sum_{X\in\Lambda_c}e^{-ip\cdot X}v(X),
 \qquad v(X)=c^{-4}\sum_p e^{ip\cdot X}\widehat v(p).
 \end{split}
 \label{eq:calibration-counting-Fourier}
\end{equation}
The sums run over the corresponding finite dual tori and over
cell types componentwise. For one fixed coarse frequency \(p\),
write a fine field in its alias sector as
\[
 u(x)=\sum_{\ell\in\mathcal A_L}e^{iq_\ell\cdot x}u_\ell,
                         \qquad v(X)=e^{ip\cdot X}v_0.
\]
Thus \(\widehat u(q_\ell)=(Lc)^4u_\ell\) and
\(\widehat v(p)=c^4v_0\). The vectors \(u_\ell,v_0\)
are plane-wave amplitudes, not counting Fourier coefficients.

Let \(w=\nu_L^{*j}\) be the original centred block weights.
Their symbol is \(\sum_\zeta w(\zeta)e^{ik\cdot\zeta}=f(k)\).
The line sum has \(L\) terms, so on a single fine mode
\begin{align*}
 Q_i(e^{ik\cdot x}a_i)(X)
 &=\sum_\zeta w(\zeta)\sum_{\nu=0}^{L-1}
                   e^{ik\cdot(LX+\zeta+\nu e_i)}a_i\\
 &=L e^{iLk\cdot X}f(k)r_L(k_i)a_i.
\end{align*}
For \(k=q_\ell\), \(e^{iLk\cdot X}=e^{ip\cdot X}\).
Therefore the alias row of \(Q\) on amplitudes is
\(Lh_1(q_\ell)\). Equivalently,
\begin{equation}
 \widehat{Qu}(p)=L^{-3}\sum_{\ell\in\mathcal A_L}
                                    h_1(q_\ell)\widehat u(q_\ell).
 \label{eq:calibration-Q-counting-row}
\end{equation}

To obtain its counting transpose, use orthogonality of the
fine alias waves. Their Gram factor is \((Lc)^4\), while
the coarse wave has Gram factor \(c^4\). If \(t_\ell\)
are the amplitudes of \(Q^{\mathsf T}v\), the adjoint identity is
\[
 c^4\sum_\ell u_\ell^*Lh_1(q_\ell)^\dagger v_0
                       =(Lc)^4\sum_\ell u_\ell^*t_\ell.
\]
Since the input amplitudes are arbitrary,
\begin{equation}
 t_\ell=L^{-3}h_1(q_\ell)^\dagger v_0,
       \qquad \widehat{Q^{\mathsf T}v}(q_\ell)
                              =Lh_1(q_\ell)^\dagger\widehat v(p).
 \label{eq:calibration-Q-transpose-factors}
\end{equation}
The star in this real-momentum calculation is the Hermitian
inner product on Fourier amplitudes. It gives the analytic
dagger already defined in C.1 when the symbols are continued.

The fine curl has symbol \(d_1(q_\ell)=L^{-1}c_\ell\),
and its kinetic matrix has symbol
\(L^{-2}c_\ell^\dagger H_fc_\ell\). Thus its product
with the first column in
\eqref{eq:calibration-alias-identities} has exactly the
counting-transpose factor \(L^{-3}\) when the fine output
amplitude is \(L^{-1}\mathcal E_\ell v_0\).
For the dual column the fourth alias identity is then
consistent with the fine amplitude \(L^{-2}\mathcal T_\ell v_0\);
the third identity is checked independently below.
Install them as
\begin{equation}
 E_r(p)=L^{-1}\sum_\ell e^{iq_\ell\cdot r}\mathcal E_\ell(p),
 \qquad T_r(p)=L^{-2}\sum_\ell e^{iq_\ell\cdot r}\mathcal T_\ell(p).
 \label{eq:calibration-counting-columns}
\end{equation}
The two representations are summarised by the following
alias entries:
\begin{equation}
 \begin{array}{c|c|c}
  &\text{plane-wave amplitudes}&\text{counting Fourier coefficients}\\ \hline
  Q&Lh_1&L^{-3}h_1\\
  Q^{\mathsf T}&L^{-3}h_1^\dagger&Lh_1^\dagger\\
  E&L^{-1}\mathcal E_\ell&L^3\mathcal E_\ell\\
  T&L^{-2}\mathcal T_\ell&L^2\mathcal T_\ell
 \end{array}
 \label{eq:calibration-alias-normalisation-table}
\end{equation}
Here \(Q\) has an alias row and the other three entries
are columns; \(h_1\) is evaluated at \(q_\ell\).
The factors follow from \eqref{eq:calibration-counting-Fourier};
no measure or Gaussian normaliser is being changed.

\smallskip\noindent
\emph{Check all four installed identities with these factors.}
Use amplitude entries. For the first identity,
\[
 (K_fE)_\ell=L^{-3}(c_\ell^\dagger H_fc_\ell)\mathcal E_\ell
               =L^{-3}h_1(q_\ell)^\dagger K=(Q^{\mathsf T}K_c)_\ell.
\]
For the second,
\[
 QE=\sum_\ell (Lh_1(q_\ell))(L^{-1}\mathcal E_\ell)
                                      =1-a^{-1}K_c.
\]
The counting transpose of \(T\) has alias row
\(L^4(L^{-2}\mathcal T_\ell)^\dagger=L^2\mathcal T_\ell^\dagger\).
The third alias identity, after taking its dagger, gives
\[
 (T^{\mathsf T}d_f)_\ell=L^2\mathcal T_\ell^\dagger L^{-1}c_\ell
                   =Lch_1(q_\ell)=(d_cQ)_\ell.
\]
Finally
\[
 (TH_cd_c)_\ell=L^{-2}\mathcal T_\ell H_cc
                    =L^{-2}H_fc_\ell\mathcal E_\ell
                                      =(H_fd_fE)_\ell.
\]
On the fine and coarse lattices the positive extensions
preserve \(\widetilde H_fd_f=H_fd_f\) and
\(\widetilde H_cd_c=H_cd_c\). The kinetic matrices are
therefore unchanged, and these are exactly the four
identities \eqref{eq:calibrated-installed-columns}.
The harmonic values \eqref{eq:calibration-harmonic-columns}
give \(E_r(0)=L^{-1}1\), \(T_r(0)=L^{-2}1\), and extend
the checks to \(p=0\). The finite Fourier transform is
invertible, so the identities hold on every compatible
pair of finite tori.

\smallskip\noindent
\emph{One phase identity includes every fine boundary crossing.}
Here the fine forward difference is in lattice units:
\(\nabla_{f,i}u(x)=u(x+e_i)-u(x)\). It does not include
a physical mesh factor. For \(r_i<(L-1)/2\), a fine step
changes \(r\) to \(r+e_i\) and multiplies an alias wave
by \(e^{iq_{\ell,i}}\). At the upper boundary write
\(r'=r+e_i-Le_i\). The new coarse base point is \(X+e_i\),
and its phase is
\begin{equation}
 e^{ip\cdot(X+e_i)}e^{iq_\ell\cdot r'}
  =e^{ip\cdot X}e^{iq_\ell\cdot r}
        e^{iq_{\ell,i}}e^{ip_i-iLq_{\ell,i}}
  =e^{ip\cdot X}e^{iq_\ell\cdot r}e^{iq_{\ell,i}}.
 \label{eq:calibration-residue-boundary-phase}
\end{equation}
The last equality is \(Lq_{\ell,i}=p_i+2\pi\ell_i\).
Thus both cases have the same multiplier
\(q_i(q_\ell)=e^{iq_{\ell,i}}-1\). Iteration gives, for
a multi-index \(\alpha=(\alpha_1,\ldots,\alpha_4)\),
\begin{align}
 (\nabla_f^\alpha E)_r(p)
   &=L^{-1}\sum_\ell e^{iq_\ell\cdot r}
                   \prod_iq_i(q_\ell)^{\alpha_i}\mathcal E_\ell(p),
                                                        \notag\\
 (\nabla_f^\alpha T)_r(p)
   &=L^{-2}\sum_\ell e^{iq_\ell\cdot r}
                   \prod_iq_i(q_\ell)^{\alpha_i}\mathcal T_\ell(p).
 \label{eq:calibration-exact-fine-differences}
\end{align}
This remains valid if successive fine steps cross several
block boundaries.

On the common complex cube of width \(\delta\), centred
residues satisfy \(|e^{iq_\ell\cdot r}|\le e^{2\delta}\).
Also
\[
 |q_i(q_\ell)|\le C L^{-1}(1+|\ell_i|).
\]
By \eqref{eq:calibration-column-alias-price}, the absolute
alias sum for \(\alpha\) is therefore bounded by
\begin{align}
 &C_jC^{|\alpha|}L^{-|\alpha|}
       \sum_{\ell\in\mathcal A_L}
                    \prod_i(1+|\ell_i|)^{-j+\alpha_i}\notag\\
 &\quad\le C_jC^{|\alpha|}L^{-|\alpha|}
                    \prod_{i=1}^4\sum_{n\in\Z}(1+|n|)^{-j+\alpha_i}.
 \label{eq:calibration-factorised-alias-sum}
\end{align}
If \(|\alpha|\le j-2\), every exponent \(j-\alpha_i\)
is at least two. Each one-dimensional sum is at most
\(1+2\int_0^\infty(1+t)^{-2}\,\dd t=3\).
The product is at most \(3^4\), independently of \(L\).
For fixed \(j\) there are only finitely many such
\(\alpha\), so their remaining constants have one common
bound. We have proved the uniform symbol estimate
\begin{equation}
 \sup_{r,p}\left(
 L^{1+|\alpha|}\|(\nabla_f^\alpha E)_r(p)\|
 +L^{2+|\alpha|}\|(\nabla_f^\alpha T)_r(p)\|\right)
                                                    \le C_j,
                      \qquad |\alpha|\le j-2.
 \label{eq:calibration-difference-symbol-bound}
\end{equation}
The supremum in \(p\) is on that same complex cube.
In particular the \(L^4\) possible aliases have been
summed with their decay, rather than counted at their
largest value. Output residues are maximised, not summed.

\smallskip\noindent
\emph{From the common cube to the coefficient norm.}
When \(p\) is translated by \(2\pi e_i\), relabel the
aliases. An alias leaving \(\mathcal A_L\) is replaced
by its representative modulo \(L\), changing its fine
momentum by \(2\pi e_i\). All fine symbols are periodic,
and \(e^{i2\pi e_i\cdot r}=1\) for integer residues.
The original formulas therefore give periodic agreement
on real overlaps away from zero. The filled formulas
extend the agreement holomorphically through zero.
Reality follows by \((p,\ell)\mapsto(-p,-\ell)\).
Thus every symbol in
\eqref{eq:calibration-exact-fine-differences} is a real-kernel,
periodic holomorphic symbol.

For an entry of \(E_r\), use the convention
\(E_r(p)=\sum_{u\in\Z^4}k_r(u)e^{ip\cdot u}\), so
\((Ev)(LX+r)=\sum_u k_r(u)v(X+u)\). Its coefficient is
\[
 k_r(u)=(2\pi)^{-4}\int_{[-\pi,\pi]^4}
                                  e^{-ip\cdot u}E_r(p)\,\dd p.
\]
Choose \(0<\sigma<\delta'<\delta\) once, and shift the
\(i\)-th contour to \(\Im p_i=-\delta'\operatorname{sgn}u_i\).
Periodicity cancels the side faces. The bound
\eqref{eq:calibration-difference-symbol-bound} gives for
the coefficients of its \(\alpha\)-th fine difference
\[
 |k_r^\alpha(u)|\le C_jL^{-1-|\alpha|}e^{-\delta'|u|_1}.
\]
For \(T\) the factor is \(L^{-2-|\alpha|}\). Summing
the fixed cell types and then the four geometric series,
\[
 \sum_{u\in\Z^4}e^{-(\delta'-\sigma)|u|_1}
      =\left(\frac{1+e^{-(\delta'-\sigma)}}
                    {1-e^{-(\delta'-\sigma)}}\right)^4,
\]
proves, in the coefficient norm used in Section~\ref{sec:calibrated-rows},
\begin{equation}
 |\nabla_f^\alpha E|_\sigma\le C_jL^{-1-|\alpha|},\qquad
 |\nabla_f^\alpha T|_\sigma\le C_jL^{-2-|\alpha|},
                              \qquad 0\le|\alpha|\le j-2.
 \label{eq:calibration-column-differences}
\end{equation}
For example \(\delta'=\delta/2\), \(\sigma=\delta/4\)
work for every history and side. Periodising these absolute
coefficient sums cannot increase the corresponding row
bound: each periodic displacement is no longer than its
native representative. This also proves uniformity in
the finite periods.

\smallskip\noindent
\emph{Why the constants do not accumulate with the history.}
All bounds used above come from the cumulative formulas
of C.1. For the extra filter masses \(m_\nu\), ordered
from the most recent index decrease backwards,
\[
 m_\nu/M_v\le3^{-\nu},\qquad
 \sum_{\nu\ge0}m_\nu/M_v\le3/2,\qquad
 \sum_{\nu\ge0}(m_\nu/M_v)^2\le9/8.
\]
The noise suffix masses likewise give
\(\sum_{u\in\mathcal P_v}M_u^{-2}\le9/8\).
These geometric bounds give the common filter and filled
metric constants. In C.2 only finitely many bounded
matrix products and inverse denominators occur, on one
common smaller cube. In C.3 the only further sums are the
factorised alias sum \eqref{eq:calibration-factorised-alias-sum}
and the displayed contour geometric series. Each depends
on the fixed profiles, not on the number of steps. This
proves \eqref{eq:calibration-column-differences} directly
for every fine history. No product of one constant per
step enters this estimate.

The installed identities also check that the same Gaussian
recursion has been used. The first two give
\((K_f+aQ^{\mathsf T}Q)E=aQ^{\mathsf T}\). Hence \(Ex\) is stationary for
the nonnegative quadratic energy
\(\frac12u^{\mathsf T}K_fu+\frac a2|x-Qu|^2\), and
\begin{align*}
 E^{\mathsf T}K_fE+a(1-QE)^{\mathsf T}(1-QE)
 &=\{(QE)^{\mathsf T}+(1-QE)^{\mathsf T}\}K_c=K_c.
\end{align*}
More explicitly, the energy at \(u=Ex+w\) is its value
\(\frac12x^{\mathsf T}K_cx\) plus
\(\frac12w^{\mathsf T}(K_f+aQ^{\mathsf T}Q)w\); its cross term vanishes
by the stationary equation. This proves minimality on
the gauge quotient as well. On the actual tree slice,
the already constructed correction \(\widehat E\) has
\(d_f\widehat E=d_fE\) and \(Q_0\widehat E=QE\).
It has the same fine and noise energy. Thus the coarse
kinetic matrix is the actual Gaussian Schur complement.
Starting at the empty history gives the literal Wilson
quadratic form; the covariance recursion appends the
same datum at each step.

The column identities, their counting factors, and the
fine difference bounds are the inputs of
\cite[Proposition~21.3 and Section~21.3, pp.~111--112]{Main}.
The displayed locators for Sections~21.1--21.3 were checked
against the current PDF on 17 November 2026.

For a sequence of indices, first take the larger prices and
smaller widths for the two possible adjacent-index steps.
Choose each allowed side \(L_j\) from these pre-ratio data,
then any required noise reserve and auxiliary width, and only
then its coupling threshold. All constants and polynomial
degrees entering that finite list of threshold conditions
are fixed before the threshold is increased.
A slowly increasing index schedule can meet any finite list of
thresholds at each index: choose integers \(m_j\) with
\(\overline\beta+b_0(m_j-1)\log3\) above that list, and keep
the index \(j\) between successive \(m_j\)'s. This uses the
same schedule for every cutoff. Comparability of the actual and
nominal couplings is still a separate nonlinear induction claim.
The present appendix proves the Gaussian history and calibration
used in the central theorem, not that later nonlinear claim.


\section{Circulation on every positive-likelihood branch}
\label{app:packet-branches}

The packet estimate uses rough configurations. We prove here the
branch-uniform bound needed there, using the actual frame and
prediction laws of Section~\ref{sec:actual-branches}. All costs
are first powers of the normalised chord \(d_G\).

\begin{lemma}[A linear mean-error estimate for synchronisation]
\label{lem:packet-frame-mean}
Let \(W_{ij,\pi}\in G\) be route matrices on a finite index set,
with \(W_{ii,\pi}=1\) and reversal symmetry
\(W_{ji,\bar\pi}=W_{ij,\pi}^{-1}\). If \(h\) minimises the
synchronisation score, then for some \(b\in G\),
\[
 \operatorname{av}_i d_G(h_i b^{-1})
              \le C_G\operatorname{av}_{i,j,\pi}d_G(W_{ij,\pi}).
\]
The constant is independent of both index cardinalities.
\end{lemma}
\begin{proof}
Put \(\epsilon=\operatorname{av}_{i,j,\pi}d_G(W_{ij,\pi})\),
\(\epsilon_i=\operatorname{av}_{j,\pi}d_G(W_{ij,\pi})\), and
\(\bar h=\operatorname{av}_i h_i\). Comparison with the constant
frame and \(d_G^2\le2d_G\) give
\[
 \operatorname{av}_{i,j,\pi}d_G(h_iW_{ij,\pi}h_j^{-1})^2\le2\epsilon.
\]
The triangle inequality therefore gives
\(\operatorname{av}_{i,j}\|h_i-h_j\|_{\rm F}^2\le8\epsilon\).
The variance identity implies
\(\operatorname{av}_i\|h_i-\bar h\|_{\rm F}^2\le4\epsilon\),
and hence \(\operatorname{dist}(\bar h,G)\le2\sqrt\epsilon\).
Rotate by a nearest group point so that \(\pi_G(\bar h)=1\).

Reversal symmetry and coordinate minimality say that \(h_i\)
is nearest to its feedback
\(A_i=\operatorname{av}_{j,\pi}h_jW_{ij,\pi}^{-1}\).
The diagonal term may be included: it adds to the comparison
with a competitor a nonnegative multiple of
\(\|h_i-v\|_{\rm F}^2\). Also
\(\|A_i-\bar h\|_{\rm F}\le\epsilon_i\).
Fix a small \(\eta>0\) for the Lipschitz nearest-point
neighbourhood of Appendix~\ref{app:projection}.
If \(\epsilon<\eta^2/64\) and \(\epsilon_i<\eta/4\),
both feedbacks lie there, so
\(d_G(h_i)\le2\epsilon_i\). At the other sites use
\(d_G(h_i)\le2\le8\epsilon_i/\eta\).
Averaging proves the bound in this case. If
\(\epsilon\ge\eta^2/64\), the trivial bound two proves it
with constant \(128/\eta^2\). Undo the common rotation.
\end{proof}

\begin{proposition}[Macro-box circulation with the actual tags]
\label{prop:packet-macro-circulation}
There is \(C\), chosen before \(L\), such that on every
positive-likelihood branch and every coarse plaquette \(p\),
\begin{equation}
 d_G(Q_t(U)_p)\le CL^{-2}\operatorname{face}_{\mathcal M_p}(U),
 \qquad
 d_G(V_p)\le CL^{-2}\operatorname{face}_{\mathcal M_p}(U)
                    +\sum_{b\in\partial p}d_G(V_bQ_t(U)_b^{-1}).
 \label{eq:packet-macro-bound}
\end{equation}
Here \(\mathcal M_p\) is the box of radius \(3(J+3)L\)
about its fine base point. A fine plaquette belongs to at most
\(C_J\) of these boxes, independently of \(L\).
\end{proposition}
\begin{proof}
On a box of side \(n\), an undirected edge occurs in at most
\(n^5\) coordinate routes with prescribed axis order: its three
transverse positions fix three endpoint coordinates, leaving
three other coordinates and the two longitudinal endpoints.
Thus the sum of route chords over endpoints is at most
\(4|\mathcal Q_R|n\operatorname{edge}(U)\).

Let \(H\) be any gauge on a frame patch. On the minimising
frame branch apply Lemma~\ref{lem:packet-frame-mean} to the
routes of \(U^H\). On fallback use
\[
 d_G(k_zk_{z+e_i}^{-1})
       \le d_G(U^{h^t}_{z,i})+d_G(U^H_{z,i}),
                  \qquad k_z=h_z^tH_z^{-1}.
\]
The route count for this pure gauge and an averaged choice of
one endpoint give a \(c\in G\) with
\[
 \sum_zd_G(k_zc^{-1})
   \le C_G n\{\operatorname{edge}_{\mathcal Q_R}(U^H)
                         +a_{\rm ax}n\operatorname{face}_{\mathcal Q_R}(U)\}.
\]
Indeed positive fallback mass imposes exactly
\eqref{eq:axial-eligible}. This comparison holds for each
retained tag, without averaging away its likelihood.

Choose one axial gauge on \(\mathcal M_p\) by the layer
argument of Lemma~\ref{lem:eligible-tag}. Its edge cost is at
most \(C_JL\operatorname{face}_{\mathcal M_p}(U)\).
For each corner frame choose the preceding constant \(c_v\).
Since \(n\le C_JL\), its summed frame error is at most
\(CL^2\operatorname{face}_{\mathcal M_p}(U)\).
The repeated average has \(\|w\|_\infty\le L^{-4}\):
convolution with a probability law cannot increase the supremum
of the first block average. Its weighted frame error is thus
\(CL^{-2}\operatorname{face}_{\mathcal M_p}(U)\).
For a straight connector each fine edge occurs for at most
\(L\) choices of its starting displacement and position, so
its weighted error is at most
\(L^{-3}\operatorname{edge}_{\mathcal M_p}(U^H)\).
Consequently the transformed connectors obey
\[
 \epsilon_e:=\sum_z w(z)d_G(c_v^{-1}Y_{e,z}c_{v+e_i})
                              \le CL^{-2}\operatorname{face}_{\mathcal M_p}(U).
\]
On the regular predictor branch its nearest group point has
chord at most \(2\epsilon_e\), by comparison with the identity.
On fallback the minimum score is at least \(c_{\rm pr}\);
testing the score at the identity gives
\(c_{\rm pr}\le\frac12\sum_z w(z)d_G(c_v^{-1}Y_{e,z}c_{v+e_i})^2
\le\epsilon_e\). Every selected connector has chord at most
two, hence at most \(2\epsilon_e/c_{\rm pr}\).
The corner constants cancel around the plaquette. Summing its
four edge errors proves the first inequality in
\eqref{eq:packet-macro-bound}. Replace its four predicted bonds
by the actual \(V_b\)'s to obtain the second.
Macro-box centres are spaced by \(L\); their radius is
\(3(J+3)L\). Their overlap count is therefore fixed before
\(L\), and all arguments use nonwrapping local boxes under
the period reserve.
\end{proof}

These are \cite[Lemma~14.5, Proposition~14.6 and
Corollary~14.7, pp.~61--62]{Main}. The proof includes both frame
branches and both predictor branches at their original thresholds.


\begin{thebibliography}{Wilson74}
\raggedright
\bibitem[Main]{Main}
\emph{Reading Notes on Four-Dimensional Yang--Mills Theory},
AI-generated working exposition, maintained by Zexi He,
publication copy of 21 November 2026, based on the research manuscript
snapshot of 17 November 2026.
\href{../4DYM_reading_notes.pdf}{Companion main exposition (PDF)}.
The mathematical text preserves the research snapshot's printed
pagination and numbered statements; the publication information page
is unnumbered. The section, equation and page locators in this rewrite
refer to that text.

\bibitem[DST23]{DST23}
W.~Dybalski, A.~Stottmeister and Y.~Tanimoto,
\emph{Lattice Green functions for pedestrians: Exponential decay},
arXiv:2303.10754v1 (2023).
Section~2.4, Lemma~2.12, p.~14, and Appendix~A, pp.~36--37,
explain the Combes--Thomas argument.
\href{https://arxiv.org/abs/2303.10754v1}{arXiv:2303.10754v1}.

\bibitem[Bal85g]{Bal85g}
T.~Ba{\l}aban,
\emph{Spaces of regular gauge field configurations on a lattice
and gauge fixing conditions},
Communications in Mathematical Physics \textbf{99} (1985), 75--102.
\href{https://doi.org/10.1007/BF01466594}{doi:10.1007/BF01466594}.
Sect.~B, Lemma~1 and (1.24)--(1.26), printed pp.~79--80,
give the axial path comparison used here.

\bibitem[Bal87]{Bal87}
T.~Ba{\l}aban,
\emph{Renormalization group approach to lattice gauge field theories.
I. Generation of effective actions in a small field approximation
and a coupling constant renormalization in four dimensions},
Communications in Mathematical Physics \textbf{109} (1987), 249--301.
\href{https://doi.org/10.1007/BF01215223}{doi:10.1007/BF01215223}.

\bibitem[Bal88a]{Bal88a}
T.~Ba{\l}aban,
\emph{Renormalization group approach to lattice gauge field theories.
II. Cluster expansions},
Communications in Mathematical Physics \textbf{116} (1988), 1--22.
\href{https://doi.org/10.1007/BF01239022}{doi:10.1007/BF01239022}.
Sect.~2, (2.11)--(2.13), printed p.~14, gives the connected
organisation of the gauge fluctuation logarithm.

\bibitem[FP07]{FP07}
R.~Fern\'andez and A.~Procacci,
\emph{Cluster expansion for abstract polymer models.
New bounds from an old approach},
Communications in Mathematical Physics \textbf{274} (2007), 123--140.
\href{https://doi.org/10.1007/s00220-007-0279-2}{doi:10.1007/s00220-007-0279-2}.
Page locators refer to arXiv:math-ph/0605041v2:
Sect.~2, (2.2)--(2.7), pp.~3--4, and Sect.~4.1,
Proposition~5, pp.~10--11.

\bibitem[Wilson74]{Wilson74}
K.~G.~Wilson, \emph{Confinement of quarks},
Physical Review D \textbf{10} (1974), 2445--2459.
\href{https://doi.org/10.1103/PhysRevD.10.2445}{doi:10.1103/PhysRevD.10.2445}.
\end{thebibliography}
\end{document}
